%\documentclass[aps,prl,singlecolumn,superscriptaddress]{revtex4-2}


\documentclass[a4paper,11pt]{article}
\usepackage{./jheppub}
%\setcounter{secnumdepth}{3}
% Useful packages
\usepackage{tikz-cd}
\usepackage{amsmath}
\usepackage{array,booktabs}
\usepackage{hyperref}
%%\usepackage{showlabels}
\usepackage{amsthm}
\newtheorem{theorem}{Theorem}
\newtheorem{lemma}[theorem]{Lemma}
\newtheorem{corollary}[theorem]{Corollary}
\newtheorem{proposition}[theorem]{Proposition}
\theoremstyle{definition}
\newtheorem{definition}{Definition}[section]
\newtheorem{prop}{Conjecture}
\newtheorem{cor}{Corollary}
%-------------------packages-------------------------
\usepackage{enumerate}
\usepackage[utf8]{inputenc}
\usepackage{amssymb,amsbsy,mathtools,amsmath}
\usepackage{braket}
\usepackage{graphicx}
\usepackage{xcolor}
%\usepackage{multirow}
\usepackage{hyperref}
%\usepackage{caption}
%\usepackage{subcaption}
%\mathcal{A}ptionsetup{compatibility=false}
\usepackage{natbib}
%\usepackage{dcolumn}% Align table columns on decimal point
\usepackage{bm}% bold math
\usepackage{soul}
%\usepackage{showlabels}
%\newtheorem{proposition}{Proposition}
\newcommand{\E}{\mathbb{E}}
\newcommand{\dd}{\,\mathrm{d}}
%\newcommand{\mathcal{A}}{\mathcal{A}}
\newcommand{\cH}{\mathcal{H}}
\newcommand{\KLL}{D_{\mathrm{KL}}}
\newcommand{\Mop}{\mathsf{M}}

%Includes "References" in the table of contents
\usepackage[nottoc]{tocbibind}

\usepackage{bbm}
\usepackage{amsthm}
%\newtheorem{theorem}{Theorem}
%\newtheorem{lemma}[theorem]{Lemma}
%\newtheorem{corollary}[theorem]{Corollary}
%\theoremstyle{definition}
%\newtheorem{definition}{Definition}[section]
\usepackage{amssymb}
\newcommand{\id}{\mathbb{I}}
\newcommand{\Tr}{\text{Tr}}
\newcommand{\Div}{\operatorname{div}}
\usepackage{amsmath}
%\usepackage{tikz}
\usepackage{braket}
%\usepackage{quantikz}
\usepackage{graphicx}
\usepackage{caption}
\newcommand{\mA}{\mathcal{A}}
\newcommand{\mB}{\mathcal{B}}
\newcommand{\mP}{\mathcal{P}}
\newcommand{\mR}{\mathcal{R}}
\newcommand{\nn}{\nonumber}
\newcommand{\mL}{\mathcal{L}}
\newcommand{\lb}{\left(}
\newcommand{\rb}{\right)}
\newcommand{\mH}{\mathcal{H}}
\usepackage{subcaption}
\newcommand{\tOO}{\tilde{\mathcal{O}}}
\newcommand{\mO}{\mathcal{O}}
\newcommand{\mK}{\mathcal{K}}
\newcommand{\ep}{\epsilon}
\newcommand{\mbR}{\mathbb{R}}
\newcommand{\mM}{\mathcal{M}}
\newcommand{\mT}{\mathcal{T}}
\newcommand{\NL}[1]{\textcolor{blue}{#1}}
\newcommand{\p}{\partial}
\newcommand{\tr}{\text{tr}}
\newcommand{\tT}{\tilde{T}}
\newcommand{\Ad}{\operatorname{Ad}}
\newcommand{\Dom}{\operatorname{Dom}}
\newcommand{\mE}{\mathcal{E}}
\newcommand{\relcomm}[2]{#1'\cap #2}

\newcommand{\SO}[1]{\textcolor{purple}{#1}}
%\newcommand{\MM}[1]{\textcolor{red}{#1}}
\newcommand{\KL}[1]{\textcolor{brown}{#1}}
\newcommand{\YC}[1]{\textcolor{teal}{#1}}
\newcommand{\cL}{\mathcal{L}}
\setlength{\emergencystretch}{2em}
\begin{document}

% Title
\title{Subsystem Flows and Many-body Ergodicity in QFT}

\author{Yidong Chen$^1$,}
\affiliation{$^1$Department of Mathematics, University of Illinois at Urbana-Champaign, Illinois, IL 61801, USA}
\author{Nima Lashkari$^2$,}
\author{Kwing Lam Leung$^2$}
\affiliation{$^2$Department of Physics and Astronomy, Purdue University, West Lafayette, IN 47907, USA}
\emailAdd{yidongc2@illinois.edu}
\emailAdd{nima@purdue.edu}
\emailAdd{leung60@purdue.edu}

% Date
%\date{\today}

% Abstract
\abstract{Recent developments in the algebraic approach to bulk reconstruction in holography have highlighted the role of half-sided modular inclusions and modular intersections in reconstructing spacetime symmetry transformations. A key ingredient in these algebraic reconstructions is the concept of subsystem flow, which composes modular evolutions associated with an algebra and that of its subsystem. In the first part of this work, we motivate the notion of coarse-grainings in the classical Hamiltonian ergodic hierarchy, distinguishing two cases: subspace and subsystem subalgebras. To identify a classical avatar of modular flow, we associate flows to states, using the Poisson bracket on phase space. We discuss the generalization to the quantum ergodic theory, and in particular to the net of subsystems associated to subregions in many-body quantum systems. We focus on {\it non-expected} and {\it split} inclusions that arise for local algebras in QFT. In the second part of this work, we compute subsystem flows corresponding to causal developments of balls in vacuum CFT and QFT and study their analytic properties. For inclusions of wedges in QFT and ball-shaped regions in CFTs, the subregion flow acts geometrically; the wedge comparison gives a translation whose displacement depends on modular time. We compute the subregion flows in conformal perturbation theory and prove bounds on the comparison of subregion flows of balls with wedge flows in vacuum QFT.}



\maketitle

% \NL{To do: 1) Merge appendix E with the end of section 2; 3) Change the title and content of Appendix C to the mathematical background for Section 5.}

\section{Introduction}


The passage from a reversible microscopic description to an apparently irreversible macroscopic one is often attributed to information loss during coarse-graining. Both in the classical and quantum settings, we have two types of coarse-graining that we call subspace and subalgebra coarse-graining.


The intuition of subspaces is a ``one-particle" system, where the particle is either in a set of states or in another set of states. Consider two separate distinguishable particles, each coming with $d_i$ states. The two-particle system has $d_1d_2$ states, whereas if we consider a one-particle subspace of it, we either have particle one in one of the $d_1$ states, or particle two in one of the $d_2$ states; therefore, in the one-particle setup, the number of states add up.

This information loss can be quantified by associating an entropy to a coarse-graining. A relative phase constructed from the fine-grained and coarse-grained densities generates this entropy and provides a first classical avatar of a coarse-graining flow. 


In modular theory, a faithful normal state $\omega$ on a von Neumann algebra $\mA$ defines the physical modular flow $\operatorname{Ad}\Delta_\omega^{-it}$, the distinguished automorphism group satisfying the upper-strip KMS condition. For a subalgebra $\mathcal B\subset\mA$ with a common cyclic and separating vector representing $\omega$, we define the {\it subalgebra flow} $D_{B\subset A}(t)=\Delta_{\omega;A}^{-it}\Delta_{\omega;B}^{it}$, which evolves backward by the intrinsic subsystem flow and then forward by the outer flow. QFT associates to an inclusion of causally complete regions $B\subset A$ an inclusion of von Neumann algebras $\mA_B\subset\mA_A$. With the vacuum as reference, we call $D_{B\subset A}(t)$ a {\it subregion flow}.

In holography, at the $G_N=0$ limit (infinite $N$), the subregion/subalgebra duality associates to every causally-complete subregion of the bulk spacetime a von Neumann subalgebra on the boundary. The subregion flows play a prominent role in the algebraic reconstruction of the bulk spacetime from the net of boundary large $N$ algebras. 

analytic properties of characterstic function imply the modular chaos bound that has been heavily used in holography \cite{}.

In the first part of this work, we motivate subregion/subalgebra flows from the point of view of classical Hamiltonian ergodic hierarchy. The comparison distinguishes conserved observables, autonomous subsystems, and generating modular orbits; Appendix~\ref{app:classical} reviews the classical ergodic hierarchy.

In the second part of this work, we study the subregion/subalgebra flows in the vacuum of QFT for the inclusion of causal developments of balls.




% \NL{Should we focus on standard inclusions on top of ergodic}

% Standard+maximal chaos (HMSI)  = PSL(2,R)

% Standard+ ergodic (cyclic-separating)

% \NL{Modular Berry phases, $[P_+,P_-]$ commuting with $K$}




\section{Subalgebra flows in quantum ergodic theory}\label{section:basics}

The observables of a physical system form a von Neumann algebra. If there are no classical observables (those that commute with all observables), this algebra is a factor; it has a trivial center. For a finite quantum system, such factors are $d\times d$ complex matrices. Infinite-dimensional factors are classified as type $I_\infty$, type $II_1$, type $II_\infty$, and type $III$ factors.

Modular theory associates to every faithful normal state a dynamical flow
$\sigma_{\omega;t}(a)=\Delta_{\omega;A}^{it}a\Delta_{\omega;A}^{-it}$.
For matrix algebras, it is $\sigma_{\omega;t}(a)=\rho_\omega^{it}a\rho_\omega^{-it}$.
Ergodicity is a property of the algebra together with the chosen state, not of the algebra alone. Type $III_1$ permits ergodic modular dynamics, but not every faithful normal state has a trivial centralizer.

A finite composite system has observable algebra $\mA_A=\mA_B\otimes\mA_R$; more generally, a subsystem is specified by an inclusion $\mA_B\subset\mA_A$. The question we carry from classical dynamics is whether the retained observables evolve autonomously. When the two algebras have a common cyclic and separating vector, we compare their modular evolutions through the {\it subsystem flow}
\begin{align}
 D_{B\subset A}(t)&=\Delta_{\omega;A}^{-it}\Delta_{\omega;B}^{it},\nn\\
 \operatorname{Ad}D_{B\subset A}(t)(b)
 &=\sigma_{\omega;A}^{-t}\bigl(\sigma_{\omega;B}^t(b)\bigr),
 \qquad b\in\mA_B.
\end{align}
We examine the relation between autonomy, conditional expectations, and generating inclusions before specializing to the vacuum subregion algebras of QFT. These notions do not by themselves imply mixing.

Appendix~\ref{app:classical} reviews the classical constructions. Two examples guide the comparison: a classical field region selects a subsystem of field observables, and a stationary marginal state need not give autonomous evolution of those observables. The latter distinction motivates comparing full and reduced state-generated flows.

\subsection{Review of quantum dynamical systems}

In this subsection, we review and fix our notation. Readers familiar with standard definitions can skip this subsection.

\paragraph{Algebra of quantum observables:}
Quantum observables are described by a 
\emph{von Neumann algebra} $\mathcal{A}$ a unital $*$-subalgebra of the bounded operators
$B(\mathcal{H})$ on some separable Hilbert space $\mathcal{H}$, closed in the
weak operator topology.
% The commutativity of the classical algebra is replaced by the requirement that $\mathcal{A}$
% be closed under the adjoint; it need not be abelian.  Indeed the classical theory is recovered
% by taking $\mathcal{A}=L^\infty(X,\Sigma,\mu)$, acting by multiplication on
% $\mathcal{H}=L^2(X,\Sigma,\mu)$.
A von Neumann algebra $\mathcal{A}$ is a \emph{factor} if its center
$\mathcal{Z}(\mathcal{A}) = \mathcal{A}\cap\mathcal{A}'$ consists only of scalar multiples
of the identity, where $\mathcal{A}' = \{b\in B(\mathcal{H}) : [b,a]=0\text{ for all }a\in\mathcal{A}\}$
is the commutant.  Physically, a factor algebra describes a system with no classical
superselection sectors \footnote{As we review in Appendix \ref{app:classical}, in classical physics, the commutator is replaced by the Poisson bracket and the commutant $\mathcal{A}'$ is the quantum analog of the Poisson
centralizer introduced in Appendix \ref{app:classical}. It consists of all observables invisible to $\mathcal{A}$. The classical analog of a factor is a non-degenerate Poisson bracket that restricts to a symplectic leaf.}.

% -------------------------------------------------------------------
\paragraph{Quantum states:}

A \emph{normal state} on $\mathcal{A}$ is a positive, normalized, $\sigma$-weakly continuous
linear functional $\varphi:\mathcal{A}\to\mathbb{C}$, with $\varphi(\id)=1$.
For a finite matrix algebra it has the form $\varphi(a)=\Tr(\rho_\varphi a)$,
where $\rho_\varphi$ is positive with unit trace. Type $III$ algebras have no
intrinsic tracial density matrix; the normal functional is the general definition.

A state $\varphi$ is \emph{faithful} if $\varphi(a^*a)=0$ implies $a=0$.
In finite dimensions this means that its density matrix has no zero eigenvalues.
Faithfulness is the counterpart of a strictly positive classical probability density
and gives a separating vector in the GNS construction.

A state $\varphi$ is \emph{tracial} if $\varphi(ab)=\varphi(ba)$ for all $a,b\in\mathcal{A}$.
On a finite algebra, a faithful normal tracial state plays the role of the reference
Liouville measure. Its modular flow is trivial, despite faithfulness. In the type $III$
algebras relevant to QFT, we use faithful normal states without a trace.

\paragraph{GNS Hilbert space and modular operators.}
Given a von Neumann algebra $\mathcal{A}$ and a faithful normal state $\varphi$, the
Gelfand--Naimark--Segal (GNS) construction produces a Hilbert space $\mathcal{H}_\varphi$,
a cyclic unit vector $|\Omega\rangle\in\mathcal{H}_\varphi$ with
$\varphi(a)=\langle\Omega|a|\Omega\rangle$, and a $*$-representation
$\pi:\mathcal{A}\to B(\mathcal{H}_\varphi)$ \footnote{The classical analog is $\mathcal{H}=L^2(X,\mu)$
with the constant function $|\id\rangle$ as the reference vector.}.

The Tomita--Takesaki Modular theory assigns to the pair $(\mathcal{A},|\Omega\rangle)$ an
antilinear operator $S_\varphi$ defined on the dense domain
$\mathcal{A}|\Omega\rangle$ by
\begin{equation}
  S_\varphi\, a|\Omega\rangle = a^*|\Omega\rangle, \qquad a\in\mathcal{A}.
  \label{eq:Tomita-S}
\end{equation}
Its polar decomposition $S_\varphi = J_\varphi\,\Delta_\varphi^{1/2}$ defines the
\emph{modular conjugation} $J_\varphi$, an antiunitary involution, and the
\emph{modular operator} $\Delta_\varphi$, a positive self-adjoint operator. 

% -------------------------------------------------------------------
\paragraph{Modular dynamics:}
The Tomita--Takesaki theorem~\cite{Takesaki1970} guarantees that the one-parameter group
of unitaries $\Delta_\varphi^{it}$, $t\in\mathbb{R}$, implements an automorphism group of
$\mathcal{A}$:
\begin{equation}
  \sigma_t^\varphi(a) = \Delta_\varphi^{it}\, a\, \Delta_\varphi^{-it},
  \qquad a\in\mathcal{A},\quad t\in\mathbb{R}.
  \label{eq:modular-flow}
\end{equation}
This is the \emph{modular flow} (or \emph{modular automorphism group}) of $\varphi$. 

The modular operator $\Delta_\varphi$ and $J_\varphi$ encode the non-commutativity of
$\mathcal A$. Appendix~\ref{subsec:classical-state-flow} describes a related
state-generated Poisson flow: for a smooth positive density $\rho$, the Hamiltonian
$K_\rho=-\log\rho$ defines a flow preserving that state. This construction uses
the Poisson bracket and differs from the trivial modular automorphism group of
an abelian von Neumann algebra. In both constructions the reference state is stationary:
$\varphi\circ\sigma_t^\varphi=\varphi$.

With the convention \eqref{eq:modular-flow}, the \emph{KMS condition} reads
\begin{equation}
 F_{a,b}(t)=\varphi\bigl(a\,\sigma_t^\varphi(b)\bigr),
 \qquad F_{a,b}(t-i)=\varphi\bigl(\sigma_t^\varphi(b)a\bigr).
 \label{eq:KMS-function}
\end{equation}
The function is bounded and analytic on $-1<\operatorname{Im}z<0$, with continuous
boundary values. Equivalently, $\beta_t^\varphi:=\sigma_{-t}^\varphi
=\operatorname{Ad}\Delta_\varphi^{-it}$ gives the upper-strip condition at inverse
temperature one. Its generator is $K_\varphi=-\log\Delta_\varphi$, while
$\partial_t\sigma_t^\varphi(b)|_{t=0}=i[\log\Delta_\varphi,b]$ on the generator domain.
The finite imaginary-time KMS shift cannot be replaced by this derivative.
The Poisson integration-by-parts identity \eqref{KMSclassical} is the classical
KMS condition; its relation to quantum KMS requires a semiclassical construction
and a specified time convention~\cite{vanDeVen2022}.

\paragraph{Ergodicity in modular dynamics:}

When the algebra $\mA$ has a non-trivial center $Z\subset \mA$, then this center is inside the centralizer of all modular flows:
\begin{align}
    \forall z\in Z,\forall t,\omega:\qquad \Delta_{\omega}^{it}z \Delta_\omega^{-it}=z
\end{align}
This means that these observables are conserved charges, and an obstruction to ergodicity in modular dynamics of arbitrary states. That is why, in quantum ergodic theory, it is natural to restrict to algebras with trivial center, i.e., factors.

For a fixed faithful normal state, modular ergodicity means that the centralizer is trivial:
\begin{align}
    \label{eq:modular-ergodicity}
    \mA^\omega\equiv \{a\in \mA:\Delta_{\omega;A}^{it}a\Delta_{\omega;A}^{-it}=a\qquad \forall t\}=\mathbb{C} 1
\end{align}

For a nontrivial factor, this requires type $III_1$, but it does not follow from the
type alone. For example, consider an infinite tensor product of identical three-level
systems with faithful site density matrix
$\rho=\sum_{i=1}^3p_i\ket{i}\bra{i}$.
If $\log(p_1/p_3)$ and $\log(p_2/p_3)$ are rationally independent, the resulting
Araki--Woods factor is type $III_1$~\cite{ArakiWoods1968}.
Nevertheless, the site spectral projections $\ket{i}\bra{i}_n$ belong to the
centralizer of the product state. Ergodicity, generation, and decay of connected
correlations must therefore be stated separately.

\subsection{Subsystem quantum coarse-grainings}
\label{subsec:quantum-coarse}


\paragraph{Quantum coarse-grainings.}
An inclusion $\mathcal B\subseteq\mathcal A$ specifies the retained observables.
When it exists, a \emph{normal conditional expectation}
$\mE_{\mathcal B}:\mathcal A\to\mathcal B$ is a completely positive, unital projection
with the bimodule property
\begin{equation}
 \mE_{\mathcal B}(bac)=b\,\mE_{\mathcal B}(a)\,c,
 \qquad a\in\mathcal A,\quad b,c\in\mathcal B.
 \label{eq:CE-bimodular}
\end{equation}
It specifies how a state on $\mathcal B$ is extended to $\mathcal A$.
Unlike the classical conditional expectation in Appendix~\ref{app:classical},
such a normal expectation need not exist. On the GNS Hilbert space, let $e_{\mathcal B}$
be the orthogonal projection onto $\overline{\mathcal B\Omega}$.
If $\mE_{\mathcal B}$ preserves the reference state $\varphi$, then
$e_{\mathcal B}a\Omega=\mE_{\mathcal B}(a)\Omega$; this is its Jones projection.
Without state preservation, the Hilbert-space projection does not represent the expectation.

\paragraph{Entropy of a coarse-graining.}
Fix a faithful normal expectation $E:\mathcal A\to\mathcal B$, independently of the
faithful normal state $\varphi$ being examined. The coarse-grained state is
$\varphi_0=\varphi|_{\mathcal B}\circ E=\varphi\circ E$.
Represent the states by their vectors $\Omega_\varphi,\Omega_{\varphi_0}$ in standard form.
The relative Tomita operator is the closure of the map below; its positive quadratic form is
\begin{align}
 S_{\varphi_0|\varphi;A}a\Omega_\varphi
 &=a^*\Omega_{\varphi_0},\qquad
 \Delta_{\varphi_0|\varphi;A}=S_{\varphi_0|\varphi;A}^\dagger S_{\varphi_0|\varphi;A},
 \nn\\
 \langle\Delta_{\varphi_0|\varphi}^{1/2}a\Omega_\varphi|
 \Delta_{\varphi_0|\varphi}^{1/2}b\Omega_\varphi\rangle
 &=\varphi_0(ba^*),\qquad a,b\in\mathcal A.
 \label{eq:relative-modular}
\end{align}
The \emph{quantum relative entropy} is
\begin{align}
 S(\varphi\|\varphi_0)
 &=-\langle\Omega_\varphi|\log\Delta_{\varphi_0|\varphi;A}|\Omega_\varphi\rangle
 \nn\\
 &=i\left.\frac{d}{dt}
 \langle\Omega_\varphi|\Delta_{\varphi_0|\varphi;A}^{it}|\Omega_\varphi\rangle
 \right|_{t=0},
 \label{eq:relative-entropy}
\end{align}
with the derivative understood when the logarithmic spectral moment is integrable.
In the matrix representation, $\Delta_{\varphi_0|\varphi}=L_{\rho_0}R_{\rho}^{-1}$.
The classical counterpart is the ratio $\rho_0/\rho$, giving
\begin{equation}
 S(\rho\|\rho_0)=\int_X\rho\log(\rho/\rho_0)\,d\mu.
 \label{eq:classical-KL}
\end{equation}
Relative entropy is nonnegative and vanishes precisely when $\varphi=\varphi\circ E$.
In particular, an expectation preserving $\varphi$ gives zero coarse-graining entropy.
The identity $S(\rho\|\rho_0)=H(\rho_0)-H(\rho)$ holds for the reference-measure-preserving
classical coarse-graining of Appendix~\ref{app:classical}, with finite entropies.
Its quantum entropy-difference counterpart requires a trace-preserving expectation.


\paragraph{Tensor-product subsystems:}
For a finite composite system, $\mH_A=\mH_B\otimes\mH_R$ and
$\mA_A=\mA_B\otimes\mA_R$. Retaining $\mA_B\otimes1_R$ gives the expectations
\begin{equation}\label{condexp}
 E_{\chi_R}(b\otimes r)=b\,\chi_R(r),
 \qquad b\in\mA_B,\quad r\in\mA_R.
\end{equation}
For this fixed expectation, $\varphi=\varphi\circ E_{\chi_R}$ precisely when
$\rho_{BR}=\rho_B\otimes\chi_R$. Thus $I_\rho(B:R)=0$ is necessary, but the hidden
marginal must also equal the chosen reference. More generally,
\begin{equation}\label{eq:quantum-coarse-entropy-decomposition}
 S(\rho_{BR}\|\rho_B\otimes\chi_R)
 =I_\rho(B:R)+S(\rho_R\|\chi_R).
\end{equation}
The two terms measure discarded correlations and discarded information about the
hidden distribution, just as in Appendix~\ref{app:classical}.
The subspace and mixed finite-dimensional inclusions are reviewed in
Appendix~\ref{app:finite-coarse}. Our main examples are subfactors, including
local QFT algebras for which this tensor factorization need not hold.

\paragraph{Dynamics reducible to a subalgebra:} Given a group action $\alpha:G\to\operatorname{Aut}(\mA)$, we say that the dynamics is reducible to $\mB\subset\mA$ if 
\begin{align}
    \forall g\in G,b\in\mB:\qquad \alpha_g(b)\in \mB\ . 
\end{align}
Then, the dynamics can be {\it reduced} to $\mB$. In dynamical systems, this is sometimes called {\it autonomous dynamics}, as well.

\paragraph{Irreducible dynamics:} Given a general dynamical system, we say it is irreducible if there is no non-trivial proper subalgebra $\mB\subset \mA$ with reducible dynamics.


\paragraph{Modular dynamics and quantum coarse-graining:}
For a faithful normal state $\omega$, autonomy of $\mB\subset\mA$ means
\begin{equation}\label{eq:quantum-coarse-autonomy}
 \mB_t:=\sigma_{\omega;A}^t(\mB)=\mB,\qquad t\in\mathbb R.
\end{equation}
Takesaki's theorem~\cite{Takesaki1972} makes this equivalent to the existence of
an $\omega$-preserving normal expectation $E_\omega:\mA\to\mB$. In that case,
\begin{align}
 \sigma_{\omega;A}^t(b)&=\sigma_{\omega;B}^t(b),\qquad b\in\mB,\nn\\
 \Delta_{\omega;A}^{it}V&=V\Delta_{\omega;B}^{it},
 \qquad V:\overline{\mB\Omega}\hookrightarrow\mH_A.
\end{align}
The intrinsic modular operator of $\mB$ acts on $\overline{\mB\Omega}$.
In the tensor-product example, autonomy is equivalent to a product state
$\omega_B\otimes\omega_R$; $E_{\omega_R}$ preserves it.
For a general inclusion, absence of correlations with $\mB'\cap\mA$ is not sufficient:
that relative commutant need not account for every discarded observable.
Appendix~\ref{app:compressed-comparison} treats the comparison compressed to
$\overline{\mB\Omega}$, separately from the characteristic operator below.

The classical Poisson analogy concerns state-generated derivations on a smooth
Poisson subalgebra. It does not identify the nonlinear assignment of a modular
Hamiltonian to a vector state with a bilinear Poisson bracket.

\begin{definition}[Generating inclusion]
We call the inclusion {\it generating} under the chosen modular flow when 
\begin{align}\label{geneating}
    \vee_{t\in \mathbb{R}}\mB_t=\mA\ .
\end{align}
\end{definition}


\paragraph{Common cyclic and separating vector:}
We now assume that $\Omega$ is cyclic and separating for both $\mB$ and $\mA$.
Then $\overline{\mB\Omega}=\mH$ and $e_B=1$. This density of vectors does not identify
the two observable algebras or imply zero information loss. In the classical
representation $L^2(X,\mu)$, a proper coarse algebra instead generates a proper
subspace from the constant reference vector. QFT vacuum cyclicity supplies the
different Hilbert-space structure used in Sections~\ref{sec:geometric-subregion-flows}--\ref{section:approx}.

We are now ready to introduce the characteristic function of the inclusion $\mathcal{B}\subset\mathcal{A}$\cite{revolutionizing}:
\begin{definition}
    Given a pair of von Neumann algebras $\mathcal{B}\subset\mathcal{A}$ and a common cyclic and separating vector $\ket{\Omega}\in\mathcal{H}$, the characteristic function is defined as:
    \begin{equation}
        D_{B\subset A}(z):=\Delta_A^{-iz}\Delta_B^{iz},\qquad z\in\mathbb R,\quad K_C=-\log\Delta_C.
    \end{equation}
\end{definition}
For $b\in\mB$, its adjoint action and the transported algebra are
\begin{equation}\label{eq:characteristic-algebra-transport}
 \operatorname{Ad}D_{B\subset A}(t)(b)=\sigma_A^{-t}\bigl(\sigma_B^t(b)\bigr),
 \qquad \operatorname{Ad}D_{B\subset A}(t)(\mB)=\sigma_A^{-t}(\mB).
\end{equation}
This compares the intrinsic subsystem motion with the outer modular motion.
It requires no conditional expectation. In general $D(t)$ is a cocycle rather than
a one-parameter group; it need not belong to $\mA$.
For a normalized vector $\Psi$ on which the logarithmic matrix elements and the
four relative entropies below are finite, it also compares information accessible
in the two algebras:
\begin{align}
 -i\left.\partial_t\langle\Psi|D_{B\subset A}(t)|\Psi\rangle\right|_{t=0}
 &=\langle\Psi|K_A-K_B|\Psi\rangle\nonumber\\
 &=S_A(\Psi\|\Omega)-S_B(\Psi\|\Omega)
   +S_{B'}(\Psi\|\Omega)-S_{A'}(\Psi\|\Omega)\geq0.
 \label{eq:characteristic-entropy-comparison}
\end{align}
The inequality uses monotonicity under $\mB\subset\mA$ and $\mA'\subset\mB'$.
It is a comparison at fixed state, not a dissipative entropy-production law.

Modular monotonicity gives bounded imaginary-time values $D_{B\subset A}(i\alpha)$ for $0\leq\alpha\leq1/2$. The characteristic function extends to the strip $S(0,1/2):=\{z\in\mathbb C:0\leq\operatorname{Im}z\leq1/2\}$ \cite{revolutionizing}:
\begin{lemma}
    Given a pair of von Neumann algebras $\mathcal{B}\subset\mathcal{A}$ with a common cyclic and separating vector $\ket{\Omega}$, the characteristic function $D_{B\subset A}(t)$ extends to the strip $\mathcal{S}(0,1/2)$. The extension is analytic in the open strip and satisfies $\|D(z)\|\leq1$ on its closure.

    Moreover, on the boundary $\partial\mathcal{S}(0,1/2) = \{z\in\mathbb{C}:\Im{z} = 0\}\sqcup\{z\in\mathbb{C}:\Im{z} = 1/2\}$, the characteristic function is strongly continuous.
\end{lemma}
For a full discussion of the characteristic function, we refer the reader to \cite{revolutionizing}. Here we list the properties of the characteristic function \cite{revolutionizing}:
\begin{enumerate}
    \item $D(0)=1$, and $D_{\mathcal{B}\subset \mathcal{A}}(t)$ and $D_{\mathcal{B}\subset \mathcal{A}}(t +\frac{i}{2})$ are unitary and strongly continuous in $t$;
    \item $D_{B\subset A}(z)$ is analytic on $0<\operatorname{Im}z<1/2$ and contractive on its closure;
    \item $D_{\mathcal{B}\subset \mathcal{A}}(t)\ket{\Omega} = \ket{\Omega}$;
    \item The cocycle condition holds: $D(t+s)=\Delta_A^{-it}D(s)\Delta_A^{it}D(t)$. Here conjugation by $\Delta_A^{it}$ acts on $B(\mH)$, since $D(s)$ need not lie in $\mA$;
    \item The intrinsic boundary symmetry is
    \[
     D(t+i/2)^*J_A D(t)=D(t)^*J_A D(t+i/2),
    \]
    independent of $t$. For the given inclusion its value is $J_B$, equivalently
    $D(t+i/2)=J_A D(t)J_B$. In particular, $D_{\mathcal{B}\subset \mathcal{A}}(\frac{i}{2}) = J_\mathcal{A}J_\mathcal{B}$. This unitary operator is crucial in subfactor theory creating a chain of algebras:
    \begin{equation}
        \mathcal{A}\supset\mathcal{B}\supset D_{\mathcal{B}\subset \mathcal{A}}(\frac{i}{2})^*\mathcal{A}D_{\mathcal{B}\subset \mathcal{A}}(\frac{i}{2})\supset D_{\mathcal{B}\subset \mathcal{A}}(\frac{i}{2})^*\mathcal{B}D_{\mathcal{B}\subset \mathcal{A}}(\frac{i}{2})\supset...
    \end{equation}
    \item For all $t\in\mathbb{R}$, $D_{\mathcal{B}\subset \mathcal{A}}(t)D_{\mathcal{B}\subset \mathcal{A}}(\frac{i}{2})^*\mathcal{A}D_{\mathcal{B}\subset \mathcal{A}}(\frac{i}{2})D_{\mathcal{B}\subset \mathcal{A}}(t)^*\subset\mathcal{A}$ holds.
\end{enumerate}
Any operator-valued function associated with a pair of von Neumann algebras that satisfy the above-listed properties is called a characteristic function \cite{revolutionizing}. The following theorem verifies the intuition that characteristic function fully captures the structure of inclusions of von Neumann algebras with common cyclic and separating vectors \cite{revolutionizing}:
\begin{theorem}
    Given a von Neumann algebra $\mathcal{A}$ and a cyclic and separating vector $\ket{\Omega}$, then each subalgebra $\mathcal{B}\subset\mathcal{A}$ such that $\ket{\Omega}$ is cyclic for $\mathcal{B}$ is uniquely associated with a characteristic function.
\end{theorem}
The converse characterization uses all the boundary, cocycle, and inclusion properties above; analyticity and vacuum invariance alone are insufficient~\cite{revolutionizing}.


\begin{lemma}[Common cyclicity implies generation]\label{lem:common-cyclicity-generating}
If $\mB\subset\mA$ has a common cyclic and separating vector $\Omega$, then
$\bigvee_{t\in\mathbb R}\sigma_A^t(\mB)=\mA$.
\end{lemma}
\begin{proof}
Let $\mP=\bigvee_t\sigma_A^t(\mB)$. It is modular invariant, so Takesaki's theorem
gives an $\Omega$-preserving expectation $E_{\mP}:\mA\to\mP$.
Since $\mB\subset\mP$ and $\Omega$ is cyclic for $\mB$, the projection onto
$\overline{\mP\Omega}$ is $1$. Hence $E_{\mP}(a)\Omega=a\Omega$ for every $a\in\mA$.
Separatingness for $\mA$ gives $E_{\mP}(a)=a$.
\end{proof}

The same argument shows that a proper common-cyclic inclusion has no
$\Omega$-preserving normal expectation onto $\mB$.
These statements do not assume ergodicity or mixing and do not establish the
classical $K$-property, which additionally requires one-sided nesting and a trivial tail.
The type $III_1$ local algebras in QFT provide the main examples studied below.


In infinite dimensions, the inclusion need not admit a tensor-product description:
\begin{enumerate}
 \item {\bf Trivial relative commutant:} A proper subfactor $\mB\subset\mA$ is
 irreducible when $\mB'\cap\mA=\mathbb C1$.
 This cannot occur for a proper inclusion of type $I$ factors.
 Null-translated wedge inclusions can provide examples in interacting QFT, subject
 to the theory's lightlike localization properties.
 \item {\bf No normal product state of $B$ and $R$:} A conormal inclusion satisfies
 $\mB\vee\mR=\mA$, with $\mR=\mB'\cap\mA$.
 If multiplication does not give a spatial tensor factorization
 $\mA\cong\mB\overline\otimes\mR$, a faithful normal expectation cannot exist,
 by Proposition~\ref{prop:conormal-factorization}.
 This obstruction occurs for adjacent local algebras in the conformal example below.
 \item {\bf Statistical independence of $B$ and $A'$:} The split property concerns
 the commuting pair $\mB,\mA'$, rather than $\mB,\mR$.
 A type $I$ factor between $\mB$ and $\mA$ gives a spatial tensor representation
 for $\mB\vee\mA'$. Under the usual split property of a QFT, a positive collar
 between nested regions supplies this structure; it is not a dynamical mixing condition.
\end{enumerate}
For finite algebras a faithful normal trace gives a trace-preserving expectation
onto every von Neumann subalgebra. A type $I$ ambient factor does not, in general,
admit a normal expectation onto an arbitrary subfactor.

\paragraph{Tensor amplification (an expected, nongenerating inclusion):}
Let $\mB$ be a type $III_1$ factor and set $\mA=\mB\overline\otimes M_n(\mathbb C)$,
with $n>1$. Then $\mR=1\otimes M_n(\mathbb C)$ and the inclusion is conormal.
The slice expectations \eqref{condexp} exist, and the minimal Jones--Kosaki index
is $n^2$. Finite index by itself neither specifies a matrix relative commutant
nor implies conormality.
For a faithful product state $\omega_A=\omega_B\otimes\omega_R$, the outer modular
flow preserves $\mB$, so its modular orbit generates only $\mB$.
The site spectral projections of $\omega_R$ lie in the centralizer of $\omega_A$.
In the product GNS representation, the relative modular comparison acts only on
the hidden factor. The reference vector is not cyclic for $\mB$ on the full Hilbert
space, consistently with Lemma~\ref{lem:common-cyclicity-generating}.

\paragraph{Local algebras of QFT (no normal conditional expectations):}
In a local conformal net $\{\mathcal{A}(I)\}$ on the circle, let $I_1,I_2$ be disjoint
open intervals with closures sharing exactly one endpoint, and set $I=\operatorname{int}\overline{I_1\cup I_2}$.
Haag duality and strong additivity give $\mathcal{A}(I_1)'\cap\mathcal{A}(I)=\mathcal{A}(I_2)$
and $\mathcal{A}(I_1)\vee\mathcal{A}(I_2)=\mathcal{A}(I)$, so the inclusion
$\mathcal{B}:=\mathcal{A}(I_1)\subset\mathcal{M}:=\mathcal{A}(I)$ is conormal with

$\mathcal R=\mathcal A(I_2)$. The adjacent pair has no spatial tensor-product
factorization in a nontrivial conformal net~\cite{Wiesbrock1993}.
Proposition~\ref{prop:conormal-factorization} therefore excludes faithful normal
expectations. The same conclusion holds for a nonfaithful normal expectation:
composing it with a normal state on $\mathcal B$ would give a normal product
state on $\mathcal B\vee\mathcal R$, whereas existence of such a state for commuting
factors entails a spatial tensor-product representation~\cite{DoplicherLongo1984}.
Thus no normal expectation $\mE:\mathcal A(I)\to\mathcal A(I_1)$ exists.
The obstruction concerns normal product states and expectations for this inclusion;
it does not exclude error correction for other channels or chosen families of states.

\paragraph{General von Neumann algebras:}
\label{sec:conormality-expectations} Here, we provide rigorous results that generalize the discussion of cases above to general von Neumann algebras.

Let $\mB\subset\mA\subset B(\mathcal H)$ be an inclusion of
$\sigma$-finite type $\mathrm{III}_1$ factors.  The inclusion is
\emph{standard} with respect to $\Omega\in\mathcal H$ if $\Omega$ is cyclic
and separating for $\mA$, $\mB$ and the relative commutant $\mR=\mB^c=\mB'\cap \mA$. It is \emph{split} if a type I factor $\mathcal F$ exists with
$\mB\subset\mathcal F\subset\mA$, and \emph{conormal} if
\begin{equation}
 \mB\vee\mB^c=\mA.
\end{equation}
Throughout this section, a conditional expectation
$\mE:\mA\to\mB$ is assumed to be \emph{faithful and normal},
unless the contrary is explicitly stated.  Tomiyama's theorem \cite{Tomiyama1957,Tomiyama1959,Takesaki1972} gives
\begin{equation}
 \mE(b_1 a b_2)=b_1\mE(a) b_2,
 \qquad b_1,b_2\in \mB,\quad a\in \mA.
\end{equation}

%\subsection{The split case}


\begin{theorem}[Split inclusions and normal expectations]\label{thm:split-no-expectation}
If $\mB$ is a type III factor and
$\mB\subset\mA$ is split, then no normal conditional expectation
$\mE:\mA\to\mB$ exists.  
\end{theorem}

\begin{proof}
Let $\mB\subset\mathcal F\subset\mA$ be an interpolating type I
factor and suppose that $\mE:\mA\to \mB$ is normal.  Its restriction
\begin{equation}
 \mE_{\mathcal F}:=\mE|_{\mathcal F}:\mathcal F\longrightarrow\mB
\end{equation}
is a normal norm-one projection with range $\mathcal B$.  Tomiyama's type
theorem states that a normally expected von Neumann subalgebra of a type I von
Neumann algebra is type I \cite{Tomiyama1959}\cite{Takesaki1972}. This contradicts $\mB$ being a type III factor.
\end{proof}

Note that in the theorem above, faithfulness and preservation of the
distinguished vector state are not needed, and 
standardness plays no role in this implication.

\begin{lemma}\label{lem:expectation-slice}
If $\mathcal B$ is a factor and $\mE:\mathcal A\to\mathcal B$ is a faithful
normal conditional expectation, there is a faithful normal state $\chi_R$ on
the relative commutant $\mB^c$ such that
\begin{equation}
 \mE(c)=\chi_R(c)\mathbbm 1,
 \qquad \mE(b c)=\chi_R(c)b,
 \qquad b\in\mathcal B,\quad c\in\mathcal R.
\end{equation}
\end{lemma}

\begin{proof}
For $b\in\mathcal B$ and $c\in\mathcal R$,
\begin{equation}
 b \mE(c)=\mE(b c)=\mE(c b)=\mE(c)b.
\end{equation}
Thus $\mE(c)\in Z(\mathcal B)=\mathbb C\mathbbm 1$, so
$\mE(c)=\chi_R(c)\mathbbm 1$ for a normal state $\chi$.  If $c\geq0$ and
$\chi_R(c)=0$, faithfulness of $\mE$ implies $c=0$; hence $\chi_R$ is faithful.  The
second identity is the bimodule property.
\end{proof}

For a faithful normal state $\varphi$ on $\mathcal B$, the state
$\varphi_0:=\varphi\circ \mE$ is faithful and normal and satisfies
\begin{equation}
 \varphi_0(b c)=\varphi_0(b)\chi_R(c).
\end{equation}

\begin{proposition}[conormality and tensor product]\label{prop:conormal-factorization}
Under the hypotheses of Lemma~\ref{lem:expectation-slice}, multiplication
extends to a normal $*$-isomorphism
\begin{equation}
 m:\mathcal B\overline\otimes\mathcal R
 \overset{\cong}{\longrightarrow}\mathcal B\vee\mathcal R,
 \qquad m(b\otimes c)=b c.
\end{equation}
If the inclusion is conormal, then
\begin{equation}\label{eq:conormal-factorization}
\mathcal A\cong
 \mathcal B\overline\otimes \mB^c .
\end{equation}
\end{proposition}

\begin{proof}
Multiplication is a $*$-homomorphism on $\mathcal B\odot\mathcal R$ because
the two algebras commute.  Moreover,
\begin{equation}
 \varphi_0\bigl(m(b\otimes c)\bigr)
 =(\varphi\otimes\chi_R)(b\otimes c).
\end{equation}
The GNS representation of the faithful normal product state
$\varphi\otimes\chi_R$ is the spatial tensor-product representation.  Hence
multiplication extends normally to $\mathcal B\overline\otimes\mathcal R$.
Faithfulness of the product state gives injectivity, while surjectivity follows
from the definition of $\mathcal B\vee\mathcal R$.  Conormality then yields
\eqref{eq:conormal-factorization}; compare
\cite{Faulkner2020,FaulknerSperanza2024}.
\end{proof}

We define the Pimsner--Popa index $\operatorname{Ind}(\mE)$ as the infimum of $\lambda\geq1$ such that, for all positive operators $x_+\in\mA$,
\begin{align}
    \mE(x_+)\geq \lambda^{-1}x_+
\end{align}
If no finite bound exists, $\operatorname{Ind}(\mE)=\infty$. In the type $III$ factor setting above, the minimal index of the inclusion is
\begin{align}
    [\mA:\mB]_0=\inf_\mE \text{Ind}(\mE)\ .
\end{align}
The index is a measure of the relative size of $\mA$ to $\mB$. When the relative commutant is infinite dimensional, or there are no normal conditional expectations, we say the index is infinite.

% If the relative commutant has a finite trace (as in type $I_d$ and $II_1$

% There is no contradiction at this abstract level.  Tensor products of type
% $\mathrm{III}_1$ factors can again be type $\mathrm{III}_1$; for example,
% $R_\infty\overline\otimes R_\infty\cong R_\infty$.  The contradiction below
% is specifically geometric: it uses the realization of $\mathcal N$ and its
% relative commutant as local algebras of adjacent AQFT regions.

% \subsection{The AQFT obstruction for adjacent regions}

% Let $\mathcal A$ be a nontrivial local conformal net, and let $I_1,I_2$ be
% disjoint intervals whose closures meet at one endpoint.  If $I$ is the interior
% of $\overline{I_1}\cup\overline{I_2}$, Haag duality and strong additivity give
% \begin{equation}\label{eq:adjacent-geometry}
%  \mathcal A(I_1)'\cap\mathcal A(I)=\mathcal A(I_2),
%  \qquad
%  \mathcal A(I_1)\vee\mathcal A(I_2)=\mathcal A(I).
% \end{equation}
% Thus $\mathcal N:=\mathcal A(I_1)\subset\mathcal M:=\mathcal A(I)$ is
% conormal and $\mathcal N'\cap\mathcal M=\mathcal A(I_2)$.

% \begin{theorem}\label{thm:adjacent-no-expectation}
% In this adjacent-region configuration, no faithful normal conditional
% expectation
% \begin{equation}
%  E:\mathcal A(I)\longrightarrow\mathcal A(I_1)
% \end{equation}
% exists.
% \end{theorem}

% \begin{proof}
% If $E$ existed, Proposition~\ref{prop:conormal-factorization} and
% \eqref{eq:adjacent-geometry} would give
% \begin{equation}\label{eq:forbidden-adjacent-factorization}
%  \mathcal A(I)\cong
%  \mathcal A(I_1)\overline\otimes\mathcal A(I_2),
%  \qquad a_1\otimes a_2\longmapsto a_1a_2.
% \end{equation}
% For commuting factors, spatiality of the multiplication map is equivalent to
% the split property \cite{DoplicherLongo1984}.  Hence the adjacent algebras
% $\mathcal A(I_1)$ and $\mathcal A(I_2)$ would form a split pair.

% This is impossible for a nontrivial chiral conformal net.
% Reference~\cite[Proposition~1.6.12]{GiorgettiThesis} proves explicitly that
% local algebras of disjoint touching intervals never satisfy the split property
% unless the net is trivial.  Its first proof reduces this to
% \cite[Lemma~13]{Wiesbrock1993}, which shows that a standard inclusion of type
% $\mathrm{III}_1$ factors with type $\mathrm{III}_1$ relative commutant cannot
% be both half-sided modular and split.  Its second proof uses the canonical
% intermediate type I factor of
% \cite[Theorems~2.1 and 4.1]{DoplicherLongo1984} and geometric modular action.
% Thus \eqref{eq:forbidden-adjacent-factorization} contradicts the
% adjacent-region structure.
% \end{proof}

% The logical chain is
% \begin{equation}
% \begin{split}
%  &\text{faithful normal expectation}+\text{conormality}
%  \Longrightarrow
%  \mathcal M\cong\mathcal N\overline\otimes
%  (\mathcal N'\cap\mathcal M)\\
%  &\Longrightarrow
%  \text{$(\mathcal N,\mathcal N'\cap\mathcal M)$ is a split commuting pair}\\
%  &\Longrightarrow
%  \text{contradiction for adjacent AQFT regions}.
% \end{split}
% \end{equation}
% The last implication is an AQFT input, not a type-classification argument.
% Rehren describes the split property as maximal decoupling for intervals with
% disjoint closures and strong additivity as maximal interaction for touching
% intervals \cite{Rehren2015}.  The failure of split decoupling when chiral
% intervals touch is also stated in \cite{LongoRehren2004}.

\subsection{Geometry from ergodic modular \texorpdfstring{$K$}{K}-systems}\label{subsec:geometry-modular-ergodic}


Half-sided modular inclusions add a one-sided nesting condition to the generating
property. A classical $K$-system also requires a trivial tail; neither this condition
nor mixing follows from common cyclicity alone.

Let $(\mathcal N\subset\mathcal M,\Omega)$ be a proper half-sided modular
inclusion and let $\omega(x)=\braket{\Omega|x\Omega}$.

\begin{proposition}\label{prop:no-vacuum-expectation}
There is no faithful normal conditional expectation
$E:\mathcal M\to\mathcal N$ satisfying $\omega\circ E=\omega$.
\end{proposition}

\begin{proof}
Takesaki's theorem would imply
$\sigma_t^{\mathcal M,\Omega}(\mathcal N)=\mathcal N$ for all $t\in\mathbb R$.
Writing $\mathcal N=T(1)\mathcal M T(-1)$, Borchers covariance gives
\begin{equation}
 \Delta_{\mathcal M}^{it}T(x)\Delta_{\mathcal M}^{-it}
 =T(e^{-2\pi t}x),
\end{equation}
and hence
\begin{equation}
 \sigma_t^{\mathcal M,\Omega}(\mathcal N)
 =T(e^{-2\pi t})\mathcal M T(-e^{-2\pi t}).
\end{equation}
Global invariance, together with the fact that $e^{-2\pi t}$ ranges through
$(0,\infty)$, forces $T(x)\mathcal M T(-x)=\mathcal M$ for every $x$ and thus
$\mathcal N=\mathcal M$, contradicting properness.
\end{proof}

This argument does not exclude an expectation preserving a different faithful
normal state or weight.  The conormal adjacent-region argument above rules out
all faithful normal expectations without requiring vacuum preservation.  For a
standard inclusion which is neither split nor conormal, however, standardness
alone does not currently yield a general no-expectation theorem.  The general
criterion remains Takesaki's theorem: a faithful normal expectation exists
precisely when there is a faithful normal semifinite weight on $\mathcal M$,
semifinite on $\mathcal N$, whose modular automorphism group globally preserves
$\mathcal N$ \cite{Takesaki1972}.

Lechner and Scotford state that half-sided modular inclusions with unique
vacuum do not admit a normal conditional expectation and separately emphasize
that such inclusions are never split \cite[Section~2]{LechnerScotford2022}.
The latter assertion is tied to the incompatibility between standard
half-sided modular and split inclusions.  The former, when cited as an
unrestricted no-expectation statement, must be distinguished from the
vacuum-preserving proof above: one must identify an additional mechanism, such
as conormal adjacent-region localization, that rules out expectations
preserving other faithful states.


%\section*{Coarse-Graining, Conditional Expectation, and Modular Invariance}

% A useful way to compare classical ergodic theory with operator-algebraic dynamics is to formulate both in terms of algebras of observables. 


% \NL{Araki-Woods lattice system} as an example of QECC: correspondence between the states of the relative commutant and the conditional algebra.


The classical comparison thus leads to transport of retained observable algebras.
For a proper local inclusion with a common cyclic and separating vacuum, its
modular orbit generates the larger algebra, while a vacuum-preserving expectation
onto the smaller one cannot exist. The characteristic operator
$D_{B\subset A}(z)=\Delta_A^{-iz}\Delta_B^{iz}$ compares the two modular evolutions
without assuming such an expectation.
Section~\ref{sec:geometric-subregion-flows} computes this transport when modular
flow is geometric; Section~\ref{section:conformalperturbation} studies its correction
away from a conformal fixed point; Section~\ref{section:approx} establishes a
nonperturbative large-region approximation.

\section{Geometric subregion flows: wedges in QFTs and balls in CFTs}
\label{sec:geometric-subregion-flows}

An inclusion of causally complete regions gives an inclusion of their observable algebras. In the vacuum, modular flow acts geometrically for wedges in QFT and for double cones associated with balls in CFT\footnote{Following terminology of AQFT, we refer to the causal development of a ball as a double cone.}. Composing these modular flows therefore gives a geometric subregion flow. For more general regions or states, the characteristic flow need not admit such a spacetime interpretation.

We use the convention of the preceding section,
\[
 D_{B\subset A}(z)=\Delta_A^{-iz}\Delta_B^{iz},
 \qquad K_A=-\log\Delta_A,
\]
and write $A,B$ for both regions and their vacuum algebras when no ambiguity arises. To study regularized matrix elements, we keep the correlator's modular-time parameter and consider
\[
 F(z)=\langle O_1|\Delta_A^{1/2}D_{B\subset A}(-z)|O_2\rangle,
 \qquad |O_j\rangle=O_j|\Omega\rangle.
\]
The normalized function $f=F/N$, where
\[
 N=\|\Delta_A^{1/4}|O_1\rangle\|\,
   \|\Delta_B^{1/4}|O_2\rangle\|>0,
\]
is bounded by one on $|\operatorname{Im}z|\leq 1/4$. This follows from the monotonicity of modular operator
\begin{align}
\Delta_A\leq \Delta_B
\end{align}
which results in the bounded analytic characteristic operator on $0<\operatorname{Im}z<1/2$ \cite{revolutionizing}; we return to it in Section~\ref{subsec:modular-chaos-bound}. In our calculations in this section, we choose point-like fields $\mO_j$; they are understood as limits of regulated or smeared insertions. We choose them to be localized strictly inside the region so that the quarter-modular norms   in $N$ corresponds to a geometrically continued finite two-point functions.

\subsection{Inclusion of Wedges}
\label{subsec:wedge-inclusion}

For a spacelike displacement between parallel wedges, Poincar\'e covariance allows us to take $B=W(a)\subset A=W(0)$, with $a>0$ and
\[
 W(a)=\{(t,x):x-a>|t|\}.
\]
We choose translation generators so that $U(b^0,b^1)=\exp[i(b^0P_0+b^1P_1)]$ acts on fields by $(t,x)\mapsto(t+b^0,x+b^1)$.

\begin{theorem}[Characteristic operator of a wedge inclusion]
\label{thm:wedge-characteristic}
In a vacuum relativistic QFT the characteristic flow is
\begin{equation}
 D_{W(a)\subset W(0)}(z)=\Delta_{W(0)}^{-iz}\Delta_{W(a)}^{iz}
 =\exp\!\left\{ia\left[\sinh(2\pi z)P_0
       +\bigl(\cosh(2\pi z)-1\bigr)P_1\right]\right\},
 \qquad z\in\mathbb R.
 \label{eq:wedge-characteristic}
\end{equation}
\end{theorem}
\begin{proof}
It follows from the Bisongano-Wichmann property that the flow $\operatorname{Ad}\Delta_{W(0)}^{-iz}$ is a boost of rapidity $2\pi z$ \cite{bisognanowichmann}; $\operatorname{Ad}\Delta_{W(0)}^{iz}$ reverses this boost. The boost about $(0,a)$ is obtained by conjugating the boost about the origin by $U(0,a)$. Their composition cancels the boost and leaves the translation $(a\sinh(2\pi z),a[\cosh(2\pi z)-1])$.
\end{proof}
For small $z$, to the first order we have $K_{W(0)}-K_{W(a)}=2\pi aP_0$ which is proportional to the Hamiltonian $P_0$; however, at order $z^2$ and higher the characteristic flow 
 also involves translations in space $P_1$. Its translation displacement traces a timelike curve;
in general it is not a one-parameter translation group generated by a fixed Hamiltonian.

\paragraph{Matrix element in 2d CFT:}

For an explicit matrix element, take a self-adjoint scalar primary of dimension $h=\bar{h}>0$ in a $1+1$ dimensional CFT. With $X=(t_1,x_1)\in W(0)$ and $Y=(t_2,x_2+a)\in W(a)$, the reflected and modularly transformed insertions are spacelike separated for real $z$. Their two-point function gives
\begin{align}
 F(z)&=\frac{C}{[\Delta s^2(z)]^{2h}},\qquad C>0
\label{eq:wedge-matrix-element}
\end{align}
where the $\Delta s^2$ is the proper distance:
\begin{align}
\Delta s^2&=-\Delta t^2+\Delta x^2
 =\bigl(u_++ae^{-2\pi z}\bigr)\bigl(u_-+ae^{2\pi z}\bigr),
 \qquad u_\pm=x_1+x_2\pm(t_1+t_2)>0\nn\\
 \Delta t&=t_1+t_2-a\sinh(2\pi z),\qquad
 \Delta x=x_1+x_2+a\cosh(2\pi z),
 \label{eq:wedge-separation}
\end{align}
Note that as long as the operators remain spacelike  on the real $z$ axis, therefore no $i\epsilon$ prescription is needed. The complex continuation is fixed by the analytic logarithm of $\Delta s^2$ that is real there.

\begin{lemma}[Wedge-inside-wedge matrix elements in CFT]
\label{lem:wedge-matrix-element}
For the scalar-primary insertions above in the CFT vacuum, $F$ is analytic and has no zeros in $|\operatorname{Im}z|<1/2$, and satisfies $F(\bar z)=\overline{F(z)}$. Its finite singularities occur at
\begin{equation}
 z=-\frac{1}{2\pi}\log\frac{u_+}{a}+i\left(n+\frac12\right),
 \qquad
 z=\frac{1}{2\pi}\log\frac{u_-}{a}+i\left(n+\frac12\right),
 \qquad n\in\mathbb Z.
 \label{eq:wedge-singularities}
\end{equation}
For real modular time, $F(z)$ decays as $e^{-4\pi h|z|}$ as $|z|\to\infty$.
\end{lemma}
\begin{proof}
The proper distance  in  \eqref{eq:wedge-separation} vanishes only at \eqref{eq:wedge-singularities}. Thus their product has a holomorphic logarithm in the open strip. Positivity on the real axis fixes its branch and gives the reflection relation. The large-$|z|$ behavior follows directly from the same factorization.
\end{proof}

These singularities correspond to null separation of the transformed insertions. At a zero of $\Delta s^2$ of multiplicity $m$, the local behavior is $F(z)\sim c(z-z_0)^{-2hm}$. It is a pole when $2hm$ is an integer and a branch point otherwise. A counterclockwise circuit multiplies $F$ by $e^{-4\pi ihm}$. The two locations in \eqref{eq:wedge-singularities} coincide when $u_+u_-=a^2$, giving $m=2$. For example, if $t_1+t_2=0$ and $x_1+x_2=a$, then
\[
 \Delta s^2=4a^2\cosh^2(\pi z),\qquad
 F(z)=\frac{C}{(2a)^{4h}}\cosh^{-4h}(\pi z).
\]
Thus $2h=3/2$ gives third-order poles despite the noninteger scaling dimension.

On $\operatorname{Im}z=\pm1/2$, the continued function is real between the two distinct singularities. In that case, cuts extending horizontally outward from them can be chosen to give an $i$-periodic continuation on the resulting cut domain. Periodicity must include this branch prescription: the denominator is always $i$-periodic, whereas its fractional power need not be. The coincident example with $2h=3/2$ is instead antiperiodic, $F(z+i)=-F(z)$. The phase of $\Delta s^2$ and the coincidence of its zeros are illustrated in Figure~\ref{fig:deltaS2_z}.

\begin{figure}[tp]
    \centering
    \begin{subfigure}{.3\textwidth}
        \label{fig:placeholder1}
        %\mathcal{A}ption{a}
    \end{subfigure}
    \begin{subfigure}{.3\textwidth}
        \includegraphics[width=1\linewidth]{images/delta_s2_plot_coincide.pdf}
        \label{fig:placeholder2}
        %\mathcal{A}ption{b}
    \end{subfigure}
    \begin{subfigure}{.3\textwidth}
        \includegraphics[width=1\linewidth]{images/delta_s2_plot_reversed.pdf}
        \label{fig:placeholder3}
        %\mathcal{A}ption{c}
    \end{subfigure}     
    \caption{The phase of $\Delta s^2$ in the complex modular-time plane. The panels show $u_+u_->a^2$ (left), $u_+u_-=a^2$ (middle), and $u_+u_-<a^2$ (right), with $u_\pm=x_1+x_2\pm(t_1+t_2)$. We take $a=1$, $x_1=0.1$, $t_1=t_2=0$, and respectively $x_2=3$, $x_2=0.9$, and $x_2=1/3$. The middle panel has double zeros of $\Delta s^2$; for $2h=3/2$ these produce third-order poles of $F$, rather than branch points.}
    \label{fig:deltaS2_z}
\end{figure}
\paragraph{Massive free fields:} The geometric operator \eqref{eq:wedge-characteristic} remains valid beyond CFT. The analytic properties of its matrix elements, however, depend on the vacuum two-point function. Lorentz covariance does not imply a zero-free complex continuation or the scalar-primary branch structure. The larger zero-free strip of Lemma~\ref{lem:wedge-matrix-element} is a property of these point-field correlators, rather than of arbitrary characteristic matrix elements. Smearing can change the singularities and introduce zeros through cancellations. The bounded-strip statement for normalized matrix elements survives whenever the required vectors exist. In a general QFT, for point-like insertions $\mO_j$ inside the region, the two-point function can only diverge when the operators become null-separated, $\Delta s^2=0$, and onlyes vanish when the proper distance goes to infinity, $\Delta s^2\to\infty$. Therefore, one expects that in a general QFT, we still have two singularities in the same locations as above, however, the branch cut structure can be more complex. 

As an example of a non-conformal theory, we consider the case of free massless boson in $d+1$ dimensions. For a wedge inclusion $W(a) \subset W(0)$ along the $x_1$-direction, let $X = (t_1, x_1, \mathbf{x}_{\perp, 1})$ and $Y = (t_2, x_2, \mathbf{x}_{\perp, 2})$ be interior points of $W(0)$ and $W(a)$ respectively. Let $\Delta s^2(z) = \Delta s_{1+1}^2(z) + |\Delta \mathbf{x}_\perp|^2$ be the analytic continuation of the squared proper distance to complex modular time $z$, where $\Delta s_{1+1}^2(z)$ is the longitudinal part defined in (3.4) and $\Delta \mathbf{x}_\perp = \mathbf{x}_{\perp, 1} - \mathbf{x}_{\perp, 2}$. We define the analytically continued proper distance as $r(z) = \sqrt{\Delta s^2(z)}$, taking the branch that is positive on the real axis. Then, we have
\begin{lemma}[Massive free-boson wedge matrix elements in $d+1$ dimensions] 
\label{lem:massive-free-boson-wedge-d-dim}
Let $\phi$ be a canonically normalized real free scalar field of mass $m > 0$ in the $(d+1)$-dimensional Minkowski vacuum. Then
\begin{equation}
    F(z) = \braket{\phi(X) | \Delta_A^{1/2} D_{B \subset A}(-z) | \phi(Y)} = \frac{1}{(2\pi)^{\frac{d+1}{2}}} \left( \frac{m}{r(z)} \right)^{\frac{d-1}{2}} K_{\frac{d-1}{2}}\!\left( m r(z) \right),
\end{equation}
where $K_{\nu}$ is the modified Bessel function of the second kind. The function $F$ is analytic and has no zeros in $|\operatorname{Im} z| < 1/2$, satisfies $F(\bar{z}) = \overline{F(z)}$, and is positive on the real axis. 

Its finite branch singularities occur at the modified null-separation points where $\Delta s^2(z) = 0$. For $d > 1$, the leading short-distance singularity is algebraic rather than logarithmic; on a fixed local branch, as $\Delta s^2(z) \to 0$,
\begin{equation}
    F(z) = \frac{\Gamma\left(\frac{d-1}{2}\right)}{4 \pi^{\frac{d+1}{2}}} \frac{1}{r(z)^{(d-1)}} \lb 1+O(r(z)^{-2})\rb.
\end{equation}
For real modular time $t$ and fixed $m>0$, $F(t)$ retains the same leading exponential suppression $e^{-m r(t)}$ as $|t| \to \infty$, where the real proper distance $\rho(t)$ grows exponentially. Thus $F$ decays faster than any exponential in $|t|$. Its quarter-modular normalization $N$ is finite and positive, and $f = F/N$ satisfies $|f(z)| \leq 1$ for $|\operatorname{Im} z| \leq 1/4$.
\end{lemma}

\begin{proof}
At spacelike separation, the $(d+1)$-dimensional vacuum mode integral evaluates to the standard Wightman propagator, yielding the expression for $F(z)$ in terms of the proper distance $r(z)$. The transverse coordinates $\mathbf{x}_\perp$ are invariant under the geometric wedge flow, so the invariant distance squared is modified only by the constant additive term $|\Delta \mathbf{x}_\perp|^2$.

For $0 < \operatorname{Im} z < 1/2$, the longitudinal part $\Delta s_{1+1}^2(z)$ maps the strip away from the negative real axis. Adding a real constant $|\Delta \mathbf{x}_\perp|^2 \geq 0$ preserves this property, so the continued squared proper distance $\Delta s^2(z) \notin (-\infty, 0]$. Reflection and positivity on the real axis give the same conclusion throughout the open strip, ensuring $\operatorname{Re} r(z) > 0$. 

The principal branch of $K_{\nu}(w)$ is analytic and zero-free in the right half-plane $\operatorname{Re}(w) > 0$, establishing the analyticity and non-vanishing of $F(z)$. The small-argument expansion of $K_{\nu}(w)$ for $\nu > 0$ yields the algebraic singularity, contrasting with the logarithmic singularity of the $d=1$ case. The large-argument asymptotic $K_{\nu}(w) \sim \sqrt{\pi / (2w)} e^{-w}$ confirms the exponential decay in modular time. The bounds on the normalized matrix element follow from the contraction property of $D_{B \subset A}(i/4 - z)$.
\end{proof}


The mass changes the long-distance decay, while the null singularity retains the logarithmic short-distance behavior of the elementary scalar in two dimensions. The factor $e^{-m r}$ expresses the suppression of correlations beyond the Compton wavelength; wedge modular flow makes $\rho$ grow exponentially in real modular time. These statements concern the elementary field, rather than arbitrary composite or smeared observables. Continuing around a null singularity changes the Bessel branch, so periodicity beyond the strip requires a branch prescription, as in the CFT example.

The geometric transformation in Lemma \ref{subsec:wedge-inclusion} means that for the inclusion of wedges the characteristic flow preserves the one-particle Hilbert space, and more generally does not result in particle production. In a general interacting QFT, the two-point admits a Lehman representation 
\begin{align}
	\braket{\phi(x)\phi(0)}=Z\int dm^2 \mu(m^2)\braket{\phi(x)\phi(0)}_{m^2}
\end{align}
where $\braket{\phi(x)\phi(0)}_{m^2}$ is the two-point function of a massive free field of mass $m$ in Lemma \ref{lem:massive-free-boson-wedge-d-dim}, $\mu(m^2)$ is the spectral density function, and $Z$ is the wave-function renormalization, which cancels out as we normalize the vectors $\phi(x)\ket{\Omega}$.


\subsection{Inclusion of Ball inside Wedge}
\label{subsec:ball-inside-wedge}

In a CFT, a vacuum-preserving conformal transformation maps an outer double cone to a wedge and its inner double cone to another double cone \cite{BrunettiGuidoLongo1993}. This relates ball-inside-ball inclusions to ball-inside-wedge inclusions; with concentric double cones being a special case. Under the usual split property, strict containment of the inner closure gives a split inclusion. Spacelike wedge inclusions are non-split in conformal theories because they share endpoints at conformal infinity. This last statement does not extend to arbitrary massive QFTs \cite{Tanimoto2014}.

Below, for simplicity, we restrict to the vacuum state of $1+1$-dimensional CFTs, however, the generalization to higher dimensions is straightforward \cite{}:
\begin{theorem}[Characteristic operator of concentric double cones]
\label{thm:concentric-characteristic}
Let $B_{u+s}\subset B_u$, with $s>0$, be concentric double cones centered at the origin, of radii $r=e^{-2\pi(u+s)}$ and $R=e^{-2\pi u}$, respectively. The geometric action of $D_{B_{u+s}\subset B_u}(z)=\Delta_{B_u}^{-iz}\Delta_{B_{u+s}}^{iz}$ on each null coordinate $x_\pm=t\pm x$ is a Mobius transformation represented by the matrix $ M_R(z)M_r(-z)$ 
with 
\begin{equation}
 M_R(z)=
 \begin{pmatrix}
  \cosh(\pi z)&R\sinh(\pi z)\\
  R^{-1}\sinh(\pi z)&\cosh(\pi z)
 \end{pmatrix}, \qquad x_\pm\to \frac{\cosh(\pi z)x_\pm+R\sinh(\pi z)}{R^{-1}\sinh(\pi z)x_\pm +\cosh(\pi z)}
 \label{eq:concentric-characteristic}
\end{equation}
The primary-field transformation includes its conformal Jacobian. The unitary lift to the conformal covering group is fixed by the product of the modular groups and continuity from $z=0$.
\end{theorem}
\begin{proof}
The vacuum modular flow of a double cone acts on $x_\pm$ as the Mobius transformation $(c x_\pm+R b)/(b x_\pm/R+c)$, with $c=\cosh(\pi z)$ and $b=\sinh(\pi z)$. Composing the outer flow with the reversed inner flow gives \eqref{eq:concentric-characteristic}.
\end{proof}

The matrix product avoids choosing a global logarithm of the conformal transformation. Its infinitesimal generator is $i(K_{B_u}-K_{B_{u+s}})$, as required by the modular-operator product.

We now take $A=W(0)$ and a double cone $B$ with tips $(t_c\pm r,x_c)$, strictly contained in the wedge:
\[
 x_c-|t_c|>r.
\]
For a scalar primary of dimension $2h$, its modular flow is
\begin{align}
 \Delta_B^{-is}O(t,x)\Delta_B^{is}
 &=\gamma^{2h}O(t_s,x_s),
 \label{eq:ball-modular-flow}\\
 t_s&=\gamma\left[(t-t_c)\cosh(2\pi s)
       +\frac{r^2-d^2}{2r}\sinh(2\pi s)\right]+t_c,
 \label{eq:ball-time}\\
 x_s&=\gamma(x-x_c)+x_c,
 \label{eq:ball-space}\\
 \gamma&=\frac{2r^2}{r^2+d^2+(r^2-d^2)\cosh(2\pi s)
                       +2r(t-t_c)\sinh(2\pi s)},
 \label{eq:ball-conformal-factor}\\
 d^2&=-(t-t_c)^2+(x-x_c)^2.
 \label{eq:ball-distance}
\end{align}
For $(t,x)\in B$ and real $s$, $\gamma>0$ and the transformed point remains in $B$. Poles of this coordinate transformation in complex modular time can represent passage through conformal infinity rather than a singularity of the correlator.

Take $X=(t_1,x_1)\in A$ and $Y=(t_2+t_c,x_2+x_c)\in B$, so that $x_1>|t_1|$ and $|t_2|+|x_2|<r$. Set
\begin{equation}
 d^2=x_2^2-t_2^2,\qquad D=r^2-d^2+2rt_2,\qquad
 p_\pm=\frac{r^2+d^2\pm2rx_2}{D}>0,
 \label{eq:ball-parameters}
\end{equation}
with $p_+p_-=(r^2-d^2-2rt_2)/D$. Writing $q=e^{2\pi z}$, the conformal factor and separations are
\begin{align}
 \gamma&=\frac{4r^2q}{D(q+p_+)(q+p_-)},
 \label{eq:ball-gamma-factorized}\\
 \Delta x+\Delta t&=-(x_1+t_1)q-(x_c+t_c)
                       -r\frac{q-p_-}{q+p_-},
 \label{eq:ball-separation-plus}\\
 \Delta x-\Delta t&=-(x_1-t_1)q^{-1}-(x_c-t_c)
                       +r\frac{q-p_+}{q+p_+},
 \label{eq:ball-separation-minus}\\
 F(z)&=C\left(\frac{\gamma}{\Delta s^2}\right)^{2h},
 \qquad \Delta s^2=(\Delta x+\Delta t)(\Delta x-\Delta t).
 \label{eq:ball-matrix-element}
\end{align}
At $z=0$, these expressions give the separation of the wedge-reflected $X$ and the original $Y$. The complex branch is fixed by continuation from real $z$, where $\gamma/\Delta s^2>0$.

To display the cancellation of coordinate poles, define $\xi_\pm=x_1\pm t_1$, $c_\pm=x_c\pm t_c$, and
\begin{align}
 P_+(q)&=\xi_+q^2+(c_++r+\xi_+p_-)q+(c_+-r)p_-,
 \label{eq:ball-polynomial-plus}\\
 P_-(q)&=(c_--r)q^2+[\xi_-+(c_-+r)p_+]q+\xi_-p_+.
 \label{eq:ball-polynomial-minus}
\end{align}
Then
\begin{equation}
 \frac{\Delta s^2}{\gamma}
 =\frac{D\,P_+(q)P_-(q)}{4r^2q^2},\qquad
 F(z)=C\left[\frac{4r^2q^2}{D\,P_+(q)P_-(q)}\right]^{2h}.
 \label{eq:ball-invariant-denominator}
\end{equation}
The factors $q+p_\pm$ have canceled. A coordinate pole is therefore harmless unless the remaining invariant denominator also vanishes.

\begin{lemma}[Ball-inside-wedge matrix elements]
\label{lem:ball-matrix-element}
Assume strict containment and insertions in the interiors of the regions, as above. Then $F$ is analytic and has no zeros in $|\operatorname{Im}z|<1/2$, and satisfies $F(\bar z)=\overline{F(z)}$. Its finite singularities are determined by the roots of $P_+P_-$. Each quadratic has two distinct negative real roots, giving generically four locations on each line $\operatorname{Im}z=n+1/2$. The local singularity at a root of multiplicity $m$ is a pole when $2hm$ is an integer and a branch point otherwise. For real modular time, $F(z)$ decays as $e^{-8\pi h|z|}$.
\end{lemma}
\begin{proof}
All coefficients in \eqref{eq:ball-polynomial-plus}--\eqref{eq:ball-polynomial-minus} are positive. Their discriminants are
\[
 [\xi_+p_--(c_++r)]^2+8r\xi_+p_->0,
 \qquad
 [\xi_--(c_-+r)p_+]^2+8r\xi_-p_+>0.
\]
Thus both roots of each quadratic are real and negative. Roots of the different quadratics may coincide. Equation~\eqref{eq:ball-invariant-denominator} gives the remaining statements, including the absence of finite zeros and the decay at $q\to0,\infty$.
\end{proof}

The extra singularities compared with the wedge case arise from the fractional linear modular flow of the inner double cone. Their interpretation is again null separation, now in the conformal compactification. Continuation beyond the strip must keep track of the logarithmic branch and root multiplicities. In a general QFT, double-cone modular flow need not be geometric, so this explicit argument does not apply. This motivates the conformal perturbation theory of the next section.

\subsection{An application: Modular Chaos Bound}
\label{subsec:modular-chaos-bound}

In this subsection, we review the modular chaos bound, as an application of the analytic property of the characteristic function. We restore the normalization and consider the strip bound underlying modular chaos \cite{deBoerLamprou2020}:
\begin{align}
 f(z)&=\frac{\langle O_1|\Delta_A^{1/2+iz}\Delta_B^{-iz}|O_2\rangle}{N}
 \nonumber\\
 &=\frac{\langle\Delta_A^{1/4}O_1|
          D_{B\subset A}(i/4-z)|\Delta_B^{1/4}O_2\rangle}{N}
 =1+i\mathcal M(z).
 \label{eq:normalized-characteristic}
\end{align}
Here $|\Delta_A^{1/4}O_1\rangle$ denotes $\Delta_A^{1/4}|O_1\rangle$, and similarly for $B$. The contraction property of $D_{B\subset A}$ gives $|f(z)|\leq1$ for $|\operatorname{Im}z|\leq1/4$, with analytic interior and continuous boundary values. Hence
\begin{equation}
 \operatorname{Im}\mathcal M(z)=1-\operatorname{Re}f(z)\geq0.
 \label{eq:positive-imaginary-part}
\end{equation}
The map $E=ie^{-2\pi z}$ sends this strip conformally to the upper half-plane. For $E=v+iw$, Schwarz--Pick gives, whenever $\operatorname{Im}\mathcal M>0$,
\begin{equation}
 \left|w\partial_w\log\operatorname{Im}\mathcal M(E)\right|
 \leq\frac{w|\partial_w\mathcal M(E)|}{\operatorname{Im}\mathcal M(E)}
 \leq1.
 \label{eq:exact-modular-chaos-bound}
\end{equation}

\begin{lemma}[Strictness of the normalized bound]
\label{lem:strict-modular-chaos-bound}
If $f=1+i\mathcal M$ is analytic and bounded by one on the upper half-plane, and $f\not\equiv1$, then $\operatorname{Im}\mathcal M>0$ and
\begin{equation}
 \left|w\partial_w\log\operatorname{Im}\mathcal M(v+iw)\right|<1
 \qquad(w>0).
 \label{eq:strict-modular-chaos-bound}
\end{equation}
The constant one is nevertheless sharp.
\end{lemma}
\begin{proof}
The image of $\mathcal M=-i(f-1)$ lies in the closed disk of radius one centered at $i$. The only point of this disk on the real axis is zero; attaining it in the interior would force $f\equiv1$. Equality in the second Schwarz--Pick inequality in \eqref{eq:exact-modular-chaos-bound} would require a nonconstant automorphism of the upper half-plane, whose image cannot lie in this disk. Constant functions give a zero derivative. Sharpness follows, for example, from $f(E)=E/(E+i\eta)$, $\eta>0$: on $E=iw$, the left side is $w/(w+\eta)$, which approaches one as $\eta/w\to0$.
\end{proof}

On the real modular-time axis, $E=iw$ with $w=e^{-2\pi t}$, so the bound reads
\[
 \left|\partial_t\log[1-\operatorname{Re}f(t)]\right|<2\pi.
\]
For the real scalar-primary correlators above, a stronger form follows by applying Schwarz--Pick directly to the disk-valued function $f$:
\begin{equation}
 |\partial_t f(t)|\leq\pi[1-|f(t)|^2],\qquad
 \left|\partial_t\log[1-f(t)]\right|\leq\pi[1+f(t)]<2\pi.
 \label{eq:real-axis-modular-chaos-bound}
\end{equation}
Exact saturation is thus excluded for a finite nonzero deficit. A rate of $2\pi$ can be approached when the inclusion is weakly deformed and the amplified deformation remains small.

\paragraph{Small geometric deformations:} For the wedge inclusion, choose identical insertions $X=Y=(0,x)$ with $x>a>0$. Equation~\eqref{eq:wedge-matrix-element}, with the quarter-modular normalization, gives
\begin{equation}
 f_a(t)=\left[
 \frac{4x(x-a)}{(2x-a)^2+a^2+2a(2x-a)\cosh(2\pi t)}
 \right]^{2h}.
 \label{eq:normalized-small-wedge}
\end{equation}
For concentric double cones with both insertions at their common center, \eqref{eq:concentric-characteristic} similarly gives
\begin{equation}
 f_\kappa(t)=\left[
 \frac{\kappa}{[1+(1-\kappa)\sinh^2(\pi t)]^2}
 \right]^{2h},\qquad \kappa=\frac rR<1.
 \label{eq:normalized-small-cones}
\end{equation}
In both cases, setting $\epsilon=a/x$ or $\epsilon=1-\kappa$, respectively, gives at fixed $t$
\begin{equation}
 f(t)=1-2h\epsilon\cosh(2\pi t)+O(\epsilon^2).
 \label{eq:small-deformation-expansion}
\end{equation}
In the joint limit $t\to+\infty$, $\epsilon\to0$ with $\epsilon e^{2\pi t}\to0$, this becomes
\begin{equation}
 1-f(t)=h\epsilon e^{2\pi t}
 +O\!\left(\epsilon e^{-2\pi t}+\epsilon^2e^{4\pi t}\right).
 \label{eq:asymptotic-saturation}
\end{equation}
The exact expressions show that $\partial_t\log[1-f(t)]\to2\pi$ in this regime. An analogous statement holds for large negative modular time.

The physical mechanism is the exponential amplification of a small geometric displacement by boost or conformal modular flow. These examples approach the modular-chaos bound in a controlled limit, rather than saturating it at finite deformation. This amplification alone does not establish many-body scrambling or the K-system property. In particular, the spacelike wedge displacement and strict concentric double-cone inclusion considered here are not half-sided modular inclusions; null-shifted wedges or suitable shared-endpoint intervals supply that structure.


\section{Subregion flows in conformal perturbation theory}\label{section:conformalperturbation}

Wedge modular flow remains geometric in the vacuum of a relativistic QFT with the Bisognano--Wichmann property. For a ball, this geometric action generally fails away from a conformal fixed point. We compute its leading correction by perturbing a CFT with a Hermitian scalar primary $O_0$ of dimension $\Delta$. We work in $1+1$ dimensions and use
\begin{equation}
 \mathcal{L}_\lambda=\mathcal{L}_{\mathrm{CFT}}-\lambda O_0,
 \qquad S_{E,\lambda}=S_{E,0}+\lambda\int d^2x_E\,O_0(x_E).
 \label{eq:cpt-deformation}
\end{equation}
The Euclidean sign fixes the signs of both the modular-Hamiltonian and vacuum corrections below. We use a common regulated description of the theories, with a fixed identification of the external operators. We assume that the regulated density matrices are faithful and their first variations exist. Continuum formulas are conditional on a renormalized distributional extension of the ordered correlators and contact terms; its existence and infrared limit depend on the theory. The expansion is in regulated correlation functions at fixed modular time; it need not be uniform at large time or in the infinite-volume limit. We use $3/2<\Delta<2$ for the pointwise kernel formulas below; the joint geometric integral has the weaker condition $\Delta>1$, as explained after Lemma~\ref{lem:cpt-ball-flow}.

\subsection{Perturbation of ball modular flow}

Let $I=(-1,1)$ at $t=0$ and let $B$ be its causal development. In a regulated description, write $k_{B,\lambda}=-\log\rho_{B,\lambda}$ for the one-sided modular Hamiltonian. Its adjoint action on $\mathcal A(B)$ agrees with that of the full generator $K_{B,\lambda}=-\log\Delta_{B,\lambda}$. We distinguish these generators because only the full modular operator acts on the vacuum Hilbert space with $\Delta_{B,\lambda}|\Omega_\lambda\rangle=|\Omega_\lambda\rangle$.

The first-order logarithmic variation is \cite{BalakrishnanParrikar2020}
\begin{align}
 \delta k_B&=\int_{-\infty+i\theta}^{\infty+i\theta}
 \frac{2\pi\,du}{4\sinh^2(\pi u)}
 e^{iu k_{B,0}}\rho_{B,0}^{-1}\delta\rho_B e^{-iu k_{B,0}},
 \qquad 0<\theta<1,\nonumber\\
 \rho_{B,0}^{-1}\delta\rho_B
 &=-\lambda\int_{-1}^1dX\int_0^1d\tau\,
 J(\tau,X)\Omega^\Delta(-i\tau,X)
 e^{\tau k_{B,0}}O_0(0,X)e^{-\tau k_{B,0}}.
 \label{eq:cpt-density-variation}
\end{align}
Here and below we suppress scalar terms from density-matrix normalization: they drop out of the adjoint flow. The contour and endpoint prescriptions are part of the distributional definition of the kernel.

For real $u$, define $\beta_C^u=\operatorname{Ad}\Delta_{C,0}^{-iu}$, so $\beta_C^u=\sigma_C^{-u}$ with the convention \eqref{eq:modular-flow}. For an insertion at $(0,X)$ in the unit ball,
\begin{gather}
 \beta_B^u(O_0(0,X))=\Omega^{-\Delta}(u,X)O_0(t_u,x_u),\nonumber\\
 t_u=\frac{(1-X^2)\sinh(2\pi u)}{2\Omega(u,X)},
 \qquad x_u=\frac{X}{\Omega(u,X)},\nonumber\\
 \Omega(u,X)=1+(1-X^2)\sinh^2(\pi u).
 \label{eq:cpt-ball-flow}
\end{gather}
The positive Euclidean Jacobian is
\begin{equation}
 J(\tau,X)=\left|\det\frac{\partial(t_E,x_E)}{\partial(\tau,X)}\right|
 =\frac{\pi(1-X^2)}{\Omega^2(-i\tau,X)},
 \qquad \Omega(-i\tau,X)=1-(1-X^2)\sin^2(\pi\tau).
 \label{eq:cpt-euclidean-jacobian}
\end{equation}
Thus $J$ is a Euclidean measure factor, rather than the Lorentzian determinant evaluated at complex time. Choosing $\theta=\tau$ in each insertion removes the Euclidean evolution from the operator and gives
\begin{align}
 \delta k_B&=\pi\lambda\int_{-1}^1dX\int_{\mathbb R}du\,
 h(u,X)\beta_B^u(O_0(0,X)),\label{eq:deltaK}\\
 h(u,X)&=-\frac{\pi(1-X^2)}{2}\int_0^1d\tau\,
 \frac{\Omega^{\Delta-2}(-i\tau,X)}{\sinh^2(\pi(u+i\tau))}.
 \label{eq:cpt-h-kernel}
\end{align}

The kernel contains a local term. With the same endpoint prescription,
\begin{equation}
 -\frac{\pi}{2}\int_0^1\frac{d\tau}{\sinh^2(\pi(u+i\tau))}=\delta(u),
 \label{eq:cpt-delta-kernel}
\end{equation}
as an identity of distributions in $u$, not a pointwise identity. Accordingly,
\begin{align}
 h(u,X)&=(1-X^2)\bigl[\delta(u)+h_{\mathrm{reg}}(u,X)\bigr],\nonumber\\
 h_{\mathrm{reg}}(u,X)&=-\frac{\pi}{2}\int_0^1d\tau\,
 \frac{\Omega^{\Delta-2}(-i\tau,X)-1}{\sinh^2(\pi(u+i\tau))}.
 \label{eq:cpt-regular-kernel}
\end{align}
The local contribution is $\pi\lambda\int_{-1}^1dX(1-X^2)O_0(0,X)$; the remaining term spreads the perturbation along the unperturbed modular orbit. At $\Delta=2$ this nonlocal kernel vanishes formally at first order. A dimension-two operator need not be exactly marginal, so this observation alone does not imply that the deformed theory is another CFT.

Write $k_{B,\lambda}=k_{B,0}+\lambda k_B^{(1)}+O(\lambda^2)$ and introduce the interaction-picture flow
\begin{equation}
 u_{s,\lambda}=e^{-is k_{B,0}}e^{is k_{B,\lambda}}
 =1+\lambda u_s^{(1)}+O(\lambda^2),
 \qquad u_s^{(1)}=i\int_0^s dw\,\beta_B^{-w}(k_B^{(1)}).
 \label{eq:cpt-interaction-flow}
\end{equation}
Duhamel's formula and a shift of the real modular integration variable give the following result.
\begin{lemma}[Perturbed modular flow of the unit ball]
\label{lem:cpt-ball-flow}
For the regulated perturbation \eqref{eq:cpt-deformation}, and real $s$,
\begin{align}
 u_{s,\lambda}&=1+i\pi\lambda\int_{-1}^1dX\int_{\mathbb R}du\,
 (1-X^2)Q(u,X;s)\Omega^{-\Delta}(u,X)O_0(t_u,x_u)
 +O(\lambda^2),\label{eq:cpt-flow-correction}\\
 Q(u,X;s)&=Q_{\mathrm{loc}}(u;s)+Q_{\mathrm{nonloc}}(u,X;s),\nonumber\\
 Q_{\mathrm{loc}}(u;s)&=\int_0^s dw\,\delta(u+w),\nonumber\\
 Q_{\mathrm{nonloc}}(u,X;s)&=-\frac{\pi}{2}\int_0^s dw\int_0^1d\tau\,
 \frac{\Omega^{\Delta-2}(-i\tau,X)-1}{\sinh^2(\pi(u+w+i\tau))}.
 \label{eq:cpt-Q-kernel}
\end{align}
For $s>0$, $Q_{\mathrm{loc}}=\mathbf 1_{(-s,0)}$; for $s<0$, it is $-\mathbf 1_{(0,-s)}$. For $3/2<\Delta<2$, integrating over $w$ yields
\begin{align}
 Q_{\mathrm{nonloc}}(u,X;s)
 &=\frac12\int_0^1d\tau\,[\Omega^{\Delta-2}(-i\tau,X)-1]\nonumber\\
 &\quad\times\bigl[\coth(\pi(u+s+i\tau))-\coth(\pi(u+i\tau))\bigr].
 \label{eq:cpt-Q-coth}
\end{align}
The integrals retain their distributional prescription at the endpoints. The remainder is understood in regulated correlation functions at fixed $s$.
\end{lemma}

The subtraction in \eqref{eq:cpt-regular-kernel} removes the endpoint singularity associated with the local term. A separate issue occurs at $X=0$, $\tau=1/2$, where the Euclidean map reaches infinity and $\Omega$ vanishes. Locally,
\[
 \Omega(-i\tau,X)\simeq X^2+\pi^2(\tau-1/2)^2.
\]
The joint $(X,\tau)$ integral of $\Omega^{\Delta-2}$ is locally integrable for $\Delta>1$, whereas the $\tau$ integral at fixed $X=0$ requires $\Delta>3/2$. We use the latter range for the pointwise kernel formulas. These conditions do not replace the contact-term prescriptions needed when the kernel is inserted in a correlation function. For a ball of radius $r$, the expansion parameter is the dimensionless coupling $\lambda r^{2-\Delta}$.

\paragraph{Almost-marginal relevant operators.}
Let $\epsilon=2-\Delta>0$ be small. The leading nonlocal correction can be evaluated explicitly:
\begin{align}
 \Omega^{-\epsilon}(-i\tau,X)-1
 &=-\epsilon\log\Omega(-i\tau,X)+O(\epsilon^2),\nonumber\\
 -\log\Omega(-i\tau,X)&=\sum_{n\in\mathbb Z}a_n^{(1)}(X)e^{2\pi in\tau},\nonumber\\
 a_0^{(1)}(X)&=-2\log\frac{1+|X|}{2},\nonumber\\
 a_{n\ne0}^{(1)}(X)&=\frac1{|n|}
 \left(-\frac{1-|X|}{1+|X|}\right)^{|n|}.
 \label{eq:cpt-near-marginal-Fourier}
\end{align}
At $X=0$ the Fourier expansion is interpreted in the integrable sense. The expansion in $\epsilon$ is not pointwise uniform near $\Omega=0$, but its logarithmic coefficients are integrable. For $0\leq\epsilon\leq\epsilon_0<1/2$, an integrable bound by $\Omega^{-\epsilon_0}\log^2\Omega$ controls the remainder in the $\tau$ integral, including $X=0$.

For real $v$, set
\begin{equation}
 L(v,X)=\log\left(\frac{1+|X|+(1-|X|)e^{-2\pi|v|}}{2}\right).
\end{equation}
Expanding $\coth$ on either side of $v=0$ and pairing its Fourier modes with \eqref{eq:cpt-near-marginal-Fourier} gives
\begin{align}
 &\int_0^1d\tau\,[\Omega^{-\epsilon}(-i\tau,X)-1]
 \coth(\pi(v+i\tau))\nonumber\\
 &\qquad=-2\epsilon\operatorname{sgn}(v)L(v,X)+O(\epsilon^2),\nonumber\\
 Q_{\mathrm{nonloc}}(u,X;s)
 &=-\epsilon\operatorname{sgn}(u+s)L(u+s,X)\nonumber\\
 &\quad+\epsilon\operatorname{sgn}(u)L(u,X)+O(\epsilon^2).
 \label{eq:cpt-near-marginal-Q}
\end{align}
This single expression covers all three real-time intervals and vanishes at $s=0$. In particular, the middle interval contains the product of the two logarithm arguments, each with its factor of $1/2$. The local term is unchanged. These real-time expressions, which contain absolute values and sign functions, are not formulas for direct substitution of complex $s$.

\subsection{First-order correction to the characteristic matrix element}

Let $B$ be the unit double cone and let $A$ be the right wedge with endpoint $(0,-x_w)$, with $x_w>1$. This guarantees $B\subset A$. Take fixed Hermitian scalar insertions $O_1(x)$ and $O_2(y)$ with $x\in A$, $y\in B$, understood as regulated or smeared fields. In the common perturbative description, write
\begin{align}
 F_\lambda(z)&=\langle O_1|\Delta_{A,\lambda}^{1/2}
 D_{B\subset A,\lambda}(-z)|O_2\rangle_\lambda\nonumber\\
 &=\langle O_1|\Delta_{A,\lambda}^{1/2+iz}
 \Delta_{B,\lambda}^{-iz}|O_2\rangle_\lambda
 =F_0(z)+\lambda F^{(1)}(z)+O(\lambda^2).
 \label{eq:cpt-characteristic-expansion}
\end{align}
Here $D_{B\subset A,\lambda}(z)=\Delta_{A,\lambda}^{-iz}\Delta_{B,\lambda}^{iz}$, and $F_\lambda$ is the unnormalized matrix element. Even when both regions are wedges, only the spacetime action remains the same as in the CFT: the vacuum correlation function generally changes with $\lambda$.

For real $t$, wedge modular conjugation sends $O_1(x)$ to its reflected operator $O_1(x_J)$, with any internal conjugation included. Since the wedge flow is geometric, Tomita's relation gives
\begin{equation}
 F_\lambda(t)=\langle\Omega_\lambda|
 \beta_A^t(O_1(x_J))\,
 \beta_B^t\bigl(u_{t,\lambda}O_2(y)u_{t,\lambda}^\dagger\bigr)
 |\Omega_\lambda\rangle.
 \label{eq:cpt-real-characteristic}
\end{equation}
We fix the perturbative identification so that the displayed wedge action and reflected external operator have no explicit coupling variation. Any additional external-operator counterterms must be included if a different identification is used.

Choose normalized vacua with $\langle\Omega_0|\Omega^{(1)}\rangle=0$. The Euclidean preparation in \eqref{eq:cpt-deformation} gives
\begin{align}
 |\Omega_\lambda\rangle&=|\Omega_0\rangle+\lambda|\Omega^{(1)}\rangle+O(\lambda^2),\nonumber\\
 |\Omega^{(1)}\rangle&=-\int_{\mathbb R}dX\int_{-\infty}^0d\tau\,
 O_{0,c}(-i\tau,X)|\Omega_0\rangle,
 \qquad O_{0,c}=O_0-\langle O_0\rangle_0.
 \label{eq:cpt-vacuum-variation}
\end{align}
This is a regulated preparation formula; an infinite-volume derivative need not exist as a Hilbert-space vector. The vacuum subtraction accounts for normalization.

Set $G_{12}(t)=\beta_A^t(O_1(x_J))\beta_B^t(O_2(y))$. There are three first-order contributions:
\begin{align}
 F^{(1)}(t)&=F^{(1)}_{\mathrm{flow}}(t)
 +F^{(1)}_{\mathrm{bra}}(t)+F^{(1)}_{\mathrm{ket}}(t),\nonumber\\
 F^{(1)}_{\mathrm{flow}}(t)
 &=\langle\Omega_0|\beta_A^t(O_1(x_J))
 \beta_B^t([u_t^{(1)},O_2(y)])|\Omega_0\rangle,\nonumber\\
 F^{(1)}_{\mathrm{bra}}(t)&=\langle\Omega^{(1)}|G_{12}(t)|\Omega_0\rangle,\nonumber\\
 F^{(1)}_{\mathrm{ket}}(t)&=\langle\Omega_0|G_{12}(t)|\Omega^{(1)}\rangle.
 \label{eq:cpt-three-corrections}
\end{align}
Using invariance of the CFT vacuum under the full unperturbed modular flow, the flow term becomes
\begin{align}
 F^{(1)}_{\mathrm{flow}}(t)
 &=i\pi\int_{-1}^1dX\int_{\mathbb R}du\,
 (1-X^2)Q(u,X;t)\Omega^{-\Delta}(u,X)\nonumber\\
 &\quad\times\langle\Omega_0|
 (\beta_B^{-t}\circ\beta_A^t)(O_1(x_J))
 [O_0(t_u,x_u),O_2(y)]|\Omega_0\rangle.
 \label{eq:cpt-flow-three-point}
\end{align}
All three terms reduce to ordered CFT three-point functions, together with the required subtractions and contact terms. For scalar primaries their separated-point form is fixed up to CFT data:
\begin{equation}
 \langle O_1(x_1)O_2(x_2)O_3(x_3)\rangle_0
 =\frac{C_{123}}
 {(x_{12}^2)^{(\Delta_1+\Delta_2-\Delta_3)/2}
  (x_{23}^2)^{(\Delta_2+\Delta_3-\Delta_1)/2}
  (x_{31}^2)^{(\Delta_3+\Delta_1-\Delta_2)/2}}.
 \label{eq:cpt-scalar-three-point}
\end{equation}
For the displayed Lorentzian ordering,
\begin{equation}
 x_{ij}^2=-(t_i-t_j-i(\epsilon_i-\epsilon_j))^2+(x_i-x_j)^2,
 \qquad \epsilon_1>\epsilon_2>\epsilon_3\to0^+.
\end{equation}
The two terms in the commutator have their respective ordering prescriptions. Conformal Jacobians follow from the modular transformations in \eqref{eq:cpt-three-corrections} and \eqref{eq:cpt-flow-three-point}. The dimensions and $C_{123}$, as well as possible contact terms, depend on the CFT; the integrated correction is therefore not theory-independent.

Microcausality restricts \eqref{eq:cpt-flow-three-point} to insertions $O_0$ that are causally related to $O_2$. Moving $O_2$ toward the entangling surface can reduce this region, suggesting a local boost approximation. Its volume alone does not bound the correction, since the kernel and correlator can become singular in the same limit. Section~\ref{section:approx} gives a separate nonperturbative estimate with specified localization and topology.

Finally, normalize as in the preceding section:
\begin{equation}
 f_\lambda=\frac{F_\lambda}{N_\lambda},
 \qquad f^{(1)}=\frac{F^{(1)}}{N_0}-\frac{F_0N^{(1)}}{N_0^2},
 \qquad N^{(1)}=\left.\partial_\lambda N_\lambda\right|_{\lambda=0}.
 \label{eq:cpt-normalized-correction}
\end{equation}
For each coupling, finite quarter-modular norms ensure that the exact $f_\lambda$ is analytic and bounded by one on the centered strip. Analyticity of its perturbative coefficient additionally requires differentiation locally uniformly in that strip. Null singularities of individual three-point integrands do not by themselves establish singularities of the integrated answer: contour pinches, ordering prescriptions, and renormalization must be analyzed. Smearing controls the external fields, but does not automatically remove every contact divergence. We therefore obtain the real-time first-order formula above without claiming an additional complex singularity classification.

\section{Nonperturbative bounds on subregion flow of balls in QFT}\label{section:approx}

We now compare subregion flows of large double cones with those of wedges. The physical idea is that a fixed local observable remains inside the large double cone for an increasingly long interval of wedge modular time. Fredenhagen's estimate \cite{fredenhagen} turns this geometric fact into a bound on time-averaged modular evolution. We first state a uniform estimate on localized vectors, then distinguish it from strong convergence on each fixed vector for an increasing exhaustion of a wedge.

Work in a vacuum representation of a Poincar\'e-covariant local net, and assume the Bisognano--Wichmann property in addition to isotony and locality. The vacuum is cyclic and separating for the double-cone and wedge algebras used below. Wedge algebras are generated by their bounded subregions. These assumptions make both geometric wedge flow and modular convergence available; geometric modular action does not follow from the Haag--Kastler axioms alone.

Let
\begin{equation}
 W_0=\{x:|x_0|<x_1\},\qquad
 W_r=\{x:|x_0|<x_1-r\},\qquad r>0.
\end{equation}
If $T_r$ implements translation by $r$ along $x_1$, covariance and vacuum invariance give
\begin{equation}
 \mathcal A(W_r)=\operatorname{Ad}_{T_r}\mathcal A(W_0),
 \qquad \Delta_{W_r}^{it}=T_r\Delta_{W_0}^{it}T_r^*.
\end{equation}
Write $B_{c,R}$ for the double cone whose $x_0=0$ section is the ball of radius $R$ centered at $(c,0,\ldots,0)$, with $c\geq R>0$. We use
\begin{equation}
 B_R:=B_{R,R},\qquad B_R^r:=T_r(B_R),\qquad
 B_{e^{\pm2\pi s}R}:=B_{e^{\pm2\pi s}R,e^{\pm2\pi s}R}.
 \label{eq:approx-ball-notation}
\end{equation}
Thus a dilation changes both the center and radius. The shrinking cones approach the wedge vertex, while the increasing family $B_{e^{2\pi s}R}$ exhausts $W_0$. A shrinking ball with its center fixed at $R$ would describe a different geometry.

\subsection{Fredenhagen's estimate and an increasing-ball version}

For a local operator $A$, abbreviate
\begin{equation}
 q(A)=\bigl(\|A\Omega\|^2+\|A^*\Omega\|^2\bigr)^{1/2}.
\end{equation}
Fredenhagen's approximation takes the following form \cite{fredenhagen}.
\begin{theorem}
For $f\in L^1(\mathbb R)$ and $s>0$, there is $c_f(s)\geq0$, with $c_f(s)\to0$ as $s\to\infty$, such that
\begin{equation}
 \left\|\int_{\mathbb R}dt\,f(t)
 (\Delta_{B_R}^{-it}-\Delta_{W_0}^{-it})A\Omega\right\|^2
 \leq c_f(s)q(A)^2,
 \qquad A\in\mathcal A(B_{e^{-2\pi s}R}).
\end{equation}
\end{theorem}
This is a Hilbert-space norm estimate for a modular-time average, not an operator-norm estimate or a pointwise-in-time bound. Shrinking the support toward the wedge vertex lengthens the interval during which its boosted images remain in $B_R$.

The input to this estimate is the following resolvent bound.
\begin{proposition}
Let $\mathcal B\subset\mathcal A$ have a common cyclic and separating vector $\Omega$. If $O\in\mathcal B$ and
\[
 \Delta_\mathcal A^{it}O\Delta_\mathcal A^{-it}\in\mathcal B
 \quad\text{for all }|t|<t_0,
\]
then
\begin{equation}
 \left\|\bigl[(1+\Delta_\mathcal A)^{-1}
 -(1+\Delta_\mathcal B)^{-1}\bigr]O\Omega\right\|
 \leq\eta(t_0)q(O),
 \qquad \eta(t_0)=\frac2\pi\operatorname{arctanh}(e^{-\pi t_0}).
 \label{eq:approx-resolvent}
\end{equation}
\end{proposition}
\begin{proof}
Equip $\operatorname{Dom}S_\mathcal A$ with the Tomita graph norm and let $I$ be its inclusion into $\mathcal H$. The subspace $\operatorname{Dom}S_\mathcal B$ is closed in this graph norm; denote its orthogonal projection by $P_\mathcal B$. The graph-projection identity is
\[
 (1+\Delta_\mathcal A)^{-1}-(1+\Delta_\mathcal B)^{-1}
 =I(1-P_\mathcal B)I^*.
\]
In the Mellin representation of the resolvent, the orbit assumption removes the part with $|t|<t_0$. The remaining modular orbit has graph norm $q(O)$, and $\|I\|\leq1$, giving
\[
 \left\|I(1-P_\mathcal B)I^*O\Omega\right\|
 \leq\frac12\int_{|t|\geq t_0}\frac{dt}{|\sinh(\pi t)|}\,q(O)
 =\eta(t_0)q(O).
\]
This is the resolvent argument of Fredenhagen; the graph identity and Mellin representation are also given in \cite[Section~3]{Hollands2019Edge}.
\end{proof}

For the remainder of the section the support stays fixed, while the containing ball grows.
\begin{corollary}
For $O\in\mathcal A(B_R)$,
\begin{equation}
 \left\|\bigl[(1+\Delta_{B_{e^{2\pi s}R}})^{-1}
 -(1+\Delta_{W_0})^{-1}\bigr]O\Omega\right\|
 \leq\eta(s)q(O).
 \label{eq:approx-growing-resolvent}
\end{equation}
\end{corollary}
\begin{proof}
The wedge boost sends $B_R$ into $B_{e^{2\pi s}R}$ for $|t|<s$. Isotony gives the orbit condition in \eqref{eq:approx-resolvent}, with $\mathcal A=\mathcal A(W_0)$ and $\mathcal B=\mathcal A(B_{e^{2\pi s}R})$.
\end{proof}

Applying the same time-averaging argument to this orbit condition gives the required version of Fredenhagen's theorem.
\begin{theorem}\label{theorem:modifiedfredenhagen}
For $f\in L^1(\mathbb R)$ and $A\in\mathcal A(B_R)$, there is $c_f(s)\geq0$, tending to zero as $s\to\infty$, such that
\begin{equation}
 \left\|\int_{\mathbb R}dt\,f(t)
 (\Delta_{B_{e^{2\pi s}R}}^{-it}-\Delta_{W_0}^{-it})A\Omega\right\|^2
 \leq c_f(s)q(A)^2.
 \label{eq:approx-growing-flow}
\end{equation}
\end{theorem}
\begin{proof}
The proof of the time-averaged estimate uses the modular orbit condition and its resolvent bound. These hold here by the preceding corollary, and $\eta(s)\to0$. Approximation of $f$ in $L^1$ extends the estimate from smooth test functions to all $L^1$ smearings \cite{fredenhagen}.
\end{proof}
Increasing the ambient region suppresses the error for fixed localized vectors and fixed time smearing. This expresses a local boost approximation without assuming a conformal symmetry of the QFT.

\subsection{A double cone inside a wedge}\label{subsection:coneinsidewedge}

Set $B_s^r=B_{e^{2\pi s}R}^r$ and compare
\begin{equation}
 D_{B_s^r\subset W_0}(t)=\Delta_{W_0}^{-it}\Delta_{B_s^r}^{it},
 \qquad D_{W_r\subset W_0}(t)=\Delta_{W_0}^{-it}\Delta_{W_r}^{it}.
\end{equation}
The characteristic operator is a contraction analytic on $0<\operatorname{Im}z<1/2$, with strongly continuous unitary boundary values. Its wedge limit is explicitly geometric, as in Section~\ref{sec:geometric-subregion-flows}.

To obtain a uniform localized bound after multiplication by the outer wedge flow, introduce the spectral family of $L=\log\Delta_{W_0}=-K_{W_0}$:
\begin{equation}
 \Delta_{W_0}^{it}=\int_{\mathbb R}e^{i\mu t}\,dE(\mu),
 \qquad E(\mu)=\mathbf1_{(-\infty,\mu]}(L),
 \qquad F(\mu_0)=E(\mu_0)-E(-\mu_0).
 \label{eq:approx-spectral-cutoff}
\end{equation}
Thus $F(\mu_0)$ projects onto $(-\mu_0,\mu_0]$. The spectrum of $L$ is not restricted to nonnegative values.

\begin{theorem}\label{theorem:mainresult}
Let $f,tf\in L^1(\mathbb R)$, $\mu_0>0$, and $A\in\mathcal A(B_R^r)$. Then
\begin{equation}
 \left\|F(\mu_0)\int_{\mathbb R}dt\,f(t)
 [D_{B_s^r\subset W_0}(t)-D_{W_r\subset W_0}(t)]A\Omega\right\|^2
 \leq b_f(\mu_0,s)q(A)^2,
 \label{eq:approx-characteristic-bound}
\end{equation}
where $b_f(\mu_0,s)\to0$ for each fixed $f$ and $\mu_0$ as $s\to\infty$. The assumptions include Schwartz functions and compactly supported bounded functions.
\end{theorem}

\begin{lemma}\label{lemma:extensionfredenhagen}
For $g\in L^1(\mathbb R)$ and $A\in\mathcal A(B_R^r)$,
\begin{equation}
 \left\|\int_{\mathbb R}dt\,g(t)
 (\Delta_{B_s^r}^{it}-\Delta_{W_r}^{it})A\Omega\right\|^2
 \leq\widetilde c_g(s)q(A)^2,
 \qquad \widetilde c_g(s)\longrightarrow0.
 \label{eq:approx-translated-flow}
\end{equation}
The constants can be chosen so that $\widetilde c_g(s)^{1/2}\leq2\|g\|_1$ and their convergence is uniform on compact subsets of $L^1(\mathbb R)$.
\end{lemma}
\begin{proof}
Translate $A$ back to $\mathcal A(B_R)$ and use
\[
 \Delta_{B_s^r}^{it}=T_r\Delta_{B_{e^{2\pi s}R}}^{it}T_r^*,
 \qquad T_r\Omega=\Omega.
\]
Translation preserves both terms in $q(A)$. Replacing $t$ by $-t$ in \eqref{eq:approx-growing-flow} supplies the positive-exponent estimate, with the reflected smearing $g(-t)$.

Choose $\widetilde c_g(s)^{1/2}$ to be the optimal bound on the ratio of the norm in \eqref{eq:approx-translated-flow} to $q(A)$. Unitarity gives, for these bounds $a_s(g)$,
\[
 a_s(g)\leq2\|g\|_1,
 \qquad |a_s(g)-a_s(h)|\leq2\|g-h\|_1.
\]
Pointwise convergence and a finite $L^1$ cover of a compact family then give uniform convergence on that family.
\end{proof}

\begin{proof}[Proof of Theorem~\ref{theorem:mainresult}]
Write $V_s(t)=\Delta_{B_s^r}^{it}-\Delta_{W_r}^{it}$ and $f_\mu(t)=e^{-i\mu t}f(t)$. Spectral integration by parts, including both endpoint terms, gives
\begin{equation}
 F(\mu_0)e^{-itL}
 =e^{-i\mu_0t}E(\mu_0)-e^{i\mu_0t}E(-\mu_0)
 +it\int_{-\mu_0}^{\mu_0}d\mu\,e^{-i\mu t}E(\mu).
 \label{eq:approx-spectral-parts}
\end{equation}
Since $\|V_s(t)\|\leq2$ and $f,tf\in L^1$, the resulting ordinary $t$ and $\mu$ integrals are absolutely bounded in norm. Fubini's theorem therefore applies without a formal interchange of two noncommuting spectral measures. Using \eqref{eq:approx-translated-flow},
\begin{align}
 &\left\|F(\mu_0)\int_{\mathbb R}dt\,f(t)e^{-itL}V_s(t)A\Omega\right\|
 \leq B_f(\mu_0,s)q(A),\nonumber\\
 B_f(\mu_0,s)
 &=\widetilde c_{f_{\mu_0}}(s)^{1/2}
 +\widetilde c_{f_{-\mu_0}}(s)^{1/2}
 +\int_{-\mu_0}^{\mu_0}d\mu\,\widetilde c_{tf_\mu}(s)^{1/2}.
 \label{eq:approx-bound-amplitude}
\end{align}
The families $\{f_\mu\}$ and $\{tf_\mu\}$ with $|\mu|\leq\mu_0$ are compact in $L^1$, by continuity in $\mu$. The lemma implies $B_f\to0$, and the squared-norm constant in \eqref{eq:approx-characteristic-bound} is $b_f=B_f^2$.
\end{proof}

For the qualitative limit on a fixed vector, a spectral cutoff is unnecessary. Indeed, the increasing algebras satisfy
\begin{equation}
 \bigvee_{s>0}\mathcal A(B_s^r)=\mathcal A(W_r).
\end{equation}
The modular convergence theorem for increasing algebras with a common cyclic and separating vector gives $\Delta_{B_s^r}^{it}\to\Delta_{W_r}^{it}$ strongly \cite{revolutionizing,takesaki2}. Consequently, for every $\psi\in\mathcal H$ and $f\in L^1(\mathbb R)$,
\begin{equation}
 \lim_{s\to\infty}\left\|\int_{\mathbb R}dt\,f(t)
 [D_{B_s^r\subset W_0}(t)-D_{W_r\subset W_0}(t)]\psi\right\|=0.
 \label{eq:approx-uncut-wedge-limit}
\end{equation}
Unitarity supplies the dominating bound $2|f(t)|\|\psi\|$. This strong limit gives no rate uniform over all vectors; the preceding theorem supplies uniform control on the specified localized vectors after the spectral cutoff. Its constants need not give an explicit decay rate; the two statements differ in their uniformity and allowed vectors.

\subsection{A double cone inside a double cone}\label{subsection:coneinsidecone}

Let $B_{s,u}=B_{e^{2\pi(s+u)}R}$ and retain $B_s^r=B_{e^{2\pi s}R}^r$. The inclusion diagram is
\begin{equation}
 \begin{tikzcd}
 B_s^r\arrow[r,"\subset"]\arrow[d,"\subset"]
 &B_{s,u}\arrow[d,"\subset"]\\
 W_r\arrow[r,"\subset"]&W_0.
 \end{tikzcd}
\end{equation}
For these balls, strict containment is ensured by
\begin{equation}
 e^{2\pi(s+u)}R>e^{2\pi s}R+\frac r2.
 \label{eq:approx-ball-containment}
\end{equation}
This follows by comparing the center separation with the difference of radii. The stronger condition with $r$ in place of $r/2$ also suffices.

We compare
\begin{equation}
 D_{B_s^r\subset B_{s,u}}(t)=\Delta_{B_{s,u}}^{-it}\Delta_{B_s^r}^{it}
 \quad\text{with}\quad D_{W_r\subset W_0}(t).
\end{equation}

\begin{theorem}\label{theorem:ballinsideball}
Under the assumptions of this section, for every $\psi\in\mathcal H$ and
$f\in L^1(\mathbb R)$,
\begin{equation}
 \lim_{\substack{s\to\infty,\ s+u\to\infty\\ B_s^r\subset B_{s,u}}}
 \left\|\int_{\mathbb R}dt\,f(t)
 [D_{B_s^r\subset B_{s,u}}(t)-D_{W_r\subset W_0}(t)]\psi\right\|=0.
 \label{eq:approx-ball-ball-limit}
\end{equation}
Thus both radii tend to infinity while the inclusion is maintained. In particular,
the iterated limit $u\to\infty$ followed by $s\to\infty$ holds.
\end{theorem}
\begin{proof}
Split the difference as
\begin{align}
 D_{B_s^r\subset B_{s,u}}(t)-D_{W_r\subset W_0}(t)
 &=(\Delta_{B_{s,u}}^{-it}-\Delta_{W_0}^{-it})\Delta_{W_r}^{it}\nonumber\\
 &\quad+\Delta_{B_{s,u}}^{-it}
 (\Delta_{B_s^r}^{it}-\Delta_{W_r}^{it}).
\end{align}
As their radii increase, the two families exhaust $W_0$ and $W_r$, respectively.
Modular strong convergence applies to the fixed vector $\Delta_{W_r}^{it}\psi$
in the first term and to $\psi$ in the second. The prefactor in the second term
is unitary. The norm of the total difference is bounded by $2\|\psi\|$,
so $2|f(t)|\|\psi\|$ gives an integrable bound for dominated convergence.
\end{proof}
This states the large-ball approximation on each fixed vector, without an
operator-norm limit or a rate uniform over the Hilbert-space unit ball.

\section{Discussion}


%\section{Characteristic function as a regulator, double cone, and Modular quasi-normal modes}




% \section{1D Lattice examples}


% Consider the one-dimensional quantum XY-model with the Hamiltonian
% \begin{eqnarray}
%     H--\frac{1}{2}\sum_i\lb \frac{1+\gamma}{2}X_i X_{i+1}+\frac{1-\gamma}{2}Y_i Y_{i+1}+\lambda Z_i\rb\ .
% \end{eqnarray}
% The cases $\gamma=0$ and $\gamma=1$ are called the XX model and the Ising model, respectively. This model can be solved using the Jordan-Wigner transform 
% that sends it to a second-order action in fermionic modes.

% We start with the case $\lambda=1$, the transverse Ising model. For $\lambda<1$ the system is in the disordered phase, and for $\lambda>1$ it is ordered. We define $k=\min(\lambda,\lambda^{-1})$ and $k'=\sqrt{1-k^2}$.

% The modular Hamiltonian of a half-line  $W_0=\{i| i\geq 0\}$ can be written explicitly as
% \begin{eqnarray}
% &&\lambda<1:\quad \hat{K}_{W(0)}=-2I(k')\sum_{j=0}^\infty \lb j X_jX_{j+1}+k (j-1/2) Z_j\rb\nn\\
% &&\lambda>1:\quad \hat{K}_{W(0)}=-2I(k')\sum_{j=0}^\infty \lb k j X_jX_{j+1}+ (j-1/2) Z_j\rb\nn
% \end{eqnarray}
% where $I(k')$ is a complete elliptic integral of the first kind. We would like to compute the characteristic function $\Delta_B^{1/4}\Delta_A^{-1/4}$ for half-spaces $A=W(a>0)\subset W(0)=B$. We can avoid the calculation of the full modular operator, by thinking in terms of the characteristic function as the following super-operator acting on the operators in the algebra:
% \begin{eqnarray}
% &&\Delta_B^{1/4}\Delta_A^{-1/4}a\ket{\Omega}=\mathcal{D}_{i/4}(a)\ket{\Omega}\nn\\
% &&\mathcal{D}_{i/4}(a)=e^{K_B/4}e^{-K_A/4} a e^{K_A/4}e^{-K_B/4}\ .
% \end{eqnarray}
% Define the operator $\tilde{K}=K/(2I(k'))$, then for $\lambda<1$ we have
% \begin{eqnarray}
% &&[\tilde{K},X_l]=-k(l-1/2)i Y_l\nn\\
% &&[\tilde{K},Z_l]=i (l Y_lX_{l+1}+(l-1)X_{l-1}Y_l)\nn\\
% &&[\tilde{K},Y_l]=-\ .
% \end{eqnarray}
% Use Jordan-Wigner transform to write it in terms of the fermions

%Similarly, the modular Hamiltonians of the commutants are
%\begin{eqnarray}
%&&\lambda<1:\quad \hat{K}_{W'(0)}=2I(k')\sum_{j=0}^{-\infty} \lb j X_jX_{j-1}+k (j+1/2) Z_j\rb\nn\\
%&&\lambda>1:\quad \hat{K}_{W'(0)}=2I(k')\sum_{j=0}^{-\infty} \lb k j X_jX_{j-1}+ (j+1/2) Z_j\rb\nn\ .
%\end{eqnarray}
%
%
% We first write down the full modular Hamiltonians for $\lambda<1$ 
%\begin{eqnarray}
%K_{W(0)}&&=\hat{K}_{W(0)}-\hat{K}_{W'(0)}=\nn\\
%&&=-2I(k')\sum_{j=\infty}^\infty \lb j X_jX_{j+\text{sgn}(j)}+k (j-\text{sgn}(j))/2) Z_j\rb\nn
%\end{eqnarray}
%and for $\lambda>1$ 
%\begin{eqnarray}
%K_{W(0)}&&=\hat{K}_{W(0)}-\hat{K}_{W'(0)}=\nn\\
%&&=-2I(k')\sum_{j=\infty}^\infty \lb k j X_jX_{j+\text{sgn}(j)}+ (j-\text{sgn}(j))/2) Z_j\rb\nn
%\end{eqnarray}
% 
%\NL{Finish the calculation; how about the XY-model?}

% As the second example, consider a 1d harmonic chain
% \begin{eqnarray}
% H=\sum_i \lb \frac{1}{2m}p_i^2+\frac{m\omega^2}{2}q_i^2+\frac{\kappa}{2}(q_i-q_{i+1})^2\rb
% \end{eqnarray}
% It is a discretization of massive scalar field in 1 dimensions. The modular Hamiltonian of a subregion $A$ of $L$ sites is
% \begin{eqnarray}
% &&\hat{K}_A=\frac{1}{2}\vec{r}^T\cdot {\bf K}_A \cdot \vec{r},\qquad \vec{r}=\begin{pmatrix}&\vec{q}\\& \vec{p}\end{pmatrix}
% \end{eqnarray}
% The matrix ${\bf K}_A$ is real, symmetric, and positive definite so that $\hat{K}_A$ is hermitian. Consider the two-point functions $Q_{ij}=\braket{q_i q_j}$ and $P_{ij}=\braket{p_i p_j}$. Since ${\bf Q}$ and ${\bf P}$ are symmetric we define ${\bf C}_A=\sqrt{{\bf Q}_A{\bf P}_A}$ and have $\sqrt{{\bf Q}_A{\bf P}_A}^T=\sqrt{{\bf P}_A{\bf Q}_A}$. Then,
% \begin{eqnarray}
% &&\hat{K}_A=({\bf P}_A\oplus {\bf Q}_A)\lb h({\bf C}_A)\oplus h({\bf C}^T_A)\rb\nn\\
% &&h(y)=\frac{1}{y}\log\lb \frac{y+1/2}{y-1/2}\rb
% \end{eqnarray}
% Therefore, we can split $\hat{K}_A$ into two terms
% \begin{eqnarray}
% &&{\bf K}_A=\hat{K}^{(M)}_A+\hat{K}^{(N)}_A\nn\\
% &&{\bf K}^{(M)}_A=L\sum_{ij}\frac{M_{ij}}{L}q_i q_j,\qquad {\bf K}_A^{(N)}=L\sum_{ij}\frac{N_{ij}}{L}\pi_i \pi_j,
% \end{eqnarray}
% where we have multiplied and dividing by a factor of $L$ to compare to the continuum limit. Note that in the continuum limit which motivates the existence of the following limits
% \begin{eqnarray}
% &&\mu_k(x_k)=\lim_{L\to \infty}\frac{M_{i,i+k}}{L},\qquad \nu_k(x_k)=\lim_{L\to \infty}\frac{N_{i,i+k}}{L}\nn\\
% &&x_k=\frac{1}{L}\lb i+k/2\rb\ .
% \end{eqnarray}

% The calculation of full modular Hamiltonian is somewhat 

% Take the transverse Ising model and compute the characteristic functions for half-space inclusions.

% Take massless free bosons and fermions and compute the characteristic function for an interval



\NL{Kubo-Ando mean for inclusions in QFT Rindler wedges, and }

\NL{Zeros of characteristic function, and a chaos bound}

%\section{Characteristic function as a regulator, double cone, and Modular quasi-normal modes}




% \section{1D Lattice examples}


% Consider the one-dimensional quantum XY-model with the Hamiltonian
% \begin{eqnarray}
%     H--\frac{1}{2}\sum_i\lb \frac{1+\gamma}{2}X_i X_{i+1}+\frac{1-\gamma}{2}Y_i Y_{i+1}+\lambda Z_i\rb\ .
% \end{eqnarray}
% The cases $\gamma=0$ and $\gamma=1$ are called the XX model and the Ising model, respectively. This model can be solved using the Jordan-Wigner transform 
% that sends it to a second-order action in fermionic modes.

% We start with the case $\lambda=1$, the transverse Ising model. For $\lambda<1$ the system is in the disordered phase, and for $\lambda>1$ it is ordered. We define $k=\min(\lambda,\lambda^{-1})$ and $k'=\sqrt{1-k^2}$.

% The modular Hamiltonian of a half-line  $W_0=\{i| i\geq 0\}$ can be written explicitly as
% \begin{eqnarray}
% &&\lambda<1:\quad \hat{K}_{W(0)}=-2I(k')\sum_{j=0}^\infty \lb j X_jX_{j+1}+k (j-1/2) Z_j\rb\nn\\
% &&\lambda>1:\quad \hat{K}_{W(0)}=-2I(k')\sum_{j=0}^\infty \lb k j X_jX_{j+1}+ (j-1/2) Z_j\rb\nn
% \end{eqnarray}
% where $I(k')$ is a complete elliptic integral of the first kind. We would like to compute the characteristic function $\Delta_B^{1/4}\Delta_A^{-1/4}$ for half-spaces $A=W(a>0)\subset W(0)=B$. We can avoid the calculation of the full modular operator, by thinking in terms of the characteristic function as the following super-operator acting on the operators in the algebra:
% \begin{eqnarray}
% &&\Delta_B^{1/4}\Delta_A^{-1/4}a\ket{\Omega}=\mathcal{D}_{i/4}(a)\ket{\Omega}\nn\\
% &&\mathcal{D}_{i/4}(a)=e^{K_B/4}e^{-K_A/4} a e^{K_A/4}e^{-K_B/4}\ .
% \end{eqnarray}
% Define the operator $\tilde{K}=K/(2I(k'))$, then for $\lambda<1$ we have
% \begin{eqnarray}
% &&[\tilde{K},X_l]=-k(l-1/2)i Y_l\nn\\
% &&[\tilde{K},Z_l]=i (l Y_lX_{l+1}+(l-1)X_{l-1}Y_l)\nn\\
% &&[\tilde{K},Y_l]=-\ .
% \end{eqnarray}
% Use Jordan-Wigner transform to write it in terms of the fermions

%Similarly, the modular Hamiltonians of the commutants are
%\begin{eqnarray}
%&&\lambda<1:\quad \hat{K}_{W'(0)}=2I(k')\sum_{j=0}^{-\infty} \lb j X_jX_{j-1}+k (j+1/2) Z_j\rb\nn\\
%&&\lambda>1:\quad \hat{K}_{W'(0)}=2I(k')\sum_{j=0}^{-\infty} \lb k j X_jX_{j-1}+ (j+1/2) Z_j\rb\nn\ .
%\end{eqnarray}
%
%
% We first write down the full modular Hamiltonians for $\lambda<1$ 
%\begin{eqnarray}
%K_{W(0)}&&=\hat{K}_{W(0)}-\hat{K}_{W'(0)}=\nn\\
%&&=-2I(k')\sum_{j=\infty}^\infty \lb j X_jX_{j+\text{sgn}(j)}+k (j-\text{sgn}(j))/2) Z_j\rb\nn
%\end{eqnarray}
%and for $\lambda>1$ 
%\begin{eqnarray}
%K_{W(0)}&&=\hat{K}_{W(0)}-\hat{K}_{W'(0)}=\nn\\
%&&=-2I(k')\sum_{j=\infty}^\infty \lb k j X_jX_{j+\text{sgn}(j)}+ (j-\text{sgn}(j))/2) Z_j\rb\nn
%\end{eqnarray}
% 
%\NL{Finish the calculation; how about the XY-model?}

% As the second example, consider a 1d harmonic chain
% \begin{eqnarray}
% H=\sum_i \lb \frac{1}{2m}p_i^2+\frac{m\omega^2}{2}q_i^2+\frac{\kappa}{2}(q_i-q_{i+1})^2\rb
% \end{eqnarray}
% It is a discretization of massive scalar field in 1 dimensions. The modular Hamiltonian of a subregion $A$ of $L$ sites is
% \begin{eqnarray}
% &&\hat{K}_A=\frac{1}{2}\vec{r}^T\cdot {\bf K}_A \cdot \vec{r},\qquad \vec{r}=\begin{pmatrix}&\vec{q}\\& \vec{p}\end{pmatrix}
% \end{eqnarray}
% The matrix ${\bf K}_A$ is real, symmetric, and positive definite so that $\hat{K}_A$ is hermitian. Consider the two-point functions $Q_{ij}=\braket{q_i q_j}$ and $P_{ij}=\braket{p_i p_j}$. Since ${\bf Q}$ and ${\bf P}$ are symmetric we define ${\bf C}_A=\sqrt{{\bf Q}_A{\bf P}_A}$ and have $\sqrt{{\bf Q}_A{\bf P}_A}^T=\sqrt{{\bf P}_A{\bf Q}_A}$. Then,
% \begin{eqnarray}
% &&\hat{K}_A=({\bf P}_A\oplus {\bf Q}_A)\lb h({\bf C}_A)\oplus h({\bf C}^T_A)\rb\nn\\
% &&h(y)=\frac{1}{y}\log\lb \frac{y+1/2}{y-1/2}\rb
% \end{eqnarray}
% Therefore, we can split $\hat{K}_A$ into two terms
% \begin{eqnarray}
% &&{\bf K}_A=\hat{K}^{(M)}_A+\hat{K}^{(N)}_A\nn\\
% &&{\bf K}^{(M)}_A=L\sum_{ij}\frac{M_{ij}}{L}q_i q_j,\qquad {\bf K}_A^{(N)}=L\sum_{ij}\frac{N_{ij}}{L}\pi_i \pi_j,
% \end{eqnarray}
% where we have multiplied and dividing by a factor of $L$ to compare to the continuum limit. Note that in the continuum limit which motivates the existence of the following limits
% \begin{eqnarray}
% &&\mu_k(x_k)=\lim_{L\to \infty}\frac{M_{i,i+k}}{L},\qquad \nu_k(x_k)=\lim_{L\to \infty}\frac{N_{i,i+k}}{L}\nn\\
% &&x_k=\frac{1}{L}\lb i+k/2\rb\ .
% \end{eqnarray}

% The calculation of full modular Hamiltonian is somewhat 

% Take the transverse Ising model and compute the characteristic functions for half-space inclusions.

% Take massless free bosons and fermions and compute the characteristic function for an interval
%\subsection{older stuff}
\NL{No zeros and recovery maps}

\NL{Characteristic Functions and Noncommutative $L_p$-Space}




% \section{CFT Hamiltonian from a ball}

% %Emphasize $T(z)T(-z)^{-1}$.

% Consider the vacuum state of a CFT in $\mathbb{R}^{d-1,1}$ conformally embedded inside the Lorentzian cylinder $\mathbb{S}^{d-1}\times \mathbb{R}_T$ with $T$ the global time on the cylinder; see figure \ref{}. The universal cover of the conformal group $\widetilde{SO(d+1,2)}$ acts on the Lorentzian cylinder. Consider concentric ball-shaped regions $B(s)$ centered at the $x^\mu=0$ and radius $e^{-2\pi s}$. 
% \begin{theorem}
%     For any real $t$ and $s>0$ we have the inclusion of ball-shaped region $B(t+s)\subset B(t)$ we have the characteristic function for $\Im (z)\in (0,1/2)$:
%     \begin{eqnarray}
%         &&\log T_{t,s}(z)=-(\beta_BK_B+\beta_H H+\beta_F K_F)\nn\\
%         &&\beta_F=\beta\cosh(\pi s)\coth(\pi z)\nn\\
%         &&\beta_B=\beta \sinh(\pi (2t+s))\nn\\
%         &&\beta_H=\beta\cosh(\pi (2t+s))\nn\\
%         &&\beta=\frac{2i\cosh^{-1}(A)}{\sqrt{\sinh^2(\pi z)\sinh^2(\pi s)-1}}\nn\\
%         &&A=1-2\sinh^2(\pi s)\sinh^2(\pi z)
%     \end{eqnarray}
% where $B=B(0)$ is a unit ball, $s=|B(t+s)|/|B(t)|$ is the relative size of the inclusion and $|B|$ is the linear size of $B$.
% \end{theorem}
% \begin{corollary}
% In this case, at small $s$ the characteristic function is
% \begin{eqnarray}
% T_{t,s}(z)&&=1-2\pi i s \lb  \frac{\sinh(\pi z)}{2\pi}\p_t K_{B(t)}+\cosh(\pi z)K_F\rb \nn\\
% &&+O(s^2)\nn\\
%  &&\frac{1}{2\pi}\p_t K_{B(t)}=\sinh(2\pi t)K_B+\cosh(2\pi t) H\nn\\
%     &&K_{B(t)}=\cosh(2\pi t)K_B+\sinh(2\pi t) H\ .
% \end{eqnarray}
% The modular S-matrix in this limit is
% \begin{eqnarray}
% S_t(z)=1-2 i s \sinh(\pi z)\p_t K_{B(t)}+O(s^2)\ .
% \end{eqnarray}
% Further expanding for large $|z|$ we obtain the scattering amplitude operator
% \begin{eqnarray}
% \tilde{M}_t=-2 s e^{\pi z}\p_t K_{B(t)}
% \end{eqnarray}
% which once again saturates the chaos bound.
% \end{corollary}


% A local observer on a unit ball $B$ has access to $K_B$. 


% Therefore, if we want to reconstruct the global Hamiltonian $H$ using only local data $K_B$ we need to remove the $K_F$ term. Under $z\to -z$ the coefficient $\beta_F$ is invariant whereas $\beta_B$ and $\beta_H$ are invariant. 

% We are particularly interested in the limit $s\ll 1$ and $s e^{\pi \Re z}\ll 1$. We define $T_{t,s}(z)=1+i M_{t,s}(z)$ so that in this limit we find
% \begin{eqnarray}
% %\beta&&\sim 4\pi i s \sinh(\pi z)\nn\\
%     &&M_{t,s}(z)=\nn\\
%     &&-4\pi s\lb \frac{\sinh(\pi z)}{2\pi}\p_t K_{B(t)}+\cosh(\pi z)K_F\rb+O(s^2)\nn\\
%     &&\frac{1}{2\pi}\p_t K_{B(t)}=\sinh(2\pi t)K_B+\cosh(2\pi t) H\nn\\
%     &&K_{B(t)}=\cosh(2\pi t)K_B+\sinh(2\pi t) H\ .
%     \end{eqnarray}
% Note that $M_{t,s}(z)$ is an operator-valued function, but its matrix elements are analogous to the scattering amplitude defined in scattering theory.

% The expansion above makes it clear that if we compute $M_{t,s}(z)$ 

% The connection to scattering theory suggests writing $z=x+iy$ and taking a further limit of large $x$ to find
% \begin{eqnarray}
%     M=-4\pi s e^{\pi z}\lb\sinh(2\pi t)K_B+\cosh(2\pi t)K_B+K_F\rb\nn
% \end{eqnarray}


% ----

% \begin{enumerate}
%     \item   The positive operator below is independent of $s$:
%     \begin{eqnarray}
%     \tilde{H}_{t}&&\equiv \frac{1}{4\pi s}T_{t,s}(-i/4)T_{t,s}(-i/4)^\dagger \nn\\
%     &&=\cosh(2\pi t)H+\sinh(2\pi t)K_B\nn\\
%     &&=\frac{1}{2\pi}\p_t K_{B(t)}\Big|_{t=0}=-R\p_R K_{B(R)}\ .
%     \end{eqnarray}
%     \item The boundary value at $z=-i/2$ is
%     \begin{eqnarray}
%         T_{t,s}(-i/2)=4\pi s (\tanh(\pi s)K_B+H)\ .
%     \end{eqnarray}
%     \item We can construct both the global Hamiltonian $J_{(-1)0}$ and the generator of rotations on the Lorentzian cylinder in terms of the local modular Hamiltonian
%        \begin{eqnarray}
%     \begin{pmatrix}
%     J_{(-1)0}\\
%         J_{0(d+1)}
%     \end{pmatrix}
%     &&=\begin{pmatrix}
%         &\cosh(2\pi t) & \sinh(2\pi t)  \\
%          &\sinh(2\pi t)& \cosh(2\pi t)
%     \end{pmatrix}
%     \begin{pmatrix}
%         & \tilde{H}_{t} \\
%         & -K_{B(t)}
%     \end{pmatrix}\nn\\
%     &&=\frac{-1}{2R}\begin{pmatrix}
%         &1+R^2 & 1-R^2  \\
%          &1-R^2& 1+R^2
%     \end{pmatrix}
%     \begin{pmatrix}
%         & R\p_R K_{B(R)}\\
%         & K_{B(R)}
%     \end{pmatrix}\ .
% \end{eqnarray}
%     \end{enumerate}




% \begin{eqnarray}
%     T_{t,s}(x)T_{t,s}(-x)^{-1}=\mathbb{I}+2\pi s\sinh(2\pi x)\begin{pmatrix}
%         &0 & -1  \\
%          &1& 0
%     \end{pmatrix}+O(s^2)\nn\ .
% \end{eqnarray}
% This means that $T_{t,s}(x) T_{t,s}(-x)^{-1}$ generates an infinitesimal step of real-time evolution $e^{i s H}$. Interestingly, this is independent of $t$.


% Our result generalize almost trivially to Conformally flat spacetimes. The results generalize to conformally flat spacetimes including double cones in de Sitter and Anti-de Sitter space \cite{frob2023modular}. In curved spacetime, ...

% Note that nuclearity implies that $e^{-\ep \tilde{H}}$ is trace class for all $\ep>0$, which in turn implies that the spectrum of $\tilde{H}$ is discrete \footnote{Note that in Minkowski space QFT, $e^{-\beta H}$ is not trace class, and the spectrum of Hamiltonian contains cuts corresponding to the continuum of multi-particle states.}. 

% Next, consider the eternal black hole in AdS in the limit of $G_N\to 0$. The boundary theory is a GFF with time-band algebras. Assuming spherical symmetry, we can view time-band algebras as time-interval algebras of $0+1$ GFF with an infinite tower of fields corresponding to the spherical harmonics. Then, the assumption of nuclearity implies that the split property, which in turn implies that the emergent type III$_1$ algebras of GFF are hyperfinite. This is related to the conjecture of Sam and Hong \cite{}.






% \section{Local approximation to global QFT Hamiltonian}

% In this section, we would like to use our reconstruction of the global Hamiltonian from the reduced state on the ball to a general QFT to define a local approximation to the global Hamiltonian of QFT. As in the case of CFT, we start with an inclusion $B_s\in B$ of balls of radius $R e^{2\pi s}\leq R$. To simplify our discussion, we assume translation invariance of the QFT state so that we can denote the characteristic function of this inclusion by $T_{R,s}(z)$, and compute $\chi_{R,s}(z)=T_{R,s}(z)T_{R,s}(-z)^{-1}$ for $\Im z\in (0,1/4)$. 
% While the CFT result holds for the inclusion of balls $B_s\in B$ with arbitrary relative size $s\sim \log(|B_s|/|B|)$. Guided by Theorem 3 \ref{} and the connection with the modular chaos bound we will discuss in the next section, in a QFT, we focus our attention on the small limit of small $s$:
% \begin{eqnarray}
%     \chi_{R,s}(z)=1+s \tilde{H}_R(z)+O(s^2)\ .
% \end{eqnarray}
% We make the following two proposals for the local approximation to the global Hamiltonian of a CFT from the inclusion of balls $B_s\in B$ for $z=t+i \alpha$. 
% \begin{enumerate}
%     \item For real $z=t$ we propose $\tilde{H}_R(t)$ for $e^{2\pi t}\ll 1/s$ is an approximation of the global Hamiltonian.
%     \item For all $\alpha\in (0,1/4)$ we propose $\tilde{H}_R(\alpha)...$ as an approximation of the global Hamiltonian.
% \end{enumerate}
% The general expectation that we will confirm in this section is that $\tilde{H}_R$ is 





% Consider QFT at finite volume; a sphere of radius $L$. Then, we propose the following operator as the local approximation to the global Hamiltonian in QFT using the local state on a sphere of radius $R$; $B(R)$:
% \begin{eqnarray}
%     H(R)=\frac{-L}{2R}\lb (1+(R/L)^2)R\p_R+(1-(R/L)^2)\rb K_{B(R)}\nn\ .
% \end{eqnarray}
% At a CFT fixed point, this operator matches the global Hamiltonian of the theory.

% In this section, we show that our proposed Hamiltonian is stable towards perturbations meaning that 1) the action of this operator on probes localized in the vicinity of the boundary $B(R)$ is close to the actual Hamiltonian of the theory. 2) In the excited states of CFT and the vacuum of deformed CFT, our Hamiltonian faithfully captures the high-energy eigenvaulues and eigenvectors of the global Hamiltonian.



% In conformal perturbation theory, one can show that the high-energy spectrum and eigenvectors of this operator are close to the actual Hamiltonian of the theory. See appendix \ref{}.


% %\subsection{New attempt}

% Consider the modular Hamiltonian of a ball of radius $R$ in conformal perturbation theory: $K_{B(R)}(\lambda)$ with $R=e^{-2\pi t}$. It is a two-parameter family of self-adjoint operators $K(t,\lambda)$, and we would like to compute
% \begin{eqnarray}
%     &&\tilde{H}_t(\lambda)\equiv \partial_s\log\lb e^{-\pi K(t,\lambda)/2} e^{\pi K(t+s,\lambda)} e^{-\pi K(t,\lambda)/2}\rb\Big|_{s=0}\ .
% \end{eqnarray}
% First, we define
% \begin{eqnarray}
%     A(s)\equiv e^{-\pi K(t,\lambda)/2} e^{\pi K(t+s,\lambda)} e^{-\pi K(t,\lambda)/2}
% \end{eqnarray}
% Then,
% \begin{eqnarray}
%     \p_s \log A(s)\Big|_{s=0}&&=\int_0^1 du \lb A(s)^u \: (\p_s A(s))\: A(s)^{-u}\rb \Big|_{s=0}= \p_sA(s)\Big|_{s=0}\ .
% \end{eqnarray}
% It follows that 
% \begin{eqnarray}\label{Hamiltonianexpress}
%     \tilde{H}_\lambda&&=\Delta_{B(t)}^{1/4}(\p_s \Delta^{-1/2}_{B(s+t)})_{s=0} \Delta_{B(t)}^{1/4}\nn\\
%     &&=\Delta_{B(t)}^{-1/4}(\p_s\Delta^{1/2}_{B(s+t)})_{s=0}\Delta_{B(t)}^{-1/4}\nn\\
%     &&=-\frac{1}{2\pi}\Delta_{B(R)}^{-1/4}(R\p_R\Delta^{1/2}_{B(R)})\Delta_{B(R)}^{-1/4}\ .
% \end{eqnarray}

% ---

% \begin{lemma}
%     We have
%     \begin{eqnarray}
%         &&\lb\Delta_s^{-\alpha}\p_s \Delta_s^{2\alpha} \Delta_s^{-\alpha}\rb_{s=0}=-\alpha\int_{-\alpha/2}^{\alpha/2} du \:\Delta^{-u}(\p_s\log \Delta)_{s=0}\Delta^u\nn\\
%         &&=\frac{-\alpha\pi}{2}\int_{-\alpha/2}^{\alpha/2}du\int_{-\infty}^\infty\frac{ dt}{\cosh^2(\pi t)}\:\Delta^{1/2+it+u}(\Delta^{-1}(\p_s\Delta))\Delta^{-(1/2+it+u)}\ .
%     \end{eqnarray}
% \end{lemma}
% \begin{proof}
%     Using the Duhammel formula, we have
%     \begin{eqnarray}
%         e^{-X/2}(\p_s e^X)e^{-X/2}=\int_0^1 du\: e^{(1/2-u)X}(\p_sX)e^{-(1/2-u)X}
%     \end{eqnarray}
%     We obtain the first equality by setting $e^X=\Delta^{-\alpha}$, and using the change of variable $u\to (1/2-u)\alpha$. To obtain the second equality, we use the variation of the logarithm
%     \begin{eqnarray}
%         \p_s\log\Delta=\frac{\pi}{2}\int_{-\infty}^\infty \frac{dt}{\cosh^2(\pi t)}\Delta^{it}(\p_s\Delta)\Delta^{-it}
%     \end{eqnarray}
% \end{proof}
% This means that for $\alpha=1/4$, and the modular Hamiltonian $K=-\frac{1}{2\pi}\log\Delta$, we have the following two expressions for our local approximation to the global Hamiltonian
% \begin{eqnarray}
%     \tilde{H}_s&&=\frac{-1}{8\pi}\int_{-\pi/4}^{\pi/4}du\:e^{i u K}(\p_s K)e^{-i u K}\label{Hs}\\
%     &&=\frac{1}{8}\int_{-1/8}^{1/8}du\:\int_{-\infty}^\infty\frac{dt}{\cosh^2(\pi t)}\Delta^{1/2+it+u}(\Delta^{-1}\p_s\Delta)\Delta^{-(1/2+it+u)}
% \end{eqnarray}
% The modular operator acts on operators in the algebra as a super-operator. 
% \begin{eqnarray}
%     \Delta^\alpha a\ket{\Omega}=\mathcal{D}_\alpha (a)\ket{\Omega}\ .
% \end{eqnarray}
% If there are density matrices, this super operator is
% \begin{eqnarray}
%     \mathcal{D}_\alpha(a)=\rho^\alpha a \rho^{-\alpha}\ .
% \end{eqnarray}
% Therefore, we can write
% \begin{eqnarray}
%     (\p_s\Delta) a\ket{\rho^{1/2}}=\p_s(\rho a \rho^{-1})\ket{\rho^{1/2}}=[(\p_s \rho)\rho^{-1},\rho a \rho^{-1}]\ket{\rho^{1/2}}\ .
% \end{eqnarray}
% Therefore, our operator $\tilde{H}$ acts on the operators in the algebra as the superoperator
% \begin{eqnarray}
%     &&\Delta^{-1}(\p_s\Delta) \mO\ket{\Omega}=[\rho^{-1}(\p_s \rho) ,\mO]\ket{\Omega}\ .
% \end{eqnarray}
% We will use the Euclidean path integrals to compute the variation
% \begin{eqnarray}
%     \rho^{-1}(\p_s \rho)
% \end{eqnarray}
% and plug it in Lemma \ref{} to compute our local approximation to the global Hamiltonian to be
% \begin{eqnarray}\label{Hsdensitymatrix}
%   H_s=\frac{-\pi}{8}\int_{-\pi/8}^{\pi/8}du \int_{-\infty}^\infty\frac{dt}{\cosh^2(\pi t)}\rho^{1/2+it+u}\lb \rho^{-1}\p_s\rho)\rb \rho^{-(1/2+it+u)}\ .
% \end{eqnarray}

% \begin{enumerate}
%     \item {\bf General QFT:} In QFT, the density matrix is given by a path integral 
% \begin{eqnarray}
%     \int_{\varphi_-}^{\varphi_+} \mathcal{D}\varphi e^{-S(\phi)}
% \end{eqnarray}
% with boundary values $\varphi_\pm$ above and below the cut on the sphere. To compute $\p_s \rho$, following \cite{faulkner}, we make a small scaling diffeomorphism with the vector field $\zeta=t \p_t+r \p_r$, and the metric changes by 
% \begin{eqnarray}
%     \eta_{\mu\nu}\to g(s)_{\mu\nu}=\eta_{\mu\nu}+s(\p_\mu \zeta_\nu+\p_\nu \zeta_\mu)+O(s^2)\ .
% \end{eqnarray}
% We find that 
% \begin{eqnarray}\label{variationrhog}
%     \rho_{B,g(s)}-\rho_{B,\eta}&&=\int d^dx (\delta g^{\mu\nu}) \rho_{B,\eta} \lb :T_{\mu\nu}:\rb+O(s^2)\nn\\
%     &&=-2s\int d^dx \zeta^\mu \rho_{B,\eta} \lb :\p^\nu T_{\mu\nu}:\rb+2s \mathcal{P}\lb \int_{\p\mathcal{M}} \zeta^\mu n^\nu \rho_{B,\eta}:T_{\mu\nu}:\rb+O(s^2) \nn\\
%     :T_{\mu\nu}:&&=T_{\mu\nu}-\braket{T_{\mu\nu}}\ .
% \end{eqnarray}
% There are three boundaries in the path integral: those above and below $B$ that we call $B^\pm$, and one that circles the entangling surface in the Euclidean plane that we refer to $\Sigma$.
% The density matrix $\rho_{B,g(s)}=U_s\rho_{B(s),\eta}U^\dagger_s$ where $B(s)$ is rescaled. To the first order in $s$, we have
% \begin{eqnarray}
%     \rho_{B,g}=\rho_{B(s),\eta}+2s \int_B \zeta^\mu[:T_{\mu} :,\rho_{B,\eta}]+O(s^2)\ .
% \end{eqnarray}
% Putting the two pieces together, we find that the term above cancels the contribution of the boundaries $B^\pm$ in (\ref{variationrhog}), and we are left with
% \begin{eqnarray}
%     \p_s\rho_{B(s),\eta}=2 \int_\Sigma  \rho_{B,\eta}(\zeta^\mu n^\nu:T_{\mu\nu}:)\ .
% \end{eqnarray}
% Therefore,
% \begin{eqnarray}
%    \rho^{-1} (\p_s\rho)=2\mathcal{P}\lb \int_\Sigma \zeta^\mu n^\nu :T_{\mu\nu}:\rb
% \end{eqnarray}
% While the expression above explicitly involves the stress tensor, the modular flows of (\ref{Hsdensitymatrix}) are highly nontrivial, and one cannot simplify the expression above any further. To make further progress, we need to resort to perturbation theory or generalized free fields.


% \item {\bf Vicinity of entangling surface}

% Consider the case that we are in a general QFT, but the subregion is almost the whole region. In this limit, we can zoom in, near the entangling surface, and the physics is well-approximated by boost, and the earlier analysis work. Therefore, in this limit, at least for operators localized in the vicinity of the entangling wedge our operator is interest is well-approximated by the Hamiltonian.

% \item {\bf Conformal perturbation theory:}

% Consider a CFT with Hamiltonian $H_0$ in $\mathbb{R}^{1,d-1}$, the vacuum state $\ket{\Omega_0}$, and the modular Hamiltonian of a ball-shaped region $B$ of radius $R$ which generates a geometric flow.
% We deform the Hamiltonian by a relevant operator $\mathcal{O}$:
% \begin{eqnarray}
%     &&H=H_0+\lambda V\nn\\
%     &&V=\int d^{d-1}x \:\mathcal{O}(\tau,x),\ .
% \end{eqnarray}
% The vacuum of the deformed theory, in the interaction picture, can be written using the Dyson formula as
% \begin{eqnarray}
%  \ket{\Omega(\lambda)}&&\sim \mathcal{T}\lb e^{-\lambda\int_{-\infty}^0  d\tau V(\tau)}\rb\ket{\Omega_0}\nn\\
%  &&=\sum_{n\geq 0}\frac{(-1)^n}{n!}\int_{-\infty}^0d\tau_1\cdots d\tau_n \:\mathcal{T}\lb V(\tau_1)\cdots V(\tau_n)\rb\ket{\Omega_0} \nn\\
%  &&V(\tau)=e^{H_0\tau}V e^{-H_0\tau}
% \end{eqnarray}
% where $\ket{\Omega_0}$ is the CFT vaccum. In terms of the density matrix, we find
% \begin{eqnarray}
%     \rho^{-1}\delta_\lambda \rho=-\int d\tau\:\mathcal{T}\lb V(\tau)\rb\ .
% \end{eqnarray}
% Plugging this into the (\ref{Hsdensitymatrix}) we find that our Hamiltonian becomes
% \begin{eqnarray}
%     \tilde{H}_s=H_0+\lambda \int \mO+O(\lambda^2)
% \end{eqnarray}
% where $H_0$ is the CFT Hamiltonian and the first-order correction takes the form of an integral with a single insertion of the relevant conformal primary operator $\mO$ in the Euclidean space and some measure that is explicitly calculated in (\ref{Hsdensitymatrix}). Alternatively, we can write this result expanded in terms of the deformation of the modular Hamiltonian in conformal perturbation theory as in (\ref{}). 

% Perturbation theory allows us to compare the eigenvalues and eigenvectors of our locally reconstructed Hamiltonian to the actual one. Schematically, both Hamiltonians take the following form
% \begin{eqnarray}
%     H=H_0+\lambda \int \mO
% \end{eqnarray}
% for some integral measure. Intuitively, one expects that continuity in $\lambda$ and the fact that $\mO$ is a light primary imply that at high energies our proposed Hamiltonian matches the actual Hamiltonian of the theory. To put this on more solid ground, we use the non-degenerate perturbation theory
% \begin{eqnarray}
%     (H_0-E_0+\lambda (H_1-E_1))\lb \ket{n^{(0)}}+\lambda\ket{n^{(1)}}\rb=O(\lambda^2)\ .
% \end{eqnarray}
% and notice that the first-order change in eigenvalues and eigenvectors is given by 
% \begin{eqnarray}
%     E_1(n)=\lambda \int\bra{n^{(0)}}\mO\ket{n^{(0)}}\nn\\
%     \ket{n^{(1)}}=\sum_{k\neq n}\frac{\bra{k^{(0)}}\mO\ket{n^{(0)}}}{(E_0(n)-E_0(k))}\ .
% \end{eqnarray}
%  In a chaotic theory, it follows from the eigenstate thermalization hypothesis that the matrix elements of light operator $\mO_p$ in these heavy energy eigenbasis, at high energies and for the vast majority of energy eigenstates, are given by 
% \begin{eqnarray}
% &&\bra{E_H}\mO_p\ket{E_{H'}}=\mO_p(E_H)\delta_{HH'}+O^p_{HH'}
% \end{eqnarray}
% where the first term is diagonal and depends only on the energy $E_H$, and the dependence of the microstate $H$ or $H'$ is all captured in the second term $O^p_{HH'}\sim e^{-S(E)}$ with $E=(E_H+E_{H'})/2$. This implies that the difference in the eigenvalues of the global Hamiltonian $H_\lambda$ and our local approximation $\tilde{H}_\lambda$ to the first order $\lambda$ is universal at high energies, and the difference in energy eigenvectors is exponentially suppressed and random to the first order in $\lambda$. In a CFT, the operator-state correspondence implies that there is a one-to-one correspondence between energy eigenvectors and there is a one-to-one correspondence between energy eigenvectors and local operators. In particular, the expectation is that, for a generic chaotic CFT, the vast majority of high-energy eigenstates correspond to heavy primary operators $H$. Then, 
% \begin{eqnarray}
%      E_1(H)\sim C_{HH}^\mO, \qquad \braket{(H')^{(0)}|H^{(1)}}\sim C_{HH'}^\mO\ .
% \end{eqnarray}
% In other words, 
% \begin{enumerate}
%     \item The first order change in energy levels of our Hamiltonian is proportional to a Heavy-Heavy-Light coefficient $C_{\mO HH}$ with the same heavy conformal primary $H$.
%     \item The first order change in energy eigenvectors is proportional to a Heavy-Heavy-Light coefficient $C_{\mO HH'}$ with differing heavy conformal primaries $H$ and $H'$.
% \end{enumerate}
% We are interested in studying the high-energy spectrum of the perturbed Hamiltonian. In the thermodynamic limit, the radius of the Lorentzian cylinder $L\to \infty$, and we look at part of the spectrum with energies that grow to infinity $E_H\sim L^d$ such that the energy density $\epsilon_H$ is kept finite. There is finite wavelength associated to $\epsilon$: 
% \begin{eqnarray}
%     l=\lb \frac{\epsilon}{N}\rb^{-1/(d+1)}\ .
% \end{eqnarray}
% where $N$ is the central charge of the CFT. Then,
% \begin{eqnarray}
%     &&C_{HH'}^\mO=\delta_{HH'}\Delta_H^{\Delta_\mO/(d+1)}f_\mO(E_H)+\mathcal{R}_{HH'}^\mO\nn\\
%     &&f_\mO(E_H)=l_H^{\Delta_\mO}(N \omega_d)^{-\Delta_\mO/(d+1)}\mO(E_H),
% \end{eqnarray}
% where $\omega_d$ is the volume of a unit sphere $S^d$, and the non-universal term suffers an exponential suppression of the form
% \begin{eqnarray}
%     \mathcal{R}^\mO_{HH'}=e^{-O(\Delta_\mO^{d/(d+1)})+\frac{\Delta_\mO}{(d+1)}\log\Delta_H}\ .
% \end{eqnarray}
% In 2d CFTs and typical holographic CFTs, the thermal one-point functions vanish in the thermodynamic limit, which implies the further restriction that the first-order perturbation $E_1(n)=0$.

% \item {\bf  Excited states of CFT:} 

% Instead of considering the vacuum of QFT, we can consider the excited states of CFT. More generally, for an inclusion $A\subset B$ and any pair of vectors $\Phi$ and $\Psi$ in the Hilbert space, we can consider the following characteristic function
% \begin{eqnarray}
%    T_{\Phi|\Psi}(t)=\Delta_{\Phi|\Psi;B}^{it}\Delta_{\Phi|\Psi;A}^{-it}
% \end{eqnarray}
% and define
% \begin{eqnarray}
%     \tilde{H}_s=\frac{-1}{8\pi}\int_{-\pi/4}^{\pi/4}du\:e^{i u K_{\Phi|\Psi}}(\p_s K_{\Phi|\Psi})e^{-i u K_{\Phi|\Psi}}\ .
% \end{eqnarray}
% The simplest cases to consider is when the excited states $U_AU_{B'}\ket{\Omega}$ and $V_AV_{B'}\ket{\Psi}$, then for either $A$ or $B$ we have
% \begin{eqnarray}
%     &&\Delta_{\Phi|\Psi;A/B}=V_{B'}U_A\Delta_{\Omega;A/B}U_A^\dagger V_{B'}^\dagger
% \end{eqnarray}
% which implies
% \begin{eqnarray}
%      &&T_{\Phi|\Psi;A\subset B}(\alpha)=\Delta_{\Phi|\Psi;B}^{\alpha}\Delta_{\Phi|\Psi;A}^{-\alpha}=V_{B'}U_A T_{\Omega;A\subset B} U_A^\dagger V_{B'}^\dagger\ .
% \end{eqnarray}
% This implies that in this case, the spectrum of our approximate Hamiltonian matches that of the actual global Hamiltonian exactly, whereas the eigenvectors differ.
% To make the example more non-trivial, we consider one-parameter families of operators
% \begin{eqnarray}
%     &&\Phi_\lambda=1+\lambda \phi^{(1)}+\frac{\lambda^2}{2}\phi^{(2)}+\cdots\nn\\
%     &&\Psi_\lambda=1+\lambda\psi^{(1)}+\frac{\lambda^2}{2}\psi^{(2)}+\cdots
% \end{eqnarray}
% and from (\ref{}) we know that 
% \begin{eqnarray}
%     &&\delta=\lambda \delta^{(10)}+\mu \delta^{(01)}+\mu\lambda \delta^{(11)}+O(\lambda^2)+O(\mu^2)\nn\\
%     &&K_{\Phi|\Psi}-K_\Omega=\lambda K^{(10)}+\mu K^{(01)}+\lambda\mu K^{(11)}+O(\lambda^2)+O(\mu^2)\nn\\
%     &&K^{(10)}=\frac{-\pi}{2}\int_{-\infty}^\infty \frac{dt}{\cosh^2(\pi t)}(\phi^{(1)}_{-i/2+t}+\phi^{(1)}_{i/2+t}),\nn\\
%     &&K^{(01)}=\pi\int_{-\infty}^\infty \frac{dt}{\cosh^2(\pi t)}J\psi_{t}^{(1)}J,\qquad K^{(11)}=0\ .
% \end{eqnarray}
% Consider the modular Hamiltonian of a ball of size $e^{-2\pi s}$ on the Lorentzian cylinder and denote it by $K(s,\lambda)$. To simplify our notation, we will use
% \begin{eqnarray}
%     &&K'(s)=\p_tK(s+t,\lambda)|_{t=0}\nn\\
%     &&\dot{K}(s)=\p_\lambda K(s,\lambda)|_{\lambda=0}\ .
% \end{eqnarray}
% We prove the following lemma
% \begin{lemma}


% \end{lemma}

% \begin{eqnarray}
%      \p_sA(s)\Big|_{s=0}= e^{-\pi K(0)/2}\lb\int_0^1 du \: e^{2\pi u K(0)}K'(0) e^{2\pi(1-u)K(0)}\rb e^{-\pi K(0)/2}
%  \end{eqnarray}
% As a result,
% \begin{eqnarray}
%     \delta\tilde{H}_t(\lambda)&&=\int_0^1 du\: e^{2\pi(u-1/2)K_{B(t)}} \lb \p_s K_{B(t+s)}\rb_{s=0} e^{-2\pi(u-1/2)K_{B(t)}}\nn\\
%     &&=-\frac{1}{\pi}\int_0^1 du\: e^{2\pi(u-1/2)K_{B(R)}} R\p_R K_{B(R)} e^{-2\pi(u-1/2)K_{B(R)}}
% \end{eqnarray}

% Next, we would like to compute the expression above in the first-order perturbation theory in $\lambda$ around $\lambda=0$:
% \begin{eqnarray}
%    \delta_\lambda (\delta\tilde{H}_t)&&=\p_\lambda\p_s\lb e^{-\pi K(t,\lambda)}e^{2\pi K(t+s,\lambda)}e^{-\pi K(t,\lambda)}\rb_{\lambda=s=0}\nn\\
%    &&=\p_\lambda\p_s A(s,\lambda)\Big|_{s=\lambda=0}\nn\\
%    &&=\int_0^1 du\: \p_\lambda\lb e^{B}K'(0) e^{-B} \rb_{\lambda=0}\nn\\
%    &&=\int_0^1 du\: \lb e^B (\p_\lambda K'(0)) e^{-B}+ (\p_\lambda e^B)K'(0)e^{-B}+e^B K'(0)(\p_\lambda e^{-B})\rb_{\lambda=0}
% \end{eqnarray}
% where we have defined $B=2\pi(u-1/2)K_{B(t)}$. We focus on the second and third terms, using 
% the Duhamel formula
% \begin{eqnarray}\label{firstorderBCH}
%     &&\delta(e^X)=\int_0^1 du\: e^X (\delta X)_{X,-u}=\int_0^1 du\:  (\delta X)_{X,u} e^X \nn\\
%     &&(\delta X)_{X,u}\equiv \lb e^{u X}(\delta X)e^{-u X}\rb
% \end{eqnarray}
% we find the integrand in the second and third terms can be written as
% \begin{eqnarray}
%   \lb(\p_\lambda e^B)K'(0)e^{-B}+e^B K'(0)(\p_\lambda e^{-B})\rb_{\lambda=0}=\int_0^1 dv\: e^{B} [ e^{-v B}\p_\lambda K(0) e^{v B},K'(0)] e^{-B}
% \end{eqnarray}
% As a result, we have
% \begin{eqnarray}
%     \delta_\lambda (\delta\tilde{H}_t)&&=\int_0^1 du\:dv\: e^B\lb  K^{(1,1)}(t,0)+ [e^{-v B}K^{(0,1)}(t,0)e^{vB},K^{(1,0)}(t,0)]\rb e^{-B}\nn\\
%     &&=\int_0^1 du e^B K^{(1,1)}(t,0)e^{-B}+\int du\:d v \: e^{(1-v)B}[K^{(0,1)}(t,0),e^{v B}K^{(1,0)}(t,0) e^{-v B}]\: e^{-(1-v)B}
% \end{eqnarray}
% where we are denoting $\p_\lambda K(s,\lambda)=K^{(0,1)}(s,\lambda)$ and $\p_sK(s,\lambda)=K^{(1,0)}$. Note that 
% \begin{eqnarray}
%     K^{(1,0)}(t,0)&&=\p_s \lb \cosh(2\pi (t+s))K_B+\sinh(2\pi (t+s) H\rb_{s=0}\nn\\
%     &&=2\pi \lb \cosh(2\pi t)H+\sinh(2\pi t)K_B\rb=-R\p_R K_{B(R)}\ .
% \end{eqnarray}
% Therefore, all we need to compute is $\delta_\lambda K_{B(R)}$ because
% \begin{eqnarray}
%     &&K^{(0,1)}(t,0)=\delta_\lambda K_{B(R)}\nn\\
%     &&K^{(1,1)}(t,0)=-\frac{1}{2\pi}R \p_R \lb \delta_\lambda K_{B(R)}\rb\ .
% \end{eqnarray}
% We will see that 
% \begin{eqnarray}
%     \delta_\lambda K_{B(R)}=\lambda R^2\int d\text{vol}_{B(t)} f(\alpha,u)\: \mO(\alpha,u)
% \end{eqnarray}
% where $\mO(\alpha,u)=e^{-\alpha K_{B(R)}}\mO(0,u) e^{\alpha K_{B(R)}}$. Then,
% \begin{eqnarray}
%     K^{(1,1)}(t,0)=-\frac{1}{\pi}K^{(0,1)}(t,0)
% \end{eqnarray}
% \NL{Can we simplify this any further?}



% In conformal perturbation theory, we can expand the Euclidean path-integral that computes the density matrix as 
% \begin{eqnarray}
%     \rho_\lambda=\int_{\varphi_-}^{\varphi_+}\mathcal{D}\varphi\: e^{-S_0(\varphi)}\lb 1+\lambda \int_{\mM} \mO\rb
% \end{eqnarray}
% Then, we have
% \begin{eqnarray}
%     \rho_\lambda=\rho_0+\lambda\mathcal{P}\lb\int_{\mM}  \mO\rb
% \end{eqnarray}





% Consider a CFT with Hamiltonian $H_0$ in $\mathbb{R}^{1,d-1}$, the vacuum state $\ket{\Omega_0}$ and the modular Hamiltonian of a ball-shaped region $B$ of radius $R$ which generates a geometric flow.
% We deform the Hamiltonian by a relevant operator $\mathcal{O}$:
% \begin{eqnarray}
%     &&H=H_0+\lambda V\nn\\
%     &&V=\int d^{d-1}x \:\mathcal{O}(\tau,x),\ .
% \end{eqnarray}
% The vacuum of the deformed theory, in the interaction picture, can be written using the Dyson formula as
% \begin{eqnarray}
%  \ket{\Omega(\lambda)}&&\sim \mathcal{T}\lb e^{-\lambda\int_{-\infty}^0  d\tau V(\tau)}\rb\ket{\Omega_0}\nn\\
%  &&=\sum_{n\geq 0}\frac{(-1)^n}{n!}\int_{-\infty}^0d\tau_1\cdots d\tau_n \:\mathcal{T}\lb V(\tau_1)\cdots V(\tau_n)\rb\ket{\Omega_0} \nn\\
%  &&V(\tau)=e^{H_0\tau}V e^{-H_0\tau}
% \end{eqnarray}
% where $\ket{\Omega_0}$ is the CFT vaccum. 


% To the first order in perturbation theory, we have
% \begin{eqnarray}
%     \ket{\Omega(\lambda)}=\lb 1-\lambda\int_{-\infty}^0d\tau\: V(\tau)\rb\ket{\Omega_0}+O(\lambda^2)
% \end{eqnarray}
% It is convenient to use polar coordinates and define the potential term at constant radius $r$ integrated on $S^{d-2}$:
% \begin{eqnarray}
% &&V=\int_0^\infty d\Omega_{d-2} r^{d-2}dr\:d\Omega_{d-2}\:\mO(r,\Omega_i)\ .
% \end{eqnarray}
% Since we are interested in spherically symmetric setups, the problem becomes effectively two-dimensional in $(\tau,r)$. 

% Consider a ball-shaped region of angular size $\theta$ at constant time slice $\tau=0$ on the Lorentzian cylinder. 

% The following change of coordinates $(\tau,r)\to (\alpha,u)$ that sends Euclidean space $\mathbb{R}^{d}$ to the $\mathbb{H}_{d-1}\times S^1$:
% \begin{eqnarray}
%     &&\tau/R=\frac{\sin(2\pi\alpha)}{\cos(2\pi\alpha)+\cosh(u)}\nn\\
%     &&r/R=\frac{\sinh(u)}{\cos(2\pi\alpha)+\cosh(u)}
% \end{eqnarray}
% with $u\in \mathbb{R}$ and $\alpha\in [0,1]$ representing the $S^1$. It is often convenient to work with rescaled coordinates $\tilde{\tau}=\tau/R$ and $\tilde{r}=r/R$.
% In these new coordinates, the first-order variation of the vacuum vector becomes
% \begin{eqnarray}\label{variationvector}
%     \delta_\lambda\ket{\Omega(\lambda)}\Big|_{\lambda=0}&&=\lb (2\pi R^2\lambda)\int_0^{1/2} d\alpha \int_{-\infty}^\infty  \frac{du\:\sinh^{d-2}(u)}{(\cos(2\pi\alpha)+\cosh(u))^{d}}V(\alpha,u)\rb\ket{\Omega(0)}=\lambda\tilde{V}\ket{\Omega(0)}\nn\\
%     \tilde{V}&&=\int_0^{1/2}\int_{\mathbb{H}_{d-1}}(d\text{vol})\mathcal{O}(\alpha,u,\Omega_i)\nn\\
%     d\text{vol}&&=2\pi R^2 \frac{d\Omega_{d-1}d\tau du\:\sinh^{d-2}(u)}{(\cos(2\pi\alpha)+\cosh(u))^{d}}=r^{d-2}drd\tau d\Omega_{d-2}\ .
% \end{eqnarray}
% It is also convenient to define the conformal factor
% \begin{eqnarray}
%     \omega(\alpha,u)=\cos(2\pi\alpha)+\cosh(u)\ .
% \end{eqnarray}
% Since the modular flow of a ball-shaped region is a shift in $\alpha$, we define the notation
% \begin{eqnarray}
%     \sigma_s(\mO)\equiv\Delta_B^{-is}\mO\Delta_B^{is}
% \end{eqnarray}
% and write 
% \begin{eqnarray}
%     \tilde{V}=\int (d\text{vol})\:\sigma_{i\alpha}(\mO)\ .
% \end{eqnarray}
% Note that the Euclidean modular flowed operator $\sigma_{i\alpha}(\mO)$ need not be inside the algebra.  However, for real $s$ we have $\sigma_s(\mO)\in \mA_B$, and for a dense set of operators $\sigma_z(\mO)$ remains inside the algebra for all complex $z$, we treat $\sigma_{i\alpha}(\mO)$ as if it belongs to the algebra. This way, we can view the vacuum of conformal perturbation theory as an ``excited state" of the CFT created by the action of an operator in the algebra $\tilde{V}$. 

% Now, we have a one-parameter family of vectors in the Hilbert space expanded around a vacuum $\ket{\Omega_0}$ that we schematically denote by
% \begin{eqnarray}
%     \ket{\Omega(\lambda)}=(1+\int\lambda \mO)\ket{\Omega}+O(\lambda^2)
% \end{eqnarray}
% with $\mO\in \mA_B$ a self-adjoint operator in the algebra of a ball-shaped region $\mA_B$. In our example, the integral has the measure (\ref{variationvector}); however, we will treat the integral schematically. 



% Then, following \cite{lashkari2023perturbation}
% \begin{eqnarray}
%    &&D(\lambda)=1-\Delta_\Omega^{-1/2} \Delta_{\Omega(\lambda)}\Delta_\Omega^{-1/2}=\lambda\int\lb \Delta^{1/2}\sigma_{i\alpha}(\mO)\Delta^{-1/2}+\Delta^{-1/2}\sigma_{i\alpha}(\mO)\Delta^{1/2}-2\sigma_{i\alpha}(\mO)\rb+O(\lambda^2)\nn\\
%    &&\frac{D(\lambda)}{1-D(\lambda)/2}=\lambda\int \delta+O(\lambda^2)\nn\\
%    &&\delta=\sigma_{i(\alpha+1/2)}(\mO)+\sigma_{i(\alpha-1/2)}(\mO)-2\sigma_{i\alpha}(\mO)
% \end{eqnarray}

% ----


% Conformal perturbation theory



% \item {\bf Coherent states of Generalized Free Fields:}


% \end{enumerate}




% Consider a one-parameter family of states 
% \begin{eqnarray}
%     \ket{\Omega(\lambda)}=1+\lambda \mO\ket{\Omega}+O(\lambda^2)
% \end{eqnarray}
% where for simplicity we will assume $\mO\in \mA$. 

% Consider an infinitesimal diffeomorphism $f:x^\mu\to x^\mu-\zeta^\mu$ so that the metric changes by
% \begin{eqnarray}
%     g_{\mu\nu}=\eta_{\mu\nu}+\p_\mu \zeta_\nu+\p_\nu  \zeta_\mu+O(\zeta^2)\ .
% \end{eqnarray}
% Under this diffeormorphism the Euclidean path-integral that computes $\Delta_{B(R)}$ changes by
% \begin{eqnarray}
%  \delta\Delta_{B(R)}&&=\int_{\tau<0}(\delta g_{\mu\nu}) \Delta_{B(R)} :T_{\mu\nu}: 
% \end{eqnarray}
% where the integral is over the lower half-cylinder in  Euclidean time; $\tau<0$.
% After integration by parts, we find
% \begin{eqnarray}
%     \delta\Delta_{B(R)}=-\int (\zeta^\nu \p^\mu T_{\mu\nu})
% \end{eqnarray}
% We choose the diffeomorphism to generate a dilatation:
% \begin{eqnarray}
%     \zeta=(r \p_r+\tau \p_\tau)\ .
% \end{eqnarray}

% Next, we define
% \begin{eqnarray}
%     K(s)\equiv K(t,\lambda)
% \end{eqnarray}
% so that 
% \begin{eqnarray}
%      \p_sA(s)\Big|_{s=0}= e^{-\pi K(0)/2}\lb\int_0^1 du \: e^{2\pi u K(0)}K'(0) e^{2\pi(1-u)K(0)}\rb e^{-\pi K(0)/2}
%  \end{eqnarray}
% As a result,
% \begin{eqnarray}
%     \delta\tilde{H}_t(\lambda)&&=\int_0^1 du\: e^{2\pi(u-1/2)K_{B(t)}} \lb \p_s K_{B(t+s)}\rb_{s=0} e^{-2\pi(u-1/2)K_{B(t)}}\nn\\
%     &&=-\frac{1}{\pi}\int_0^1 du\: e^{2\pi(u-1/2)K_{B(R)}} R\p_R K_{B(R)} e^{-2\pi(u-1/2)K_{B(R)}}
% \end{eqnarray}

% Next, we would like to compute the expression above in the first-order perturbation theory in $\lambda$ around $\lambda=0$:
% \begin{eqnarray}
%    \delta_\lambda (\delta\tilde{H}_t)&&=\p_\lambda\p_s\lb e^{-\pi K(t,\lambda)}e^{2\pi K(t+s,\lambda)}e^{-\pi K(t,\lambda)}\rb_{\lambda=s=0}\nn\\
%    &&=\p_\lambda\p_s A(s,\lambda)\Big|_{s=\lambda=0}\nn\\
%    &&=\int_0^1 du\: \p_\lambda\lb e^{B}K'(0) e^{-B} \rb_{\lambda=0}\nn\\
%    &&=\int_0^1 du\: \lb e^B (\p_\lambda K'(0)) e^{-B}+ (\p_\lambda e^B)K'(0)e^{-B}+e^B K'(0)(\p_\lambda e^{-B})\rb_{\lambda=0}
% \end{eqnarray}
% where we have defined $B=2\pi(u-1/2)K_{B(t)}$. We focus on the second and third terms, using 
% the Duhamel formula
% \begin{eqnarray}\label{firstorderBCH}
%     &&\delta(e^X)=\int_0^1 du\: e^X (\delta X)_{X,-u}=\int_0^1 du\:  (\delta X)_{X,u} e^X \nn\\
%     &&(\delta X)_{X,u}\equiv \lb e^{u X}(\delta X)e^{-u X}\rb
% \end{eqnarray}
% we find the integrand in the second and third terms can be written as
% \begin{eqnarray}
%   \lb(\p_\lambda e^B)K'(0)e^{-B}+e^B K'(0)(\p_\lambda e^{-B})\rb_{\lambda=0}=\int_0^1 dv\: e^{B} [ e^{-v B}\p_\lambda K(0) e^{v B},K'(0)] e^{-B}
% \end{eqnarray}
% As a result, we have
% \begin{eqnarray}
%     \delta_\lambda (\delta\tilde{H}_t)&&=\int_0^1 du\:dv\: e^B\lb  K^{(1,1)}(t,0)+ [e^{-v B}K^{(0,1)}(t,0)e^{vB},K^{(1,0)}(t,0)]\rb e^{-B}\nn\\
%     &&=\int_0^1 du e^B K^{(1,1)}(t,0)e^{-B}+\int du\:d v \: e^{(1-v)B}[K^{(0,1)}(t,0),e^{v B}K^{(1,0)}(t,0) e^{-v B}]\: e^{-(1-v)B}
% \end{eqnarray}
% where we are denoting $\p_\lambda K(s,\lambda)=K^{(0,1)}(s,\lambda)$ and $\p_sK(s,\lambda)=K^{(1,0)}$. Note that 
% \begin{eqnarray}
%     K^{(1,0)}(t,0)&&=\p_s \lb \cosh(2\pi (t+s))K_B+\sinh(2\pi (t+s) H\rb_{s=0}\nn\\
%     &&=2\pi \lb \cosh(2\pi t)H+\sinh(2\pi t)K_B\rb=-R\p_R K_{B(R)}\ .
% \end{eqnarray}
% Therefore, all we need to compute is $\delta_\lambda K_{B(R)}$ because
% \begin{eqnarray}
%     &&K^{(0,1)}(t,0)=\delta_\lambda K_{B(R)}\nn\\
%     &&K^{(1,1)}(t,0)=-\frac{1}{2\pi}R \p_R \lb \delta_\lambda K_{B(R)}\rb\ .
% \end{eqnarray}
% We will see that 
% \begin{eqnarray}
%     \delta_\lambda K_{B(R)}=\lambda R^2\int d\text{vol}_{B(t)} f(\alpha,u)\: \mO(\alpha,u)
% \end{eqnarray}
% where $\mO(\alpha,u)=e^{-\alpha K_{B(R)}}\mO(0,u) e^{\alpha K_{B(R)}}$. Then,
% \begin{eqnarray}
%     K^{(1,1)}(t,0)=-\frac{1}{\pi}K^{(0,1)}(t,0)
% \end{eqnarray}
% \NL{Can we simplify this any further?}
% % and 
% % \begin{eqnarray}
% %     e^{vB}K^{(1,0)}(t,\lambda)e^{-v B}=e^{2\pi v (u-1/2)K_{B(t)}}K^{(1,0)}(t,\lambda) e^{-2\pi v(u-1/2)K_{B(t)}}
% % \end{eqnarray}



\section*{Acknowledgements}
We thank Keiichiro Furuya, Mudassir Moosa, Shoy Ouseph, Laimei Nie, 

% Acknowledgments
%\begin{acknowledgments}
%\end{acknowledgments}

\appendix

%\begin{widetext}
% \section{Proof of Proposition \ref{prop:density}}\label{app:proof}
% In this appendix, we prove Proposition \ref{prop:density}. To prove this proposition, recall that we need to answer two questions:
% \begin{enumerate}
%     \item if there exists a common cyclic and separating vector;
%     \item if there exists a dense subset of common cyclic and separating vectors.
% \end{enumerate}
% The existence question is answered by the following observation due to Dixmier and Marechal\cite{dixmiermarechal}:
% \begin{lemma}
%     For a sequence of von Neumann algebras $(A_1,A_2,...)$ acting on the same Hilbert space $\mathcal{H}$, if each von Neumann algebra has a cyclic and separating vector, then there exists a common cyclic and separating vector for all $A_i$'s. 
% \end{lemma}
% This lemma directly shows that $\mathcal{C}_\mathcal{A}\mathcal{A}p\mathcal{C}_\mathcal{B} \neq \emptyset$. The density question explores whether such common cyclic and separating vectors are hard to find. The answer to this question follows from two simple observations:
% \begin{lemma}
%     For any von Neumann algebra $\mathcal{A}$ in its standard form acting on a Hilbert space $\mathcal{H}$, the set of cyclic vectors is dense in $\mathcal{H}$ and the set of separating vectors is dense in $\mathcal{H}$.
% \end{lemma}
% \begin{proof}
%     Since $\mathcal{A}$ is in its standard form, there exists a cyclic vector $\ket{\xi}\in\mathcal{H}$. For any invertible operator $\mathcal{O}\in\mathcal{A}$, the vector $\mathcal{O}\ket{\xi}\in\mathcal{H}$ is also a cyclic vector:
%     \begin{equation}
%         \overline{A\mathcal{O}\ket{\xi}} = \overline{A\mathcal{O}^{-1}\mathcal{O}\ket{\xi}} = \overline{A\ket{\xi}} = \mathcal{H}
%     \end{equation}
    
%     Finally by an observation due to Dixmier and Marechal\cite{dixmiermarechal}, any element in $\mathcal{A}$ is a strong-operator limit of a sequence of invertible elements in $\mathcal{A}$. Therefore the set $\{\mathcal{O}\ket{\xi}: \mathcal{O}\in \text{Inv}(\mathcal{A})\}$ is dense in $\mathcal{H}$ where $\text{Inv}(\mathcal{A})$ is the set of invertible operators in $\mathcal{A}$.

%     Since $\mathcal{A}$ is in its standard form, its commutant $\mathcal{A}'$ acting on $\mathcal{H}$ is also in its standard form. Hence, the set of cyclic vectors of $\mathcal{A}'$ is dense in $\mathcal{H}$. Thus, the set of separating vectors of $\mathcal{A}$ is dense in $\mathcal{H}$.
% \end{proof}
% The next observation is again due to Dixmier and Marachel\cite{dixmiermarechal}. The statement applies to general von Neumann algebras.
% \begin{lemma}
%     For any von Neumann algebra $A$ acting on a Hilbert space $\mathcal{H}$ the set of cyclic vector forms a $G_\delta$ subset of $\mathcal{H}$ and the set of separating vector forms a $G_\delta$ subset of $\mathcal{H}$
% \end{lemma}
% The proof of this lemma is elementary and short. So we include it here. Recall a $G_\delta$ set of a topological space is the intersection of countably infinite open subsets. 
% \begin{proof}
%     An empty set is a $G_\delta$ set by definition.

%     Suppose the set of cyclic vectors of $A$ is non-empty, i.e. there exists $\ket{\xi_0}\in \mathcal{H}$ such that $\overline{A\ket{\xi_0}}= \mathcal{H}$. Then for any operator $\mathcal{O}\in A$, we consider the countably infinite collection of open sets:
%     \begin{equation}
%         U(\mathcal{O},n):=\{\ket{\xi}\in\mathcal{H}:||\mathcal{O}\ket{\xi} - \ket{\xi_0}|| < \frac{1}{n}\}
%     \end{equation}
%     If $\ket{\xi}\in\mathcal{H}$ is a cyclic vector of $A$, then for all $n > 0$ there exists an operator $\mathcal{O}_n$ such that $||\mathcal{O}_n\ket{\xi} - \ket{\xi_0}|| < \frac{1}{n}$. Hence we have:
%     \begin{equation}
%         \ket{\xi}\in \mathcal{A}p_{n > 0}\cup_{\mathcal{O}\in A}U(\mathcal{O}, n)
%     \end{equation}
%     Note that because $\cup_{\mathcal{O}\in A} U(\mathcal{O}, n)$ is open, the intersection $\mathcal{A}p_{n > 0}\cup_{\mathcal{O}\in A}U(\mathcal{O}, n)$ is a $G_\delta$ set. 
    
%     Conversely, any vector in this intersection is a cyclic vector of $A$. This is because the reference vector $\ket{\xi_0}$ is a cyclic vector. Thus for an arbitrary vector $\ket{\xi}\in\mathcal{H}$ and for any $\epsilon  > 0$, there exists an operator $\mathcal{O}\in A$ such that:
%     \begin{equation}
%         ||\mathcal{O}\ket{\xi_0} - \ket{\xi}|| < \frac{\epsilon}{2}
%     \end{equation}
%     For any vector $\ket{\xi'}$ in the aforementioned $G_\delta$ set, there exists another operator $\mathcal{O}'$ such that:
%     \begin{equation}
%         ||\mathcal{O}'\ket{\xi'
%         } - \ket{\xi_0}|| < \frac{\epsilon}{2 ||\mathcal{O}||}
%     \end{equation}
%     Then we have:
%     \begin{equation}
%         ||\mathcal{O}\mathcal{O}'\ket{\xi'} - \ket{\xi}||\leq||\mathcal{O}||\cdot||\mathcal{O}\ket{\xi'} - \ket{\xi_0}|| + ||\mathcal{O}\ket{\xi_0} - \ket{\xi}|| <\epsilon
%     \end{equation}
%     Hence $\overline{\mathcal{A}\ket{\xi'}} = \mathcal{H}$. 

%     Thus it follows that the set of cyclic vectors is exactly the $G_\delta$ set $\mathcal{A}p_{n >0}\cup_{\mathcal{O}\in A}U(\mathcal{O},n)$. Apply the same argument to the commutant $A'$ and since the set of separating vectors of $A$ is the set of cyclic vectors of $A'$, we arrive at the conclusion.
% \end{proof}
% Based on these two observations, the density of common cyclic and separating vectors follows from the Baire category theorem. 
% \begin{corollary}
%     For two type $III$ von Neumann algebras $\mathcal{A}, \mathcal{B}$ acting on a common separable Hilbert space $\mathcal{H}$, then the set of common cyclic and separating vectors is a dense $G_\delta$ subset of $\mathcal{H}$.
% \end{corollary}
% \begin{proof}
%     Following the previous two lemmas, for the algebra $\mathcal{A}$, the set of cyclic vectors is a dense $G_\delta$ subset of $\mathcal{H}$ and the set of separating vectors is also a dense $G_\delta$ subset of $\mathcal{H}$. By Baire category theorem, the intersection of these two sets is also a dense $G_\delta$ subset of $\mathcal{H}$, i.e. $\mathcal{C}_\mathcal{A}$ is a dense $G_\delta$ subset of $\mathcal{H}$. Similarly, $\mathcal{C}_\mathcal{B}$ is a dense $G_\delta$ subset of $\mathcal{H}$.

%     Finally, by another application of Baire category theorem, the intersection $\mathcal{C}_\mathcal{A}\mathcal{A}p\mathcal{C}_\mathcal{B}$ is a dense $G_\delta$ subset of $\mathcal{H}$.
% \end{proof}

%\subsection{$L^2$-Nuclearity Condition}\label{subsection:l2nuclearity}
%We are now ready to present the definition of $L^2$-nuclearity.
%\begin{definition}
%    Given a pair of von Neumann algebras $\mathcal{B}\subset\mathcal{A}$ with a common cyclic and separating vector $\ket{\Omega}\in\mathcal{H}$, $D_{\mathcal{B}\subset \mathcal{A}}(i/4) = \Delta_\mathcal{A}^{1/4}\Delta_\mathcal{B}^{-1/4}$ is a well-defined map between the standard Hilbert spaces:
%    \begin{equation}
%        \Delta_\mathcal{A}^{1/4}\Delta_\mathcal{B}^{-1/4}: L_2(\mathcal{B})\rightarrow L_2(\mathcal{A})
%    \end{equation} 
%    If $\Delta_\mathcal{A}^{1/4}\Delta_\mathcal{B}^{-1/4}$ is a nuclear operator, then we say the triplet $(\mathcal{B}\subset\mathcal{A}, \ket{\Omega})$ satisfies the \textit{$L^2$-nuclearity condition}.
%\end{definition}
%For a map between two separable Hilbert spaces, it is nuclear if and only if it is trace-class. For maps between general Banach spaces, the notion of a nuclear operator generalizes the notion of a trace-class operator \cite{nuclearandsumming}.

%For algebraic quantum field theories, one of the most important consequences of $L^2$-nuclearity is the following observation \cite{buchholz2007nuclearity}:
%\begin{theorem}
%    Given a pair of von Neumann algebras $\mathcal{B}\subset\mathcal{A}$ with a common cyclic and separating vector $\ket{\Omega}$, if it satisfies $L^2$-nuclearity condition, then the inclusion is split. More specifically, we have isomorphism:
%    \begin{equation}
%        \mathcal{B}'\vee\mathcal{A}\cong\mathcal{B}'\otimes\mathcal{A}
%    \end{equation}
%\end{theorem}
%For other important consequences of $L^2$-nuclearity, we refer the readers to \cite{buchholz2007nuclearity}. 

%In the context of algebraic quantum field theory, $L^2$-nuclearity is usually proved using representation theory. For example, in the context of massless Lorentzian conformal field theories of dimension $d + 1$, the global conformal group $SO(d + 1, 2)$ admits discrete series representations (i.e. positive energy representations). In this case, any operator representing an element in the maximal compact sub Lie algebra has discrete spectrum that is bounded from below. This provides ample examples of nuclear operators. The calculations in the main text reflect this fact.

%\section{Modular Scattering matrix}

%In scattering theory, we start with an interacting Hamiltonian that splits as 
%\begin{eqnarray}
%H=H_0+V
%\end{eqnarray}
%where $H_0$ is some reference (free) Hamiltonian with multi-particle states $\ket{\Phi_\alpha}$ and eigen-energies $E_\alpha$:
%\begin{eqnarray}
%H_0\ket{\Phi_\alpha}=E_\alpha \ket{\Phi_\alpha}\ .
%\end{eqnarray}
%The asymptotic states of the interacting theory are defined as two sets of eigenstates of the interacting Hamiltonian $H$ with the same eigenvalues $E_\alpha$ as those of $H_0$:
%\begin{eqnarray}
%H\ket{\Psi_\alpha^\pm}=E_\alpha \ket{\Psi_\alpha^\pm}
%\end{eqnarray}
%with the property that coarse-grained multi-particle states $\ket{\Phi_\alpha^\pm}$ behave like free multi-particle eigenstates $\ket{\Phi_\alpha}$. To make this coarse-graining rigorous, we consider states of the interacting theory that belong to a microcanonical ensemble of small width $\Delta E$ with a smooth wave-packet $g(\alpha)$ and evolve them for long times $t\gg 1/(\Delta E)$. Then, the requirement is that 
%\begin{eqnarray}
%\|\int e^{-i E_\alpha t}g_\alpha (\ket{\Psi^\pm_\alpha}-\ket{\Phi_\alpha})\|
%\end{eqnarray}
%can be made arbitrarily small for large $t$. The intuition is that while the free multi-particle wave-packet $\ket{\Phi_\alpha}$ can never be an eigenket of the interacting theory, as we separate excitations by a large amount, the interaction falls off so fast that they can be approximated arbitrarily well by suitably coarse-grained eigenkets $\ket{\Psi_\alpha^\pm}$. The asymptotic states can be defined as 
%\begin{eqnarray}
%&&\ket{\Psi_\alpha^\pm}=\Omega(\mp \infty)\ket{\Phi_\alpha}\nn\\
%&&\Omega(t)=e^{i Ht}e^{-i H_0t}\ .
%\end{eqnarray}
%Of course, this is only meaningful after coarse-graining, and not as an exact equality. The S-matrix is defined to be the overlap of the in and out asymptotic states
%\begin{eqnarray}
%\braket{\Phi_\beta|S \Phi_\alpha}&&=\braket{\Psi_\beta^+|\Psi_\alpha^-}=\lim_{t\to \infty}\braket{\Phi_\beta|\Omega(t)^\dagger \Omega(-t)|\Phi_\alpha}\ .
%\end{eqnarray}
%In other words, we have
%\begin{eqnarray}
%S=\lim_{t\to \infty} e^{i H_0t} e^{-2i Ht}e^{i H_0t}\ .
%\end{eqnarray} 

%In direct analogy, we define a modular S-matrix for any pair of modular Hamiltonians $K$ and $K_0$ as 
%\begin{eqnarray}
%S=\lim_{t\to \infty} e^{i K_0t} e^{-2i Kt} e^{i K_0t}\ .
%\end{eqnarray}
%There are two obvious choices: 
%\begin{enumerate}
%\item {\bf Two algebras and one state:} For a general pair of algebras $\mA$ and $\mB$ and a single state $\ket{\Omega}$, we define the modular S-matrix as 
%\begin{eqnarray}
%S_{A,B}(z)&&=e^{2\pi i z K_B }e^{-4\pi i z K_A}e^{2\pi i z K_B}\nn\\
%&&=\Delta_B^{-iz}\Delta_A^{2iz}\Delta_B^{-iz}\ .
%\end{eqnarray}
%In the special case, $\mA\subset \mB$ the operator $T(z)$ and admits an analytic in the complex $z$-domain $\Im(z)\in (0,1/2)$: 
%\begin{eqnarray}
%    S_{A\subset B}(z)=T(z)T(-z)^{-1}\ .
%\end{eqnarray}
% In particular, for imaginary values of $z=i\alpha$, the modular S-matrix is the norm of the characteristic function
% \begin{eqnarray}
%     S(i\alpha)=|\Delta^{\alpha}\Delta^{-\alpha}|\ .
% \end{eqnarray}
%In the remainder of this work, we focus on the case of inclusion, because this analyticity allows us to prove a modular chaos bound \NL{cite Maldacena et al, Simon, Faulkner}. In the special case of a common cyclic and separating vector for the inclusion of algebras, $T(z)$ is a characteristic function which involves more conditions than just analyticity such as unitarity on the other boundary (see appendix \ref{}).

%\item {\bf One algebra and two states:} For a pair of states $\psi$ and $\phi$ of the algebra $\mA$ we define
%\begin{eqnarray}
%S_{\psi,\phi}(z)=e^{i z K_\phi} e^{-2i z K_\psi}e^{i z K_\phi}=u_{\psi|\phi}(z)u_{\psi|\phi}(-z)^{-1}
%\end{eqnarray}
%where $u_{\psi|\phi}$ is a cycle which is analytic in the complex $z$-domain $\Im(z)\in (0,1/2)$. In the special case, $\phi$ majorizes $\psi$ we obtain further properties on the boundaries.

%\NL{cite de Boer et al.} \NL{Is there a chaos bound here?}
%\end{enumerate}

%When the modular S-matrix admits an analytic extension to the complex $z$ strip, for imaginary values $z=i\alpha$, the modular S-matrix is the norm of the characteristic function or the cocycle:
%\begin{eqnarray}
%&&S_{A\subset B}(z)=|T(z)|^2\nn\\
%&&S_{\psi,\phi}(z)=|u_{\psi|\phi}(i\alpha)|^2\ .
%\end{eqnarray}

%The modular S-matrix expressions can be viewed as an operator algebraic regulator for modular Hamiltonian. In the previous section, we regulated by the operator using the intuition of Euclidean path-integrals and the geometric action of $P$ and $H$ we introduced the regulator 
%\begin{eqnarray}
%&&\Delta^{1/2+it}\to e^{i \ep Ht} \Delta^{1/2+it}e^{i \ep Ht}\nn\\
%&&\Delta^{1/2+it}\to e^{i \ep Pt}\Delta^{1/2+it} e^{-i \ep Pt}\ .
%\end{eqnarray}
%Our modular $S-matrix$ can be interpreted as a regulator procedure in the absence of operators $P$ and $H$. However, in the absence of Euclidean path-integral picture, it is not apriori clear that the modular S-matrix has discrete spectrum. The assumption of nuclearity is the operator algebraic avatar of our intuition of locality coming from the Euclidean path-integrals. In the case of inclusion of algebras $\mA\subset \mB$, $L^2$-nuclearity posits that for purely the analytic continuation of the modular S-matrix to purely imaginary $z=i\alpha$ is trace-class for all $\alpha\in (0,1/2)$:
%\begin{eqnarray}
%&&\forall \alpha\in (0,1/2):\quad \tr(Z(i\alpha)) <\infty\ .
%\end{eqnarray}
%In analogy with scattering theory, we define the scattering amplitude operator $M$ using
%\begin{eqnarray}
%S=1+i M\ .
%\end{eqnarray}
%Consider a sequence of interpolating algebras $\mA\subset \mA_{n-1}\subset \mA_n\subset \mB$ such that $\lim_{n\to -\infty}\mA_n=\mA$ and $\lim_{n\to \infty}\mA_n=\mB$. Then,  we have
%\begin{eqnarray}
%&&\lim_{n\to \infty} K_{\mA_n}=K_B\nn\\
%&&\lim_{n\to -\infty} K_{\mA_n}=K_A\ .
%\end{eqnarray}
%For the characteristic function we find
%\begin{eqnarray}
%&&T_{A\subset B}=T_{A_n\subset B}T_{A\subset A_n}\nn\\
%&&\lim_{n\to \infty}T_{A_n\subset B}=1\nn\\
%&&\lim_{n\to -\infty}T_{A\subset A_n}=1\ .
%\end{eqnarray}
%Therefore, we can intuit that the relative size of the inclusion $A\subset B$ controls the size of the regulator, and we should keep it small.   
%If the inclusion of algebras has finite index, one can use it as a measure of relative size of an inclusion. In our application to QFT, we will be interested in inclusion of type III$_1$ algebras with infinite index that split. In these cases, we will use the space volume of regions to define a measure of relative size. For instance, in the inclusion of Rindler wedges, we need to take the proper distance between entangling surfaces to be small, whereas for the inclusion of balls, we take the conformal cross ratio as a measure of relative size.

%We often refer to the relative size of inclusion as the cut-off parameter $\ep$. We first expand the S-martix in small $\ep$ to obtain non-trivial scattering:
%\begin{eqnarray}
%S(t)=1+\ep i \tilde{M}(t)+O(\ep^2)
%\end{eqnarray}
%Then, take the large $t$ limit. Analyticity of $S(t)$ results in a modular chaos bound 
%\begin{eqnarray}
%\lim_{t\to \infty}\frac{\log |M(t)|}{2\pi}\leq 1\ .
%\end{eqnarray}
%Therefore, we will have to work in the approximation $\ep\ll 1$ and $1\gg 1$ but $\ep e^{2\pi t}\ll 1$. In the standard chaos bound of \cite{}, the gravitational constant $G_N$ or $\hbar$ plays the role of $\ep$, and the interpretation is that we are quantizing a classical theory corresponding to $G_N=0$ and $\hbar=0$. In contrast, our definition of chaos is genuinely quantum, and there is no clear classical limit.\NL{make a milder statement}


\section{Subsystem flows in classical ergodic theory}\label{app:classical}

%\subsection{Coarse-gainings and Classical net of subregion algebras}

We begin by reviewing the measure-theoretic definition of a classical dynamical system in subsection \ref{subsec:classical-definitions}, the subspaces and subsystems of observables in subsection \ref{subsec:classical-region-subsystem}, and their role in classical ergodic theory in subsection \ref{subsec:classical-ergodic-hierarchy}. 

In subsection \ref{subsec:classical-state-flow}, we argue that when the configuration space $X$ is equipped with a Poisson bracket (as suggested by the phase space intuition from physics) {\it the Poisson flow of a state} is a classical avatar of modular dynamics associated to every state. Similar to quantum modular flow, the Poisson flow of a state generates an automorphism of the algebra leaving the state invariant, and satisfies a classical avatar of the KMS condition \footnote{Recall that in the quantum case, these two properties uniquely fix the modular flow of a state~\cite{Takesaki1970}.}. However, crucially for us, as opposed to the quantum modular Hamiltonian, the Poisson Hamiltonian of a state does not satisfy a monotonicity condition under coarse-graining as in the quantum case \footnote{Recall that in the quantum case, this monotonicity originates from the fact that the Tomita operator of the algebra is an extension of its restriction to the subalgebra. In the classical case, there is no analog of the Tomita operator.}. This monotonicity is at the heart of non-commutative results on the reconstruction of geometry from a net of subsystems, such as half-sided modular inclusions and intersections as we review in Section~\ref{subsec:geometry-modular-ergodic}. 

For each coarse-graining we define a notion of a coarse-graining flow. In the case of subspaces, the {\it subspace flow} has the interpretation of a generating function for entropy fluctuations. However, the focus of this work is {\it subsystem flows} which are richer.
Classical many-body systems and field theory come equipped with a net of subsystems, and the quantum generalization of subsystem flows plays a crucial role in the emergence of geometry from a net of abstract subsystems. 

\subsection{Review of classical dynamical systems}\label{subsec:classical-definitions}
% We will see that it is the failure of conditional expectation to intertwine the microscopic phase action with its coarse-grained counterpart that is of interest to us.  The same relative density that measures this failure generates a cocycle whose expectation is exactly the Kullback--Leibler divergence, and this divergence is exactly the entropy increase produced by coarse-graining.  These results provide a rigorous kinematic account of information loss, but they do not by themselves derive an irreversible physical time evolution.

\paragraph{Algebra of observables:} In classical mechanics, one often begins with a phase space that we take to be an orientable smooth manifold \(X\) that admits a strictly positive volume measure. Points \(x\in X\) represent microstates of the system. But from the algebraic viewpoint, the primary object is not the set of points but the commutative von Neumann algebra of observables (complex measurable functions)
\[
\mA=L^\infty(X,\Sigma,\mu),
\]
where \((X,\Sigma,\mu)\) is a standard probability space. When $X$ is a smooth phase space of finite volume, we take $\mu$ to be the normalized volume measure\footnote{$\Sigma$ is the completion of the Borel $\sigma$-algebra with respect to $\mu$. Elements of $L^\infty(X,\Sigma,\mu)$ are equivalence classes of essentially bounded measurable functions that agree $\mu$-almost everywhere.}.


\paragraph{States:} The normalized measure $\mu$ above is interpreted as a distinguished reference state\footnote{A $\sigma$-finite volume measure of infinite total volume instead defines a distinguished weight. We return to this extension in later sections.}. General normal states correspond to probability measures $\nu$ on $X$ that are absolutely continuous with respect to $\mu$ (denoted $\nu\ll \mu$), meaning that for every measurable set $E\subseteq X$, if $\mu(E)=0$ then $\nu(E)=0$.
There exists a $\mu$-measurable function $\rho_\nu:X\to [0,\infty)$ such that $\nu(E)=\int_E \rho_\nu d\mu$ \footnote{In mathematical terminology, $\rho_\nu$ is called a Radon-Nikodym derivative.}. In other words, every normal state can be viewed as a probability density function $\rho_\nu$ on $X$, with expectation values of an observable $f$ in this state captured by integrals on $X$: 
\begin{align}
\braket{f}_\nu=\int_{x\in X}f(x) d\nu(x)=\int_{x\in X} \rho_\nu(x) f(x) d\mu(x)\ .
\end{align}

\paragraph{GNS Hilbert space:} The GNS Hilbert space of the reference state $\mu$ is
\[
  \cH=L^2(X,\Sigma,\mu),
  \qquad
  \mA=L^\infty(X,\Sigma,\mu).
\]
Observables act by multiplication, and the reference state $\mu$ is represented by the constant unit vector $\ket{1}$. A general normal state with density $\rho$ is represented by
\[
  \ket{\sqrt{\rho}}\in\cH,
  \qquad
  \|\ket{\sqrt{\rho}}\|_2^2=\int_X\rho\,\dd\mu=1.
\]
The GNS construction builds a connection between observables and states.


\paragraph{Dynamics:} Dynamics will be described by a measurable one-parameter group of invertible, measure-preserving transformations $\Phi_t:X\to X$. General dynamics need not be smooth. In physics, we are often interested in smooth flows described by a one-parameter family of volume-preserving diffeomorphisms of $X$. Taking a $t$-derivative of how a point $x\in X$ moves under this flow defines a vector field $V$, which preserves the volume measure if and only if $\mathcal{L}_V\mu=0$; i.e., it is divergence-free with respect to $\mu$. 


% \paragraph{Coarse-graining:} Let $\Sigma_0\subseteq\Sigma$ be the sub-$\sigma$-algebra of macroscopically resolved events. We define a coarse-graining as 
% \[
%   Pf:=\E_\mu[f\mid\Sigma_0].
% \]
% On the Hilbert space $\mH$, $P$ is represented by the orthogonal projection onto $L^2(X,\Sigma_0,\mu)$ \footnote{Note that $P$ acts on functions and not measurable sets.}. Coarse-graining by construction preserves the volume measure $\mu$; however, general states are coarse-grained as $\rho_0=P \rho$. A coarse-graining can be viewed as a conditional expectation $E:L^\infty(X,\Sigma,\mu)\to L^\infty(X,\Sigma_0,\mu\mid \Sigma_0)$ that projects to the algebra of coarse-grained observables.
% It is a defining property of this projection that the fine-grained and coarse-grained states are observationally indistinguishable on the macroscopic algebra, even though they may differ on the full $\sigma$-algebra.


% \paragraph{Net of subregion (subspace) coarse-grainings:} In our measure-theoretic setting above, there is an important class (net) of coarse-grainings that are provided geometrically by restricting to subregions $O\subset X$. For a measurable subregion $O \subset X$, we wish to associate a true unital von Neumann subalgebra $\mathcal{A}_O \subset \mathcal{A}=L^\infty(X,\Sigma,\mu)$ that captures the observables strictly localized to $O$. 

% To ensure that $\mathcal{A}_O$ shares the global identity element $\mathbf{1}_X$, this subalgebra is the space of essentially bounded functions that are arbitrarily measurable within $O$ but strictly constant almost everywhere outside of $O$:
% \[
% \mathcal{A}_O = \{ f \in L^\infty(X, \Sigma, \mu) \mid f = \chi_O f + c \chi_{O^c} \text{ for some } c \in \mathbb{C} \}
% \]
% where $\chi_O$ and $\chi_{O^c}$ are the indicator functions of the region and its complement, respectively. Similarly, there is a subalgebra associated to the complement $\mA_{O^c}$ \footnote{Note that in this construction, the regions $O$ have the same dimension as the phase space $X$.}

% The projections from the global algebra $\mathcal{A}$ into the local subalgebra $\mathcal{A}_O$ are normal conditional expectation $E: \mathcal{A} \to \mathcal{A}_O$ \footnote{A conditional expectation is a unit-preserving, positive, ultraweakly continuous linear map satisfying the module property $E(f_1 g f_2) = f_1 E(g) f_2$ for all $f_1, f_2 \in \mathcal{A}_O$ and $g \in \mathcal{A}$.}. Given a state $\omega_{O^c}$ on $\mA_{O^c}$, these projections (normal conditional expectations) are defined as
% \[
% E_\omega(f)(x) = \chi_O(x) f(x) + \chi_{O^c}(1) \left( \frac{1}{\omega_{O^c}(1)} \int_{O^c} f(y) \rho_\omega(y) \, d\mu(y) \right)\ .
% \]
% In fact, there is a one-to-one correspondence between the states of $\mA_{O^c}$ and the normal-condition expectations for $\mA_O$. It is straightforward to see that this conditional expectation preserves any state for which the probability density $\rho$ restricted to $O^c$ is proportional to $\rho_\omega$. However, if the volume measure of $O^c$ is infinite, there does not exist any volume-preserving conditional expectation.


% \paragraph{Subspace algebras versus subsystem algebras:} Note that, in this construction, we have a direct sum $\mA_O\oplus \mA_{O^c}$ as opposed to a subsystem where we expect a tensor product structure. For a measurable partition \(X=O\sqcup O^c\),
% \[
% L^\infty(X)
% \simeq L^\infty(O)\oplus L^\infty(O^c).
% \]These summands represent mutually exclusive alternatives: the microstate belongs to one region or the other. Their units are \(\chi_O\) and \(\chi_{O^c}\), and their products vanish.
% However, your unital region algebras are
% \[
% \mathcal A(O)=L^\infty(O)\oplus\mathbb C\chi_{O^c},
% \qquad
% \mathcal A(O^c)=\mathbb C\chi_O\oplus L^\infty(O^c).
% \]They are not themselves the two direct summands. They overlap:
% \[
% \mathcal A(O)\cap\mathcal A(O^c)
% =\operatorname{span}\{\chi_O,\chi_{O^c}\}.
% \]The direct summands are the supported algebras \(\chi_OL^\infty(X)\) and \(\chi_{O^c}L^\infty(X)\).
% For a global product \(X=Y\times Z\), equipped with a product measure,
% \[
% L^\infty(X)
% \simeq L^\infty(Y)\,\bar\otimes\,L^\infty(Z).
% \]The base and hidden-variable algebras are
% \[
% \mathcal N=L^\infty(Y)\otimes1,
% \qquad
% \mathcal R=1\otimes L^\infty(Z).
% \]They represent simultaneously existing degrees of freedom: a microstate specifies both \(y\) and \(z\). Their intersection is \(\mathbb C1\), and together they generate the full algebra.
% For the Hamiltonian interpretation, require a product Poisson structure as well:
% \[
% \{f(y),g(z)\}=0.
% \]
% This factorization does not require the state to factorize; \(\rho(y,z)\) may contain correlations.
% The second qualification is that a general Poisson fibration \(X\to Y\) need not admit such a global product decomposition or a globally defined complementary fiber algebra.
% A precise formulation for the paper would be:
% A partition of phase space gives a direct-sum decomposition into mutually exclusive sectors. A product decomposition of phase space gives a tensor-product decomposition into simultaneous subsystems. General Poisson coarse-grainings extend the subsystem picture without necessarily admitting a global tensor-product factorization.


\paragraph{Coarse-grainings:} Classical coarse-graining retains a subalgebra of observables and discards the information needed to distinguish states that agree on this subalgebra. Two useful constructions arise from a subregion of the space of microstates and from a subsystem of observed variables. Both can be formulated in measure theory, before introducing a Hamiltonian or a Poisson bracket. We first describe these constructions and then ask when the retained observables evolve autonomously.

In the measure-theoretic setting above, a coarse description is specified by a sub-$\sigma$-algebra $\Sigma_0\subseteq\Sigma$, with associated unital von Neumann subalgebra
\begin{align}
    \mB=L^\infty(X,\Sigma_0,\mu)\subseteq\mA\ .
\end{align}
The reference measure $\mu$ determines the coarse-graining map
\begin{align}
    E_{\mB}(f)=\mathbb E_\mu[f\mid\Sigma_0]\ .
\end{align}
This is the unique $\mu$-preserving normal conditional expectation onto $\mB$. On $L^2(X,\mu)$ it acts as the orthogonal projection onto $L^2(X,\Sigma_0,\mu)$. A general normal state $\omega_\rho(f)=\int_X\rho f\,d\mu$ is coarse-grained to $\omega_0=\omega_\rho\circ E_{\mB}$, whose density is $\rho_0=E_{\mB}(\rho)$, using the $L^1$ extension of conditional expectation. The two states agree on $\mB$. Thus the subalgebra specifies the retained information, while the chosen expectation specifies how a state on this subalgebra is extended back to the full algebra.


\paragraph{Entropy of a coarse-graining:} Consider a $\mu$-preserving coarse-graining $E$, a normal state with density $\rho$, and its coarse-grained density $\rho_0=E(\rho)$. Then, the relative entropy (KL-divergence) of these two states associates an entropy to this coarse-graining:
\begin{align}
    \label{eq:KL-classical}
    S(\rho\|\rho_0)=\int_X \rho \log(\rho/\rho_0)d\mu=H(\rho_0)-H(\rho)
\end{align}
where $H(\rho)=-\int_X\rho\log\rho\,d\mu$ is the entropy relative to the chosen reference measure, and the equality assumes that the two entropies are finite. The nonnegativity of relative entropy implies that coarse-graining cannot decrease this entropy. This compares a state with its coarse-grained version; it does not by itself imply monotonic entropy growth under time evolution.

\paragraph{Generating function for entropy:} For faithful densities $\rho,\rho_0>0$ almost everywhere, the multiplicative relative phase 
\[
u_t
=
\left(\frac{\rho}{\rho_0}\right)^{it}
\]
is the generating function for coarse-graining entropy because 
\begin{align}
&-i\left.\frac{\dd}{\dd t}u_t\right|_{t=0}
=
\log\frac{\rho}{\rho_0},\nn\\
&S(\rho\|\rho_0)=-i \int_X \rho (\partial_t u_t\big|_{t=0})d\mu\ .
\end{align}
Here the differentiation under the integral is justified when $\rho\log(\rho/\rho_0)$ is integrable. More generally, for every faithful density $\rho$ we can define the strongly continuous unitary multiplication group
\[
  U_t^\rho:=\Mop_{\rho^{it}}
  =\exp\!\bigl(it\Mop_{\log\rho}\bigr),
  \qquad
  (U_t^\rho\psi)(x)=e^{it\log\rho(x)}\psi(x).
\]
Similarly, the coarse-grained density defines
\[
  U_t^{\rho_0}:=\Mop_{\rho_0^{it}}.
\]
The parameter $t$ is an auxiliary phase parameter. Nothing in this construction identifies it with laboratory time. In particular, $U_t^\rho$ does not transport points of $X$; it only rotates the phase of a Hilbert space vector pointwise\footnote{For faithful normal states on the same abelian algebra, $u_t=(\rho/\rho_0)^{it}$ is the Connes Radon-Nikodym cocycle. Relative modular operators remain nontrivial, even though the modular automorphism group of each state is trivial.}.

\subsection{Classical subspace (single-particle) and subsystem (multi-particle) coarse-grainings}
\label{subsec:classical-region-subsystem}

Subspace and subsystem coarse-grainings organize incomplete information differently. An example of a subspace coarse-graining resolves microstates within a selected region of the phase space and identifies those outside it. A microstate is either inside the coarse-graining region or outside it. The set of microstates adds up in this case; that is why, in anticipation of the generalization to Hilbert space, we refer to them as {\it subspaces} so that addition is a direct sum. As we will see, there is a net of subspace algebras on the phase space of a single particle. As an example of coarse-graining to a subspace, consider the reduced state to the part of the phase space of a pendulum where the angle $\theta$ is in some limited range. 

Next, consider a microstate described using several (potentially independent) degrees of freedom. In this case, the set of microstates multiplies, and in anticipation of generalization to Hilbert spaces, we refer to this case as {\it subsystems}, so that multiplication is tensor product. As a motivating example, consider two particles with the phase space $(q_1,p_1,q_2,p_2)$, and the restriction to the first particle as a coarse-graining on the second particle degree of freedom. 

Note that since subspaces could have the same dimension as the entire phase space, they can be discussed even in a single-particle setup. A subsystem is instead described by a coarse space $Y$ and a map $\pi:X\to Y$; $Y$ need not be an embedded submanifold of $X$. This construction naturally describes many-body systems and field theory. 


% Evolution transports the resolved region in the first case and the fibers of indistinguishable microstates in the second. Their invariance has correspondingly different meanings: an invariant region defines a conserved event, while an invariant subsystem defines an autonomous factor dynamics. Ergodicity excludes nontrivial conserved events but permits nontrivial autonomous factors.

\paragraph{Subspace coarse-grainings:}
For a measurable region $O\subset X$, define the algebra generated by observables that are constant outside of $O$:
\begin{align}
    \mA_O=\{f\in\mA:\ f=f \cdot 1_O +c 1_{O^c}\text{ for some }c\in\mathbb C\}
\end{align}
where $1_O$ is the indicator function on $O$, and $(f\cdot 1_O)(x)=f(x)1_O(x)$ is point by point multiplication.
This algebra resolves the microstate completely within $O$ and retains only the event ``outside $O$'' on its complement. Assuming $0<\mu(O)<1$, its conditional expectation is
\begin{align}
    E_O(f)=1_O f+
    1_{O^c}\mathbb{E}_{O^c}(f)
\end{align}
where 
\begin{align}
    \mathbb{E}_{O^c}(f)=\frac{1}{\mu(1_{O^c})}\int_{O^c} f d\mu
\end{align}
is the average of $f$ on $O^c$.
For $O_1\subseteq O_2$, we have $\mA_{O_1}\subseteq\mA_{O_2}$, so this construction gives a net of coarse observable algebras. Since a microstate $x\in X$ is either in $O$ or $O^c$ we have 
\begin{align}
    \mA_X=1_O\mA\oplus 1_{O^c}\mA
    \simeq L^\infty(O,\mu|_O)\oplus L^\infty(O^c,\mu|_{O^c})\ .
\end{align}
Our unital subspace algebras have the form $\mA_O=1_O\mA\oplus\mathbb C1_{O^c}$, so $\mA_O$ and $\mA_{O^c}$ are not themselves the two direct summands. In fact,
\begin{align}
    \mA_O\cap\mA_{O^c}=\operatorname{span}\{1_O,1_{O^c}\}\ .
\end{align}

\paragraph{Subsystem coarse-grainings:}
Suppose instead that the probability space has a product decomposition
\begin{align}
    (X,\Sigma,\mu)=(Y\times Z,\Sigma_Y\otimes\Sigma_Z,\mu_Y\otimes\mu_Z)\ .
\end{align}
A microstate now specifies both an observed variable $y$ and a hidden variable $z$. The observable algebra factorizes as
\begin{align}
    \mA\simeq\mA_Y\,\bar\otimes\,\mA_Z,
    \qquad \mA_Y=L^\infty(Y,\mu_Y),\quad
    \mA_Z=L^\infty(Z,\mu_Z)\ .
\end{align}
Retaining the subsystem $Y$ means retaining $\mB_Y=\mA_Y\otimes 1$. The corresponding conditional expectation averages over the hidden variables:
\begin{align}
    (\mathcal{E}_Y f)(y,z)=\int_Z f(y,z')\,d\mu_Z(z')\ .
\end{align}
This operation forgets the hidden variables at every microstate. The tensor factors describe simultaneously specified degrees of freedom. Although the reference measure is a product, a general state with density $\rho(y,z)$ can have correlations between them. 

More generally, a measurable map $\pi:X\to Y$ between standard probability spaces specifies coarse observables through
\begin{align}
    &\mB_Y=\mA_Y\circ \pi\subset \mA,
    \qquad \mu_Y=\mu\circ \pi^{-1}\nn\\
    &\mA_Y=L^\infty(Y,\mu_Y)
\end{align}

Here $\mB_\pi=\pi^*\mA_Y$ is the pullback of $\mA_Y$, and $\mu_Y=\pi_*\mu$ is the
push-forward measure. Under the standard Borel assumptions, disintegration gives
conditional probability measures $\mu_y$ concentrated on the fibers $\pi^{-1}(y)$.

The conditional expectation is
\begin{align}
    (\mathcal{E}_\pi f)(x)=\int_{\pi^{-1}(\pi(x))}f(x')\,d\mu_{\pi(x)}(x')\ .
\end{align}
The fiber measures may depend on $y$, and no global product decomposition is required. 

Note that both subspace and subsystem conditional expectations are special cases of a common measure-theoretic notion of coarse-graining.

\paragraph{Dynamically closed and invariant coarse-grainings:}
Let $\Phi_t:X\to X$ be a measurable one-parameter group of invertible, measure-preserving transformations. It evolves observables according to
\begin{align}
    \alpha_t(f)=f\circ\Phi_{-t}\ .
\end{align}
For any coarse algebra $\mB\subseteq\mA$, we can follow the family $\mB_t=\alpha_t(\mB)$. We call the coarse description dynamically closed when
\begin{align}\label{eq:classical-coarse-invariance}
    \alpha_t(\mB)=\mB,\qquad t\in\mathbb R\ .
\end{align}
For nonatomic $\mu$ and $0<\mu(O)<1$, invariance of the subspace algebra $\mA_O$ violates ergodicity, as we show below.
For the subsystem $Y$ in a product space $Y\times Z$, dynamical closure means that the $Y$-component of $\Phi_t(y,z)$ depends only on $y$ \footnote{Note that this does not follow from preservation of the product measure.}. For a coarse-graining specified by $\pi$, this means that there is a measure-preserving evolution $\Psi_t$ on the coarse space such that
\begin{align}\label{eq:classical-factor-intertwining}
    \pi\circ\Phi_t=\Psi_t\circ\pi\ ,
\end{align}
understood modulo null sets. Evolving the microstate and then retaining its coarse variables gives the same result as evolving the coarse variables themselves. In ergodic theory, such a coarse description is called a {\it factor system} in ergodic theory, and we say the coarse-grained variables evolve autonomously. Intuitively, this means that the future of the macroscopic (or observed) variables can be determined from the present macroscopic state, without the need to know the hidden microscopic details. In other words, there is an effective theory that describes the coarse-grained variables.


% For a subregion, evolution transports its associated algebra:
% \begin{align}
%     \alpha_t(\mA_O)=\mA_{\Phi_t(O)}\ .
% \end{align}
% An invariant region therefore defines an invariant coarse algebra. For nonatomic $\mu$ and $0<\mu(O)<1$, the converse holds as well: the single atom representing the unresolved complement identifies $O^c$ from $\mA_O$. Ergodicity excludes nontrivial invariant regions, but it does not exclude nontrivial factor systems. Invariance of a subalgebra permits its observables to evolve into one another; they need not be individually fixed.

%  For example, on the two-torus with uniform product measure, the flow
% \begin{align}
%     \Phi_t(y,z)=\bigl(y+t\sin(2\pi z)\!\!\pmod 1,\ z\bigr)
% \end{align}
% preserves the measure, but the future observed coordinate depends on the hidden coordinate. Hence $\mB_Y$ is not invariant. An autonomous observed subsystem is also weaker than a product dynamics: the hidden variables may evolve in a way that depends on the observed variables.

\paragraph{Effective theory (Compatibility of coarse-graining and evolution):}
For the measure-preserving dynamics above, invariance of $\mB$ is equivalent to
\begin{align}\label{eq:classical-expectation-equivariance}
    \mE_{\pi}\circ\alpha_t=\alpha_t\circ \mE_{\pi},
    \qquad t\in\mathbb R\ .
\end{align}
Indeed, when $\mB_Y$ is invariant, $\alpha_t\circ \mE_{\pi}\circ\alpha_{-t}$ is another $\mu$-preserving conditional expectation onto $\mB_Y$, so uniqueness gives the equality. Conversely, this equality ensures that $\alpha_t$ maps the range of $\mE_{\pi}$ onto itself. The physical interpretation is macroscopic predictability: the future observed variables can be determined from their present values. This compatibility supplies an effective dynamics on the coarse algebra.


\paragraph{Correlation functions of observables and subalgebras:}
For complex functions $f,g\in\mA$, define the correlation function and its connected part as
\begin{align}
    &C_{f,g}(t)=\mu\bigl(f^*\alpha_t(g)\bigr)\nn\\
    &C^{con}_{f,g}(t)
    :=\mu\bigl(f^*\alpha_t(g)\bigr)-\mu(f^*)\mu(g),
    \qquad \mu(f)=\int_X f\,d\mu\ .
\end{align}
Ergodicity is equivalent to the vanishing of the time-averaged connected correlation for every pair of observables:
\begin{align}
    \lim_{T\to\infty}\frac{1}{T}\int_0^T C^{con}_{f,g}(t)\,dt=0\ .
\end{align}
A stronger condition, {\it weak mixing}, requires
\begin{align}
    \lim_{T\to\infty}\frac{1}{T}\int_0^T |C^{con}_{f,g}(t)|^2\,dt=0\ .
\end{align}
The flow is {\it mixing} if $C^{con}_{f,g}(t)\to0$ as $t\to+\infty$. These conditions distinguish cancellation after time averaging from actual loss of correlation at late times. Mixing does not specify a decay rate and does not imply that there are no autonomous subsystems.

\paragraph{Correlations between subspaces:} Write $C(f,g)=C_{f,g}(0)$ and similarly for its connected part. Of course, the definition of correlation above changes if we multiply the function by a constant. To account this and define a measure of correlations between two subalgebras, we define
\begin{align}
    C(O_1,O_2)\equiv \sup_{\text{supp}(f_i)\in O_i,|f_i|_\infty\leq 1}|C(f_1,f_2)|
\end{align}
where $|f|_\infty=\sup_{x\in X}|f(x)|$. It follows from the definition above that 
\begin{align}
    C(O_1,O_2)\leq \int_{O_1\cap O_2} 1=\mu(O_1\cap O_2)=C(1_{O_1},1_{O_2})
\end{align}
Therefore, we have
\begin{align}
    C(O_1,O_2)=C(1_{O_1},1_{O_2})\ .
\end{align}

For connected correlations, define the event coefficient
\begin{equation}
 \alpha(O_1,O_2)=
 \sup_{\substack{A\subset O_1\\B\subset O_2}}
 |\mu(A\cap B)-\mu(A)\mu(B)|.
\end{equation}
In particular,
\begin{equation}
 C^{con}_{1_{O_1},1_{O_2}}(t)
 =\mu\bigl(O_1\cap\Phi_t(O_2)\bigr)-\mu(O_1)\mu(O_2).
\end{equation}
This measures the correlation of an event with a transported event.

\paragraph{Correlations between subsystems:}
Let $X=Y\times Z$ carry a joint probability measure $\mu_{YZ}$ with marginals
$\mu_Y,\mu_Z$. The unconnected correlation has supremum one over functions bounded
by one, because constant functions attain it. For connected correlations, distinguish
\begin{align}
 \alpha(Y,Z)&=\sup_{\substack{A\subset Y\\B\subset Z}}
 |\mu_{YZ}(A\times B)-\mu_Y(A)\mu_Z(B)|,\nn\\
 \Gamma(Y,Z)&=\sup_{\substack{f,g\ {\rm real}\\\|f\|_\infty,\|g\|_\infty\leq1}}
 |\mu_{YZ}(fg)-\mu_Y(f)\mu_Z(g)|=4\alpha(Y,Z).
\end{align}
The factor four follows by rescaling real functions from $[-1,1]$ to $[0,1]$
and writing the latter as integrals of indicator functions. With
$\nu=\mu_{YZ}-\mu_Y\otimes\mu_Z$, total variation satisfies
\begin{equation}
 \alpha(Y,Z)\leq\operatorname{TV}(\mu_{YZ},\mu_Y\otimes\mu_Z)
 :=\sup_E|\nu(E)|=\tfrac12\|\nu\|_1,
 \qquad \Gamma(Y,Z)\leq2\operatorname{TV}.
\end{equation}
Here $\|\nu\|_1$ is the total variation norm of the signed measure. All these
coefficients vanish precisely when the joint measure is a product.

\paragraph{Subsystem coarse-grainings of a free relativistic field theory:} Classical field theory provides a natural setting for a net of subsystem algebras with infinitely many degrees of freedom. For an unconstrained real scalar field on a Cauchy surface $\Sigma$, a phase-space point is a pair of initial data $(\phi,\pi)$. An open spatial region $U\subset\Sigma$ selects the unital polynomial algebra
\begin{align}
    \mathcal B_U
    :=\operatorname{Alg}_{\mathbb R}\{1,\phi(f),\pi(g):f,g\in C_c^\infty(U)\},
    \qquad
    \{\phi(f),\pi(g)\}=\int_\Sigma fg\,d\Sigma\ .
\end{align}
Here $\phi(f)=\int_\Sigma f\phi\,d\Sigma$, and similarly for $\pi(g)$; the other brackets between generators vanish. Thus $\mathcal B_U$ is a Poisson subalgebra, and algebras associated with disjoint spatial regions Poisson commute. Retaining $\mathcal B_U$ resolves the field data in $U$ while discarding the exterior data. Although labeled by a spatial region, this is a subsystem coarse-graining: $U$ selects degrees of freedom, whereas a region of phase space selects entire field configurations.


Choose an ambient Poisson algebra $\mathcal A$ containing these observables and stable under a complete Hamiltonian evolution. With our convention for the pullback, the transported subsystem is
\begin{align}
    \alpha_t^H(F)=F\circ\Phi_{-t}^H,
    \qquad \partial_t\alpha_t^H(F)=\{H,\alpha_t^H(F)\},
    \qquad \mathcal B_{U,t}=\alpha_t^H(\mathcal B_U)\ .
\end{align}
Hamiltonian evolution preserves Poisson brackets, so every $\mathcal B_{U,t}$ remains a Poisson subalgebra. Invariance for all $t$ gives an autonomous reversible subsystem. More generally, propagation can make the transported observables depend on initially discarded data. Whether their time translates generate $\mathcal A$, or instead form a nested family, is a property of the chosen subsystem and dynamics together.


\subsection{Ergodic hierarchy, subspaces and subalgebras}
\label{subsec:classical-ergodic-hierarchy}

The ergodic hierarchy describes how much information about a state remains accessible at late times. In the algebraic language, it can be studied both through correlations of observables and through the evolution of coarse observable algebras. For our purposes, the relevant distinction is between conserved events, autonomous subsystems, and subalgebras that are preserved only in one time direction.

\paragraph{Ergodicity, conserved subspaces and conserved subsystems:}
Let $(\mA,\mu,\alpha_t)$ be the measure-preserving dynamical system introduced above. The flow is {\it ergodic} if every measurable set invariant under all $\Phi_t$ has measure zero or one. Equivalently, its fixed-point algebra is trivial:
\begin{align}
    \mA^\alpha
    :=\{f\in\mA:\alpha_t(f)=f\text{ for all }t\in\mathbb R\}
    =\mathbb C1\ .
\end{align}
For nonatomic $\mu$ and $0<\mu(O)<1$, ergodicity forbids invariant subspace algebras. If 
\begin{align}
    \alpha_t(\mA_O)=\mA_O,\qquad t\in\mathbb R,
\end{align}
then $\Phi_t(O)=O$ modulo null sets and $1_O\in\mA^\alpha$ (it is a {\it conserved event}), which contradicts ergodicity. However, ergodicity does not exclude an invariant subsystem $\mB\subset\mA$: the condition $\alpha_t(\mB)=\mB$ permits observables in $\mB$ to evolve into one another. Such a subsystem is an autonomous factor of the dynamics (conserved subsystem). Taking two identical independent copies of a classical dynamical system always results in autonomous factors; however, a standard result in ergodic theory implies that if the original system is weakly mixing, the combined dynamical system is ergodic \footnote{For an invertible discrete-time example, take two independent bilateral Bernoulli shifts on $\{0,1\}^{\mathbb Z}$ with their product probability measures. Each shift, and their product, is mixing.}.

\paragraph{Information recovered by subalgebra evolution:}
Given a coarse algebra $\mB$, the family $\mB_t=\alpha_t(\mB)$ describes the observables accessible as this coarse description is transported by the dynamics. If $\mB_t=\mB$, evolution remains within the original description. Otherwise, the transported observables can depend on variables that were initially hidden. The algebra
\begin{align}
    \mA_{\mB}:=\bigvee_{t\in\mathbb R}\alpha_t(\mB)
\end{align}
contains the information accessible from all time translates of the coarse observables. Here $\bigvee$ denotes the generated von Neumann algebra. The condition $\mA_{\mB}=\mA$ says that these translates generate the full observable algebra. This is a property of the chosen coarse description, distinct from ergodicity of the dynamics.

%{\paragraph Correlations among subalgebras:} We can also define 


\paragraph{Half-sided subalgebra translations and the Kolmogorov property:}
Of particular interest are subalgebras satisfying
\begin{align}\label{eq:classical-one-sided-inclusion}
    \alpha_t(\mB_+)\subseteq\mB_+,\qquad t\geq0\ .
\end{align}
The family $\mB_t=\alpha_t(\mB_+)$ is then decreasing: $\mB_{t+s}\subseteq\mB_t$ for $s\geq0$. To interpret this inclusion, one can take $\mB_+$ to be generated by observations on a half-infinite time interval, with the orientation chosen so that positive $t$ removes the earliest part of the record.

The Kolmogorov, or $K$, property can be formulated through the existence of such a one-sided algebra whose translates generate the full algebra and whose tail is trivial:
\begin{align}\label{eq:classical-K-algebra}
    \bigvee_{t\in\mathbb R}\mB_t=\mA,
    \qquad
    \bigcap_{t\geq0}\mB_t=\mathbb C1\ .
\end{align}
The first condition says that the complete record resolves all observables. The second says that no nonconstant observable remains accessible arbitrarily far into the retained half of the record. The analogous conditions with integer time describe a $K$-automorphism. They provide an inclusion-based formulation of loss of memory and motivate the future and past algebras used in quantum ergodic theory~\cite{ouseph2024local}.

\paragraph{Example: a Bernoulli shift:}
Consider independent binary variables $x_n$, indexed by $n\in\mathbb Z$, with the uniform product measure on $X=\{0,1\}^{\mathbb Z}$. Let $\alpha(x_n)=x_{n+1}$, and let $\mB_+$ be the von Neumann algebra generated by $x_n$ with $n\geq0$. Then
\begin{align}
    \alpha^k(\mB_+)
    =\operatorname{vN}\{x_n:n\geq k\},
    \qquad
    \alpha(\mB_+)\subsetneq\mB_+\ .
\end{align}
All integer translates generate $\mA$, and the tail algebra is trivial by the independence of the coordinates and the Kolmogorov zero-one law. At each step the retained description forgets one more coordinate, although the shift on the full system is invertible and preserves the measure. This is a subsystem coarse-graining: the retained and discarded coordinates give tensor factors of the observable algebra.

This example also distinguishes one-sided subsystem inclusions from shrinking phase-space regions. In a finite measure space, $\Phi_t(O)\subseteq O$ and measure preservation imply $\mu(O\setminus\Phi_t(O))=0$. Thus a region cannot shrink strictly modulo null sets under a measure-preserving flow. In the nonatomic setting of Section~\ref{subsec:classical-region-subsystem}, the same obstruction applies to proper one-sided inclusions of the algebras $\mA_O$. It does not apply to the subsystem algebras of the Bernoulli example.

\paragraph{Coarse-graining along a decreasing family:}
Each $\mB_t$ has a $\mu$-preserving conditional expectation $\mathcal{E}_t:\mA\to\mB_t$. For $s\geq t$, the tower property gives
\begin{align}
    \mathcal{E}_s\mathcal{E}_t=\mathcal{E}_s\ .
\end{align}
For a fixed state with density $\rho$, its coarse-grained densities $\rho_t=\mathcal{E}_t(\rho)$ therefore satisfy
\begin{align}
    H(\rho_s)\geq H(\rho_t),
    \qquad
    S(\rho\|\rho_s)\geq S(\rho\|\rho_t),
    \qquad s\geq t,
\end{align}
whenever the entropies are finite. Here the loss of information follows from the nesting of the retained algebras. It does not require irreversible microscopic dynamics. For a trivial tail, the reverse martingale theorem further gives $E_t(f)\to\mu(f)1$ in $L^2$ for every $f\in L^2(X,\mu)$.

% \paragraph{Towards state-dependent subalgebra flows:}
% The structures above depend on a specified dynamics $\alpha_t$. Measure theory alone does not assign a nontrivial flow to a state, and all the coarse algebras considered here admit $\mu$-preserving conditional expectations. In the quantum setting, modular theory supplies a distinguished dynamics, and compatibility with that dynamics controls the existence of state-preserving expectations. The one-sided relation \eqref{eq:classical-one-sided-inclusion} anticipates the inclusion structure of a half-sided modular flow, without by itself supplying its quantum modular properties.

% To make a classical comparison with state-generated modular evolution, we next introduce a Poisson bracket and the Hamiltonian associated with $-\log\rho$. The ergodic properties of this dynamics must be assessed on its appropriate invariant sectors: on a symplectic phase space it conserves $\rho$ pointwise, so a nonconstant smooth density already supplies nontrivial conserved events on the full phase space.

\subsection{Classical avatar of modular flow in Hamiltonian systems}
\label{subsec:classical-state-flow}

So far, the state and the dynamics have been independent inputs. A probability density assigns weights to microstates, but measure theory alone does not turn these weights into a motion of the microstates. A Poisson bracket supplies the missing structure: it converts a function into a Hamiltonian vector field. Applied to the logarithm of a density, it gives a distinguished state-dependent dynamics. This construction will let us compare the dynamics associated with a state and with its coarse description.

\paragraph{Poisson brackets and Hamiltonian motion:}
A Poisson structure on a smooth manifold $X$ is an antisymmetric bivector $\Pi$ whose bracket
\begin{align}
    \{f,g\}=\Pi(df,dg),\qquad f,g\in C^\infty(X),
\end{align}
satisfies the Jacobi identity. The bracket obeys the Leibniz rule in each argument, so fixing a smooth real Hamiltonian $H$ defines a vector field. We use the convention
\begin{align}\label{eq:classical-Hamiltonian-convention}
    X_H(f)&=\{f,H\},\qquad
    \alpha_t^H(f)=f\circ\Phi_{-t}^H,\nn\\
    \delta_H(f)&:=\left.\partial_t\alpha_t^H(f)\right|_{t=0}=\{H,f\},
\end{align}
where $\Phi_t^H$ is the flow of $X_H$. Thus the generator on observables has the opposite sign from the vector field transporting points. The Jacobi identity ensures that Hamiltonian evolution preserves the Poisson bracket. In canonical coordinates, $\{q_i,p_j\}=\delta_{ij}$ and the point trajectories satisfy the usual Hamilton equations.

\paragraph{Casimirs and symplectic leaves:}
A Casimir is a function $c$ with $\{c,f\}=0$ for every smooth $f$. It is conserved by every Hamiltonian, rather than by one chosen dynamics. The points accessible from a given microstate by successive Hamiltonian motions form a symplectic leaf. Every Casimir is constant on each leaf, although global Casimirs need not distinguish all leaves. A leaf is therefore a classical analogue of an irreducible sector for Hamiltonian motion. This analogy concerns the Poisson bracket: the pointwise algebra $L^\infty(X,\mu)$ remains abelian, and restricting it to a leaf does not turn it into a nonabelian factor.

\paragraph{Liouville measure and the state Hamiltonian:}
For the explicit constructions below, take $X$ to be a connected compact symplectic manifold without boundary, with symplectic form $\omega$ and normalized Liouville measure
\begin{align}
    d\mu=\frac{\omega^n/n!}{\int_X\omega^n/n!},\qquad \dim X=2n.
\end{align}
Hamiltonian flows preserve this measure. Equivalently,
\begin{align}\label{statepreserve}
    \int_X\{f,g\}\,d\mu=0
\end{align}
for smooth $f,g$. Compactness makes smooth Hamiltonian flows complete and avoids boundary terms. On a noncompact phase space these properties, together with the required integrability, must be checked separately.

Let $\rho$ be a smooth strictly positive probability density relative to $\mu$. Define its {\it state Hamiltonian} and the associated evolution by
\begin{align}\label{eq:classical-state-Hamiltonian}
    K_\rho=-\log\rho,\qquad
    \Phi_t^\rho:=\Phi_t^{K_\rho},\qquad
    \alpha_t^\rho:=\alpha_t^{K_\rho}.
\end{align}
The infinitesimal evolution of an observable is
\begin{align}\label{eq:classical-state-derivation}
    \delta_\rho(f)=\{K_\rho,f\}=\{f,\log\rho\}.
\end{align}
The flow moves along level sets of the density, since $X_{K_\rho}(\rho)=0$. It preserves both $\mu$ and the state $\omega_\rho(f)=\int_X\rho f\,d\mu$. For a Gibbs density $\rho=Z^{-1}e^{-\beta H}$, we have $K_\rho=\beta H+\log Z$ and $\Phi_t^\rho=\Phi_{\beta t}^H$. Thus the state recovers the Hamiltonian motion, with a rescaled parameter; for a general density there is no prior identification of this parameter with physical time.

\paragraph{The classical KMS identity:}
The logarithm is distinguished by an integration-by-parts identity. Using \eqref{statepreserve},
\begin{align}
    0=\int_X\{g,f\rho\}\,d\mu
    =-\omega_\rho(\{f,g\})+
    \omega_\rho\bigl(f\{g,\log\rho\}\bigr).
\end{align}
Consequently,
\begin{align}\label{KMSclassical}
    \omega_\rho(\{f,g\})=\omega_\rho(f\delta_\rho(g)).
\end{align}
For a Gibbs state this reads $\omega_\beta(\{f,g\})=\beta\omega_\beta(f\delta_H(g))$. This is the classical, or static, KMS condition, written with the time orientation in \eqref{eq:classical-Hamiltonian-convention}~\cite{vanDeVen2022}. It is stronger than stationarity: for a given smooth positive density, testing against all smooth $f$ determines $\delta_\rho(g)$ for every $g$. Conversely, on a connected symplectic manifold, imposing this identity for a specified Hamiltonian dynamics determines a smooth positive density to be Gibbs, up to normalization. The identity also motivates formulations on algebras of local observables in infinite systems, where existence of the dynamics and equilibrium states requires additional hypotheses; a global Gibbs density need not exist.

\paragraph{Poisson modular flow:}
There is a geometric formulation of the same construction. For a smooth positive measure $\nu$, define its Poisson modular vector field by
\begin{align}
    \mathcal M_\nu(f):=\operatorname{div}_\nu X_f.
\end{align}
This measures the failure of Hamiltonian motions to preserve $\nu$~\cite{Weinstein1997}. For $d\nu=\rho\,d\mu$ and Liouville $\mu$, the divergence formula gives
\begin{align}
    \mathcal M_\nu(f)=X_f(\log\rho)
    =\{f,-\log\rho\}=X_{K_\rho}(f).
\end{align}
Thus $\Phi_t^\rho$ is precisely this Poisson modular flow; its action on observables uses the inverse pullback, so $\delta_\rho=-\mathcal M_\nu$. On a general Poisson manifold and for an arbitrary smooth reference measure, the formula instead reads $\mathcal M_{\rho\mu}=\mathcal M_\mu+X_{-\log\rho}$. A reference measure invariant under every Hamiltonian need not exist. Our symplectic setting removes this additional obstruction.

The construction uses the Poisson structure in an essential way. It is distinct from the multiplication group $U_t^\rho$ introduced above and from the trivial modular automorphism group of the abelian von Neumann algebra. It also supplies no automatic ergodic or entropy-production property. A nonconstant $\rho$ is itself a conserved observable, so the flow cannot be ergodic on the full probability space. Questions about ergodicity and mixing must be posed on invariant sectors with suitable invariant measures. The comparison of full and coarse state Hamiltonians, however, is meaningful without assuming ergodicity.

\subsection{Subsystem Poisson flows in classical many-body systems and field theory}
\label{subsec:classical-subalgebra-Hamiltonian}

A coarse description determines a reduced state and hence, when the coarse space is Poisson, its own state-generated dynamics. We can then ask whether the reduced dynamics agrees with what the full dynamics does to the observed variables. Comparing the two motions gives a classical model of the subalgebra flows studied in this paper. The comparison separates evolution already visible within the retained algebra from evolution involving the discarded variables.

\paragraph{Poisson subsystems:} We say a subalgebra $\mB_Y\subset \mA$ is a Poisson subsystem if for all $b_1,b_2\in \mB_Y$, we have $\{b_1,b_2\}\in \mB_Y$. In the geometric setting below, $\pi:X\to Y$ is a smooth surjective Poisson submersion:
\begin{align}\label{eq:classical-Poisson-subsystem}
    \{\pi^*f,\pi^*g\}_X=\pi^*\{f,g\}_Y,
    \qquad f,g\in C^\infty(Y).
\end{align}
The smooth coarse observables $\mathcal B=\pi^*C^\infty(Y)$ are closed under the bracket. This condition allows us to speak of Hamiltonian dynamics within the coarse description.

The simplest setting makes both the geometric and probabilistic structures explicit. Take a product $X=Y\times Z$ of connected compact symplectic manifolds, with product symplectic structure and normalized Liouville measure $\mu=\mu_Y\otimes\mu_Z$. The projection to $Y$ is Poisson, and the associated von Neumann subalgebra is $\mB_Y=L^\infty(Y,\mu_Y)\otimes1$. For a smooth positive density $\rho$, restriction of the state gives the marginal density
\begin{align}\label{eq:classical-marginal-Hamiltonian}
    \rho_Y(y)=\int_Z\rho(y,z)\,d\mu_Z(z),\qquad
    K_Y=-\log\rho_Y.
\end{align}
The expectation $\mE_Y$ from Section~\ref{subsec:classical-region-subsystem} extends this marginal back to $X$ as $\rho_0(y,z)=\rho_Y(y)$. Here, $\rho_Y=\pi_*\rho$ is the classical analog of the reduced state on $Y$, and $\rho_0=\pi^*\rho_Y=\pi^*\pi_* \rho$ is the pullback of $\rho_Y$ to $X$, which is the coarse-grained version of $\rho$. Its Hamiltonian is $K_0=\pi^*K_Y$, whose flow is
\begin{align}
    \Phi_t^0(y,z)=(\Phi_t^Y(y),z).
\end{align}
This provides a definite lift of the reduced motion to the full phase space. For a general Poisson submersion, a specified coarse Hamiltonian can likewise be pulled back, but the marginal measure need not be Liouville and a complementary subsystem need not exist. We use the product setting below to keep the marginal, its reference measure, and its lifted Hamiltonian unambiguous.


% \NL{Include discussion of singular inclusions, if they exist}

\paragraph{Agreement between reduced and full Poisson dynamics:}
Agreement means that the observed trajectory is determined by its initial coarse state and follows the flow generated by $K_Y$:
\begin{align}\label{eq:classical-state-intertwining}
    \pi\circ\Phi_t^\rho=\Phi_t^Y\circ\pi.
\end{align}
For our complete smooth flows, this is equivalent to the infinitesimal condition
\begin{align}\label{eq:classical-state-defect}
    \{\pi^*g,K_\rho\}_X-\pi^*\{g,K_Y\}_Y
    =\{\pi^*g,K_\rho-K_0\}_X=0,
    \qquad g\in C^\infty(Y).
\end{align}
The left side measures the part of the observed velocity missed by the reduced state Hamiltonian. A Hamiltonian invisible to all coarse observables belongs to the Poisson centralizer
\begin{align}
    \mathcal B^{\mathrm c}_{\mathrm P}
    :=\{h\in C^\infty(X):\{h,b\}=0\text{ for every }b\in\mathcal B\}.
\end{align}
Thus agreement is precisely the condition $K_\rho-K_0\in\mathcal B^{\mathrm c}_{\mathrm P}$. This is an analogue of a relative commutant defined using Poisson brackets. The ordinary operator commutant of a classical subalgebra contains the whole abelian observable algebra and cannot express this condition.

In the product setting, nondegeneracy of the symplectic form and connectedness of $Y$ imply that this centralizer consists of functions of $z$ alone. Hence \eqref{eq:classical-state-intertwining} is equivalent to
\begin{align}\label{eq:classical-product-state-closure}
    K_\rho(y,z)=K_Y(y)+K_Z(z),
    \qquad
    \rho(y,z)=\rho_Y(y)\rho_Z(z).
\end{align}
Normalization of the marginal fixes $\int_Z\rho_Z\,d\mu_Z=1$. The condition \eqref{eq:classical-state-defect} implies that the coarse observables have the same Poisson evolution under $\rho_0$ and $\rho$. In other words:
\begin{align}
    \forall g\in C^\infty(Y):\qquad \alpha^\rho_t\circ\alpha_{-t}^0(\pi^*g)=\pi^*g\ .
\end{align}
This compares full and reduced observable evolution, as in Section~\ref{subsec:quantum-coarse}; the state-preserving expectation criterion itself differs in the classical and quantum theories.

\paragraph{Comparing the two flows:}
Define the relative transformation by first undoing the lifted coarse motion and then applying the full motion:
\begin{align}\label{eq:classical-relative-transformation}
    R_t=\Phi_t^\rho\circ\Phi_{-t}^0,
    \qquad
    \mathcal T_t(f):=f\circ R_t^{-1}
    =\alpha_t^\rho\circ\alpha_{-t}^0(f).
\end{align}
For each $t$, this is a symplectic transformation and a measure-preserving automorphism, respectively. Since the lifted coarse flow preserves $\mB_Y$,
\begin{align}\label{eq:classical-relative-algebra-flow}
    \mathcal T_t(\mB_Y)=\alpha_t^\rho(\mB_Y).
\end{align}
The relative transformation therefore follows the same moving subalgebra as the full dynamics, while canceling the chosen internal coarse evolution. It fixes every coarse observable exactly when \eqref{eq:classical-state-intertwining} holds. For a product state, it can still move hidden variables (the relative commutant): $R_t(y,z)=(y,\Phi_t^Z(z))$.

The infinitesimal generator of $R_t$ is the difference of the state Hamiltonians:
\begin{align}\label{eq:classical-relative-Hamiltonian}
    K_{\mathrm{rel}}&=K_\rho-K_0=-\log\frac{\rho}{\rho_0},\nn\\
    \left.\partial_tR_t\right|_{t=0}&=X_{K_{\mathrm{rel}}},\qquad
    \left.\partial_t\mathcal T_t(f)\right|_{t=0}=\{K_{\mathrm{rel}},f\}.
\end{align}

Every normal state of the hidden factor, with density $\rho_Z$ relative to $\mu_Z$,
gives an expectation $E_{\rho_Z}:\mA\to\mB_Y$:
\begin{equation}
 (E_{\rho_Z}f)(y)=\int_Z\rho_Z(z)f(y,z)\,d\mu_Z(z).
\end{equation}
These states parameterize the expectations with $y$-independent averaging.
General normal expectations allow measurable probability kernels $\nu_y(dz)$
depending on $y$, with $\mu_Y(dy)\nu_y(dz)\ll\mu_Y\otimes\mu_Z$:
$E_\nu f(y)=\int_Z f(y,z)\nu_y(dz)$.

\paragraph{Example: a uniform marginal with nonautonomous observables:}
Take $Y$ and $Z$ to be two-dimensional tori with angular coordinates $(q,p)$ and $(Q,P)$, each of period $2\pi$, and the product symplectic form $dq\wedge dp+dQ\wedge dP$. Consider
\begin{align}
    \rho(q,p,Q,P)=\frac{e^{\epsilon\cos(q-Q)}}{\mathcal Z_\epsilon},
    \qquad
    \mathcal Z_\epsilon=\int_X e^{\epsilon\cos(q-Q)}\,d\mu,
    \qquad \epsilon\ne0.
\end{align}
Translation invariance of the angular integral gives $\rho_Y=1$, so $K_0=0$ and the reduced state predicts no motion. The full state Hamiltonian is $K_\rho=-\epsilon\cos(q-Q)+\log\mathcal Z_\epsilon$, giving
\begin{align}
    \dot q=\dot Q=0,\qquad
    \dot p=-\epsilon\sin(q-Q),\qquad
    \dot P=\epsilon\sin(q-Q).
\end{align}
For example, the bounded coarse observable $e^{ip}$ evolves as
\begin{align}
    \alpha_t^\rho(e^{ip})
    =\exp\bigl(i[p+\epsilon t\sin(q-Q)]\bigr).
\end{align}
Its future value depends on the hidden angle $Q$, so it leaves $\mB_Y$ for $t\ne0$. The full state and its uniform marginal are stationary, but the retained observables do not evolve autonomously. Here $R_t=\Phi_t^\rho$ directly displays the motion that the marginal state Hamiltonian misses.

\paragraph{Relative entropy and the information discarded:}
The same relative logarithm appeared in the phase construction at the beginning of this section. Here it generates a vector field, while its expectation gives
\begin{align}\label{eq:classical-relative-entropy-Hamiltonian}
    S(\rho\|\rho_0)=-\omega_\rho(K_{\mathrm{rel}})
    =H(\rho_0)-H(\rho).
\end{align}
The vector field alone does not determine this number: on a connected symplectic manifold it determines its Hamiltonian only up to a constant. The normalized densities fix that constant. On a general Poisson manifold, Casimirs give a larger ambiguity.

Information loss and failure of autonomous coarse dynamics are also distinct. For a product state, the relative transformation acts only on $Z$, yet
\begin{align}
    S(\rho_Y\rho_Z\|\rho_Y\otimes1)
    =\int_Z\rho_Z\log\rho_Z\,d\mu_Z
\end{align}
can be positive. Coarse-graining has discarded the hidden distribution even without correlations. More generally,
\begin{align}
    S(\rho\|\rho_Y\otimes1)
    =S(\rho\|\rho_Y\rho_Z)+S(\rho_Z\|1),
\end{align}
where $\rho_Z$ is the other marginal. The first term is the mutual information and the second measures the hidden marginal relative to the chosen reference. Neither this entropy nor the relative transformation supplies a monotone entropy-production law without an additional coarse-graining or dynamical assumption.

\paragraph{Subregions and the passage to field theory:}
The measurable region algebras $\mA_O$ remain valid coarse descriptions. What can fail is the smoothness needed to assign their coarse densities a Hamiltonian flow. For the region expectation $E_O$,
\begin{align}
    \rho_0=E_O(\rho)=\chi_O\rho+c_O\chi_{O^c},\qquad
    c_O=\frac{1}{\mu(O^c)}\int_{O^c}\rho\,d\mu.
\end{align}
For an open region with smooth boundary this density is generally discontinuous there. In the interior of $O$, $K_0=K_\rho$; in the interior of the complement, $K_0$ is constant. These local statements do not specify a smooth global Hamiltonian flow across the boundary. Additional regularization or boundary conditions are needed. This is a limitation of the Hamiltonian construction, not of the measure-theoretic net.

Finally, every classical coarse algebra still admits a state-preserving conditional expectation for a faithful probability state, even when its Poisson state flow fails to preserve that algebra. In the product example this expectation is explicitly
\begin{align}
    E_{\rho,Y}(f)(y)=\int_Z f(y,z)\frac{\rho(y,z)}{\rho_Y(y)}\,d\mu_Z(z).
\end{align}
Its existence does not imply \eqref{eq:classical-state-intertwining}. In quantum modular theory the relation changes: for a faithful normal state $\varphi$ on $\mA$, Takesaki's theorem identifies invariance of a unital von Neumann subalgebra $\mB$ under $\sigma_t^\varphi$ with the existence of a $\varphi$-preserving normal conditional expectation onto $\mB$~\cite{Takesaki1972}. Failure of invariance rules out an expectation preserving this state; it does not by itself rule out expectations preserving other states. This is the additional structure that the quantum theory contributes to the comparison of full and coarse dynamics.

The classical construction thus motivates studying how relative state evolutions transport retained observables. In QFT a spacetime region selects local field observables, whereas a region $O\subset X$ above selects possible microstates. The connection to Section~\ref{sec:geometric-subregion-flows} is through inclusions of observable algebras and their state-dependent flows; it does not identify these two notions of region.

% Bibliography addition for the new citation key (insert into modularscattering.bib):
% @misc{vanDeVen2022,
%   author = {van de Ven, Christiaan J. F.},
%   title = {Gibbs states and their classical limit},
%   year = {2022},
%   eprint = {2211.01755},
%   archivePrefix = {arXiv},
%   primaryClass = {math-ph},
%   url = {https://arxiv.org/abs/2211.01755}
% }

A simple example of strict nesting is supplied by a classical chiral current $j(x)$ on the line. On rapidly decreasing smooth configurations, choose
\begin{align}
    \{j(x),j(y)\}=-\partial_x\delta(x-y),
    \qquad P=\frac12\int_{\mathbb R}j(x)^2\,dx,
    \qquad \tau_a j(x)=j(x+a)\ .
\end{align}
The nonnegative Hamiltonian $P$ generates $\tau_a$ through $\{P,j(x)\}=\partial_xj(x)$. Let $\mathcal B_b$ be generated by $1$ and $j(f)$ with $f\in C_c^\infty((b,\infty))$. Since $\tau_a j(f)=j(f(\,\cdot-a))$,
\begin{align}
    \tau_a(\mathcal B_b)=\mathcal B_{b+a}\subsetneq\mathcal B_b\quad(a>0),
    \qquad
    \bigcup_{b\in\mathbb R}\mathcal B_b=\mathcal A_{\mathrm{loc}},
    \qquad
    \bigcap_{b\geq0}\mathcal B_b=\mathbb R1\ .
\end{align}
Here $\mathcal A_{\mathrm{loc}}$ is the polynomial algebra of compactly smeared currents; $P$ is a global functional generating its automorphisms. The full translation group is reversible, but its positive half acts on $\mathcal B_0$ by injective, nonsurjective Poisson endomorphisms. The decreasing family progressively discards accessible field degrees of freedom.

The dilation Hamiltonian $D=\frac12\int_{\mathbb R}xj(x)^2\,dx$ gives a further connection with half-sided modular inclusions. Its flow $\delta_s$ satisfies
\begin{align}
    \{D,P\}=P,
    \qquad \delta_s(\mathcal B_b)=\mathcal B_{e^s b},
    \qquad \delta_s\tau_a\delta_{-s}=\tau_{e^s a}\ .
\end{align}
Thus dilations preserve $\mathcal B_0$ and act half-sidedly on $\mathcal B_b$ for $b>0$. These are classical Hamiltonian realizations of the translation--dilation relations. Their construction uses the Poisson structure; identifying the dilation with a state-generated modular evolution requires additional input.


\section{Continuity of Characteristic Functions}\label{section:continuity}
In this appendix, we consider the question of continuity of characteristic function. First we recall the following key property of modular operators:
\begin{theorem}\label{theorem:continuityofmodular}
    Consider a sequence of normal faithful states $\{\omega_n\}$ on a von Neumann algebra $\mathcal{A}$ converging in norm to a normal faithful state $\omega$. Then the modular operators converge in the strong resolvent sense. Or equivalently, the unitary groups $\{\Delta_{\omega_n}^{-it}\}$ converges (point-wise in time) strongly to $\Delta_\omega^{-it}$. And on any bounded domain in time, the one-parameter unitary groups $\{\Delta_{\omega_n}^{-it}\}$ converges uniformly to $\Delta_\omega^{-it}$.
\end{theorem}
Using the theorem and the Uniform Boundedness Principle, the following corollary is obvious:
\begin{corollary}\label{corollary:firstcorollary}
    Given a pair of type $III$ von Neumann algebras $\mathcal{B}\subset\mathcal{A}$ acting on an common separable Hilbert space $\mathcal{H}$. Consider a sequence of normal faithful states $\{\omega_n:=\bra{\Omega_n}\cdot\ket{\Omega_n}\}$ where each vector $\ket{\Omega_n}$ is a common cyclic and separating vector of $\mathcal{A}$ and $\mathcal{B}$. Assume this sequence converges in norm to a normal faithful state $\omega:=\bra{\Omega}\cdot\ket{\Omega}$ where $\ket{\Omega}$ is also a common cyclic and separating vector. Then the characteristic functions $\{D_{\mathcal{B}\subset \mathcal{A};n}(t) := \Delta_{\omega_n,\mathcal{A}}^{-it}\Delta_{\omega_n, \mathcal{B}}^{it}\}$ converges (point-wise in time) strongly to $\Delta_{\omega,\mathcal{A}}^{-it}\Delta_{\omega,\mathcal{B}}^{it}$. In addition, on any bounded domain in time, the sequence of characteristic functions converges strongly uniformly.
\end{corollary}
Notice that the assumptions in this corollary is not restrictive. By completeness of Hilbert space and the fact that the set of common cyclic and separating vectors is a dense $G_\delta$ subset of the Hilbert space (cf. Proposition~\ref{prop:dense-common}), such sequences of normal faithful states as required in the corollary are easy to find. One can simply take any convergent sequence of vectors in the $G_\delta$ set. The only caveat is that the limit needs to be inside the same $G_\delta$ set as well. 

Moreover, notice that because both $\mathcal{A}$ and $\mathcal{B}$ acting on $\mathcal{H}$ are in standard forms, all vectors $\ket{\Omega_n}$ representing states $\omega_n$ are in the same positive cone. Hence the modular conjugations of all of these states are the same. Unambiguously, we denote the modular conjugation associated with $\mathcal{A}$ as $J_\mathcal{A}$ and the modular conjugation associated with $\mathcal{B}$ as $J_\mathcal{B}$. We can now make the following simple calculation:
\begin{align}
    \begin{split}
        J_\mathcal{A}\Delta_{\omega_n,\mathcal{A}}^{-it + 1/2}\Delta_{\omega_n,\mathcal{B}}^{it - 1/2}J_\mathcal{B} &= \Delta_{\omega_n, \mathcal{A}}^{-it}J_\mathcal{A}\Delta_{\omega_n, \mathcal{A}}^{1/2}\Delta_{\omega_n,\mathcal{B}}^{-1/2}J_\mathcal{B}\Delta_{\omega_n,\mathcal{B}}^{it}
        \\&=\Delta_{\omega_n, \mathcal{A}}^{-it}S_{\omega_n,\mathcal{A}}S_{\omega_n,\mathcal{B}}\Delta_{\omega_n,\mathcal{B}}^{it}
        \\& = \Delta_{\omega_n,\mathcal{A}}^{-it}\Delta_{\omega_n,\mathcal{B}}^{it}
    \end{split}
\end{align}
where we have used the fact:
\begin{equation}
    S_{\omega_n,\mathcal{A}}S_{\omega_n,\mathcal{B}}\mathcal{O}_\mathcal{B}\ket{\Omega_n} = S_{\omega_n,\mathcal{A}}\mathcal{O}^*_\mathcal{B}\ket{\Omega_n} = \mathcal{O}_\mathcal{B}\ket{\Omega_n}
\end{equation}
This identity holds for the dense subspace of vectors $\{\mathcal{O}_\mathcal{B}\ket{\Omega_n}:\mathcal{O}_\mathcal{B}\in\mathcal{B}\}$. Hence the identity extends to the entire Hilbert space by closability.

This calculation shows immediately that:
\begin{corollary}
    Under the same assumptions as Corollary \ref{corollary:firstcorollary}, the characteristic functions $\{D_{\mathcal{B}\subset \mathcal{A};n}(t+i/2)\}$ converges (point-wise in time) strongly to $\Delta_{\omega,\mathcal{A}}^{-it + 1/2}\Delta_{\omega, \mathcal{B}}^{it - 1/2}$. In addition, on any bounded domain in time, the sequence converges strongly uniformly.
\end{corollary}
Now for any vectors $\ket{\eta}, \ket{\xi}\in\mathcal{H}$, we consider the function:
\begin{equation}
    f^{(n)}_{\eta,\xi}(z) := \bra{\eta}\Delta^{-iz}_{\omega_n,\mathcal{A}}\Delta^{iz}_{\omega_n, \mathcal{B}} - \Delta^{-iz}_{\omega,\mathcal{A}}\Delta^{iz}_{\omega,\mathcal{B}}\ket{\xi}
\end{equation}
where $z\in\mathcal{S}(0,1/2)$. By the basic properties of characteristic functions, $f^{(n)}_{\eta,\xi}(z)$'s are bounded analytic functions in the strip $\mathcal{S}(0, 1/2)$. In addition, the corollaries show that on the boundary we have limits:
\begin{equation}
    \lim_{n\rightarrow\infty}f^{(n)}_{\eta,\xi}(t) = 0
\end{equation}
\begin{equation}
    \lim_{n\rightarrow\infty}f^{(n)}_{\eta, \xi}(t +i/2) = 0
\end{equation}
Hence by Vitali's theorem, on any compact subset of the strip $\mathcal{S}(0, 1/2)$ (in particular, point-wise), $\{f^{(n)}_{\eta,\xi}(z)\}$ converges uniformly to $0$. Because this holds for any vectors $\ket{\eta}, \ket{\xi}\in\mathcal{H}$, we arrive at the conclusion:
\begin{theorem}
    Under the same assumptions as Corollary \ref{corollary:firstcorollary}, the characteristic functions $\{D_{\mathcal{B}\subset \mathcal{A};n}(z):z\in\mathcal{S}(0, 1/2)\}$ converges strongly (point-wise) to $D_{\mathcal{B}\subset \mathcal{A}}(z):=\Delta_{\omega,\mathcal{A}}^{-iz}\Delta_{\omega,\mathcal{B}}^{iz}$. In addition, on any compact subset of the strip, the convergence is uniform.
\end{theorem}
%In particular in the context of $L^2$-nuclearity, we have strong convergence:
%\begin{equation}
   % \lim_{n\rightarrow\infty}\Delta_{\omega_n,\mathcal{A}}^{1/4}\Delta_{\omega_n,\mathcal{B}}^{-1/4} = \Delta_{\omega,\mathcal{A}}^{1/4}\Delta_{\omega,\mathcal{B}}^{-1/4}
%\end{equation}
%If the limiting map is nuclear, then we can decompose:
%\begin{equation}
%    \Delta_{\omega,\mathcal{A}}^{1/4}\Delta_{\omega,\mathcal{B}}^{-1/4} = \sum_k\lambda_k\ket{\xi_k}\bra{\eta_k}
%\end{equation}
%By strong continuity, we know immediately that:
%\begin{equation}
    %\lim_{n\rightarrow\infty}\Delta_{\omega_n,\mathcal{A}}^{1/4}\Delta_{\omega_n,\mathcal{B}}^{-1/4}\ket{\eta_k} = \lambda_k\ket{\xi_k}
%\end{equation}
\subsection{Characteristic Functions of Relative Modular Operators}\label{app:relcharacter}
We take a slight digression in this subsection to briefly discuss a slight generalization of characteristic function. Instead of modular operator, we can study the \textit{characteristic function of relative modular operators}:
\begin{definition}
    Given a pair of von Neumann algebras $\mathcal{B}\subset\mathcal{A}$ and two common cyclic and separating vector $\ket{\Omega_\varphi}, \ket{\Omega_\psi}\in\mathcal{H}$, the \textit{characteristic function of relative modular operator} is defined as:
    \begin{equation}
        D_{\mathcal{B}\subset \mathcal{A};\varphi|\psi}(t) := \Delta_{\varphi|\psi;\mathcal{A}}^{-it}\Delta_{\varphi|\psi;\mathcal{B}}^{it}
    \end{equation}
    where $\varphi$ (and $\psi$ respectively) is the state corresponding to the vector $\ket{\Omega_\varphi}$ (and $\ket{\Omega_\psi}$ respectively) and $\Delta_{\varphi|\psi;\mathcal{A}}$ (and $\Delta_{\varphi|\psi;\mathcal{B}}$) is the relative modular operator of the two states $\varphi,\psi$ restricted to the algebra $\mathcal{A}$ (and $\mathcal{B}$ respectively).
\end{definition}
The main result of this appendix is the following continuity theorem:
\begin{theorem}
    Let $B\subset A$ be a standard inclusion. Consider two sequences of normal faithful states $\{\varphi_n\}$ and $\{\psi_n\}$ such that the corresponding representing vectors $\{\xi_{\varphi_n}\}$ and $\{\xi_{\psi_n}\}$ are common cyclic and separating vectors of $A$ and $B$. Assume both sequences converge in norm $\lim_n\varphi_n = \varphi$ and $\lim_n\psi_n = \psi$ such that the limiting states are normal faithful state and their representing vectors are also common cyclic and separating vectors for $A$ and $B$. In this case, for all complex numbers $z$ in the strip $\Re{z}\in[0, 1/2]$, we have strong convergence:
    \begin{equation}
        \lim_n \Delta_{\varphi_n|\psi_n;A}^z \Delta_{\varphi_n|\psi_n;B}^{-z} = \Delta^z_{\varphi|\psi;A}\Delta^{-z}_{\varphi|\psi; B}
    \end{equation}
\end{theorem}
Before proving the theorem, we comment on the conditions on the theorem. When $B\subset A$ is a standard inclusion, by definition the set of common cyclic and separating vectors is not empty. In this case, it turns out that this set is actually a dense $G_\delta$ set in the representing Hilbert space. Thus, sequences like $\{\varphi_n\}$ and $\{\psi_n\}$ can be easily found within this dense $G_\delta$ set. However, it is necessary to check that the norm limit $\varphi$ and $\psi$ can still be represented by common cyclic and separating vectors. A priori, it is possible to have limiting states that cannot be represented by common cyclic and separating vectors.

The proof is based on a few simple observations.
\begin{lemma}
    Under the same conditions as the theorem above, $\{\Delta_{\varphi_n|\psi_n;A}\}_n$ converges in strong resolvent sense to $\Delta_{\varphi|\psi;A}$. Same for $\{\Delta_{\varphi_n|\psi_n;B}\}$
\end{lemma}
\begin{proof}
    Let $\mathcal{H}$ be the standard representing Hilbert space. Then by definition, we have:
    \begin{equation}
        J_A\Delta^{1/2}_{\varphi|\psi;A}x\xi_{\psi} = x^*\xi_{\varphi}
    \end{equation}
    For any $n$, we have:
    \begin{equation}
        ||\Delta^{1/2}_{\varphi_n|\psi_n;A}x\xi_{\psi} - \Delta^{1/2}_{\varphi|\psi;A}x\xi_{\psi}|| \leq ||\Delta^{1/2}_{\varphi_n|\psi_n;A}x\xi_\psi - \Delta^{1/2}_{\varphi_n|\psi;A}x\xi_\psi|| + ||\Delta^{1/2}_{\varphi_n|\psi;A}x\xi_\psi - \Delta^{1/2}_{\varphi|\psi;A}x\xi_\psi|| 
    \end{equation}
    In addition, for any $n$, we have:
    \begin{equation}
        ||\Delta^{1/2}_{\varphi_n|\psi_n;A}x\xi_\psi - \Delta^{1/2}_{\varphi_n|\psi;A}x\xi_\psi|| \leq ||\Delta^{1/2}_{\varphi_n|\psi_n;A}(x\xi_\psi - x\xi_{\psi_n})|| + ||\Delta^{1/2}_{\varphi_n|\psi_n;A}x\xi_{\psi_n} - \Delta^{1/2}_{\varphi_n|\psi;A}x\xi_{\psi}||
    \end{equation}
    Since $J_A\Delta^{1/2}_{\varphi_n|\psi_n;A}x\xi_{\psi_n} = x^*\xi_{\varphi_n}$ and $J_A\Delta^{1/2}_{\varphi_n|\psi;A}x\xi_\psi = x^*\xi_{\varphi_n}$, we have:
    \begin{align}
        \begin{split}
            ||\Delta^{1/2}_{\varphi_n|\psi_n;A}x\xi_{\psi_n} - \Delta^{1/2}_{\varphi_n|\psi;A}x\xi_{\psi}|| = ||J_A\Delta^{1/2}_{\varphi_n|\psi_n;A}x\xi_{\psi_n} - J_A\Delta^{1/2}_{\varphi_n|\psi;A}x\xi_{\psi}|| = 0
        \end{split}
    \end{align}
    From Powers-Stormer inequality, we have:
    \begin{align}
        ||x\xi_\psi -x\xi_{\psi_n}||^2 \leq ||x||^2||\psi-\psi_n||
    \end{align}
    Hence, we have:
    \begin{align}
        \begin{split}
            ||\Delta^{1/2}_{\varphi_n|\psi_n;A}(x\xi_\psi -x\xi_{\psi_n})|| &= ||J_A\Delta^{1/2}_{\varphi_n|\psi_n;A}(x\xi_\psi - x\xi_{\psi_n})||
            \\&
            \leq ||x^*(\xi_\psi - \xi_{\psi_n})|| \leq ||x||||\psi - \psi_n||\rightarrow 0
        \end{split}
    \end{align}
    Finally, by modular theory we have:
    \begin{equation}
        \lim_n||\Delta^{1/2}_{\varphi_n|\psi;A}x\xi_\psi - \Delta^{1/2}_{\varphi|\psi;A}x\xi_\psi|| = 0
    \end{equation}
    Therefore, combining all these observations, we have:
    \begin{equation}
        \lim_n||\Delta^{1/2}_{\varphi_n|\psi_n;A}x\xi_\psi - \Delta^{1/2}_{\varphi|\psi;A}x\xi_\psi|| = 0
    \end{equation}
    By a theorem in \cite{takesaki2}, this means that $\{\Delta_{\varphi_n|\psi_n;A}\}$ converges to $\Delta_{\varphi|\psi;A}$ in the strong resolvent sense. 

    Because $B\subset A$ is a standard inclusion, $B$ is also in the standard form on $\mathcal{H}$. Replacing all modular conjugation from $J_A$ to $J_B$, the proof above can be used verbatim to show that $\{\Delta_{\varphi_n|\psi_n;B}\}$ converges to $\Delta_{\varphi|\psi;B}$ in the strong resolvent sense as well.
\end{proof}
\begin{corollary}
    Under the same conditions as in the theorem, $\{\Delta^{-it}_{\varphi_n|\psi_n;A}\}$ converges strongly to $\Delta^{-it}_{\varphi|\psi;A}$. Same for $\{\Delta^{-it}_{\varphi_n|\psi_n;B}\}$.

    In addition, the product $\{\Delta^{-it}_{\varphi_n|\psi_n;A}\Delta^{it}_{\varphi_n|\psi_n;B}\}$ converges strongly to $\Delta^{-it}_{\varphi|\psi;A}\Delta^{it}_{\varphi|\psi;B}$.
\end{corollary}
\begin{proof}
    The first statement follows from the characterization of convergence in strong resolvent sense \cite{takesaki2}.

    The second statement follows from the Uniform Boundedness Principle and the first statement.
\end{proof}
Notice that in the proof above, we have used the fact that both $A$ and $B$ represented on $\mathcal{H}$ are in their standard forms. Therefore, the modular conjugations do not depend on the states under consideration. That is why we have used the notations $J_A$ and $J_B$ to emphasize that the modular conjugations only depend on the algebra not the states.

Next we recall some simple relations regarding relative modular operator:
\begin{equation}
    J_A\Delta_{\varphi|\psi;A}J_A = \Delta^{-1}_{\psi|\varphi;A}
\end{equation}
where $\varphi,\psi$ are both normal faithful states.

Next, we make a simple calculation:
\begin{align}
    \begin{split}
        J_A \Delta^{-it + 1/2}_{\varphi|\psi;A}\Delta^{it-1/2}_{\varphi|\psi;B} J_B &= \Delta^{-it}_{\psi|\varphi;A} J_A\Delta^{1/2}_{\varphi|\psi;A}J_B\Delta^{1/2}_{\psi|\varphi;B}\Delta_{\psi|\varphi;B}^{it}
        \\&=\Delta^{-it}_{\psi|\varphi;A}S_{\varphi|\psi;A}S_{\psi|\varphi;B}\Delta^{it}_{\psi|\varphi;B}
    \end{split}
\end{align}
where $S_{\varphi|\psi;A}$ and $S_{\psi|\varphi;B}$ are the relative modular operators. A simple calculation shows:
\begin{equation}
    S_{\varphi|\psi;A}S_{\psi|\varphi;B}x\xi_{\varphi} = S_{\varphi|\psi;A}x^*\xi_{\psi} = x\xi_{\varphi}
\end{equation}
Hence, we have the identity:
\begin{equation}
    J_A\Delta^{-it+1/2}_{\varphi|\psi;A}\Delta^{it-1/2}_{\varphi|\psi;B}J_B = \Delta^{-it}_{\psi|\varphi;A}\Delta^{it}_{\psi|\varphi;B}
\end{equation}
Using this identity and the previous corollary, it is straight-forward to see:
\begin{corollary}
    Under the same condition as the theorem, we have strong convergence:
    \begin{equation}
        \lim_n\Delta^{-it + 1/2}_{\varphi_n|\psi_n;A}\Delta^{it-1/2}_{\varphi_n|\psi_n;B} = \Delta^{-it + 1/2}_{\varphi|\psi;A}\Delta^{it-1/2}_{\varphi|\psi;B}
    \end{equation}
\end{corollary}
%For any pairs of vectors $\xi,\eta\in\mathcal{H}$, we consider the functions:
%\begin{equation}
%    f^{(n)}_{\xi,\eta}(z) := \bra\eta, \Delta^{z}_{\varphi_n|\psi_n;A}\Delta^{-z}_{\varphi_n|\psi_n;B} - \Delta^z_{\varphi|\psi;A}\Delta^{-z}_{\varphi|\psi;B}\ket{\xi} 
%\end{equation}
%For each given $n$, the function is analytic inside the strip $\Re{z}\in(0,1/2)$
%Consider the relative modular operator $\Delta_{\phi|\psi}$. For the special case, $\phi=\psi$ this is simply the modular operator of $\psi$. The first observation is that a sequence of states $\phi_n\to \phi$ and $\psi_n\to \psi$ and any bounded continuous function $f$ we have
%\begin{eqnarray}
%    \lim_n f(\Delta_{\phi_n|\psi_n})=f(\Delta_{\phi|\psi})\ .
%\end{eqnarray}
%In particular this implies for all $\beta>0$ we have
%\begin{eqnarray}
%    &&\lim_n \frac{1}{\beta+\Delta_{\phi_n|\psi_n}} =\frac{1}{\beta+\Delta_{\phi|\psi}}
    %\nn\\
    %&&\lim_n \frac{\Delta_{\phi_n|\psi_n}}{\beta(\beta+\Delta_{\phi_n|\psi_n})}=\frac{1}{\beta}-\lim_n\frac{1}{\beta+\Delta_{\phi_n|\psi_n}} =\frac{\Delta_{\phi_n|\psi_n}}{\beta(\beta+\Delta_{\phi|\psi})}
%\end{eqnarray}
%in operator norm topology. Note that for the sequence $(\beta+\Delta_{\phi_n|\psi_n})^{-1}$ is uniformly bounded by $1/\beta$. 

%If $A_n\to A$ and $B_n\to B$ in operator norm topology and $A_n$ are uniformly bounded meaning that $\sup_n\|A_n\|<\infty$ then we have $\lim_n A_nB_n=A B$. The proof is straightforward:
%\begin{eqnarray}
%    \lim_n \|A_nB_n-A B\|&=&\lim_n\|A_n (B_n-B)\|+\lim_n \|(A_n-A)B\|\nn\\
%    & \leq &\sup_n \|A_n\|\lim_n\|B_n-B\|+\|B\|\lim_n\|A_n-A\|=0\nn\ .
%\end{eqnarray}
%Now, consider two algebras $\mA$ and $\mB$ (need not satisfy inclusion relation), then for all $\beta,\beta'>0$ we have
%\begin{eqnarray}\label{limitmultiply}
%    \lim_n \lb\frac{1}{\beta+\Delta_{\phi_n|\psi_n;\mA}}\rb\lb \frac{1}{\beta'+\Delta_{\phi_n|\psi_n;\mB}}\rb=\lb\frac{1}{\beta+\Delta_{\phi|\psi;\mA}}\rb\lb \frac{1}{\beta'+\Delta_{\phi|\psi;\mB}}\rb
%\end{eqnarray}
%Next, we use the following expansions for $\alpha\in (0,1)$ and positive operator $\Delta$:
%\begin{eqnarray}
%    &&\Delta^{-\alpha}=\frac{\sin(\pi\alpha)}{\pi}\int_0^\infty d\beta \:\beta^{-\alpha}\lb \frac{1}{\beta+\Delta}\rb\nn\\
%    &&\Delta^{\alpha}=\frac{\sin(\pi\alpha)}{\pi}\int_0^\infty d\beta \:\beta^{\alpha}\lb \frac{1}{\beta}-\frac{1}{\beta+\Delta}\rb
%\end{eqnarray}
%As a result, we obtain
%\begin{eqnarray}
%    \Delta_{\phi_n|\psi_n;\mA}^\alpha\Delta_{\phi|\psi;\mB}^{-\alpha}=\frac{\sin^2(\pi \alpha)}{\pi^2}\int_0^\infty d\beta d\beta' \lb\frac{\beta}{\beta'}\rb^\alpha \lb \frac{1}{\beta}-\frac{1}{\beta+\Delta_{\phi_n
%|\psi_n;\mA}}\rb\lb \frac{1}{\beta'+\Delta_{\phi_n|\psi_n;\mB}}\rb 
%\end{eqnarray}
%To commute the limit and integral, we can use the dominated convergence them.
%It follows from (\ref{limitmultiply}) that the integral above converges pointwise, but apriori we do not know the integral for all $n$ is dominated by a single integrable operator. If $\mA\subset \mB$ from the monotonicity condition we know that for all $z$ in the complex strip $\Re(z)\in (0,1/2)$ we have
%\begin{eqnarray}
%    |\Delta_{\phi|\psi;\mA}^z \Delta_{\phi|\psi;\mB}^{-z}|\leq 1\ .
%\end{eqnarray}
%Therefore, for inclusion of algebra, we can use the dominated convergence theorem to find
%\begin{eqnarray}
%    \lim_n \Delta^\alpha_{\phi_n|\psi_n;\mA}\Delta_{\phi_n|\psi_n;\mB}^{-\alpha}=\Delta^\alpha_{\phi|\psi;\mA}\Delta_{\phi|\psi;\mB}^{-\alpha}\ .
%\end{eqnarray}
%\NL{Does this justify continuity?}


% Wrong, don't use it \KL{We consider $f_{ab}(z) \equiv \bra{a}\Delta_A^{1/2+iz}\Delta_B^{-iz}\ket{b} = \bra{\Delta_A^{1/4}a} T_{B\subset A}(z-i/4)\ket{\Delta_B^{1/4} b} $ with some special choice of $a,b$. $f_{ab}(z)$ is holomorphic on the stripe Im$(z) \in (-\frac{1}{4}, \frac{1}{4})$ and continuous in the closure. The characteristic function has the boundary value that, for real $t$,
% \begin{align}
%     T_{B\subset A}(t-i/2) = J_AT_{B\subset A}(t)J_B.
% \end{align}
% By choosing $a\in A$ and $b\in B \subset A$ as positive operators, we have
% \begin{align}
%     J_A \ket{\Delta_A^{1/4}a} = \ket{\Delta_A^{1/4}a},\nonumber\\
%     J_B \ket{\Delta_B^{1/4}b} = \ket{\Delta_B^{1/4}b},
% \end{align}
% because both vectors are in the corresponding natural cones. Thus, the boundary value of $f_{ab}$ satisfies, for real $t$,
% \begin{align}
%     f_{ab}(t+i/4) = f_{ab}(t-i/4).
% \end{align}
% \textbf{Mistake, check the anti-unitary steps}
% The claim is that, with this choice of $a$ and $b$, we can extend $f_{ab}(z)$ to an entire function, i.e. a holomorphic function on the whole complex plane. In fact, $f_{ab}(z)$ is a constant function. 
% \begin{lemma}
%     If $a\in A$ and $b\in B \subset A$ are both positive operators, $f_{ab}(z)$ is a constant function.
% \end{lemma}
% To prove it, we need the following result.
% \begin{lemma}
%     If $f(z)$ is holomorphic on the stripe Im$(z) \in (0,1)$ and continuous in the closure with boundary values
%     \begin{align}
%         f(t+i) = f(t)
%     \end{align}
%     for real $t$, there exists an entire function that is an extension of $f(z)$ to the whole complex plane.
% \end{lemma}
% \begin{proof}
%     The function $\tilde{g}(z) = \overline{f(\overline{z})}$ is holomorphic on the stripe Im$(z) \in (-1,0)$ and continuous in the closure with boundary values
%     \begin{align}
%         \tilde{g}(t-i) = \tilde{g}(t) = \overline{f(t)}
%     \end{align}  
%     for real $t$. We define $g(z) = \tilde{g}(z-i)$ which shifts the domain to Im$(z) \in (0,1)$. Let $F_+(z) = \frac{f+g}{2}$ and $F_-(z) = \frac{f-g}{2i}$. $F_\pm(z)$ are both holomorphic on the stripe Im$(z) \in (0,1)$ and continuous in the closure with real boundary values. By the Schwarz reflection principle, we can extend $F_\pm(z)$ to Im$(z) \in (-1,1)$ with $F_\pm(z) = \overline{F_\pm(\overline{z})}$. Since the boundary values on two sides of the stripe are real, we can shift the stripe and apply the Schwarz reflection principle repeatedly to extend $F_\pm(z)$ to the whole complex plane. Thus, $f(z) = F_+(z) + i F_-(z)$ can be extended to the whole complex plane.
% \end{proof}
% If we have $|f(z)|\leq 1$ on the stripe Im$z \in [0,1]$, then $|g(z)|$ and $|F_\pm(z)|$ are also bounded by $1$ on the stripe. The extension through the Schwarz reflection principle does not change the bound of $F_\pm(z)$. A bounded entire function must be a constant function. Therefore, $F_\pm(z)$ are constant functions, and $f(z)$ is also a constant function. This finishes the proof, as $f_{ab}(z)$ is a bounded function. }

% \section{Titchmarsh theorem}\label{app:titchmarsh}

% Titchmarsh theorems are a class of Paley-Weiner type theorems that relate the behavior of complex-valued square-integrable functions $S(t)\in L^2(\mathbb{R})$ to the analytic domain of their Fourier transform
% \begin{equation}
%     \tilde{S}(\omega) = \int_{-\infty}^{\infty} S(t) e^{i\omega t} \, dt\ .
% \end{equation}
% Theorem \ref{titchmarshUHP} gives a condition that is equivalent to $\tilde{S}(\omega)$ being analytic on the upper-half-plane. Theorem \ref{titchmarshStrip} generalizes it to the condition equivalent to the analyticity of $\tilde{S}(\omega)$ in the strip $\Im(\omega)\in (-A,B)$.

% \subsection{Analyticity in frequency upper-half-plane:}
% \begin{theorem}[Titchmarsh]\label{titchmarshUHP}
% The following three statements are mutually equivalent:

% \begin{enumerate}
%     \item \textbf{Causality:} $S(t) = 0$ for almost every $t < 0$: $S(t)=\text{Sgn}(t)S(t)$.
%     \item \textbf{Analyticity:} The function $\tilde{S}(\omega)$ is the boundary limit of a function $G(z)$ analytic in the upper half-plane $\mathbb{C}^+$ (belonging to Hardy space $H^2$).
%     \item \textbf{Hilbert Transform Pair:}
%     The real and imaginary parts of the spectrum $\tilde{S}(\omega) = U(\omega) + i V(\omega)$ are Hilbert transforms of each other:
%     \begin{align}
%         V(\omega) &= -\mathcal{H}[U](\omega) \\
%         U(\omega) &= \mathcal{H}[V](\omega)
%     \end{align}
%     where the Hilbert transform is defined as:
%     \begin{equation}
%         \mathcal{H}[f](\omega) = \frac{1}{\pi} \mathcal{P} \int_{-\infty}^{\infty} \frac{f(\omega')}{\omega' - \omega} \, d\omega'
%     \end{equation}
% \end{enumerate}
% \end{theorem}

% \begin{proof}
% \textbf{1) $\implies$ 2) (Causality to Analyticity)} \\
% Assume $S(t) = 0$ for $t < 0$. We define the analytic extension $G(z)$ for $z = \omega + i\eta$ with $\eta > 0$:
% \begin{equation}
%     G(z) = \int_{0}^{\infty} S(t) e^{i(\omega + i\eta)t} \, dt = \int_{0}^{\infty} S(t) e^{-\eta t} e^{i\omega t} \, dt
% \end{equation}
% Because $S(t) \in L^2$ and $e^{-\eta t} \le 1$ for $t > 0$, the integral converges absolutely and defines a holomorphic function in the upper half-plane. By Plancherel's theorem, it satisfies the Hardy space condition. As $\eta \to 0^+$, $G(z)$ converges to $\tilde{S}(\omega)$ in the $L^2$ norm.

% \vspace{0.5cm}
% \textbf{2) $\implies$ 3) (Analyticity to Hilbert Transform)} \\
% By Cauchy's Integral Formula for $z \in \mathbb{C}^+$:
% \begin{equation}
%     G(z) = \frac{1}{2\pi i} \int_{-\infty}^{\infty} \frac{\tilde{S}(\omega')}{\omega' - z} \, d\omega'
% \end{equation}
% Taking the limit $z \to \omega + i0^+$ using the Plemelj-Sokhotski relations:
% \begin{equation}
%     \tilde{S}(\omega) = \frac{1}{2}\tilde{S}(\omega) + \frac{1}{2\pi i} \mathcal{P} \int_{-\infty}^{\infty} \frac{\tilde{S}(\omega')}{\omega' - \omega} \, d\omega'
% \end{equation}
% Rearranging to isolate $\tilde{S}(\omega)$:
% \begin{equation}\label{contourint}
%     \tilde{S}(\omega) = \frac{1}{\pi i} \mathcal{P} \int_{-\infty}^{\infty} \frac{\tilde{S}(\omega')}{\omega' - \omega} \, d\omega' = -i \mathcal{H}[\tilde{S}](\omega)
% \end{equation}
% Substituting $\tilde{S} = U + iV$:
% \begin{equation}
%     U + iV = -i (\mathcal{H}[U] + i\mathcal{H}[V]) = \mathcal{H}[V] - i\mathcal{H}[U]
% \end{equation}
% Comparing real and imaginary parts yields the relations: $U = \mathcal{H}[V]$ and $V = -\mathcal{H}[U]$.

% \vspace{0.5cm}
% \textbf{3) $\implies$ 1) (Hilbert Transform to Causality)} \\
% From the derivation in step 2, condition (3) implies that the full complex spectrum satisfies:
% \begin{equation}
%     \tilde{S}(\omega) = -i \mathcal{H}[\tilde{S}](\omega)
% \end{equation}
% We apply the inverse Fourier transform, $\mathcal{F}^{-1}$, to both sides. We utilize the property that the Hilbert transform in the frequency domain corresponds to multiplication by the signum function in the time domain. Specifically, for the transform definition used here ($\int e^{i\omega t}$):
% \begin{equation}
%     \mathcal{F}^{-1} \big\{ -i \mathcal{H}[\tilde{S}] \big\}(t) = \text{sgn}(t) S(t)
% \end{equation}
% Substituting this into the equation:
% \begin{equation}
%     S(t) = \text{sgn}(t) S(t)
% \end{equation}
% which implies $S(t) = 0$ for almost all $t < 0$. 
% \end{proof}
% Note that the Hilbert transform can be viewed as the convolution of the function with $f(\omega)=-\frac{1}{\pi\omega}$:
% \begin{eqnarray}
%     U(\omega)=(f*V)(\omega)\ .
% \end{eqnarray}
% Since $\mathcal{F}^{-1}(\frac{-1}{\pi \omega})=\frac{i}{2\pi}\text{sgn}(t)$ we find
% \begin{eqnarray}
%     \tilde{U}(t)=i\text{sgn}(t) \tilde{V}(t)\ .
% \end{eqnarray}
% However, note that both $\tilde{U}(t)$ and $\tilde{V}(t)$ are complex, as they are defined as the inverse Fourier transform of $U(\omega)$ and $V(\omega)$. 

% \begin{corollary}
%     If $S(t)$ that satisfies the conditions of Theorem \ref{titchmarshUHP} is a real function on $\mathbb{R}$, then we have
%     \begin{eqnarray}
%         &&U(\omega)=\frac{2}{\pi}\mathcal{P}\int_0^\infty \frac{\omega' V(\omega')}{\omega^2-(\omega')^2}\nn\\
%         &&V(\omega)=\frac{-2\omega}{\pi}\mathcal{P}\int_0^\infty \frac{U(\omega')}{\omega^2-(\omega')^2}
%     \end{eqnarray}
% \end{corollary}
% \begin{proof}
%     Since $S(t)$ is real for all $t\in\mathbb{R}$, then $S(-\omega)=S(\omega)^*$ which implies that $U(\omega)$ and $V(\omega)$ are even and odd functions of $\omega$, respectively. Then, the corollary follows from multiplying the numerator and denominator of the integrand inside the Contour integral in (\ref{contourint}) by $\omega+\omega'$.
% \end{proof}



% \subsection{Analyticity in the frequency strip:}

% We would like to generalize the Titchmarsh theorem to functions that are analytic in a complex frequency strip. The key intuition is the following: If a complex function of a real parameter $t$ decays exponentially fast in time its Fourier transform is analytic on a strip, and decays exponentially fast in real frequency space. More precisely, we have the following generalized Titchmarsh theorem


% \begin{theorem}[Generalized Titchmarsh for a Strip]\label{titchmarshStrip}
% Let $S(t)$ be a complex-valued function and $\tilde{S}(\omega)$ its Fourier transform. Let $A \ge 0$ and $B \ge 0$ (with $A+B > 0$) be real constants defining a horizontal strip $\mathcal{D}$ in the complex frequency plane:
% \begin{equation}
%     \mathcal{D} = \{ \omega \in \mathbb{C} \mid -A < \text{Im}(\omega) < B \}
% \end{equation}

% The following three statements are mutually equivalent:

% \begin{enumerate}
%     \item \textbf{Two-sided exponential decay:}
%     The function $S(t)$ satisfies the asymptotic decay conditions:
%     \begin{align}
%         S(t) &\sim e^{-B|t|} \quad (t \to -\infty) \\
%         S(t) &\sim e^{-At} \quad (t \to +\infty)
%     \end{align}
%     In other words, for any $\epsilon > 0$, the weighted norms are finite:
%     \begin{equation}
%         e^{(B-\epsilon)t} S(t) \in L^2(-\infty, 0] \quad \text{and} \quad e^{-(A-\epsilon)t} S(t) \in L^2[0, \infty)\ .
%     \end{equation}

%     \item \textbf{Strip Analyticity:}
%     The function $\tilde{S}(\omega)$ is the boundary limit of a function $G(z)$ holomorphic inside the strip $-A < \text{Im}(z) < B$, satisfying the uniform Hardy condition:
%     \begin{equation}
%         \sup_{-A < y < B} \int_{-\infty}^{\infty} |G(x + iy)|^2 \, dx < \infty
%     \end{equation}

%     \item \textbf{Spectral Decomposition and shifted dispersion Relations:}
%     The spectrum $\tilde{S}(\omega)$ uniquely decomposes into:
%     \begin{equation}
%         \tilde{S}(\omega) = \tilde{S}_+(\omega) + \tilde{S}_-(\omega)
%     \end{equation}
%     where $\tilde{S}_+(\omega)$ is analytic for $\text{Im}(\omega) > -A$ and $\tilde{S}_-(\omega)$ is analytic for $\text{Im}(\omega) < B$. \\
    
%     Crucially, the real and imaginary parts of these components satisfy \textbf{Shifted Hilbert Transform} relations on the boundary lines:
    
%     \begin{itemize}
%         \item \textbf{For $\tilde{S}_+$ (on the line $\omega = \nu - iA$):} \\
%         Let $\tilde{S}_+(\nu - iA) = U_+(\nu) + i V_+(\nu)$. Then:
%         \begin{equation}
%             U_+(\nu) = \mathcal{H}[V_+](\nu), \quad V_+(\nu) = -\mathcal{H}[U_+](\nu)
%         \end{equation}
        
%         \item \textbf{For $\tilde{S}_-$ (on the line $\omega = \nu + iB$):} \\
%         Let $\tilde{S}_-(\nu + iB) = U_-(\nu) + i V_-(\nu)$. Then:
%         \begin{equation}
%             U_-(\nu) = -\mathcal{H}[V_-](\nu), \quad V_-(\nu) = \mathcal{H}[U_-](\nu)
%         \end{equation}
%     \end{itemize}
% \end{enumerate}
% \end{theorem}

% \begin{proof}
% \textbf{1) $\implies$ 2) (Decay to Analyticity)} \\
% We define the complex integral $G(z) = \int_{-\infty}^{\infty} S(t) e^{i z t} \, dt$ for $z = x + iy$.
% Splitting the integral:
% \begin{equation}
%     G(z) = \int_{-\infty}^{0} S(t) e^{(B-y)t} e^{-(B-i x)t} \, dt + \int_{0}^{\infty} S(t) e^{-(A+y)t} e^{(A+ix)t} \, dt
% \end{equation}
% Convergence requires $y < B$ (for the first term) and $y > -A$ (for the second term). Thus, $G(z)$ is analytic in $-A < y < B$. The Hardy condition follows from Plancherel's theorem applied to $S(t)e^{-yt}$.

% \vspace{0.5cm}
% \textbf{2) $\implies$ 3) (Analyticity to Decomposition)} \\
% Using Cauchy's Integral Formula on a rectangular contour inside the strip and sending the horizontal width to infinity, we obtain:
% \begin{equation}
%     \tilde{S}(z) = \frac{1}{2\pi i} \int_{-\infty - iA}^{\infty - iA} \frac{\tilde{S}(\zeta)}{\zeta - z} \, d\zeta - \frac{1}{2\pi i} \int_{-\infty + iB}^{\infty + iB} \frac{\tilde{S}(\zeta)}{\zeta - z} \, d\zeta
% \end{equation}
% We define the first integral as $\tilde{S}_+(z)$ and the second as $\tilde{S}_-(z)$. The analyticity domains and shifted Hilbert relations follow directly from the standard Titchmarsh theorem applied to the shifted functions $\tilde{S}_+(z-iA)$ and $\tilde{S}_-(z+iB)$.

% \vspace{0.5cm}
% \textbf{3) $\implies$ 1) (Decomposition to Decay)} \\
% Inverse transforming $\tilde{S}_+$ yields a function supported on $t>0$ with weight $e^{-At}$. Inverse transforming $\tilde{S}_-$ yields a function supported on $t<0$ with weight $e^{Bt}$. The sum satisfies the two-sided decay.
% \end{proof}

% \subsection*{Discussion of Limits}

% \paragraph{1. The Causal Limit ($A \to 0, B \to \infty$)}
% The strip becomes the Upper Half-Plane. The condition $e^{Bt}S(t) \in L^2$ as $B \to \infty$ forces $S(t) = 0$ for $t<0$. This recovers the standard Kramers-Kronig relations on the real axis.

% \paragraph{2. The Anti-Causal Limit ($A \to \infty, B \to 0$)}
% The strip becomes the Lower Half-Plane. The condition $e^{-At}S(t) \in L^2$ as $A \to \infty$ forces $S(t) = 0$ for $t>0$.

% \paragraph{3. The Degenerate Limit ($A \to 0, B \to 0$)}
% The strip collapses to the real axis. The region of analyticity vanishes. The theorem reduces to the Plancherel Theorem ($S \in L^2 \iff \tilde{S} \in L^2$). The real and imaginary parts of the spectrum become independent, and no dispersion relations exist.
% %\section{Coherent Perturbations of GFF}



\section{Approximation of Modular Scattering Matrix}\label{app:modularSmatrix}

The results on the approximations of characteristic functions in the previous subsection provide the necessary technical tools for us to prove corresponding results on the approximations of modular scattering matrices. Recall the following definition of modular scattering matrix:
\begin{definition}
    Given a pair of von Neumann algebras $\mathcal{B}\subset\mathcal{A}$ and a common cyclic and separating vector $\ket{\Omega}\in\mathcal{H}$, the modular scattering matrix is defined as:
    \begin{equation}
        \mathcal{M}_{\mathcal{A},\mathcal{B}}(t):=\Delta_\mathcal{A}^{-it}\Delta_\mathcal{B}^{2it}\Delta_\mathcal{A}^{-it}
    \end{equation}
\end{definition}
Using the same set-up as Theorem \ref{theorem:mainresult}, we can state the main approximation results on the modular scattering matrices:
\begin{theorem}\label{theorem:approxscattering}
    Consider a net of local algebras $\{\mathcal{A}(\mathcal{O})\}_{\mathcal{O}\in\mathbb{R}^{d,1}}$ that satisfies the Haag-Kastler axioms. Let $W_0$ be the right Rindler wedge and let $W_r$ be a translated right Rindler wedge shifted to the right by $r > 0$ along the $x_1$ direction. Consider a one-parameter family of shifted double cones inside the shifted wedge $\{ B^r_{e^{2\pi s}R}\subset W_r\}_{s > 0}$. 

    In addition, consider the following modular operators associated with the local algebras:
    \begin{enumerate}
        \item Let $\Delta_{W_0}$ be the modular operator of $\mathcal{A}(W_0)$;
        \item Let $\Delta_{W_r}$ be the modular operator of $\mathcal{A}(W_r)$;
        \item Let $\Delta_{B^r_{e^{2\pi s}R}}$ be the modular operator of $\mathcal{A}(B^r_{e^{2\pi s}R})$
    \end{enumerate}
    Finally consider the spectral decomposition of the modular flow of the right Rindler wedge:
    \begin{equation}
        \Delta_{W_0}^{it} = \int_0^{W_0}e^{i\mu t}dE(\mu)
    \end{equation}
    Then for any spectral cut-off $\mu_0 > 0$, for any Schwartz function $f\in \mathcal{D}(\mathbb{R})$, for any operator $A\in\mathcal{A}(B^r_R)$ and for any constant $s > 0$, there exists a constant $b_f(\mu_0, s) > 0$ such that:
    \begin{align}
        \begin{split}
            ||&E(\mu_0)\lb\int_\mathbb{R} dt f(t)\lb \mathcal{M}_{W_0, B^r_{e^{2\pi s}R}}(t) - \mathcal{M}_{W_0, W_r}(t) \rb E(\mu_0)A\ket{\Omega}\rb||^2 \\&\leq b_f(\mu_0, s)\lb||A\ket{\Omega}||^2 + ||A^*\ket{\Omega}||^2\rb
        \end{split}
    \end{align}
    In addition, for any spectral cut-off, we have a well-defined limit:
    \begin{equation}
        \lim_{s\rightarrow\infty}||E(\mu_0)\lb\int_\mathbb{R} dt f(t)\lb \mathcal{M}_{W_0, B^r_{e^{2\pi s}R}}(t) - \mathcal{M}_{W_0, W_r}(t) \rb E(\mu_0)A\ket{\Omega}\rb||^2 = 0
    \end{equation}
\end{theorem}
Before presenting the proof, the physical intuition behind this theorem is simple. For operators localized to a fixed
\begin{proof}
    The proof is similar to the proof of Theorem \ref{theorem:mainresult}. Using the spectral decomposition we can rewrite the norm as follows:
    \begin{align}
        \begin{split}
            ||&E(\mu_0)\lb\int_\mathbb{R} dt f(t)\lb\mathcal{M}_{W_0, B^r_{e^{2\pi s}R}}(t) - \mathcal{M}_{W_0, W_r}(t)\rb E(\mu_0)A\ket{\Omega}\rb||
            \\&
            =||\int_\mathbb{R} dt f(t)\int_0^{\mu_0} dE(\mu) e^{-i\mu t}\lb\Delta^{2it}_{B^r_{e^{2\pi s}R}} - \Delta_{W_r}^{2it}\rb \int_0^{\mu_0}dE(\nu) e^{-i\nu t} A\ket{\Omega}||
            \\&
            =\sup_{\substack{B\in\mathcal{A}(W_0)\\s.t.||B\ket{\Omega}||\leq 1}}|\langle B\Omega\text{ , }\int_\mathbb{R} dt f(t)\int_0^{\mu_0} dE(\mu)e^{-i\mu t}\lb\Delta_{B^r_{e^{2\pi s}R}}^{2it} - \Delta_{W_r}^{2it}\rb\int_0^{\mu_0}dE(\nu)e^{-i\nu t}A\Omega\rangle|
            \\&
            =\sup_{\substack{B\in\mathcal{A}(W_0)\\s.t.||B\ket{\Omega}||\leq 1}}|\int_\mathbb{R} dt f(t)\int_0^{\mu_0}e^{-i\mu t}\int_0^{\mu_0}e^{-i\nu t}d_{\mu\nu}\langle E(\mu)B\Omega\text{ , }\lb\Delta^{2it}_{B^r_{e^{2\pi s}R}} - \Delta^{2it}_{W_r}\rb E(\nu)A\Omega\rangle|
        \end{split}
    \end{align}
    To simplify notation, we denote the $(t,\mu,\nu)$-dependent measurable function as:
    \begin{equation}
        g_t(\mu,\nu):=\langle E(\mu)B\Omega\text{ , }\lb\Delta^{2it}_{B^r_{e^{2\pi s}R}} - \Delta^{2it}_{W_r}\rb E(\nu)A\Omega \rangle
    \end{equation}
    Notice that $g_t(0,\nu) = g_t(\mu,0) = 0$ because $E(0) = 0$ is the zero projection. To apply integration by parts, we have the following formula:
    \begin{align}
        \begin{split}
            e^{-i\mu t}&e^{-i\nu t}d_{\mu\nu}g_t(\mu,\nu) = e^{-i\mu t}e^{-i\nu t}d_{\mu\nu}\langle E(\mu)B\Omega\text{ , }\lb\Delta^{2it}_{B^r_{e^{2\pi s}R}} - \Delta^{2it}_{W_r}\rb E(\nu)A\Omega\rangle
            \\&
            = d_{\mu\nu}\lb e^{-i\mu t}e^{-i\nu t}g_t(\mu,\nu) \rb + ite^{-i\mu t}e^{-i\nu t} d_\nu g_t(\mu,\nu) d\mu+ ite^{-i\mu t}e^{-i\nu t} d_\mu g_t(\mu,\nu)d\nu \\&+ t^2 e^{-i\mu t}e^{-i\nu t} g_t(\mu,\nu)d\mu d\nu
        \end{split}
    \end{align}
    Then by integration by parts, we have:
    \begin{align}
        \begin{split}
            ||&E(\mu_0)\lb\int_\mathbb{R} dt f(t)\lb\mathcal{M}_{W_0,B^r_{e^{2\pi s}R}}(t) - \mathcal{M}_{W_0, W_r}(t)\rb E(\mu_0) A\ket{\Omega}\rb||
            \\&
            = \sup_{\substack{B\in\mathcal{A}(W_0)\\s.t.||B\ket{\Omega}||\leq 1}}|\int_\mathbb{R} dt f(t)\int_0^{\mu_0}\int_0^{\mu_0}d_{\mu\nu}\lb e^{-i\mu t}e^{-i\nu t}g_t(\mu,\nu)\rb 
            \\&
            + i\int_\mathbb{R} dt \lb tf(t)\rb\int_0^{\mu_0}\int_0^{\mu_0}\lb e^{-i\mu t}d_\nu \lb e^{-i\nu t}g_t(\mu,\nu) \rb d\mu\rb
            \\&
            +i\int_\mathbb{R} dt \lb tf(t)\rb\int_0^{\mu_0}\int_0^{\mu_0}\lb e^{-i\nu t} d_\mu\lb e^{-i\mu t} g_t(\mu,\nu)\rb d\nu\rb 
            \\&
            - \int_\mathbb{R} dt \lb t^2f(t) \rb\int_0^{\mu_0}\int_0^{\mu_0} e^{-i\mu t}e^{-i\mu t}g_t(\mu,\nu)d\mu d\nu|
        \end{split}
    \end{align}
    We estimate each term separately. For the first term, we have the following direct calculation:
    {\small
\begin{align}
        \begin{split}
            &\sup_{\substack{B\in\mathcal{A}(W_0)\\s.t.||B\ket{\Omega}||\leq 1}}|\int_\mathbb{R} dt f(t)\int_0^{\mu_0}\int_0^{\mu_0} d_{\mu\nu}\lb e^{-i\mu t}e^{-i\nu t}g_t(\mu,\nu) \rb| = \sup_{\substack{B\in\mathcal{A}(W_0)\\s.t.||B\ket{\Omega}||\leq 1}}|\int_\mathbb{R} dt f(t) e^{-2i\mu_0 t}g_t(\mu_0, \mu_0)|
            \\&
            =\sup_{\substack{B\in\mathcal{A}(W_0)\\s.t.||B\ket{\Omega}||\leq 1}}|\frac{1}{2}\int_\mathbb{R} dt f(\frac{t}{2}) e^{-i\mu_0 t}\langle E(\mu_0)B\Omega\text{ , }\lb\Delta_{B^r_{e^{2\pi s}R}}^{it} - \Delta_{W_r}^{it}\rb E(\mu_0)A\Omega\rangle|\\&
            \leq \frac{1}{2}c_{f_{\mu_0}(t/2)}^{1/2}(s)\lb||A\ket{\Omega}||^2 + ||A^*\ket{\Omega}||^2\rb^{1/2}
        \end{split}
    \end{align}
}
    where the last inequality follows from Lemma \ref{lemma:extensionfredenhagen} and $c_{f_{\mu_0}(t/2)}$ is Fredenhagen's constant associated with the function $e^{-i\mu_0 t}f_(\frac{t}{2})$.

    For the second term, we have the following estimate:
    {\small
\begin{align}
        \begin{split}
            &\sup_{\substack{B\in\mathcal{A}(W_0)\\s.t.||B\ket{\Omega}||\leq 1}}|\int_\mathbb{R} dt \lb tf(t) \rb\int_0^{\mu_0}\int_0^{\mu_0}\lb e^{-i\mu t}d_\nu\lb e^{-i\nu t}g_t(\mu,\nu) \rb d\mu \rb|
            \\&
            =\sup_{\substack{B\in\mathcal{A}(W_0)\\s.t.||B\ket{\Omega}||\leq 1}}|\int_\mathbb{R}dt \lb tf(t)\rb\int_0^{\mu_0}e^{-i\mu t} e^{-i\mu_0 t}g_t(\mu,\mu_0) d\mu|
            \\&
            \leq \sup_{\substack{B\in\mathcal{A}(W_0)\\s.t.||B\ket{\Omega}||\leq 1}}\sup_{0\leq\mu\leq\mu_0}\mu_0|\int_\mathbb{R} dt\lb tf(t) \rb e^{-i(\mu + \mu_0) t}\langle E(\mu)B\Omega\text{ , }\lb\Delta^{2it}_{B^r_{e^{2\pi s}R}} - \Delta^{2it}_{W_r}\rb E(\mu_0)A\Omega \rangle|
            \\&
            =\sup_{\substack{B\in\mathcal{A}(W_0)\\s.t.||B\ket{\Omega}||\leq 1}}\sup_{0\leq \mu\leq\mu_0}\mu_0|\frac{1}{4}\int_\mathbb{R} dt \lb tf(\frac{t}{2}) \rb e^{-i\frac{\mu + \mu_0}{2}t}\langle E(\mu)B\Omega\text{ , }\lb\Delta^{it}_{B^r_{e^{2\pi s}R}} - \Delta^{it}_{W_r}\rb E(\mu_0)A\Omega \rangle|
            \\&
            \leq \sup_{0\leq \mu\leq\mu_0}\frac{\mu_0}{4} c^{1/2}_{tf_{\frac{\mu + \mu_0}{2}}(\frac{t}{2})}(s)\lb||A\ket{\Omega}||^2 + ||A^*\ket{\Omega}||^2\rb^{1/2}
        \end{split}
    \end{align}
}
    where the last inequality follows from Lemma \ref{lemma:extensionfredenhagen} and $c_{tf_{\frac{\mu + \mu_0}{2}}(\frac{t}{2})}$ is Fredenhagen's constant associated with the function $e^{-i\frac{\mu + \mu_0}{2}t}tf(\frac{t}{2})$.

    For the third term, we have the same estimate:
    \begin{align}
        \begin{split}
            &\sup_{\substack{B\in\mathcal{A}(W_0)\\s.t.||B\ket{\Omega}||\leq 1}}|\int_\mathbb{R} dt \lb tf(t) \rb\int_0^{\mu_0}\int_0^{\mu_0}\lb e^{-i\nu t}d_\mu\lb e^{-i\mu t}g_t(\mu,\nu) \rb d\nu \rb|\\&
            \leq\sup_{0\leq \mu\leq\mu_0}\frac{\mu_0}{4}c^{1/2}_{tf_{\frac{\mu + \mu_0}{2}}(\frac{t}{2})}(s)\lb||A\ket{\Omega}||^2 + ||A^*\ket{\Omega}||^2\rb^{1/2}
        \end{split}
    \end{align}

    For the last term, we directly apply Lemma \ref{lemma:extensionfredenhagen}:
    \begin{align}
        \begin{split}
            &\sup_{\substack{B\in\mathcal{A}(W_0)\\s.t.||B\ket{\Omega}||\leq 1}}|\int_\mathbb{R} dt \lb t^2f(t) \rb\int_0^{\mu_0}\int_0^{\mu_0}e^{-i(\mu + \nu)t}g_t(\mu,\nu)d\mu d\nu|
            \\&
            \leq \frac{1}{8}c_{t^2f_{\frac{\mu + \nu}{2}}(\frac{t}{2})}^{1/2}(s)\lb||A\ket{\Omega}||^2 + ||A^*\ket{\Omega}||^2\rb^{1/2}
        \end{split}
    \end{align}
    Combining these four estimates, we have:
    {\small
\begin{align}
        \begin{split}
            ||&E(\mu_0)\lb\int_\mathbb{R} dt f(t)\lb \mathcal{M}_{W_0, B^r_{e^{2\pi s}R}}(t) - \mathcal{M}_{W_0, W_r}(t) \rb E(\mu_0)A\ket{\Omega}\rb||
            \\&
            \leq \lb \frac{1}{2}c^{1/2}_{f_{\mu_0}(\frac{t}{2})}(s) + \sup_{0\leq\mu\leq\mu_0}\frac{\mu_0}{2} c^{1/2}_{tf_{\frac{\mu + \mu_0}{2}}(\frac{t}{2})}(s) + \frac{1}{8}c^{1/2}_{t^2f_{\frac{\mu + \nu}{2}}(\frac{t}{2})}(s)\rb\cdot\lb||A\ket{\Omega}||^2 + ||A^*\ket{\Omega}||^2\rb^{1/2}
        \end{split}
    \end{align}
}
    By Lemma \ref{lemma:shiftedfredenhagenconstant}, we have the limit:
    \begin{equation}
        \lim_{s\rightarrow\infty}\frac{1}{2}c^{1/2}_{f_{\mu_0}(\frac{t}{2})}(s) + \sup_{0\leq\mu\leq\mu_0}\frac{\mu_0}{2} c^{1/2}_{tf_{\frac{\mu + \mu_0}{2}}(\frac{t}{2})}(s) + \frac{1}{8}c^{1/2}_{t^2f_{\frac{\mu + \nu}{2}}(\frac{t}{2})}(s) = 0\qedhere
    \end{equation}
\end{proof}

\subsection{Calculation of Fredenhagen's Constant}\label{subsection:gaussiansmear}
In this subsection, we collect several results (including explicit calculations) of Fredenhagen's constant \cite{fredenhagen}. Recall the Fredenhagen's original estimate depends on the constant:
\begin{equation}
    c_f(s):=\inf_{\substack{g\in\mathcal{D}(\mathbb{R})\\ c > 0}}\{2||f - g||_1 + 2\int_{|u|\geq c}du|\varphi_g(u)| + \frac{2}{\pi}\tanh^{-1}(e^{-\pi s})e^{c/2}||\varphi_g||_1\}
\end{equation}
where $g\in\mathcal{D}(\mathbb{R})$ is a Schwartz function and $\varphi_g(u)$ is  a twisted Fourier transform of $g$:
\begin{equation}
    \varphi_g(u):=\frac{1}{2\pi}\int_\mathbb{R} dt g(t)\lb\frac{2}{i}\sinh(\pi t)\rb e^{itu}
\end{equation}
We first make the following observation on Fredenhagen's constant for shifted function $f_\mu(t) = e^{-i\mu t}f(t)$:
\begin{lemma}\label{lemma:shiftedfredenhagenconstant}
    For a Schwartz function $f\in \mathcal{D}(\mathbb{R})$, the following limits hold:
    \begin{equation}
        \lim_{s\rightarrow\infty}c_{f_\mu}(s) = 0
    \end{equation}
    \begin{equation}
        \lim_{s\rightarrow\infty}\sup_{-\mu_0\leq \mu\leq\mu_0}c_{tf_\mu}(s) = 0
    \end{equation}
\end{lemma}
\begin{proof}
    The proof is based on an explicit calculation of $c_{f_\mu}(s)$. Using the definition of $c_f(s)$, we have the following:
    \begin{align}
        \begin{split}
            c_{f_\mu}(s) &= \inf_{\substack{g\in\mathcal{D}(\mathbb{R})\\ c>0}}\{2||f_\mu - g|| + 2\int_{|u|\geq c}du|\varphi_g(u)| + \frac{2}{\pi}\tanh^{-1}(e^{-\pi s})e^{c/2}||\varphi_g||_1\}
            \\&
            =\inf_{\substack{g_\mu\in\mathcal{D}(\mathbb{R})\\ c > 0}}\{2||f_\mu - g_\mu||_1 + 2\int_{|u|\geq c}du|\varphi_{g_\mu}(u)| + \frac{2}{\pi}\tanh^{-1}(e^{-\pi s})e^{c/2}||\varphi_{g_\mu}||_1\}
            \\&
            =\inf_{\substack{g_\mu\in\mathcal{D}(\mathbb{R})\\ c > 0}}\{2||f-g||_1 + 2\int_{|u| \geq c}du|\varphi_g(u + \mu)| + \frac{2}{\pi}\tanh^{-1}(e^{-\pi s})e^{c/2}||\varphi_g||_1\}
            \\&
            =\inf_{\substack{g\in\mathcal{D}(\mathbb{R})\\ c > 0}}\{2||f-g||_1 + 2\int_{|u - \mu|\geq c}du|\varphi_g(u)| + \frac{2}{\pi}\tanh^{-1}(e^{-\pi s})e^{c/2}||\varphi_g||_1\}
        \end{split}
    \end{align}
    where $g_\mu(t) := e^{-i\mu t}g(t)\in\mathcal{D}(\mathbb{R})$. In addition, by definition of $\varphi_g$, we have:
    {\small
\begin{align}
        \begin{split}
            \varphi_{g_\mu}(u)&=\frac{1}{2\pi}\int_\mathbb{R} dt g_\mu(t)\lb\frac{2}{i}\sinh(\pi t)\rb e^{-itu} = \frac{1}{2\pi}\int_\mathbb{R} dt g(t)\lb\frac{2}{i}\sinh(\pi t)\rb e^{-i(\mu + u)t} = \varphi_g(u + \mu)
        \end{split}
    \end{align}
}
    In particular by translation invariance of the Lebesgue measure, we have $||\varphi_g||_1 = ||\varphi_{g_\mu}||_1$.

    Since the function $f$ is a Schwartz function, we have the following upper bound by taking $g = f$:
    \begin{align}
        \begin{split}
            c_{f_\mu}(s) &\leq \inf_{c > 0}\{2\int_{|u - \mu| > c}du|\varphi_f(u)| + \frac{2}{\pi}\tanh^{-1}(e^{-\pi s})e^{c/2}||\varphi_f||_1\}
        \end{split}
    \end{align}
    In addition, because $f$ is a Schwartz function, the twisted Fourier transform $\varphi_f$ is integrable. Hence, for any small constant $\epsilon > 0$, there exists $c_\epsilon > 0$ such that:
    \begin{equation}
        2\int_{|u - \mu| > c_\epsilon}du|\varphi_f(u)| < \epsilon
    \end{equation}

    Hence for any $\epsilon > 0$, we have:
    \begin{align}
        \begin{split}
            \lim_{s\rightarrow\infty}c_{f_\mu}(s) &\leq \lim_{s\rightarrow\infty}2\int_{|u-\mu|>c_\epsilon}du|\varphi_f(u)| + \frac{2}{\pi}\tanh^{-1}(e^{-\pi s})e^{c_\epsilon/2}||\varphi_f||_1
            \\&
            <\lim_{s\rightarrow\infty}2\epsilon + \frac{2}{\pi}\tanh^{-1}(e^{-\pi s})e^{c_\epsilon/2}||\varphi_f||_1 = 2\epsilon
        \end{split}
    \end{align}
    Since the bound holds for any $\epsilon > 0$, we have $\lim_{s\rightarrow\infty}c_{f_\mu}(s) = 0$.

    To calculate the limit $\lim_{s\rightarrow\infty}c_{tf_\mu}(s)$, first notice that because $f$ is a Schwartz function, $tf_\mu(t) = tf(t)e^{-i\mu t}$ is also a Schwartz function. Using the same arguments as above, we have the following upper bound:
    \begin{equation}
        \sup_{-\mu_0\leq \mu\leq\mu_0}c_{tf_\mu}(s) \leq \sup_{-\mu_0\leq \mu\leq\mu_0}\inf_{c > 0}\{2\int_{|u - \mu| > c}du|\varphi_{tf}(u)|+ \frac{2}{\pi}\tanh^{-1}(e^{-\pi s})e^{c/2}||\varphi_{tf}||_1\}
    \end{equation}
    Consider $\int_{|u-\mu| > c}du|\varphi_{tf}(u)|$ as a family of functions $\{F_\mu(c)\}_{-\mu_0\leq\mu\leq\mu_0}$. For any $\mu \in [-\mu_0, \mu_0]$, the function $F_\mu$ is Lipschitz with the same constant:
    \begin{align}
        \begin{split}
            ||\frac{d}{dc}F_\mu(c)||_\infty &= ||\frac{d}{dc}\int_{|u| > c}|\varphi_{tf}(u + \mu)|du||_\infty \leq 2||\varphi_{tf}||_\infty
        \end{split}
    \end{align}
    Hence the family of functions $\{F_\mu\}_{0\leq \mu\leq \mu_0}$ is equicontinuous. In particular, since for any $\mu\in [-\mu_0, \mu_0]$ the limit $\lim_{c\rightarrow \infty}F_\mu(c) = 0$ (by the integrability of $\varphi_{tf}$), then for any $\epsilon > 0$ there exists $c_\epsilon > 0$ such that for all $\mu \in [-\mu_0, \mu_0]$, we have $|F_\mu(c_\epsilon)| < \epsilon$. Thus for any $\epsilon > 0$, we have:
    \begin{align}
        \begin{split}
            &\sup_{-\mu_0\leq \mu\leq \mu_0}c_{tf_\mu}(s) \leq \sup_{-\mu_0\leq \mu\leq\mu_0}\{2\int_{|u - \mu| > c_\epsilon} du |\varphi_{tf}(u)| + \frac{2}{\pi}\tanh^{-1}(e^{-\pi s})e^{c_\epsilon / 2}||\varphi_{tf}||_1\}
            \\&
            < 2\epsilon + \frac{2}{\pi}\tanh^{-1}(e^{-\pi s})e^{c_\epsilon / 2}||\varphi_{tf}||_1
        \end{split}    
    \end{align}
    Taking the limit $s\rightarrow\infty$, for any $\epsilon > 0$ we have the limit:
    \begin{align}
        \begin{split}
            \lim_{s\rightarrow\infty}\sup_{-\mu_0\leq \mu\leq \mu_0}c_{tf_\mu}(s) \leq \lim_{s\rightarrow\infty}2\epsilon + \frac{2}{\pi}\tanh^{-1}(e^{-\pi s})e^{c_\epsilon / 2}||\varphi_{tf}||_1 = 2\epsilon
        \end{split}
    \end{align}
    Since the bound holds for any $\epsilon > 0$, we have the limit $\lim_{s\rightarrow\infty}\sup_{0\leq \mu\leq \mu_0}c_{tf_\mu}(s) = 0$.
\end{proof}
In the remainder of this section, we provide explicit upper bound on Fredenhagen's constant when the smearing function is Gaussian. The purpose of this calculation is to confirm several physical intuitions.
\begin{lemma}
    For Gaussian $f_{\sigma,\mu}(t) := \frac{1}{\sigma\sqrt{2\pi}}e^{-\frac{(t - \mu)^2}{2\sigma^2}}$, there exists $s_0 > 0$ and positive constants $K_3, K_4, K_5 > 0$ such that for $s > s_0$, to the first order in $s$ the following bound holds:
    \begin{equation}
        c_{f_{\sigma,\mu}}(s) \leq K_3e^{-\frac{K_4}{2}s^2} + K_5e^{-K_4 s^2}
    \end{equation}
    These positive constants depend on $\sigma,\mu$.
\end{lemma}
\begin{proof}
    Since Gaussian function is a Schwartz function, we have the following bound:
    \begin{equation}\label{equation:bound}
        c_{f_{\sigma,\mu}}(s) \leq \inf_{c > 0}\{2\int_{|u|\geq c}du |\varphi_{f_{\sigma,\mu}}(u)| + \frac{2}{\pi}\tanh^{-1}(e^{-\pi s})e^{c/2}||\varphi_{f_{\sigma,\mu}}||_1\}
    \end{equation}
    The modified Fourier transform of $f_{\sigma,\mu}$ can be calculated explicitly:
    \begin{align}
        \begin{split}
            \varphi_{f_{\sigma,\mu}}(u) &= \frac{1}{2\pi}\int_\mathbb{R} dt \lb \frac{1}{\sigma\sqrt{2\pi}}e^{-\frac{(t - \mu)^2}{2\sigma^2}} \rb\lb \frac{2}{i}\sinh(\pi t) \rb e^{-it u}
            \\&
            =\frac{1}{\sigma\sqrt{2\pi }2\pi i}\int_\mathbb{R} dt e^{-\frac{(t-  \mu)^2}{2\sigma^2} + (\pi - i u) t} - e^{-\frac{(t - \mu)^2}{2\sigma^2} - (\pi + i u) t}
            \\&
            =\frac{1}{2\pi i}\lb e^{\frac{(\pi - iu)(2\mu + \sigma^2(\pi - iu))}{2}} - e^{-\frac{(\pi + iu)(2\mu - \sigma^2(\pi + iu))}{2}}\rb
            \\&
            =\frac{1}{\pi i}e^{\frac{\sigma^2}{2}(\pi^2 - u^2) - i\mu u}\sinh(\pi \mu  - \pi i u\sigma^2)
        \end{split}
    \end{align}
    This $L^1$-norm of $\varphi_{f_{\sigma,\mu}}$ can be bounded above by:
    \begin{align}
        \begin{split}
            ||\varphi_{f_{\sigma,\mu}}||_1 &= \int_\mathbb{R} du |\varphi_{f_{\sigma,\mu}}(u)| = \int_\mathbb{R} du \frac{e^{\frac{\sigma^2}{2}(\pi^2 - u^2)}}{\pi}|\sinh(\pi\mu - \pi i u \sigma^2)|
            \\&
            \leq \int_\mathbb{R}du\frac{e^{\frac{\sigma^2}{2}(\pi^2 - u^2)}}{\pi}e^{\pi\mu}|\sin(\pi u \sigma^2)|\leq \frac{e^{\pi\mu + \frac{\sigma^2\pi^2}{2}}}{\pi}\frac{\sqrt{2\pi}}{\sigma} = \frac{e^{\pi\mu}\sqrt{2}}{\sqrt{\pi}}\frac{e^{\frac{\sigma^2\pi^2}{2}}}{\sigma}
        \end{split}
    \end{align}
    To calculate the infimum on the right hand side of Equation \ref{equation:bound}, notice that the function on the right hand side is differentiable in $c$. Hence the infimum can be calculated using elementary methods. The infimum is achieved at $c_0 > 0$ where $c_0$ is the solution to:
    \begin{equation}
        -|\varphi_{f_{\sigma,\mu}}(c)| - |\varphi_{f_{\sigma,\mu}}(-c)| + \frac{1}{\pi}\tanh^{-1}(e^{-\pi s})e^{c/2}||\varphi_{f_{\sigma,\mu}}||_1 = 0
    \end{equation}
    Notice that $|\varphi_{f_{\sigma,\mu}}(c)| = |\varphi_{f_{\sigma,\mu}}(-c)|$:
    \begin{align}
        \begin{split}
            |\varphi_{f_{\sigma,\mu}}(c)| &= \frac{e^{\frac{\sigma^2}{2}(\pi^2 - c^2)}}{\pi}|\sinh(\pi\mu - \pi i c\sigma^2)| \\&= \frac{e^{\frac{\sigma^2}{2}(\pi^2 - c^2)}}{\pi}\sqrt{\sin^2(\pi c\sigma^2)\sinh^2(\pi\mu) + \cos^2(\pi c \sigma^2)\cosh^2(\pi\mu)}= |\varphi_{f_{\sigma,\mu}}(-c)|
        \end{split}
    \end{align}
    Hence the infimum is achieved at $c_0$ where $c_0$ is the solution to:
    \begin{equation}
        |\varphi_{f_{\sigma,\mu}}(c_0)| = \frac{1}{2\pi}\tanh^{-1}(e^{-\pi s})e^{c_0/2}||\varphi_{f_{\sigma,\mu}}||_1
    \end{equation}
    We can provide a more explicit lower bound on $c_0$:
    {\small
\begin{align}
        \begin{split}
            \frac{1}{2\pi}\tanh^{-1}(e^{-\pi s})e^{c_0 / 2}||\varphi_{f_{\sigma,\mu}}||_1 &=\frac{e^{\frac{\sigma^2}{2}(\pi^2 - c_0^2)}}{\pi}\sqrt{\sin^2(\pi c_0\sigma^2)\sinh^2(\pi\mu) + \cos^2(\pi c_0 \sigma^2)\cosh^2(\pi\mu)}
            \\&
            = \frac{e^{\frac{\sigma^2}{2}(\pi^2 - c_0^2)}}{\pi}\sqrt{\cosh^2(\pi\mu) - \sin^2(\pi c_0\sigma^2)\lb\cosh^2(\pi\mu) - \sinh^2(\pi\mu)\rb}
            \\&
            =\frac{e^{\frac{\sigma^2}{2}(\pi^2 - c_0^2)}}{\pi}\sqrt{\cosh^2(\pi\mu) - \sin^2(\pi c_0\sigma^2)}
            \\&\geq\frac{e^{\frac{\sigma^2}{2}(\pi^2 - c_0^2)}}{\pi}\sqrt{\cosh^2(\pi\mu) - 1}
            =\frac{e^{\frac{\sigma^2}{2}(\pi^2 - c_0^2)}}{\pi}\sinh(\pi\mu)
        \end{split}
    \end{align}
}
    Therefore, we have:
    \begin{equation}
        e^{-\frac{\sigma^2}{2}c_0^2 - \frac{c_0}{2}} \leq \frac{\tanh^{-1}(e^{-\pi s})||\varphi_{f_{\sigma,\mu}}||_1 e^{-\frac{\pi^2\sigma^2}{2}}}{2\sinh(\pi\mu)}
    \end{equation}
    Since $c_0 > 0$, using this formula we can solve for a lower bound of $c_0$:
    \begin{equation}
        c_0 \geq c_0(\sigma, \mu, s) := -\frac{1}{2\sigma^2} + \sqrt{\frac{1}{4} - \frac{2}{\sigma^2}\log\lb\frac{\tanh^{-1}(e^{-\pi s})||\varphi_{f_{\sigma,\mu}}||_1 e^{-\frac{\pi^2\sigma^2}{2}}}{2\sinh(\pi\mu)}\rb}
    \end{equation}
    Notice that $c_0(\sigma,\mu, s)$ monotonely increases as $s$ increases. In particular, there exists $s_0 > 0$ such that:
    \begin{equation}
        c_0(\sigma,\mu,s_0) \geq \frac{1}{\sinh^2(\pi\mu)}
    \end{equation}
    Notice that for $u \geq c_0(\sigma,\mu,s_0)\geq \frac{1}{\sinh^2(\pi\mu)}$, the norm $|\varphi_{f_{\sigma,\mu}}|$ is monotone decreasing:
    {\small
\begin{align}
        \begin{split}
            \frac{d}{du}|\varphi_{f_{\sigma,\mu}}(u)| 
            &=\frac{d}{du}\frac{e^{\frac{\sigma^2}{2}(\pi^2 - u^2)}}{\pi}\sqrt{\cosh^2(\pi\mu) - \sin^2(\pi u\sigma^2)}
            \\&= -\sigma^2 u e^{\frac{\sigma^2}{2}(\pi^2 - u^2)}\sqrt{\cosh^2(\pi\mu) - \sin^2(\pi u \sigma^2)} - \frac{2\pi\sigma^2e^{\frac{\sigma^2}{2}(\pi^2 - u^2)}\sin(\pi u\sigma^2)\cos(\pi u \sigma^2)}{2\pi\sqrt{\cosh^2(\pi\mu) - \sin^2(\pi u \sigma^2)}}
            \\&
            =\frac{\sigma^2 e^{\frac{\sigma^2}{2}(\pi^2 - u^2)}}{\sqrt{\cosh^2(\pi\mu) - \sin^2(\pi u\sigma^2)}}\lb -u\lb \cosh^2(\pi\mu) - \sin^2(\pi u \sigma^2) \rb - \sin(\pi u \sigma^2)\cos(\pi u \sigma^2) \rb
            \\&
            \leq \frac{\sigma^2 e^{\frac{\sigma^2}{2}(\pi^2 - u^2)}}{\sqrt{\cosh^2(\pi\mu) - \sin^2(\pi u \sigma^2)}}\lb -u\lb \cosh^2(\pi\mu) - 1 \rb + 1 \rb
            \\&
            =\frac{\sigma^2 e^{\frac{\sigma^2}{2}(\pi^2 - u^2)}}{\sqrt{\cosh^2(\pi\mu) - \sin^2(\pi u \sigma^2)}}\lb  -u\sinh^2(\pi\mu) + 1\rb
        \end{split}
    \end{align}
}
    Hence for $u \geq c_0(\sigma,\mu,s_0) \geq \frac{1}{\sinh^2(\pi\mu)}$, $\frac{d}{du}|\varphi_{f_{\sigma,\mu}}(u)| \leq 0$ and thus $|\varphi_{f_{\sigma,\mu}}(c_0)| \geq |\varphi_{f_{\sigma,\mu}}(u)|$. And since $|\varphi_{f_{\sigma,\mu}}(u)|$ is even, for $u \leq -c_0(\sigma,\mu, s_0)\leq -\frac{1}{\sinh^2(\pi\mu)}$, the function $|\varphi_{f_{\sigma,\mu}}(u)|$ is monotonely increasing and thus $|\varphi_{f_{\sigma,\mu}}(-c_0)|\geq |\varphi_{f_{\sigma,\mu}}(u)|$. Moreover, to the leading order in $s$, $c_0(\sigma,\mu, s)$ is approximated by:
    \begin{equation}
        c_0(\sigma,\mu,s) = -\frac{1}{2\sigma^2} + \sqrt{\frac{1}{4} - \frac{2}{\sigma^2}\log\lb\frac{||\varphi_{f_{\sigma,\mu}}||_1 e^{-\frac{\pi^2\sigma^2}{2}}}{2\sinh(\pi \mu)}\rb + \frac{2\pi s }{\sigma^2}} \sim K_1 + K_2 s
    \end{equation}
    where $K_1, K_2$ are two constants which depend on $\sigma, \mu$ and $K_2 > 0$. Finally using the formula for $|\varphi_{f_{\sigma,\mu}}(c_0)|$, we can provide an upper bound:
    \begin{align}
        \begin{split}
            |\varphi_{f_{\sigma,\mu}}(c_0)| &=\frac{e^{\frac{\sigma^2}{2}(\pi^2 - c_0^2)}}{\pi}|\sinh(\pi\mu - \pi i c_0\sigma^2)| \leq \frac{e^{\pi\mu}e^{\frac{\sigma^2}{2}(\pi^2 - c_0^2)}}{\pi}
        \end{split}
    \end{align}
    Using these observations, we can calculate the following bound for all $s\geq s_0$ (i.e. $c_0(\sigma,\mu,s_0)\geq\frac{1}{\sinh^2(\pi\mu)}$):
    {\small
\begin{align}
        \begin{split}
            c_{f_{\sigma,\mu}}(s) &\leq 2\int_{|u| \geq c_0}du |\varphi_{f_{\sigma,\mu}}(u)| + \frac{2}{\pi}\tanh^{-1}(e^{-\pi s})e^{c_0 / 2}||\varphi_{f_{\sigma,\mu}}||_1
            \\&
            \leq 2\lb|\varphi_{f_{\sigma,\mu}}(c_0)|^{1/2} + |\varphi_{f_{\sigma,\mu}}(-c_0)|^{1/2}\rb\int_\mathbb{R} du |\varphi_{f_{\sigma,\mu}}(u)|^{1/2} + \frac{2}{\pi}\tanh^{-1}(e^{-\pi s })e^{c_0 / 2}||\varphi_{f_{\sigma,\mu}}||_1
            \\&
            \leq 4|\varphi_{f_{\sigma,\mu}}(c_0)|^{1/2}\frac{e^{\frac{\sigma^2\pi^2}{4}}}{\sqrt{\pi}}e^{\frac{\pi\mu}{2}}\int_\mathbb{R}du e^{-\frac{\sigma^2}{4}u^2} + \frac{2}{\pi}\tanh^{-1}(e^{-\pi s})e^{c_0/2}||\varphi_{f_{\sigma,\mu}}||_1
            \\&
            =\frac{8e^{\frac{\sigma^2\pi^2}{4}}e^{\frac{\pi\mu}{2}}}{\sigma}|\varphi_{f_{\sigma,\mu}}(c_0)|^{1/2} + 4|\varphi_{f_{\sigma,\mu}}(c_0)|
            \\&
            \leq\frac{8e^{\frac{\sigma^2\pi^2}{4}}e^{\frac{\pi\mu}{2}}}{\sigma}\frac{e^{\frac{\pi\mu}{2}}e^{\frac{\sigma^2}{4}(\pi^2 - c_0^2)}}{\sqrt{\pi}} + 4\frac{e^{\pi\mu}e^{\frac{\sigma^2}{2}(\pi^2 - c_0^2)}}{\pi}
            \\&
            \leq\frac{8e^{\frac{\sigma^2\pi^2}{4}}e^{\frac{\pi\mu}{2}}}{\sigma}\frac{e^{\frac{\pi\mu}{2}}e^{\frac{\sigma^2}{4}(\pi^2 - c_0(\sigma,\mu,s)^2)}}{\sqrt{\pi}} + 4\frac{e^{\pi\mu}e^{\frac{\sigma^2}{2}(\pi^2 - c_0(\sigma,\mu,s)^2)}}{\pi}
        \end{split}
    \end{align}
}
    As $s$ increases, the upper bound on $c_{f_{\sigma,\mu}}(s)$ is monotonely decreasing. And to the leading order in $s$, the upper bound can be approximated by:
    \begin{equation}
        c_{f_{\sigma,\mu}}(s) \leq K_3e^{-\frac{K_4 s^2}{2}} + K_5e^{-K_4 s^2}
    \end{equation}
    where $K_3, K_4, K_5 > 0$ are positive constants which depend on $\sigma, \mu$. 
\end{proof}
In addition, as $\sigma \rightarrow 0$, the upper bound of $c_{f_{\sigma, \mu}}$ approaches:
\begin{align}
    \begin{split}
        c_{f_{\sigma,\mu}}(s) &\leq C_1 \frac{e^{-C_2/\sigma^2}e^{C_3\sigma^2}}{\sigma}e^{-C_4s^2}
    \end{split}
\end{align}
where $C_1,C_2,C_3, C_4 > 0$ are positive constants that depend on $\mu$. The right hand side as a well-defined limit as $\sigma\rightarrow 0^+$ and the limit is $0$ (uniformly in $s$).


\subsection{Example in 1+1d CFT}
We consider the setup in a 1+1d CFT with $R = e^{2\pi s}$ and compute an asymptotic upper bound of $||\int_\mathbb{R} dt f(t)(\Delta_{B_R}^{-it} - \Delta_{W_0}^{-it})A\ket{\Omega}||^2$ with $f_{\sigma,\mu}(t)$ the normalized Gaussian distribution. For simplicity, We choose self-adjoint $A = \mathcal{O}(g) = \int dtdx \;g(t,x) \mathcal{O}(t,x)$ with a well-behave smearing function $g(t,x)$ sported in $B_{r}$. From the expression in Section~\ref{subsec:ball-inside-wedge}, we have 
\begin{gather}
    \Delta_{B_R}^{-iz} \mathcal{O}(g) \Delta_{B_R}^{iz} = \mathcal{O}(\tilde{g}),\\
    \tilde{g}(X) = \left|\text{det}\frac{d\alpha_{B_R}^{-z}(X)}{dX}\right|\gamma^{2h}(\alpha_{B_R}^{-z}(X)) g(\alpha_{B_R}^{-z}(X)),
\end{gather}
where we denote $(t,x)$ as $X$ and $\alpha_{B_R}^{z}$ is the geometric action of $\text{Ad}_{ \Delta_{B_R}^{-iz}}$. The parameters of the action of $\text{Ad}_{\Delta_{B_R}^{-iz}}$ are
\begin{gather}
    t_c = 0, \qquad x_c  = R, \qquad 
    |-t^2 + (x-r)^2| < r\\
    d^2 = -t^2 +(x-R)^2 = -t^2 + (x-r)^2 +(R-r)^2 - 2(x-r)(R-r).
\end{gather}
For large $R=re^{2\pi s}\gg r$, we have the approximation
\begin{align}
    d^2 &= R^2\left[1 -2x/R +O({(r/R)}^2)\right],\\
    \gamma(t,x)&= \left(1 + \frac{x}{R}(\cosh(2\pi z) -1)+\frac{t}{R}\sinh (2\pi z) + O((r/R)^2)\right)^{-1}\nonumber\\
    &= 1 - \frac{x}{R}(\cosh(2\pi z) -1)-\frac{t}{R}\sinh (2\pi z) + O((r/R)^2),\\
    t_z &= t\cosh (2\pi z) + x\sinh (2\pi z) + r O(re^{4\pi|z|}/R), \quad \text{for } \;re^{4\pi|z|}\ll R,\\
    x_z &= t\sinh (2\pi z) + x\cosh (2\pi z) + r O(re^{2\pi|z|}/R), \quad \text{for } \;re^{2\pi|z|}\ll R.
\end{align}
For $re^{4\pi|z|}\ll R$, We have the approximation of $g(X)$ as
\begin{gather}
    g(\alpha_{B_R}^{-z}(X)) = g(\alpha_{W}^{-z}(X)) + O(re^{4\pi|z|}/R),\\
    \tilde{g}(X) = g(\alpha_{W}^{-z}(X)) + O(re^{4\pi|z|}/R).
\end{gather}
Notice that $O(re^{4\pi|z|}/R)$ is in the form of
\begin{align}
    O(re^{4\pi|z|}/R) = \sum_{n=-2}^2\frac{r}{R} e^{2n\pi z} \; a_n(X),
\end{align}
for some functions $a_n(X)$. This expression comes from the function $\sinh(2\pi z)$ and $\cosh(2\pi z)$ in $\gamma$, $t_z$, and $x_z$.

We break $||\int_\mathbb{R} dz f(t)(\Delta_{B_R}^{-iz} - \Delta_{W_0}^{-iz})A\ket{\Omega}||$ into
\begin{align}
    &||\int_\mathbb{R} dz f_{\sigma,\mu}(z)(\Delta_{B_R}^{-iz} - \Delta_{W_0}^{-iz})A\ket{\Omega}|| \nn\\
    &=\left|\left|\int_{|z|<s_0} dz f_{\sigma,\mu}(z) (\Delta_{B_R}^{-iz} - \Delta_{W_0}^{-iz})A\ket{\Omega}+\int_{|z|\geq s_0} dz f_{\sigma,\mu}(z)(\Delta_{B_R}^{-iz} - \Delta_{W_0}^{-iz})A\ket{\Omega}\right|\right|\\
    & \leq \left|\left|\int_{|z|<s_0} dz f_{\sigma,\mu}(z) (\Delta_{B_R}^{-iz} - \Delta_{W_0}^{-iz})A\ket{\Omega}\right|\right| + \left|\left| \int_{|z|\geq s_0} dz f_{\sigma,\mu}(z)(\Delta_{B_R}^{-iz} - \Delta_{W_0}^{-iz})A\ket{\Omega}\right|\right|,
\end{align}
for some $s_0$ such that $s_0\ll s$. The second term has is bounded by
\begin{align}
    \left|\left|\int_{|z|\geq s_0} dz f_{\sigma,\mu}(z)(\Delta_{B_R}^{-iz} - \Delta_{W_0}^{-iz})A\ket{\Omega}\right|\right| &\leq \int_{|z|\geq s_0} dz |f_{\sigma,\mu}(z)|\;\left|\left|(\Delta_{B_R}^{-iz} - \Delta_{W_0}^{-iz})A\ket{\Omega}\right|\right|\\
    &\leq 2||A\ket{\Omega}|| \int_{|z|\geq s_0} dz |f_{\sigma,\mu}(z)| .
\end{align}
This integral decays like a rescaled complementary error function of $s_0$. For the first term, we have $|z|< s_0\ll s$ and we can use the approximation for $re^{4\pi|z|}\ll R$. We have
\begin{align}
    &\left|\left|\int_{|z|<s_0} dz f_{\sigma,\mu}(z) (\Delta_{B_R}^{-iz} - \Delta_{W_0}^{-iz})A\ket{\Omega}\right|\right| \leq \int_{|z|<s_0} dz |f_{\sigma,\mu}(z)| \left|\left|(\Delta_{B_R}^{-iz} - \Delta_{W_0}^{-iz})A\ket{\Omega}\right|\right| \\
    &= \int_{|z|<s_0} dz |f_{\sigma,\mu}(z)| \sqrt{2||A\ket{\Omega}||^2 - 2\text{Re}\bra{\Omega}A \Delta_{W_0}^{+iz}\Delta_{B_R}^{-iz} A\ket{\Omega}}.
\end{align}
From the approximation we have
\begin{align}
    \text{Ad}_{\Delta_{W_0}^{+iz}\Delta_{B_R}^{-iz}} A &= \text{Ad}_{\Delta_{W_0}^{+iz}\Delta_{B_R}^{-iz}} \mathcal{O}(g(X))\\
    & = \text{Ad}_{\Delta_{W_0}^{+iz}} \left(\mathcal{O}(g(\alpha_{W}^{-z}(X))) + \mathcal{O}(O(re^{4\pi|z|}/R))\right)\\
    & = \mathcal{O}(g(X)) + \mathcal{O}(O(re^{4\pi|z|}/R)).
\end{align}
Thus, we can write 
\begin{align}
    2||A\ket{\Omega}||^2 - 2\bra{\Omega}A \Delta_{W_0}^{+iz}\Delta_{B_R}^{-iz} A\ket{\Omega} = \frac{r}{R} \sum_{n=-2}^2e^{2n\pi z} \bra{\Omega} \mathcal{O}(g) \tilde{\mathcal{O}}_n\ket{\Omega},
\end{align}
for some $\tilde{\mathcal{O}}_n$ that depends on the explicit form of $g(X)$. This gives 
\begin{align}
    \left|\left|\int_{|z|<s_0} dz f_{\sigma,\mu}(z) (\Delta_{B_R}^{-iz} - \Delta_{W_0}^{-iz})A\ket{\Omega}\right|\right| 
    \leq \sqrt{\frac{r}{R}}\int_{\mathbb{R}} dz |f_{\sigma,\mu}(z)| \sqrt{ \sum_{n=-2}^2e^{2n\pi z} \bra{\Omega} \mathcal{O}(g) \tilde{\mathcal{O}}_n\ket{\Omega}}.
\end{align}
$f_{\sigma,\mu}(z)$ is a Gaussian function so that the integral is a finite constant. $\sqrt{r/R}$ decays as $e^{-\pi s}$. Therefore, we have the bound for large $s$
\begin{align}
    \left|\left|\int_\mathbb{R} dz f_{\sigma,\mu}(z)(\Delta_{B_R}^{-iz} - \Delta_{W_0}^{-iz})A\ket{\Omega}\right|\right| \lesssim  Ke^{-\pi s} + 2||A\ket{\Omega}|| \int_{|z|\geq s_0} dz |f_{\sigma,\mu}(z)|,
\end{align}
for some constant $K$. We can take $s_0/s = \epsilon\ll 1$. The integral $\int_{|z|\geq \epsilon s} dz |f_{\sigma,\mu}(z)|$ decays much faster than $e^{-\pi s}$. Thus, the asymptotic behaviour of the difference is
\begin{align}
    \left|\left|\int_\mathbb{R} dz f_{\sigma,\mu}(z)(\Delta_{B_R}^{-iz} - \Delta_{W_0}^{-iz})A\ket{\Omega}\right|\right|^2 \lesssim  K'e^{-2\pi s},
\end{align}
for some constant $K'$.

The calculation for the characteristic operators is similar in this case. Without loss of generality, we compare the characteristic operator of the inclusion $W_0\subset W_{-a}$ and that of $B_R\subset W_{-a}$, with $a>0$. We have
{\small
\begin{align}
    &\left|\left|\int_\mathbb{R} dz f_{\sigma,\mu}(z)(D_{W_0\subset W_{-a}}(z) - D_{B_R\subset W_{-a}}(z))A\ket{\Omega}\right|\right| = \left|\left|\int_\mathbb{R} dz f_{\sigma,\mu}(z)\Delta_{W_{-a}}^{-iz}(\Delta_{B_R}^{iz} - \Delta_{W_0}^{iz})A\ket{\Omega}\right|\right| \\
    & \leq \left|\left|\int_{|z|<s_0} dz f_{\sigma,\mu}(z)\Delta_{W_{-a}}^{-iz}(\Delta_{B_R}^{iz} - \Delta_{W_0}^{iz})A\ket{\Omega}\right|\right| + \left|\left|\int_{|z|\geq s_0} dz f_{\sigma,\mu}(z)\Delta_{W_{-a}}^{-iz}(\Delta_{B_R}^{iz} - \Delta_{W_0}^{iz})A\ket{\Omega}\right|\right|\\
    &\leq \int_{|z|<s_0} dz |f_{\sigma,\mu}(z)|\;\left|\left|(\Delta_{B_R}^{iz} - \Delta_{W_0}^{iz})A\ket{\Omega}\right|\right| + 2 \left|\left|A\ket{\Omega}\right|\right|\int_{|z|\geq s_0} dz |f_{\sigma,\mu}(z)|.
\end{align}
}
This is the same expression as in the case of modular flow, but with $z\to-z$ in the first term because of the sign convention. Thus, the difference in the characteristic operator has a similar large $s$ asymptotic behaviour
\begin{align}
    \left|\left|\int_\mathbb{R} dz f_{\sigma,\mu}(z)(D_{W_0\subset W_{-a}}(z) - D_{B_R\subset W_{-a}}(z))A\ket{\Omega}\right|\right|^2 \lesssim K''e^{-2\pi s}.
\end{align}
We observe that the asymptotic behaviour decays much faster than the general bound. This is expected as the general bound uses neither the details of the QFT nor the extra symmetry in a CFT. In addition, it is straightforward to derive the bound for the characteristic operator case in this example. The reason is that we use the triangle inequality to bound the norm of an integral with an integral of the norm. In general QFT, $\mathcal{A}_{B_{R}} \subset \mathcal{A}_{B_{R'}}$ for $R<R'$ and $(\cup_R \mathcal{A}_{B_{R}})'' = \mathcal{A}_{W_0}$ implies $\text{s-}\lim_{R\to\infty} \Delta^{-iz}_{B_{R}} = \Delta^{-iz}_{W_0}$ in the strong operator topology. We have
\begin{align}
    \text{s-}\lim_{R\to \infty}\left|\left|(\Delta_{B_{R}}^{-iz} - \Delta_{W_0}^{-iz})A\ket{\Omega}\right|\right| = 0
\end{align}
for all $z$. It is a question of whether we also have
\begin{align}
    \lim_{R\to\infty}\int_{|z|<s_0} dz |f_{\sigma,\mu}(z)|\;\left|\left|(\Delta_{B_R}^{iz} - \Delta_{W_0}^{iz})A\ket{\Omega}\right|\right| = 0.
\end{align} In other words, we need $\left|\left|(\Delta_{B_{R}}^{-iz} - \Delta_{W_0}^{-iz})A\ket{\Omega}\right|\right| \to 0$ uniformly for the range of $z$. In this CFT example, the integral of the norm indeed converges to zero. The calculation for the characteristic operator case is thus reduced to the same calculation as in the modular flow case. The geometric action of the modular flows allows us to bound all $\left|\left|(\Delta_{B_{R}}^{-iz} - \Delta_{W_0}^{-iz})A\ket{\Omega}\right|\right|$ in the truncated range of $z$.


% \section{Anosov Cocycle Derivative and Borchers Characteristic Function of Standard Subspaces}\label{app:classicalborchers}


% In this appendix, we consider characteristic functions of nested standard subspaces. These characteristic functions can be understood as the one-particle analog of the characteristic functions of local algebras of observables in an algebraic quantum field theory. Simplest examples of nested standard subspaces come from anti-unitary representation of $\mathbb{Z}_2$ extension of certain real semi-simple Lie groups. These examples are nothing but slight generalization of the familiar examples of chiral conformal nets. These nets can be constructed out of nets of standard subspaces using representation theory of $PSL(2,\mathbb{R})$. However, the new observation here is that the first order variation of characteristic function is naturally given by the Anosov cocycle derivative. In the classical theory of deterministic chaos, Anosov cocycle derivative is a key quantity that can be used to calculate Lyapnov exponents, metric entropy, and to study mixing properties of the dynamical system. To keep the discussion precise, we choose to work with $SO_0(1,n)$. However, the same discussion can be generalized to other suitable real semi-simple Lie groups and their homogeneous spaces.

% Let
% \[
% G=SO_0(1,n),\qquad K=SO(n),
% \]
% and let
% \[
% \mathbb H^n\simeq G/K
% \]
% be real hyperbolic \(n\)-space.  We choose a one-dimensional split Cartan subgroup
% \[
% A=\{a_r=e^{rh}:r\in\mathbb R\}
% \]
% generated by an element \(h\in\mathfrak g=\mathfrak{so}(1,n)\) normalized so that
% \(\operatorname{ad}h\) has eigenvalues \(-1,0,1\).  Thus \(h\) is an Euler element and
% \[
% \mathfrak g
% =
% \mathfrak g_{-1}(h)\oplus \mathfrak g_0(h)\oplus \mathfrak g_1(h)
% =
% \mathfrak n_-\oplus(\mathfrak m\oplus\mathfrak a)\oplus\mathfrak n_+,
% \]
% where
% \[
% \dim \mathfrak n_+=\dim\mathfrak n_-=n-1,\qquad
% \mathfrak m\simeq\mathfrak{so}(n-1),\qquad
% \mathfrak a=\mathbb Rh.
% \]
% For \(X_\pm\in\mathfrak n_\pm\),
% \begin{equation}
% [h,X_+]=X_+,\qquad [h,X_-]=-X_-,
% \end{equation}
% and therefore
% \begin{equation}
% \Ad(a_r)X_+=e^rX_+,\qquad
% \Ad(a_r)X_-=e^{-r}X_-.
% \label{eq:root-scaling}
% \end{equation}
% For real hyperbolic space the positive nilpotent algebra \(\mathfrak n_+\) is abelian.
% Writing
% \[
% N_+=\exp\mathfrak n_+\simeq\mathbb R^{n-1},
% \qquad
% n_v=\exp X_v,\qquad v\in\mathbb R^{n-1},
% \]
% one obtains
% \begin{equation}
% a_r n_v a_{-r}=n_{e^rv}.
% \label{eq:horospherical-dilation}
% \end{equation}
% This identity is the basic algebraic input for both the modular and the Anosov parts of
% the construction.

% The recent homogeneous-net viewpoint of Neeb makes it natural to regard the standard
% subspace associated with \(h\) not as an isolated object, but as the wedge value of a
% covariant net of real subspaces.  An Euler element determines the involution
% \[
% \tau_h=e^{\pi i\,\operatorname{ad}h},
% \]
% which acts by \(-1\) on \(\mathfrak g_{\pm1}(h)\) and by \(+1\) on
% \(\mathfrak g_0(h)\).  Let
% \[
% U:G_{\tau_h}\longrightarrow \operatorname{AU}(\cH)
% \]
% be an antiunitary representation of the corresponding \(\mathbb Z_2\)-extension.
% Set
% \[
% J:=U(\tau_h),
% \qquad
% U(a_r)=e^{irH},
% \]
% with \(H=H^*\).  We normalize the modular operator by
% \begin{equation}
% \Delta=e^{-2\pi H},
% \qquad
% \Delta^{it}=U(a_{-2\pi t}).
% \label{eq:modular-boost}
% \end{equation}
% Since \(JHJ=-H\), one has
% \[
% J\Delta J=\Delta^{-1}.
% \]
% Hence
% \[
% S:=J\Delta^{1/2}=Je^{-\pi H}
% \]
% is a closed antilinear involution on its natural domain, and
% \begin{equation}
% V
% =
% \operatorname{Fix}(S)
% =
% \left\{
% \xi\in\Dom(e^{-\pi H}):Je^{-\pi H}\xi=\xi
% \right\}
% \label{eq:standard-subspace}
% \end{equation}
% is a standard real subspace:
% \[
% V\cap iV=\{0\},\qquad \overline{V+iV}=\cH.
% \]
% Its modular objects are \(J_V=J\) and \(\Delta_V=\Delta\).

% If \(H\) is diagonalized as
% \[
% \cH\simeq L^2(\mathbb R,\dd\nu(\lambda);\mathcal K_\lambda),
% \qquad
% (Hf)(\lambda)=\lambda f(\lambda),
% \]
% and the spectral identification is chosen compatibly with \(J\) so that, schematically,
% \[
% (Jf)(\lambda)=\overline{f(-\lambda)},
% \]
% then the Tomita fixed-point condition becomes
% \begin{equation}
% f(-\lambda)=e^{-\pi\lambda}\overline{f(\lambda)}.
% \label{eq:spectral-standard}
% \end{equation}
% Thus the real standard subspace is explicitly determined by a modular spectral weight.

% Neeb's recent formulation places this construction on a homogeneous causal space.  The
% canonical geometric partner of an Euler element in the semisimple setting is a
% non-compactly causal symmetric space \(M_{\rm c}=G/H_{\rm c}\).  For
% \(G=SO_0(1,n)\), the model relevant here is de Sitter space
% \begin{equation}
% M_{\rm c}=dS^n
% \simeq
% SO_0(1,n)/SO_0(1,n-1),
% \label{eq:de-sitter}
% \end{equation}
% which should be distinguished from the Riemannian symmetric space
% \[
% \mathbb H^n=SO_0(1,n)/SO(n).
% \]
% The point is that the same pair \((G,h)\) acts on both spaces.  On the causal side \(h\)
% generates the modular flow of a wedge domain; on the Riemannian side it generates the
% geodesic flow after passing to the unit tangent bundle.

% Let \(W\subset M_{\rm c}\) denote the distinguished wedge domain associated with \(h\).
% In Neeb's homogeneous-net framework one seeks a net
% \[
% \mathcal O\longmapsto \cH_M(\mathcal O)
% \]
% of closed real subspaces satisfying isotony, covariance and the Reeh--Schlieder
% property, with the Bisognano--Wichmann identification
% \begin{equation}
% \cH_M(W)=V,\qquad
% \Delta_{\cH_M(W)}^{-it/2\pi}=U(e^{th})=U(a_t).
% \label{eq:BW}
% \end{equation}
% For semisimple non-compactly causal symmetric spaces, the maximal net constructed from
% the wedge standard subspace satisfies isotony, covariance, Reeh--Schlieder and
% Bisognano--Wichmann under the hypotheses established in
% \cite{Neeb2025,FNO2025}.  Concretely, if
% \[
% A_{\mathcal O}:=
% \{g\in G:g^{-1}\mathcal O\subseteq W\},
% \]
% then the maximal net can be written in the form
% \begin{equation}
% \cH_M^{\max}(\mathcal O)
% =
% \bigcap_{g\in A_{\mathcal O}}U(g)V.
% \label{eq:maximal-net}
% \end{equation}
% Thus the standard subspace \(V\) appearing in \eqref{eq:standard-subspace} is naturally
% interpreted as a wedge-localized one-particle space, and its group translates are
% geometrically localized standard subspaces.

% This observation gives a geometric meaning to the horospherical translates used in the
% Borchers calculation.  For \(v\in\mathbb R^{n-1}\), put
% \[
% V_v:=U(n_v)V.
% \]
% Covariance identifies this as
% \[
% V_v=\cH_M(n_vW),
% \]
% whenever the net is evaluated on the translated wedge.  Covariance of modular objects
% gives
% \[
% \Delta_{V_v}^{it}
% =
% U(n_v)\Delta_V^{it}U(n_v)^{-1}.
% \]
% Using \eqref{eq:modular-boost} and \eqref{eq:horospherical-dilation},
% \begin{align}
% \Delta_V^{it}U(n_v)\Delta_V^{-it}
% &=
% U(a_{-2\pi t}n_va_{2\pi t})
% \nn\\
% &=
% U(n_{e^{-2\pi t}v}).
% \label{eq:borchers-scaling}
% \end{align}
% Define the relative modular product
% \begin{equation}
% D_v(t):=\Delta_V^{-it}\Delta_{V_v}^{it}.
% \label{eq:relative-modular}
% \end{equation}
% Then
% \begin{align}
% D_v(t)
% &=
% \Delta_V^{-it}U(n_v)\Delta_V^{it}U(n_v)^{-1}
% \nn\\
% &=
% U(a_{2\pi t}n_va_{-2\pi t}n_{-v})
% \nn\\
% &=
% U(n_{(e^{2\pi t}-1)v}).
% \label{eq:borchers-characteristic}
% \end{align}
% This has exactly the functional form of Borchers' characteristic function.  When the
% horospherical generator has the positivity needed for a half-sided modular inclusion,
% \eqref{eq:borchers-characteristic} is the characteristic function of that inclusion.
% Without that additional positivity hypothesis, it remains an exact relative-modular
% identity for the geometrically related wedge subspaces \(V\) and \(V_v\).

% Writing
% \[
% U(n_{sv})=e^{isP_v},
% \]
% and differentiating \eqref{eq:borchers-scaling} at \(s=0\), one obtains
% \begin{equation}
% \Delta_V^{-it}P_v\Delta_V^{it}
% =
% e^{2\pi t}P_v.
% \label{eq:modular-generator-scaling}
% \end{equation}
% Thus the adjoint modular action has a positive weight \(2\pi\) in modular time.  If
% geometric time is defined by
% \[
% r=2\pi t,
% \]
% the expansion becomes \(P_v\mapsto e^rP_v\).  This is precisely the same restricted-root
% weight that will appear as the unstable Lyapunov exponent of the geodesic flow.

% The relation between the causal-net and hyperbolic-dynamical realizations can be displayed
% as
% \[
% \begin{tikzcd}[column sep=large,row sep=large]
% & (G,h,U) \arrow[dl] \arrow[dr] & \\
% dS^n=G/H_{\rm c} \arrow[d,"\text{wedge net}"]
% &&
% \mathbb H^n=G/K \arrow[d,"\text{compact quotient}"]\\
% W\mapsto \cH_M(W)=V
% \arrow[d,"\Delta_V^{-it/2\pi}=U(a_t)"]
% &&
% T^1(\Gamma\backslash\mathbb H^n)\simeq\Gamma\backslash G/M_0
% \arrow[d,"\phi_t"]\\
% \Ad(\Delta_V^{-it})|_{\mathcal P^u}
% \arrow[dr]
% &&
% D\phi_{2\pi t}|_{E^u}\arrow[dl]\\
% &
% e^{2\pi t\,\operatorname{ad}h}|_{\mathfrak g_1(h)}
% &
% \end{tikzcd}
% \]
% where \(M_0=Z_K(A)\simeq SO(n-1)\).  The two lower arrows are therefore not merely
% analogous: they are two representations of the same Euler-element dilation on
% \(\mathfrak g_1(h)\).

% To make the dynamical side explicit, let \(\Gamma<G\) be a torsion-free cocompact lattice.
% Then
% \[
% X_\Gamma=\Gamma\backslash\mathbb H^n
% \]
% is a closed hyperbolic \(n\)-manifold and
% \begin{equation}
% T^1X_\Gamma
% \simeq
% \Gamma\backslash G/M_0.
% \label{eq:unit-tangent}
% \end{equation}
% Choose the convention
% \begin{equation}
% \phi_r(\Gamma gM_0)=\Gamma ga_{-r}M_0.
% \label{eq:geodesic-flow}
% \end{equation}
% At the tangent level,
% \[
% T_{\Gamma gM_0}(\Gamma\backslash G/M_0)\simeq \mathfrak g/\mathfrak m,
% \]
% and a tangent vector represented by
% \[
% s\longmapsto \Gamma ge^{sX}M_0
% \]
% is carried to
% \[
% \phi_r(\Gamma ge^{sX}M_0)
% =
% \Gamma ga_{-r}\left(a_re^{sX}a_{-r}\right)M_0.
% \]
% Consequently,
% \begin{equation}
% D\phi_r[X]=[\Ad(a_r)X].
% \label{eq:derivative-adjoint}
% \end{equation}
% Using \eqref{eq:root-scaling},
% \begin{equation}
% E^u\simeq\mathfrak n_+=\mathfrak g_1(h),\qquad
% E^s\simeq\mathfrak n_-=\mathfrak g_{-1}(h),\qquad
% E^0\simeq\mathfrak a,
% \label{eq:splitting}
% \end{equation}
% and
% \begin{equation}
% D\phi_r|_{E^u}=e^r I_{n-1},
% \qquad
% D\phi_r|_{E^s}=e^{-r}I_{n-1}.
% \label{eq:anosov-derivative}
% \end{equation}
% The geodesic flow is therefore uniformly Anosov.

% Combining \eqref{eq:modular-generator-scaling} and
% \eqref{eq:anosov-derivative}, with \(r=2\pi t\), yields the direct dictionary
% \begin{equation}
% \Ad(\Delta_V^{-it})P_v=e^{2\pi t}P_v,
% \qquad
% D\phi_{2\pi t}(v)=e^{2\pi t}v.
% \label{eq:dictionary}
% \end{equation}
% At the group level the Borchers-type relative modular product can be rewritten as
% \begin{equation}
% D_v(t)
% =
% U\!\left(
% n_{\bigl(D\phi_{2\pi t}|_{E^u}-I\bigr)v}
% \right).
% \label{eq:char-vs-anosov}
% \end{equation}
% Equation \eqref{eq:char-vs-anosov} should be interpreted carefully: the relative modular
% operator is not literally the tangent cocycle.  Rather, the displacement appearing in
% the characteristic function is obtained by applying \(D\phi_{2\pi t}-I\) to the same
% root coordinate \(v\).  Both structures are generated by
% \[
% \Ad(A)|_{\mathfrak g_1(h)}.
% \]

% The Lyapunov exponents now follow immediately.  For \(0\neq v\in E^u\),
% \begin{equation}
% \lambda_+(v)
% =
% \lim_{r\to\infty}
% \frac{1}{r}
% \log\|D\phi_rv\|
% =
% 1,
% \end{equation}
% and similarly the stable exponent is \(-1\).  Since
% \(\dim E^u=\dim E^s=n-1\), the Lyapunov spectrum is
% \begin{equation}
% \underbrace{1,\ldots,1}_{n-1},
% \qquad 0,\qquad
% \underbrace{-1,\ldots,-1}_{n-1}.
% \label{eq:lyapunov-spectrum}
% \end{equation}
% The positive exponent can be read directly from the modular action:
% \begin{equation}
% \lambda_+
% =
% \frac{1}{2\pi}
% \frac{\dd}{\dd t}
% \log
% \left\|
% \Ad(\Delta_V^{-it})P_v
% \right\|.
% \label{eq:modular-lyapunov}
% \end{equation}
% For \(SO_0(1,n)\), the unstable Jacobian is
% \begin{equation}
% J_r^u(x)
% =
% \left|
% \det D\phi_r|_{E^u(x)}
% \right|
% =
% e^{(n-1)r}.
% \label{eq:unstable-jacobian}
% \end{equation}
% Thus
% \[
% \frac{\dd}{\dd r}\log J_r^u=n-1.
% \]
% For normalized Liouville measure \(\mu_L\), Pesin's entropy formula gives
% \begin{align}
% h_{\mu_L}(\phi_1)
% &=
% \int_{T^1X_\Gamma}
% \sum_{\lambda_j(x)>0}m_j(x)\lambda_j(x)\,
% \dd\mu_L(x)
% \nn\\
% &=
% n-1.
% \label{eq:pesin}
% \end{align}
% Equivalently,
% \[
% h_{\mu_L}(\phi_r)=|r|(n-1).
% \]
% The same quantity can be reconstructed from the modular representation.  Let
% \[
% \mathcal P^u
% =
% \operatorname{span}\{P_1,\ldots,P_{n-1}\}
% \]
% be the space of unstable horospherical generators.  Since
% \[
% \Ad(\Delta_V^{-it})|_{\mathcal P^u}
% =
% e^{2\pi t}I_{n-1},
% \]
% one has
% \[
% \det\left(
% \Ad(\Delta_V^{-it})|_{\mathcal P^u}
% \right)
% =
% e^{2\pi(n-1)t},
% \]
% and hence
% \begin{equation}
% h_{\mu_L}(\phi_1)
% =
% \frac{1}{2\pi}
% \left.
% \frac{\dd}{\dd t}
% \log
% \det\left(
% \Ad(\Delta_V^{-it})|_{\mathcal P^u}
% \right)
% \right|_{t=0}.
% \label{eq:modular-entropy}
% \end{equation}
% Thus, in this constant-curvature rank-one model, the Pesin entropy of the Riemannian
% dynamics is the logarithmic modular Jacobian of the causal wedge net on the unstable
% Euler subspace.  The equality does not identify the two homogeneous spaces; rather, it
% expresses the fact that both invariants are determined by the same triple \((G,h,U)\)
% together with the same root space \(\mathfrak g_1(h)\).

% This perspective clarifies the role of the compact quotient.  The Euler element and its
% root dilation already contain the local hyperbolicity:
% \[
% \Ad(a_r)|_{\mathfrak g_{\pm1}(h)}=e^{\pm r}.
% \]
% Passing to \(\Gamma\backslash G/M_0\) does not create the instability.  Instead it turns
% the homogeneous dynamics into a finite-measure system on which global invariants such as
% metric entropy and mixing are defined.  In particular, the right regular representation
% \[
% (R(g)F)(\Gamma x)=F(\Gamma xg)
% \]
% acts on \(L^2(\Gamma\backslash G)\), and the geodesic flow is represented by
% \(R(a_{-r})\) on the \(M_0\)-invariant subspace
% \[
% L^2(\Gamma\backslash G/M_0)
% \simeq
% L^2(\Gamma\backslash G)^{M_0}.
% \]
% For \(F,H\in L^2_0(\Gamma\backslash G)\),
% \[
% C_{F,H}(r)
% =
% \langle F,R(a_{-r})H\rangle.
% \]
% Since \(a_r\to\infty\) in \(G\) as \(|r|\to\infty\), the Howe--Moore theorem implies
% \begin{equation}
% \langle F,R(a_{-r})H\rangle\longrightarrow0.
% \label{eq:howe-moore}
% \end{equation}
% Therefore for arbitrary \(F,H\in L^2(\Gamma\backslash G/M_0)\),
% \begin{equation}
% \int F(x)H(\phi_rx)\,\dd\mu_L(x)
% \longrightarrow
% \left(\int F\,\dd\mu_L\right)
% \left(\int H\,\dd\mu_L\right),
% \label{eq:mixing}
% \end{equation}
% which proves mixing of the geodesic flow.

% The construction can therefore be summarized as follows.  The Euler element \(h\)
% produces the three-grading
% \[
% \mathfrak g_{-1}(h)\oplus\mathfrak g_0(h)\oplus\mathfrak g_1(h).
% \]
% Neeb's homogeneous-net theory interprets this grading on the causal symmetric space
% \(dS^n=G/H_{\rm c}\): the wedge \(W\) carries the standard subspace \(V\) and the
% Bisognano--Wichmann modular group is \(U(e^{th})\).  Horospherical wedge translates
% \(n_vW\) give the relative modular product
% \[
% \Delta_V^{-it}\Delta_{V_v}^{it}
% =
% U(n_{(e^{2\pi t}-1)v}).
% \]
% The Riemannian dual geometry \(G/K=\mathbb H^n\) uses the same \(A\)-action, and after a
% compact quotient its derivative on the unstable bundle is
% \[
% D\phi_r|_{E^u}
% =
% e^rI.
% \]
% Thus
% \[
% \text{Euler weight}
% \longrightarrow
% \text{modular dilation}
% \longrightarrow
% \text{Borchers displacement}
% \]
% and
% \[
% \text{Euler weight}
% \longrightarrow
% \text{Anosov derivative}
% \longrightarrow
% \text{Lyapunov spectrum}
% \longrightarrow
% \text{Pesin entropy}
% \]
% are two branches of the same Lie-theoretic mechanism.  The final identity
% \[
% h_{\mu_L}(\phi_1)
% =
% \frac{1}{2\pi}
% \left.
% \frac{\dd}{\dd t}
% \log
% \det\left(
% \Ad(\Delta_V^{-it})|_{\mathcal P^u}
% \right)
% \right|_{t=0}
% =
% n-1
% \]
% is therefore not an accidental agreement of exponents: it is the common Jacobian of the
% Euler dilation as seen in the modular localization representation and in the unstable
% tangent representation.

% This formulation also suggests the natural generalizations.  For other rank-one groups
% one replaces the single root space by
% \[
% \mathfrak n_+
% =
% \mathfrak g_\alpha\oplus\mathfrak g_{2\alpha},
% \]
% so that
% \[
% \Ad(a_r)X_\alpha=e^rX_\alpha,\qquad
% \Ad(a_r)X_{2\alpha}=e^{2r}X_{2\alpha},
% \]
% and the entropy becomes
% \[
% h_\mu(\phi_1)
% =
% m_\alpha+2m_{2\alpha}
% =
% 2\rho(h).
% \]
% For a general semisimple group and regular \(H_0\in\mathfrak a\), the same calculation
% formally gives
% \[
% h_\mu(a_{H_0})
% =
% \sum_{\alpha(H_0)>0}
% m_\alpha\alpha(H_0),
% \]
% while on the modular side
% \[
% \Ad(\Delta^{-it})P_\alpha
% =
% e^{2\pi t\alpha(H_0)}P_\alpha.
% \]
% The more ambitious extension is to variable negative curvature, where the constant
% restricted-root factor \(e^r\) is replaced by the point-dependent unstable cocycle
% \[
% A^u(r,x)=D\phi_r|_{E^u(x)}.
% \]
% A corresponding family of relative modular data would amount to a modular analogue of an
% Oseledets cocycle.  The homogeneous \(SO_0(1,n)\) calculation above should then be viewed
% as the constant-coefficient model in which the modular and Anosov cocycles are both
% representations of the same Euler-element action.

% % \section{Characteristic flows in classical ergodic theory}

% % The goal of this section is to argue for a classical version of Connes cocycle as $(\rho/\rho_0)^{it}$ for probability measures $\rho$ and $\rho_0=\rho\circ E$ a coarse-graining associated to a conditional expectation $E$ that projects to a sub-$\sigma$-algebra. From the classical ergodic theory point of view, this cocycle has no natural interpretation because $t$ here has nothing to do with the Hamiltonian flow. From our point of view, the cocycle above is simply an object whose time derivative captures the increase in entropy. In this sense, $t$ is the conjugate variable to entropy production. 


% % The passage from a reversible microscopic description to an apparently irreversible macroscopic one is often attributed to a loss of information under coarse-graining.  This idea admits a precise classical formulation using conditional expectation, relative entropy, and the commutative form of a relative Connes cocycle.  We will see that it is the failure of conditional expectation to intertwine the microscopic phase action with its coarse-grained counterpart that is of interest to us.  The same relative density that measures this failure generates a cocycle whose expectation is exactly the Kullback--Leibler divergence, and this divergence is exactly the entropy increase produced by coarse-graining.  These results provide a rigorous kinematic account of information loss, but they do not by themselves derive an irreversible physical time evolution.

% % \paragraph{Measure-theoretic setting.}
% % Let $(\Gamma,\Sigma,\Omega)$ be a probability space.  In a statistical-mechanical application, $\Omega$ may be a normalized Liouville measure on a finite accessible region of phase space.  Let
% % \[
% %   \mu\ll\Omega,
% %   \qquad
% %   \rho=\frac{\dd\mu}{\dd\Omega},
% %   \qquad
% %   \int_\Gamma \rho\,\dd\Omega=1.
% % \]
% % For the simplest operator-theoretic presentation, assume for now that $\rho$ is faithful and bounded away from both zero and infinity: there exist constants $0<c<C<\infty$ such that\footnote{This assumption can be weakened to suitable support and integrability conditions, but they eliminate domain issues associated with logarithms and unbounded multiplication operators.}
% % \[
% %   c\leq \rho(x)\leq C
% %   \quad\text{for $\Omega$-almost every }x.
% % \]


% % Set
% % \[
% %   \cH=L^2(\Gamma,\Sigma,\Omega),
% %   \qquad
% %   \mA=L^\infty(\Gamma,\Sigma,\Omega),
% % \]
% % where $\mA$ acts on $\cH$ by multiplication.  The state $\mu$ is represented by the unit vector
% % \[
% %   \xi_\mu=\sqrt{\rho}\in\cH,
% %   \qquad
% %   \|\xi_\mu\|_2^2=\int_\Gamma\rho\,\dd\Omega=1.
% % \]
% % Let $\Sigma_0\subseteq\Sigma$ be the sub-$\sigma$-algebra of macroscopically resolved events, and define
% % \[
% %   Pf:=\E_\Omega[f\mid\Sigma_0].
% % \]
% % On $L^2$, $P$ is the orthogonal projection onto $L^2(\Gamma,\Sigma_0,\Omega)$ \footnote{Note that $P$ acts on functions and not measurable sets.}.  The coarse-grained density and its associated probability measure are
% % \[
% %   \rho_0:=P\rho,
% %   \qquad
% %   \mu_0(A):=\int_A\rho_0\,\dd\Omega,
% %   \quad A\in\Sigma.
% % \]
% % For every $A\in\Sigma_0$, the defining property of conditional expectation gives
% % \[
% %   \mu_0(A)=\int_A P\rho\,\dd\Omega
% %   =\int_A\rho\,\dd\Omega
% %   =\mu(A).
% % \]
% % Thus the fine-grained and coarse-grained states are observationally indistinguishable on the macroscopic algebra, even though they may differ on the full $\sigma$-algebra.

% % \paragraph{Information phase groups and their interpretation.}
% % The self-information associated with $\rho$ is
% % \[
% %   I_\rho(x)=-\log\rho(x).
% % \]
% % It defines the strongly continuous unitary multiplication group
% % \[
% %   U_t^\rho:=\Mop_{\rho^{it}}
% %   =\exp\!\bigl(it\Mop_{\log\rho}\bigr),
% %   \qquad
% %   (U_t^\rho\psi)(x)=e^{it\log\rho(x)}\psi(x).
% % \]
% % Similarly, the coarse-grained density defines
% % \[
% %   U_t^{\rho_0}:=\Mop_{\rho_0^{it}}.
% % \]
% % The parameter $t$ is an auxiliary modular or phase parameter.  Nothing in this construction identifies it with laboratory time.  In particular, $U_t^\rho$ does not transport points of $\Gamma$; it only rotates the phase of an $L^2$ vector pointwise.

% % If $\Gamma$ additionally carries a symplectic structure and $\rho$ is sufficiently smooth, then $-\log\rho$ also defines the Hamiltonian vector field
% % \[
% %   X_{-\log\rho}(f)=\{f,-\log\rho\}.
% % \]
% % The resulting geometric flow uses the Poisson structure and is mathematically distinct from the multiplication group $U_t^\rho$.  Any connection with thermal-time ideas requires this additional geometric structure and cannot be inferred from the measure-theoretic $L^2$ representation alone.


% % \paragraph{Relative modular formulation.}
% % Define the reference trace and the state induced by $\rho$ by
% % \[
% %   \tau(f)=\int_\Gamma f\,\dd\Omega,
% %   \qquad
% %   \varphi_\rho(f)=\int_\Gamma f\rho\,\dd\Omega,
% %   \quad f\in\mA.
% % \]
% % Because $\mA$ is abelian, the modular automorphism group of any faithful normal state acts trivially on $\mA$.  Consequently, $U_t^\rho$ should not be described as a nontrivial Tomita--Takesaki time evolution of the classical observable algebra.  There is, nevertheless, an exact relative modular statement.  Relative to $\tau$, the Connes cocycle derivative is represented by
% % \[
% %   [D\varphi_\rho:D\tau]_t=\rho^{it},
% % \]
% % and for the fine- and coarse-grained states,
% % \[
% %   [D\varphi_\rho:D\varphi_{\rho_0}]_t
% %   =\left(\frac{\rho}{\rho_0}\right)^{it}.
% % \]
% % Since the underlying algebra is abelian, the modular automorphism groups are trivial and this relative cocycle is an ordinary one-parameter unitary group.

% % Define
% % \[
% %   u_t:=U_t^\rho\bigl(U_t^{\rho_0}\bigr)^*
% %   =\Mop_{(\rho/\rho_0)^{it}}.
% % \]
% % Because both $U_t^\rho$ and $U_s^{\rho_0}$ are multiplication operators,
% % \[
% %   U_t^\rho U_s^{\rho_0}=U_s^{\rho_0}U_t^\rho
% %   \qquad\text{for every }s,t\in\mathbb{R}.
% % \]
% % Hence a nontrivial $u_t$ is not evidence that the two phase groups fail to commute.  It records that their generators differ:
% % \[
% %   U_t^\rho=u_tU_t^{\rho_0}=U_t^{\rho_0}u_t.
% % \]
% % Moreover, $u_t=I$ for every $t$ if and only if $\rho=\rho_0$ almost everywhere, equivalently, if and only if the original density is already $\Sigma_0$-measurable.

% % \paragraph{The actual intertwining defect.}
% % The nontrivial relation involves the conditional expectation $P$, not a commutator between the two multiplication groups.

% % \begin{proposition}[Exact intertwining identity]
% % For every $t\in\mathbb{R}$,
% % \[
% %   PU_t^\rho=U_t^{\rho_0}Pu_t.
% % \]
% % Equivalently,
% % \[
% %   PU_t^\rho-U_t^{\rho_0}P
% %   =U_t^{\rho_0}P(u_t-I).
% % \]
% % \end{proposition}

% % \begin{proof}
% % Because $\rho_0^{it}$ is $\Sigma_0$-measurable, conditional expectation satisfies the module property
% % \[
% %   P(gf)=gP(f)
% %   \qquad
% %   \text{for }g\in L^\infty(\Sigma_0).
% % \]
% % Thus, for every $\psi\in\cH$,
% % \begin{align*}
% %   PU_t^\rho\psi
% %   &=P\!\left(\rho^{it}\psi\right)\\
% %   &=P\!\left(\rho_0^{it}
% %        \left(\frac{\rho}{\rho_0}\right)^{it}\psi\right)\\
% %   &=\rho_0^{it}
% %        P\!\left(\left(\frac{\rho}{\rho_0}\right)^{it}\psi\right)\\
% %   &=U_t^{\rho_0}Pu_t\psi.
% % \end{align*}
% % Subtracting $U_t^{\rho_0}P\psi$ yields the second identity.
% % \end{proof}

% % The relative phase $u_t$ therefore controls the failure of ``apply the microscopic information phase and then coarse-grain'' to agree with ``coarse-grain and then apply the effective information phase.''  This is a genuine non-intertwining statement, but it concerns $P$ and $U_t^\rho$, not the pair $U_t^\rho$ and $U_t^{\rho_0}$.

% % \paragraph{Cocycle generator, relative entropy, and entropy increase.}
% % Let
% % \[
% %   K:=\Mop_{\log(\rho/\rho_0)},
% %   \qquad
% %   u_t=e^{itK}.
% % \]
% % Under the standing boundedness assumptions,
% % \[
% %   -i\left.\frac{\dd}{\dd t}u_t\right|_{t=0}=K
% % \]
% % in operator norm.  Taking its expectation in the fine-grained state gives
% % \begin{align*}
% %   \langle\xi_\mu,K\xi_\mu\rangle
% %   &=\int_\Gamma \rho
% %        \log\left(\frac{\rho}{\rho_0}\right)\dd\Omega\\
% %   &=\KLL(\mu\,\|\,\mu_0).
% % \end{align*}
% % Thus the Kullback--Leibler divergence is exactly the state expectation of the generator of the relative cocycle.  This statement is algebraic rather than thermodynamically dynamical: it expresses a fixed information difference through differentiation with respect to the cocycle parameter, not with respect to an independently specified physical time.

% % Define the entropy relative to the chosen reference measure $\Omega$ by
% % \[
% %   S_\Omega(\rho):=-\int_\Gamma\rho\log\rho\,\dd\Omega.
% % \]
% % For continuous phase spaces this differential entropy depends on the reference measure, but the following identity is exact once $\Omega$ is fixed.  Since $\log\rho_0$ is $\Sigma_0$-measurable,
% % \[
% %   \int_\Gamma\rho\log\rho_0\,\dd\Omega
% %   =\int_\Gamma\rho_0\log\rho_0\,\dd\Omega.
% % \]
% % Therefore,
% % \begin{align*}
% %   \KLL(\mu\,\|\,\mu_0)
% %   &=\int_\Gamma\rho\log\rho\,\dd\Omega
% %     -\int_\Gamma\rho\log\rho_0\,\dd\Omega\\
% %   &=\int_\Gamma\rho\log\rho\,\dd\Omega
% %     -\int_\Gamma\rho_0\log\rho_0\,\dd\Omega\\
% %   &=S_\Omega(\rho_0)-S_\Omega(\rho).
% % \end{align*}
% % Because relative entropy is nonnegative,
% % \[
% %   S_\Omega(\rho_0)\geq S_\Omega(\rho),
% % \]
% % with equality if and only if $\rho=\rho_0$ almost everywhere.  The same inequality follows from conditional Jensen applied to the convex function $x\mapsto x\log x$:
% % \[
% %   \rho_0\log\rho_0
% %   =(P\rho)\log(P\rho)
% %   \leq P(\rho\log\rho).
% % \]
% % After integration, this is precisely the entropy inequality above.

% % \paragraph{Meaning and limitation of the result.}
% % The construction establishes four exact facts.  First, $\mu$ and $\mu_0$ agree on every event in $\Sigma_0$.  Second, the information discarded by coarse-graining is
% % \[
% %   \KLL(\mu\,\|\,\mu_0)
% %   =S_\Omega(\rho_0)-S_\Omega(\rho)\geq0.
% % \]
% % Third, the same quantity is encoded by the relative Connes cocycle
% % \[
% %   u_t=[D\varphi_\rho:D\varphi_{\rho_0}]_t
% %   =\left(\frac{\rho}{\rho_0}\right)^{it},
% % \]
% % whose generator has expectation $\KLL(\mu\,\|\,\mu_0)$ in the state $\xi_\mu$.  Fourth, coarse-graining generally fails to intertwine the microscopic and effective phase groups, with exact defect
% % \[
% %   PU_t^\rho-U_t^{\rho_0}P
% %   =U_t^{\rho_0}P(u_t-I).
% % \]

% % These facts do not, by themselves, prove an arrow of physical time.  To discuss irreversibility, one must introduce an independent physical evolution, such as a measure-preserving flow $\Phi_s$ or a Markov semigroup, and then study the time-dependent densities $\rho_s$ and their coarse-grainings $P\rho_s$.  Reversible Hamiltonian evolution preserves fine-grained entropy, whereas the coarse-grained entropy $S_\Omega(P\rho_s)$ need not be monotone without additional assumptions.  Monotonicity may arise from repeated coarse-graining, dissipative or Markovian dynamics, a data-processing structure, suitable mixing limits, or another mechanism that effectively breaks the exact reversibility of the microscopic description.

% % The rigorous conclusion is therefore that conditional expectation produces an exact, nonnegative entropy gap and that a relative modular cocycle encodes this gap algebraically.  This gives a precise kinematic account of information loss under coarse-graining.  A thermodynamic arrow of time requires additional dynamical input.


% % \section{Poisson Modular Flow and the Geometry of Relative Densities}

% % The phrase \emph{classical modular flow} has a precise meaning in Poisson geometry.  It does not refer merely to the pointwise Hilbert-space phase multiplication $\rho^{it}$, nor is it identical to the Tomita--Takesaki modular automorphism group of a commutative von Neumann algebra, whose adjoint action is trivial.  Rather, it is a Poisson-geometric construction introduced by Weinstein: a positive density on a Poisson manifold determines a distinguished Poisson vector field, and the corresponding flow is a classical, first-order analogue of modular dynamics \cite{Weinstein1997,Kosmann2008}.  The relative-flow construction obtained from two densities fits naturally into this theory at the level of modular vector fields and their Hamiltonian representatives.

% % \paragraph{Poisson manifolds and Hamiltonian dynamics.}
% % Let $(M,\pi)$ be a smooth Poisson manifold, with Poisson bracket
% % \[
% % \{f,g\}=\pi(\dd f,\dd g),
% % \qquad f,g\in C^{\infty}(M).
% % \]
% % We use the convention
% % \[
% % X_f(g)=\{g,f\}.
% % \]
% % Thus $X_f$ is the Hamiltonian vector field generated by $f$.  Let $\mu$ be a smooth positive density on $M$.  The divergence of a vector field $Y$ relative to $\mu$ is defined by
% % \[
% % \cL_Y\mu=(\Div_\mu Y)\mu.
% % \]
% % The \emph{modular vector field} associated with $\mu$ is the derivation
% % \[
% % \mathcal{M}_\mu(f):=\Div_\mu(X_f).
% % \]
% % A fundamental result is that $\mathcal{M}_\mu$ is itself a Poisson vector field:
% % \[
% % \cL_{\mathcal{M}_\mu}\pi=0.
% % \]
% % Consequently, whenever it is complete, it integrates to a one-parameter group of Poisson diffeomorphisms.  Without completeness, one obtains a local flow.  This is the Poisson modular flow determined by $\mu$.

% % The choice of density changes the representative but not its Poisson-cohomology class.  If
% % \[
% % \mu'=e^h\mu,
% % \]
% % then, with the convention above,
% % \[
% % \mathcal{M}_{\mu'}=\mathcal{M}_\mu-X_h.
% % \]
% % Hence the class
% % \[
% % \operatorname{mod}(M):=[\mathcal{M}_\mu]\in H^1_\pi(M)
% % \]
% % is independent of $\mu$.  It is called the \emph{modular class} of the Poisson manifold.  Its vanishing is the obstruction-free case: it means that one can choose a density invariant under every Hamiltonian flow.  In this sense the modular class measures the failure of a global Liouville theorem on a general Poisson manifold \cite{Weinstein1997,Kosmann2008}.

% % \paragraph{The symplectic case and the appearance of $-\log\rho$.}
% % Suppose now that $(M,\omega)$ is symplectic and let
% % \[
% % \Omega=\frac{\omega^n}{n!}
% % \]
% % be the Liouville density.  Every Hamiltonian vector field preserves $\Omega$, so
% % \[
% % \mathcal{M}_\Omega=0.
% % \]
% % Given a strictly positive density function $\rho$, set
% % \[
% % \mu_\rho=\rho\Omega,
% % \qquad
% % K_\rho=-\log\rho.
% % \]
% % Using the change-of-density formula with $h=\log\rho$ gives
% % \[
% % \boxed{
% % \mathcal{M}_{\mu_\rho}=X_{-\log\rho}=X_{K_\rho}.
% % }
% % \]
% % Equivalently, for every observable $f$,
% % \[
% % \mathcal{M}_{\mu_\rho}(f)
% % =
% % \{f,-\log\rho\}.
% % \]
% % Thus the vector field discussed earlier is not only a suggestive state-dependent Hamiltonian flow: on a symplectic manifold it is exactly the modular vector field associated with the density $\rho\Omega$.

% % This statement requires a careful interpretation.  Since a symplectic manifold possesses the invariant Liouville density $\Omega$, its modular \emph{class} vanishes.  The nonzero field $X_{-\log\rho}$ is therefore not a new intrinsic cohomology class; it is the density-dependent Hamiltonian representative obtained after replacing $\Omega$ by $\rho\Omega$.  Its flow moves points along level sets of $\rho$ because
% % \[
% % X_{-\log\rho}(\rho)
% % =
% % \{\rho,-\log\rho\}
% % =0.
% % \]
% % It is reversible and symplectic.  It reorganizes points among equally probable microstates but does not by itself create dissipation or thermodynamic irreversibility.

% % \paragraph{Why this is a classical analogue of modular flow.}
% % In a noncommutative von Neumann algebra, a faithful state produces a modular automorphism group whose generator is related formally to the logarithm of the state.  For an abelian algebra such as $L^\infty(M)$, pointwise multiplication by
% % \[
% % U_t^\rho=\rho^{it}=e^{it\log\rho}
% % \]
% % defines a unitary group on an appropriate $L^2$ space, but conjugation is trivial:
% % \[
% % U_t^\rho M_f(U_t^\rho)^{-1}=M_f.
% % \]
% % The phase evolution therefore does not move classical observables or phase-space points.  Its infinitesimal generator is the multiplication operator
% % \[
% % -i\left.\frac{\dd}{\dd t}U_t^\rho\right|_{t=0}
% % =
% % \log\rho.
% % \]

% % Poisson geometry retains the first-order remnant of noncommutativity through the bracket.  The same logarithm now generates the derivation
% % \[
% % \left.\frac{\dd}{\dd t}(f\circ\Phi_t^\rho)\right|_{t=0}
% % =
% % \{f,-\log\rho\},
% % \]
% % where $\Phi_t^\rho$ is the flow of $X_{-\log\rho}$.  The two constructions therefore encode complementary aspects of the same logarithmic data:
% % \[
% % \begin{array}{ccl}
% % \rho^{it} &:& \text{pointwise phase; retains the value of }\log\rho,\\[1mm]
% % \Phi_t^\rho &:& \text{Poisson flow; retains the Hamiltonian derivative of }-\log\rho.
% % \end{array}
% % \]
% % Weinstein's terminology is justified in precisely this semiclassical sense: the Poisson modular flow is the classical limit of modular automorphism theory after the commutator has been replaced by the Poisson bracket \cite{Weinstein1997}.  It should not be confused with the trivial modular automorphism group of the abelian algebra considered without its Poisson structure.

% % \paragraph{Two densities and a relative modular vector field.}
% % Let
% % \[
% % \mu_\rho=\rho\Omega,
% % \qquad
% % \mu_{\rho_0}=\rho_0\Omega,
% % \]
% % where $\rho$ and $\rho_0$ are smooth and strictly positive.  Their modular vector fields are
% % \[
% % \mathcal{M}_{\mu_\rho}=X_{K_\rho},
% % \qquad
% % \mathcal{M}_{\mu_{\rho_0}}=X_{K_{\rho_0}},
% % \]
% % with $K_\rho=-\log\rho$ and $K_{\rho_0}=-\log\rho_0$.  Their difference is
% % \[
% % \boxed{
% % Y_{\rho\mid\rho_0}
% % :=
% % \mathcal{M}_{\mu_\rho}-\mathcal{M}_{\mu_{\rho_0}}
% % =
% % X_{K_\rho-K_{\rho_0}}
% % =
% % X_{-\log(\rho/\rho_0)}.
% % }
% % \]
% % This is an exact and natural relative modular object on the same Poisson manifold.  It is generated by the negative log-likelihood ratio.

% % There is, however, a terminological distinction.  The established phrase \emph{relative modular class} usually refers to a cohomological invariant of a Lie-algebroid morphism or a Poisson map, measuring the obstruction to choosing modular representatives that are related by that map \cite{KosmannWeinstein2005,CaseiroFernandes2013}.  The field $Y_{\rho\mid\rho_0}$ above is instead the difference of two modular representatives associated with two densities on the same manifold.  It is closely analogous and fits naturally into the same geometry, but it should not automatically be identified with the relative modular class of a Poisson map.

% % \paragraph{The relative echo map.}
% % Let $\Phi_t^\rho$ and $\Phi_t^{\rho_0}$ be the modular flows generated by $X_{K_\rho}$ and $X_{K_{\rho_0}}$.  Define
% % \[
% % R_t
% % =
% % (\Phi_t^{\rho_0})^{-1}\circ\Phi_t^\rho
% % =
% % \Phi_{-t}^{\rho_0}\circ\Phi_t^\rho.
% % \]
% % The map $R_t$ compares the motion generated by the microscopic density with the motion generated by the reference or coarse-grained density.  It is a symplectic diffeomorphism for every $t$ for which both flows exist.  Its infinitesimal generator at the identity is
% % \[
% % \boxed{
% % \left.\frac{\dd}{\dd t}R_t\right|_{t=0}
% % =
% % X_{K_\rho}-X_{K_{\rho_0}}
% % =
% % Y_{\rho\mid\rho_0}.
% % }
% % \]
% % The equality is an equality of tangent vectors to the diffeomorphism group, hence of vector fields on $M$.  The scalar Hamiltonian generating that tangent vector is
% % \[
% % K_{\mathrm{rel}}
% % =
% % K_\rho-K_{\rho_0}
% % =
% % -\log\frac{\rho}{\rho_0}.
% % \]

% % In general, $R_t$ is not itself a one-parameter group, because the two modular vector fields need not commute.  It is more accurately a \emph{relative modular isotopy} or an echo map.  Its right-trivialized time-dependent generator is
% % \[
% % \dot R_t\circ R_t^{-1}
% % =
% % X_{H_t},
% % \qquad
% % H_t
% % =
% % K_\rho\circ\Phi_t^{\rho_0}-K_{\rho_0}.
% % \]
% % Only in the commuting case does the finite relative transformation reduce to the autonomous flow
% % \[
% % R_t=\exp\bigl(tX_{K_{\mathrm{rel}}}\bigr).
% % \]
% % Thus the log-density ratio controls the first-order discrepancy exactly, while higher-time behavior also depends on the noncommuting geometry of the two modular flows.

% % \paragraph{Relative entropy as the mean relative modular Hamiltonian.}
% % The Kullback--Leibler divergence is
% % \[
% % \KLL(\rho\Vert\rho_0)
% % =
% % \int_M \rho\log\frac{\rho}{\rho_0}\,\Omega.
% % \]
% % Since $K_{\mathrm{rel}}=-\log(\rho/\rho_0)$,
% % \[
% % \boxed{
% % \KLL(\rho\Vert\rho_0)
% % =
% % -\int_M K_{\mathrm{rel}}\,\dd\mu_\rho
% % =
% % -\mathbb E_{\mu_\rho}[K_{\mathrm{rel}}].
% % }
% % \]
% % This is the clean bridge between the relative modular geometry and information theory.  The infinitesimal relative map determines the Hamiltonian vector field $X_{K_{\mathrm{rel}}}$, while relative entropy is the negative expectation of its Hamiltonian potential.

% % The distinction between the vector field and its Hamiltonian is essential.  On a connected symplectic manifold, $X_{K_{\mathrm{rel}}}$ determines $K_{\mathrm{rel}}$ only up to an additive constant.  The density relation
% % \[
% % \rho=\rho_0e^{-K_{\mathrm{rel}}}
% % \]
% % and the normalization
% % \[
% % \int_M\rho_0e^{-K_{\mathrm{rel}}}\Omega=1
% % \]
% % fix that constant.  On a general Poisson manifold the ambiguity is larger: two Hamiltonians generate the same vector field when they differ by a Casimir, that is, a function constant on every symplectic leaf.  Reconstructing the full KL divergence from the relative vector field alone then also requires the missing transverse or leafwise normalization data.

% % The multiplicative relative phase
% % \[
% % u_t
% % =
% % \left(\frac{\rho}{\rho_0}\right)^{it}
% % \]
% % retains the scalar generator directly:
% % \[
% % -i\left.\frac{\dd}{\dd t}u_t\right|_{t=0}
% % =
% % \log\frac{\rho}{\rho_0},
% % \]
% % and therefore
% % \[
% % \KLL(\rho\Vert\rho_0)
% % =
% % \mathbb E_{\mu_\rho}\!\left[
% % -i\left.\frac{\dd}{\dd t}u_t\right|_{t=0}
% % \right].
% % \]
% % The geometric relative flow instead retains the Poisson derivative of the same generator:
% % \[
% % \left.\frac{\dd}{\dd t}R_t\right|_{t=0}
% % =
% % X_{-\log(\rho/\rho_0)}.
% % \]
% % This explains why the phase cocycle is the more direct carrier of the numerical KL divergence, whereas the Poisson relative flow is the carrier of its Hamiltonian geometry.

% % \paragraph{A rate identity and its limitation.}
% % Introduce a second parameter $\tau$ describing actual physical evolution.  Suppose that $\rho_\tau$ is transported by a Hamiltonian $G_\tau$:
% % \[
% % \partial_\tau\rho_\tau+\{\rho_\tau,G_\tau\}=0,
% % \]
% % and let $\rho_0$ be fixed.  Define
% % \[
% % Y_\tau
% % =
% % X_{-\log(\rho_\tau/\rho_0)}
% % =
% % \left.\frac{\partial}{\partial t}R_t^{(\tau)}\right|_{t=0}.
% % \]
% % Under the usual integrability and boundary assumptions, differentiation and Hamiltonian integration by parts give
% % \[
% % \boxed{
% % \frac{\dd}{\dd\tau}
% % \KLL(\rho_\tau\Vert\rho_0)
% % =
% % \int_M \rho_\tau\,Y_\tau(G_\tau)\,\Omega.
% % }
% % \]
% % Equivalently, if $\mathcal R_t^{(\tau)}f=f\circ R_t^{(\tau)}$, then
% % \[
% % \frac{\dd}{\dd\tau}
% % \KLL(\rho_\tau\Vert\rho_0)
% % =
% % \int_M\rho_\tau
% % \left.
% % \frac{\partial}{\partial t}
% % \mathcal R_t^{(\tau)}G_\tau
% % \right|_{t=0}
% % \Omega.
% % \]
% % The rate of relative entropy is therefore the mean action of the infinitesimal relative modular flow on the physical Hamiltonian.

% % If $\rho_{0,\tau}=P\rho_\tau$ is obtained from a fixed conditional expectation $P$ onto a coarse sub-$\sigma$-algebra, the additional derivative of the reference density cancels under the conditional-expectation identity, provided the relevant functions are integrable and $\rho_{0,\tau}>0$.  One then obtains the same formula with $\rho_0$ replaced by $\rho_{0,\tau}$.  Moreover,
% % \[
% % \KLL(\rho_\tau\Vert\rho_{0,\tau})
% % =
% % S(\rho_{0,\tau})-S(\rho_\tau).
% % \]
% % For Hamiltonian physical evolution $S(\rho_\tau)$ is constant, so the KL rate equals the rate of change of the coarse-grained entropy.

% % This identity is kinematic, not a second-law theorem.  The right-hand side has no definite sign for a generic reversible Hamiltonian evolution.  In particular, if the physical Hamiltonian coincides with the relative modular Hamiltonian, then
% % \[
% % Y_\tau(G_\tau)
% % =
% % \{G_\tau,G_\tau\}=0.
% % \]
% % A nonnegative entropy-production law requires extra irreversible structure, such as a Markov semigroup, repeated projection, stochastic noise, coupling to reservoirs, or dissipative phase-space contraction.  Poisson modular geometry supplies a rigorous language connecting densities, relative logarithms, and infinitesimal flows; it does not by itself generate an arrow of time.

% % \paragraph{Example: two Gibbs densities.}
% % Let
% % \[
% % \rho_\beta=Z_\beta^{-1}e^{-\beta H},
% % \qquad
% % \rho_{\beta_0}=Z_{\beta_0}^{-1}e^{-\beta_0H}.
% % \]
% % Then
% % \[
% % K_{\rho_\beta}=\beta H+\log Z_\beta,
% % \qquad
% % \mathcal{M}_{\rho_\beta\Omega}=\beta X_H.
% % \]
% % The relative modular vector field is
% % \[
% % Y_{\beta\mid\beta_0}
% % =(\beta-\beta_0)X_H,
% % \]
% % and the two flows commute, so
% % \[
% % R_t=\Phi_H^{(\beta-\beta_0)t}.
% % \]
% % The relative map sees only the difference in dynamical clock rates.  By contrast,
% % \[
% % \KLL(\rho_\beta\Vert\rho_{\beta_0})
% % =
% % (\beta_0-\beta)\langle H\rangle_\beta
% % +
% % \log\frac{Z_{\beta_0}}{Z_\beta},
% % \]
% % which also contains the partition-function term.  That term is an additive constant in the relative Hamiltonian and is therefore invisible to the vector field.  This example displays the general principle: the relative modular flow captures the geometry of the statistical mismatch, while the KL divergence additionally remembers its normalization.

% % \paragraph{Conclusion.}
% % The established Poisson modular-flow theory supports the following precise synthesis:
% % \[
% % \boxed{
% % \begin{gathered}
% % \rho\Omega
% % \longmapsto
% % \mathcal{M}_{\rho\Omega}=X_{-\log\rho},\\[1mm]
% % (\rho,\rho_0)
% % \longmapsto
% % Y_{\rho\mid\rho_0}=X_{-\log(\rho/\rho_0)},\\[1mm]
% % R_t=(\Phi_t^{\rho_0})^{-1}\circ\Phi_t^\rho,
% % \qquad
% % \dot R_0=Y_{\rho\mid\rho_0},\\[1mm]
% % \KLL(\rho\Vert\rho_0)
% % =-\mathbb E_{\rho}[K_{\mathrm{rel}}],
% % \qquad
% % K_{\mathrm{rel}}=-\log(\rho/\rho_0).
% % \end{gathered}
% % }
% % \]
% % The first line is standard Poisson modular geometry.  The second and third lines give a natural relative modular representative and its finite echo construction.  The fourth line connects the Hamiltonian potential of that representative to relative entropy.  Together they provide a rigorous classical analogue of modular dynamics, while clearly separating reversible Poisson motion from genuine irreversible entropy production.

% \section{Coarse-graining entropy versus Kolmogorov-Sinai entropy}

% Kolmogorov-Sinai involves an infimum over various coarse-grainings...


\section{Modular theory for standard subspaces}


Let
\[
    \mathcal{A}=L^\infty(X,\mu)
\]
act by multiplication on
\[
    \cH=L^2(X,\mu;\mathbb{C}).
\]
The intrinsic Tomita--Takesaki modular dynamics of the abelian algebra \(\mathcal{A}\), relative to any faithful normal state, is trivial. Nevertheless, the same complex Hilbert space \(\cH\) can carry nontrivial standard real subspaces, and each such subspace has its own modular operator. This produces a meaningful nontrivial unitary dynamics on \(\cH\). The essential point is that this dynamics belongs to the enriched pair \((\cH,H)\), not to the abelian algebra \(\mathcal{A}\) by itself.

\paragraph{Standard real subspaces.}
A closed real-linear subspace \(H\subset\cH\) is called \emph{standard} when
\[
    H\mathcal{A}p iH=\{0\},
    \qquad
    \overline{H+iH}=\cH.
\]
On the dense domain \(H+iH\), one defines the closed antilinear Tomita operator
\[
    S_H(h_1+ih_2)=h_1-ih_2,
    \qquad h_1,h_2\in H.
\]
Its polar decomposition is
\[
    S_H=J_H\Delta_H^{1/2},
\]
where \(J_H\) is antiunitary and \(\Delta_H\) is a positive self-adjoint operator. The associated one-particle modular flow is the unitary group
\[
    t\longmapsto \Delta_H^{it}.
\]
This is the real-subspace form of modular theory. Under bosonic second quantization, the real subspace \(H\) determines the Weyl von Neumann algebra
\[
    \mathcal{R}(H)=\{W(h):h\in H\}'',
\]
and the one-particle modular data second-quantize to the ordinary Tomita--Takesaki modular data of \(\mathcal{R}(H)\).

\paragraph{The canonical abelian subspace gives no new flow.}
To see what the intrinsic abelian theory produces, represent a faithful state by a strictly positive density \(\rho\) and the vector
\[
    \Omega_\rho=\sqrt{\rho}.
\]
The canonical real subspace associated with \((\mathcal{A},\Omega_\rho)\) is
\[
    H_{\mathcal{A},\rho}
    =\overline{\{a\Omega_\rho:a=a^*\in\mathcal{A}\}}.
\]
After passing to the equivalent GNS realization \(L^2(X,\rho\,d\mu)\), with cyclic vector \(1\), this becomes
\[
    H_{\mathcal{A},\rho}=L^2(X,\rho\,d\mu;\mathbb{R}).
\]
Its Tomita operator is complex conjugation,
\[
    S_{H_{\mathcal{A},\rho}}\psi=\overline{\psi},
\]
so that
\[
    J_{H_{\mathcal{A},\rho}}=C,
    \qquad
    \Delta_{H_{\mathcal{A},\rho}}=I,
    \qquad
    \Delta_{H_{\mathcal{A},\rho}}^{it}=I.
\]
This agrees exactly with ordinary modular theory. Since \(\mathcal{A}\) is abelian, every faithful state on \(\mathcal{A}\) is tracial in the algebraic sense,
\[
    \varphi(ab)=\varphi(ba),
\]
and therefore its modular automorphism group acts trivially on \(\mathcal{A}\).

\begin{proposition}
If the standard subspace is the canonical real subspace obtained from an abelian von Neumann algebra and a faithful cyclic separating vector, then the standard-subspace modular operator is the identity.
\end{proposition}

Thus Longo's one-particle formalism does not alter the intrinsic modular theory of an abelian algebra: when applied canonically, it reproduces the same trivial flow.

\paragraph{Nontrivial standard subspaces in \(L^2(X,\mu)\).}
The Hilbert space \(L^2(X,\mu)\) contains many standard real subspaces other than the canonical real form \(L^2(X,\mu;\mathbb{R})\). Choose an independently specified standard subspace
\[
    H\subset L^2(X,\mu;\mathbb{C})
\]
with
\[
    \Delta_H\neq I.
\]
Then \(\Delta_H^{it}\) is a nontrivial modular flow on the one-particle Hilbert space. This construction is mathematically legitimate, but it uses additional data:
\[
    (X,\mu)
    \quad\rightsquigarrow\quad
    \bigl(L^2(X,\mu),H\bigr).
\]
The measure space determines the complex Hilbert space and its multiplication algebra, whereas the real subspace \(H\) supplies an extra real-symplectic or localization structure. The nontrivial modular operator is determined by that added structure.

Consequently, it is inaccurate to call \(\Delta_H^{it}\) the intrinsic modular flow of \(L^\infty(X,\mu)\). A more precise description is
\[
    \boxed{\begin{gathered}
    \text{standard-subspace modular dynamics carried by an }L^2\\
    \text{representation of an abelian algebra}.
    \end{gathered}}
\]

\paragraph{When does the flow act on the classical observables?}
A unitary group on \(\cH\) need not preserve the multiplication algebra. To obtain an actual flow of the classical observables one must impose
\[
    \Delta_H^{it}\mathcal{A}\Delta_H^{-it}=\mathcal{A}
    \qquad (t\in\mathbb{R}).
\]
When this normalization condition holds, conjugation defines a one-parameter automorphism group
\[
    \alpha_t(a)=\Delta_H^{it}a\Delta_H^{-it},
    \qquad a\in\mathcal{A}.
\]
Automorphisms of an abelian von Neumann algebra are implemented, modulo standard measure-theoretic qualifications, by nonsingular measurable transformations \(T_t:X\to X\). Hence one has schematically
\[
    \Delta_H^{it}M_f\Delta_H^{-it}
    =M_{f\circ T_{-t}},
\]
where \(M_f\) denotes multiplication by \(f\). In such a compatible system, the standard-subspace modular group induces a genuine Koopman-type classical dynamics.

This does not contradict the triviality of intrinsic abelian modular theory. The measurable flow \(T_t\) is selected by the extra standard subspace \(H\), not by the pair \((\mathcal{A},\varphi)\) alone. Without the normalization condition, \(\Delta_H^{it}\) may carry multiplication operators into non-multiplication operators and therefore leave the classical observable algebra altogether.

\paragraph{Why nontriviality signals hidden symplectic structure.}
For a real subspace \(H\), define its symplectic complement by
\[
    H'=\{\xi\in\cH:\operatorname{Im}\langle\xi,h\rangle=0
    \text{ for all }h\in H\}.
\]
If \(H\) is symplectically isotropic, so that \(H\subset H'\), and is also standard, then \(H''=H\) implies \(H=H'\). In this abelian component the modular operator is trivial. Therefore a nontrivial standard-subspace modular operator requires a nontrivial symplectic relation between vectors in \(H\).

The same fact appears after second quantization. Weyl operators satisfy
\[
    W(h)W(k)
    =e^{-\frac{i}{2}\operatorname{Im}\langle h,k\rangle}W(h+k).
\]
If the symplectic form is nonzero on \(H\), then the algebra \(\mathcal{R}(H)\) is noncommutative. Thus
\begin{center}
\fbox{\parbox{0.84\linewidth}{\centering
A nontrivial modular operator \(\Delta_H\neq I\) typically reflects
symplectic structure that becomes noncommutativity after second quantization.}}
\end{center}
The original multiplication algebra may remain abelian, but the additional modular geometry is no longer encoded by that algebra alone.

\paragraph{Relation to coarse-graining.}
Classical coarse-graining associated with a sub-\(\sigma\)-algebra \(\Sigma_0\subset\Sigma\) is represented by the conditional-expectation projection
\[
    P=E[\,\cdot\mid\Sigma_0]:L^2(X,\Sigma,\mu)\to
    L^2(X,\Sigma_0,\mu).
\]
This is a complex-linear orthogonal projection. Its range is a complex Hilbert subspace and is therefore not a standard real subspace of the original \(L^2\)-space. The projection is naturally interpreted as the abelian version of a Jones projection implementing a conditional expectation, rather than as a standard-subspace cutting projection.

One may nevertheless enrich the coarse-graining setup by specifying a family of standard subspaces \(H_0\subset H\), or by requiring compatibility relations between \(P\) and the modular dynamics, for example
\[
    P\Delta_H^{it}=\Delta_{H_0}^{it}P
\]
or studying the failure of this relation. Such constructions can compare localization, information loss, and modular dynamics, but they involve more than classical conditional expectation alone.

\paragraph{Relative modular data remain nontrivial even in the abelian theory.}
Although the ordinary modular automorphism group of \(\mathcal{A}\) is trivial, comparison of two faithful densities \(\rho\) and \(\rho_0\) gives the relative modular operator
\[
    \Delta_{\rho\mid\rho_0}=M_{\rho/\rho_0},
\]
up to the conventional order of the two states. Its unitary powers act by pointwise phase multiplication,
\[
    \Delta_{\rho\mid\rho_0}^{it}\psi
    =\left(\frac{\rho}{\rho_0}\right)^{it}\psi.
\]
This action is nontrivial on Hilbert-space vectors, even though its adjoint action on the multiplication algebra remains trivial. The logarithm of the relative modular operator contains the classical log-likelihood ratio,
\[
    \log\Delta_{\rho\mid\rho_0}
    =M_{\log(\rho/\rho_0)},
\]
and its expectation yields the Kullback--Leibler divergence,
\[
    D_{\mathrm{KL}}(\rho\Vert\rho_0)
    =\left\langle\sqrt\rho,
    \log\Delta_{\rho\mid\rho_0}\sqrt\rho\right\rangle.
\]
Thus the canonical abelian theory has trivial absolute modular flow but nontrivial relative modular data.

\paragraph{A precise enriched classical framework.}
A potentially useful construction is an enriched quadruple
\[
    \bigl(X,\mu,\mathcal{A},H\bigr),
    \qquad
    \mathcal{A}=L^\infty(X,\mu),
    \quad
    H\subset L^2(X,\mu)
    \text{ standard},
\]
subject to the compatibility condition
\[
    \Delta_H^{it}\mathcal{A}\Delta_H^{-it}=\mathcal{A}.
\]
Then
\begin{enumerate}
    \item \(\mathcal{A}\) supplies the abelian classical observables;
    \item \(H\) supplies additional real-symplectic or localization data;
    \item \(\Delta_H^{it}\) supplies a distinguished one-parameter unitary group;
    \item normalization of \(\mathcal{A}\) turns this group into a measurable classical flow;
    \item inclusions of standard subspaces can produce relative modular cocycles;
    \item conditional expectations can be added independently to model coarse-graining.
\end{enumerate}
This may support a meaningful modular formulation of classical dynamics, but its modular flow is not intrinsic to the abelian algebra. It is modular dynamics of an enriched structure compatible with that algebra.

\subsection*{Standard Subspace Analog of an Inclusion of von Neumann Algebras}

Let
\[
\mathcal{N}\subset \mathcal{M}\subset B(\cH)
\]
be an inclusion of von Neumann algebras, and suppose that a vector
$\Omega\in\cH$ is cyclic and separating for both algebras.  The natural
one-particle objects associated with the two pairs $(\mathcal{N},\Omega)$ and
$(\mathcal{M},\Omega)$ are the closed real subspaces
\[
H_{\mathcal{N},\Omega}
   :=\overline{\mathcal{N}_{\mathrm{sa}}\Omega},
\qquad
H_{\mathcal{M},\Omega}
   :=\overline{\mathcal{M}_{\mathrm{sa}}\Omega}.
\]
Because $\mathcal{N}\subset\mathcal{M}$, one has
\[
H_{\mathcal{N},\Omega}\subset H_{\mathcal{M},\Omega}.
\]
Cyclicity and separatingness of $\Omega$ imply that both real subspaces
are standard:
\[
H\cap iH=\{0\},
\qquad
\overline{H+iH}=\cH.
\]
Thus the direct one-particle analog of an algebra inclusion with a common
cyclic and separating vector is an inclusion
\[
\boxed{H_1\subset H_2\subset\cH}
\]
in which both $H_1$ and $H_2$ are standard in the same ambient complex
Hilbert space.

The common vector $\Omega$ is not displayed separately after passing to
the real subspaces.  Its modular information is encoded in the real
embeddings of $H_1$ and $H_2$ in $\cH$.  Conversely, after Bose second
quantization, an inclusion of standard subspaces gives an inclusion of
Weyl von Neumann algebras
\[
\mathcal{R}(H_1)\subset\mathcal{R}(H_2),
\]
with the Fock vacuum cyclic and separating for both algebras.

\paragraph{Relative commutants.}
For a real subspace $H\subset\cH$, its symplectic complement is
\[
H'
   :=\left\{\xi\in\cH:
      \operatorname{Im}\langle \xi,h\rangle=0
      \text{ for all }h\in H
      \right\}.
\]
Accordingly, the geometric analog of the relative commutant
$\mathcal{N}'\cap\mathcal{M}$ is
\[
\boxed{H_1'\cap H_2.}
\]
The standardness of $H_1'\cap H_2$ is an additional condition; it does
not follow merely from $H_1\subset H_2$.  It records how much localization
or independent symplectic information remains between the two standard
subspaces.

\paragraph{Tomita and modular data.}
For each standard subspace $H_j$, define the Tomita operator
\[
S_j(h+ik)=h-ik,
\qquad h,k\in H_j,
\]
with polar decomposition
\[
S_j=J_j\Delta_j^{1/2}.
\]
Here $J_j$ is an antiunitary involution and $\Delta_j$ is a positive,
injective self-adjoint operator.  The one-particle modular flow is
\[
t\longmapsto \Delta_j^{it}.
\]

\subsection*{The relative modular cocycle}

The natural operator comparing the two modular flows is
\[
\boxed{
   D_{H_1\subset H_2}(t):=\Delta_2^{-it}\Delta_1^{it}.
}
\]
We write $D(t)=D_{H_1\subset H_2}(t)$ below. This family is strongly
continuous and unitary. It composes modular evolution with $H_1$ and
reversed modular evolution with $H_2$.  The residual unitary therefore measures the mismatch between the
two modular notions of localization.

Define the modular automorphism groups of $B(\cH)$ by
\[
\alpha_t^{(j)}(A)
   :=\Delta_j^{it}A\Delta_j^{-it}.
\]
Then $D(t)$ satisfies the cocycle identity
\[
\boxed{
   D(t+s)=D(t)\,\alpha_{-t}^{(1)}(D(s)).
}
\]
Indeed,
\begin{align*}
 D(t)\alpha_{-t}^{(1)}(D(s))
 &=\Delta_2^{-it}\Delta_1^{it}
   \Delta_1^{-it}
   \bigl(\Delta_2^{-is}\Delta_1^{is}\bigr)
   \Delta_1^{it} \\
 &=\Delta_2^{-i(t+s)}\Delta_1^{i(t+s)}
 =D(t+s).
\end{align*}
It also intertwines the two modular dynamics:
\[
\boxed{
   \alpha_{-t}^{(2)}
   =\text{Ad}(D(t))\circ\alpha_{-t}^{(1)}.
}
\]
Thus $D(t)$ is a genuine relative modular intertwining cocycle.

In general, it is not a one-parameter group.  One has
\[
 D(t+s)\neq D(t)D(s)
\]
unless the modular operators strongly commute.  If they do commute, then
\[
 D(t)
 =\exp\!\left[it\bigl(\log\Delta_1-\log\Delta_2\bigr)\right].
\]
For a generic inclusion, however, the cocycle law rather than the group
law is the correct structure.

\paragraph{Infinitesimal relative modular generator.}
On the common domain
\[
\Dom(\log\Delta_1)\cap\Dom(\log\Delta_2),
\]
or more generally in the quadratic-form sense,
\[
\boxed{
   \left.\frac{d}{dt}D(t)\right|_{t=0}
   =i\bigl(\log\Delta_1-\log\Delta_2\bigr).
}
\]
The formal relative modular Hamiltonian is therefore
\[
K_{\mathrm{rel}}
   :=\log\Delta_1-\log\Delta_2.
\]
For a vector $\xi$ in the appropriate form domain,
\[
-i\left.\frac{d}{dt}
  \langle\xi,D(t)\xi\rangle\right|_{t=0}
=
\langle\xi,K_{\mathrm{rel}}\xi\rangle.
\]
This is a relative modular-energy expectation.  It is not automatically a
relative entropy: an entropy requires additional state or vector data and,
in standard-subspace formulations, generally also the cutting projection
that encodes localization.

\paragraph{Relation to a Connes cocycle.}
The expression $\Delta_2^{-it}\Delta_1^{it}$ resembles a Connes cocycle,
but the two notions should not be identified without further hypotheses.
A Connes cocycle compares two faithful states or weights on the same von
Neumann algebra and is affiliated with, or belongs to, that algebra.  The
operator $D(t)$ above compares the modular structures of two real
subspaces and a priori lies only in $B(\cH)$.  It is therefore more precise
to call it a relative modular intertwining cocycle.  After second
quantization one obtains the unitary $\Gamma(D(t))$ on Fock space, but it
becomes a genuine Connes cocycle only when suitable algebraic and state
structures are supplied.

\subsection*{Half-sided modular inclusions}

The relative cocycle acquires much stronger meaning when the inclusion is
half-sided modular.  With one common sign convention, this means
\[
H_1\subset H_2,
\qquad
\Delta_2^{-it}H_1\subset H_1
\quad (t\geq 0).
\]
Under the standard hypotheses, the operator
\[
\boxed{
P
   :=\frac{1}{2\pi}
      \overline{\log\Delta_1-\log\Delta_2}
   \geq 0
}
\]
is the positive generator of a translation group
\[
U(a)=e^{iaP}.
\]
After normalization, the translation satisfies
\[
U(1)H_2=H_1,
\]
and the modular group of $H_2$ acts by dilations on $U$:
\[
\Delta_2^{it}U(a)\Delta_2^{-it}
   =U(e^{-2\pi t}a),
\]
up to the convention used for the modular parameter and the orientation
of the inclusion.  The modular and translation groups therefore generate
a positive-energy representation of the affine group of the line.

In this setting the derivative of the relative modular cocycle is no
longer an arbitrary formal difference:
\[
\left.\frac{d}{dt}D(t)\right|_{t=0}
   =2\pi iP
\]
on a suitable common core.  The mismatch between the two modular clocks
organizes itself into a positive translation generator.  This is the
one-particle form of the Borchers--Wiesbrock mechanism.

\subsection*{Summary dictionary}

\renewcommand{\arraystretch}{1.45}
\setlength{\tabcolsep}{4pt}
\begin{center}
\begin{tabular}{>{\raggedright\arraybackslash}p{0.40\textwidth}
                >{\centering\arraybackslash}p{0.07\textwidth}
                >{\raggedright\arraybackslash}p{0.40\textwidth}}
\toprule
\textbf{von Neumann algebra picture}
& &
\textbf{standard-subspace picture}\\
\midrule
$\mathcal{N}\subset\mathcal{M}$ with a common cyclic and separating vector
& $\longleftrightarrow$ &
$H_1\subset H_2$, both standard in the same $\cH$\\
$\mathcal{N}'\cap\mathcal{M}$
& $\longleftrightarrow$ &
$H_1'\cap H_2$\\
Modular groups of the two inclusions
& $\longleftrightarrow$ &
$\Delta_1^{it}$ and $\Delta_2^{it}$\\
Relative modular comparison
& $\longleftrightarrow$ &
$D_{H_1\subset H_2}(t)=\Delta_2^{-it}\Delta_1^{it}$\\
Infinitesimal relative generator
& $\longleftrightarrow$ &
$\log\Delta_1-\log\Delta_2$\\
Half-sided modular inclusion
& $\longleftrightarrow$ &
$\Delta_2^{-it}H_1\subset H_1$ for $t\geq0$\\
Positive translation generator
& $\longleftrightarrow$ &
$\frac{1}{2\pi}\overline{\log\Delta_1-\log\Delta_2}$\\
\bottomrule
\end{tabular}
\end{center}

The principal conclusion is therefore
\[
\boxed{
\begin{minipage}{0.88\textwidth}
The one-particle analog of an inclusion of von Neumann algebras with a
common cyclic and separating vector is an inclusion of two standard real
subspaces in one complex Hilbert space.  The direct analog of the relative
modular comparison is the cocycle
$D_{H_1\subset H_2}(t)=\Delta_{H_2}^{-it}\Delta_{H_1}^{it}$.  For a generic inclusion its
derivative gives a relative modular-energy generator.  For a half-sided
modular inclusion, the same difference of modular Hamiltonians closes to
a positive translation generator and yields a representation of the
affine group.
\end{minipage}
}
\]


 \section{Extras (old stuff)}



\section{Quantum coarse-grainings and modular dynamics}\label{app:quantum-coarse}
\subsection{Modular dynamics and conditional expectations}\label{subsec:quantum-ergodic}

Appendix~\ref{subsec:classical-ergodic-hierarchy} distinguishes ergodicity, mixing,
autonomous subsystems, and the $K$-property. Their quantum comparison uses modular
dynamics of a chosen faithful normal state. Takesaki's theorem supplies an additional
relation between autonomy and state-preserving expectations. Common cyclicity and
the split property give further algebraic structure, but neither is a definition of mixing.


\paragraph{State-preserving expectations and recovery.}
Let $\mathcal N\subset\mathcal M$ be an inclusion of factors and
$\mathcal R=\mathcal N'\cap\mathcal M$ its relative commutant.
A fixed faithful normal expectation $E:\mathcal M\to\mathcal N$ gives the extension
$\omega_{\mathcal M}\circ E$ of the retained state. Thus
\begin{equation}
 S(\omega_{\mathcal M}\|\omega_{\mathcal M}\circ E)=0
 \quad\Longleftrightarrow\quad \omega_{\mathcal M}=\omega_{\mathcal M}\circ E.
 \label{eq:QEC-entropy}
\end{equation}
This is equality of a state with its specified extension. Exact error correction
requires a channel, a recovery map, and a family of states or observables on which
their composition is the identity; a single preserved state does not establish a code.
For a tensor-product erasure and a fixed hidden state, appending that state
recovers the family of product states, as in \eqref{condexp}.

Takesaki's theorem makes existence of an $\omega_{\mathcal M}$-preserving normal
expectation equivalent to
\begin{equation}
 \sigma_t^{\omega_{\mathcal M}}(\mathcal N)=\mathcal N,\qquad t\in\mathbb R.
 \label{eq:Takesaki-QEC}
\end{equation}
This expresses autonomous modular dynamics, not that every observable of $\mathcal N$
is a conserved charge. The flows intertwine on the retained algebra:
\begin{equation}
 \sigma_t^{\omega_{\mathcal M}}(n)=\sigma_t^{\omega_{\mathcal N}}(n),
 \qquad n\in\mathcal N,\quad\omega_{\mathcal N}=\omega_{\mathcal M}|_{\mathcal N}.
 \label{eq:modular-intertwining}
\end{equation}
The intrinsic modular operator on $\overline{\mathcal N\Omega}$ is understood
through its isometric inclusion into the full GNS space.
For a tracial reference state, the trace-preserving expectation has this property.
Arbitrary expectations need not preserve the trace. Neither a Jones projection
nor a product of modular conjugations is, by itself, a Petz recovery channel.


For a faithful normal expectation onto a factor, the bimodule property gives
a faithful normal state $\chi$ on the relative commutant such that
\begin{equation}
 E(nr)=\chi(r)n,\qquad n\in\mathcal N,\quad r\in\mathcal R.
 \label{eq:CE-state-correspondence}
\end{equation}
For a spatial tensor factorization $\mathcal M=\mathcal N\overline\otimes\mathcal R$,
this is the slice expectation, and every choice of $\chi$ gives one.
Without such a factorization, the displayed slice formula only specifies the
expectation on $\mathcal N\vee\mathcal R$.
Finite index forces a finite-dimensional relative commutant but neither a single
matrix block nor conormality. Infinite index does not force the relative commutant
to be infinite-dimensional.
If an algebra has a nontrivial center, those central observables are fixed by every
modular flow. This is the obstruction to modular ergodicity used in Section~\ref{section:basics}.


\paragraph{Product states, autonomy, and mixing.}
In the spatial tensor-product setting, a faithful state
\begin{equation}
 \omega_{\mathcal M}=\omega_{\mathcal N}\otimes\omega_{\mathcal R}
 \label{eq:product-state}
\end{equation}
has autonomous modular dynamics on $\mathcal N$ and is preserved by
$E_{\omega_{\mathcal R}}$. For a different fixed hidden reference, its
coarse-graining entropy can still be positive; see
\eqref{eq:quantum-coarse-entropy-decomposition}.
Without conormality, correlations with the relative commutant alone do not test autonomy.
Locality gives commuting algebras, while normal product states require additional
spatial tensor-product structure. Strong additivity does not itself produce an expectation.

For a common cyclic and separating vector, write the relative observable comparison as
\begin{equation}
 \gamma_t=\sigma_{-t}^{\omega_{\mathcal M}}\circ\sigma_t^{\omega_{\mathcal N}}
 \quad\hbox{on }\mathcal N.
 \label{eq:coarse-graining-flow}
\end{equation}
It need not be a one-parameter group or remain within $\mathcal N$.
Modular mixing instead tests connected correlations:
\begin{equation}
 C_{a,b}^{\omega}(t)=\omega_{\mathcal M}(a^*\sigma_t^{\omega_{\mathcal M}}(b))
 -\omega_{\mathcal M}(a^*)\omega_{\mathcal M}(b),\qquad a,b\in\mathcal M.
 \label{eq:quantum-correlation}
\end{equation}
Mixing requires their decay for every pair. The separate generating condition is
\begin{equation}
 \bigvee_{t\in\mathbb R}\sigma_t^{\omega_{\mathcal M}}(\mathcal N)=\mathcal M.
 \label{eq:topological-transitivity}
\end{equation}
Common cyclicity implies this equality by Lemma~\ref{lem:common-cyclicity-generating};
it does not supply a decay rate or a trivial one-sided tail.

\paragraph{Common cyclic and separating vectors.}
A vector $|\Omega\rangle\in\mathcal{H}$ is \emph{cyclic and separating} for a von Neumann
algebra $\mathcal{A}$ if $\mathcal{A}|\Omega\rangle$ is dense in $\mathcal{H}$ (cyclicity)
and $a|\Omega\rangle=0$ implies $a=0$ for $a\in\mathcal{A}$ (separating).  For an
inclusion $\mathcal{B}\subset\mathcal{A}$, a vector that is simultaneously cyclic and
separating for both algebras is called a \emph{common cyclic and separating vector}.

The first question is whether such a vector exists at all.  For type~$III$ algebras,
the answer is affirmative and robust, by the following two results.

\begin{lemma}[Blackadar]\label{lem:typeIII-standard}
Every type~$III$ von Neumann algebra acting on a separable Hilbert space is in
standard form.
\end{lemma}

\noindent Since both $\mathcal{A}$ and $\mathcal{B}$ are in standard form, their individual
sets of cyclic and separating vectors, $\mathcal{C}_\mathcal{A}$ and $\mathcal{C}_\mathcal{B}$,
are non-empty.  The non-trivial question is whether their intersection
$\mathcal{C}_\mathcal{A}\cap\mathcal{C}_\mathcal{B}$ is non-empty and, if so, how large it is.

\begin{proposition}[Dixmier--Mar\'echal]\label{prop:dense-common}
For two type~$III$ von Neumann algebras $\mathcal{A}$, $\mathcal{B}$ acting on a common
separable Hilbert space $\mathcal{H}$, the set of common cyclic and separating vectors
is a dense $G_\delta$ subset of $\mathcal{H}$.
\end{proposition}

\noindent The density of common cyclic and separating vectors has a strong implication for
entanglement.  Fix any such vector $|\Omega\rangle\in\mathcal{C}_\mathcal{A}\cap\mathcal{C}_\mathcal{B}$
and consider the inclusion $\mathcal{N}\subset\mathcal{M}$.  The Tomita operator satisfies
\begin{equation}
  S_\mathcal{M}\,\mathcal{O}_\mathcal{N}|\Omega\rangle
  = J_\mathcal{M}\Delta_\mathcal{M}^{1/2}\mathcal{O}_\mathcal{N}|\Omega\rangle
  = \mathcal{O}_\mathcal{N}^*|\Omega\rangle
  = J_\mathcal{N}\Delta_\mathcal{N}^{1/2}\mathcal{O}_\mathcal{N}|\Omega\rangle
  = S_\mathcal{N}\,\mathcal{O}_\mathcal{N}|\Omega\rangle
  \label{eq:Tomita-extension}
\end{equation}
for every $\mathcal{O}_\mathcal{N}\in\mathcal{N}$.  Since $S_\mathcal{M}$ is an extension
of $S_\mathcal{N}$, the associated closed quadratic forms give the form inequality:
\begin{equation}
  \Delta_\mathcal{N} \geq \Delta_\mathcal{M}.
  \label{eq:modular-inequality}
\end{equation}
Because $f(z)=z^\alpha$ is operator monotone for $0\leq\alpha\leq 1$, this extends to
\begin{equation}
  \Delta_\mathcal{N}^\alpha \geq \Delta_\mathcal{M}^\alpha, \qquad 0\leq\alpha\leq 1.
  \label{eq:modular-inequality-alpha}
\end{equation}
For $0\leq\alpha\leq1/2$, the quadratic-form order gives $\mathcal{D}(\Delta_\mathcal{N}^\alpha)\subseteq\mathcal{D}(\Delta_\mathcal{M}^\alpha)$,
so the operator $\Delta_\mathcal{M}^\alpha\Delta_\mathcal{N}^{-\alpha}$ is well-defined
and closable for $0\leq\alpha\leq 1/2$, and its closure is bounded:
\begin{equation}
  \|\Delta_\mathcal{M}^\alpha\Delta_\mathcal{N}^{-\alpha}\|^2
  = \|\Delta_\mathcal{N}^{-\alpha}\Delta_\mathcal{M}^{2\alpha}\Delta_\mathcal{N}^{-\alpha}\|
  \leq 1.
  \label{eq:characteristic-bounded}
\end{equation}
This boundedness, which uses $2\alpha\leq 1$, is the foundation for the characteristic
function of the inclusion.


For a common cyclic and separating vector, the projection onto
$\overline{\mathcal N\Omega}$ is $e_{\mathcal N}=1$.
This is a statement about Hilbert-space cyclicity, not universal entanglement
of every normal state or the index of the inclusion.

\paragraph{Proper inclusions and reference-state preservation.}
The following result concerns preservation of this particular vector state.

\begin{proposition}\label{prop:common-css-no-invariance}
Let $\mathcal{N}\subset\mathcal{M}$ be an inclusion of factors with a common cyclic and
separating vector $|\Omega\rangle$.  If $\sigma_t^{\omega_\mathcal{M}}(\mathcal{N})=\mathcal{N}$
for all $t\in\mathbb{R}$, then $\mathcal{N}=\mathcal{M}$.
\end{proposition}


\begin{proof}
Takesaki's theorem gives an $\Omega$-preserving normal expectation
$E:\mathcal M\to\mathcal N$. Its GNS projection is $e_{\mathcal N}=1$ by common
cyclicity. Thus $E(m)\Omega=e_{\mathcal N}m\Omega=m\Omega$.
Separatingness for $\mathcal M$ gives $E(m)=m$ and $\mathcal N=\mathcal M$.
\end{proof}

For a proper inclusion this rules out expectations preserving $\Omega$.
It does not exclude expectations preserving another state, finite index,
or error correction for a specified channel and state family.
Generation \eqref{eq:topological-transitivity} is compatible with this distinction.


\paragraph{Split inclusions and statistical independence.}
The split property concerns the commuting algebras $\mathcal N$ and $\mathcal M'$,
whereas conormal factorization concerns $\mathcal N$ and $\mathcal N'\cap\mathcal M$.
It is a condition of statistical independence, separate from decay under a chosen dynamics.

\begin{definition}[Split inclusion]\label{def:split}
An inclusion $\mathcal{N}\subset\mathcal{M}$ is \emph{split} if there exists a type~$I$
factor $\mathcal{F}$ such that $\mathcal{N}\subset\mathcal{F}\subset\mathcal{M}$.
\end{definition}


\noindent Physically, the intermediate type $I$ factor supplies a tensor-product
representation of $\mathcal N\vee\mathcal M'$~\cite{standardinclusion}.
In QFT, the usual split property applies to nested bounded regions separated
by a positive collar. It neither asserts that the vacuum is a product state
nor supplies an exponential correlation decay rate.

For a type $III$ range, the split property rules out every normal expectation:

% TODO: Verify the original attribution "Theorem 30 in ChLaLe2026".
\begin{theorem}[Split inclusions and normal conditional expectations]\label{thm:split-no-CE}
If $\mathcal{N}$ has a non-zero type~$III$ summand and $\mathcal{N}\subset\mathcal{M}$
is split, then no normal conditional expectation $E:\mathcal{M}\to\mathcal{N}$ exists.
\end{theorem}

\begin{proof}[Proof sketch]
Let $\mathcal{N}\subset\mathcal{F}\subset\mathcal{M}$ be the interpolating type~$I$
factor, and suppose $E:\mathcal{M}\to\mathcal{N}$ is normal.  Its restriction
$E|_\mathcal{F}:\mathcal{F}\to\mathcal{N}$ is a normal norm-one projection with range
$\mathcal{N}$.  Tomiyama's theorem then implies that $\mathcal{N}$ must be type~$I$ as a
normally expected subalgebra of the type~$I$ factor $\mathcal{F}$, contradicting the
hypothesis that $\mathcal{N}$ has a non-zero type~$III$ summand.
\end{proof}

\noindent Hence the logical implication is:
\begin{equation}
  \text{split inclusion with type~}III\text{ range}
  \;\Longrightarrow\;
  \text{no normal conditional expectation.}
  \label{eq:split-no-CE}
\end{equation}


The distinctions relevant to this work can be recorded without identifying
the split property with mixing or common cyclicity with the classical $K$-property:
\begin{center}
\begin{tabular}{ll}
 \textit{Property} & \textit{Modular formulation}\\[4pt]
 Ergodicity & $\mathcal M^\omega=\mathbb C1$\\
 Mixing & $C_{a,b}^\omega(t)\to0$ for all $a,b\in\mathcal M$\\
 Generation & $\bigvee_t\sigma_t^\omega(\mathcal N)=\mathcal M$\\
 Half-sided nesting & $\sigma_t^\omega(\mathcal N)\subseteq\mathcal N$ for one sign of $t$
\end{tabular}
\end{center}
The classical $K$-property additionally requires a trivial tail. Section~\ref{subsec:geometry-modular-ergodic}
discusses the extra geometric structure supplied by half-sided modular inclusions.

\subsection{Finite-dimensional coarse-grainings}\label{app:finite-coarse}

The tensor-product example \eqref{condexp} is one of three useful finite-dimensional
inclusions. Here the ambient algebra is $B(\mathcal H)$.
\begin{enumerate}
 \item {\bf Subspace inclusions:} Let $\mathcal H=P\mathcal H\oplus(1-P)\mathcal H$.
 The retained unital algebra is
 \[
  \mathcal B_P=P B(\mathcal H)P\oplus\mathbb C(1-P).
 \]
 Its center is $\mathbb CP\oplus\mathbb C(1-P)$. For a fixed state $\chi_{1-P}$
 on the complementary corner, the expectation is
 \[
  E_{\chi_{1-P}}(a)=PaP+
  \chi_{1-P}\bigl((1-P)a(1-P)\bigr)(1-P).
 \]
 It preserves precisely the block-diagonal states whose complementary block
 is proportional to the density matrix of $\chi_{1-P}$.
 Block diagonality alone does not suffice for a fixed expectation.
 For $P_i\leq P_{i+1}$, $\mathcal B_{P_i}\subseteq\mathcal B_{P_{i+1}}$.
 \item {\bf Subsystem inclusions:} If $\mathcal H=\mathcal H_B\otimes\mathcal H_R$,
 the retained algebra is $B(\mathcal H_B)\otimes1_R$.
 Its expectations and preserved states are those of \eqref{condexp}.
 \item {\bf Mixed inclusions:} A general unital finite-dimensional subalgebra has
 \[
  \mathcal H=\bigoplus_a\mathcal H_B^{(a)}\otimes\mathcal H_R^{(a)},\qquad
  \mathcal B=\bigoplus_a B(\mathcal H_B^{(a)})\otimes1_R^{(a)}.
 \]
 The center records the block label $a$. Bimodularity removes off-diagonal blocks,
 and a choice of normalized hidden state $\chi_R^{(a)}$ specifies a slice
 expectation in each block. The preserved states are
 \[
  \rho=\bigoplus_a p_a\,\rho_B^{(a)}\otimes\chi_R^{(a)}.
 \]
 With the label stored in a classical register $Z$, these have $I(B:R|Z)=0$.
 They maximize entropy at fixed retained state only when each hidden reference
 is maximally mixed; for a general reference, they are fixed states of that expectation.
\end{enumerate}
The subspace construction resolves alternatives among microstates, while a tensor
subsystem retains simultaneously specified degrees of freedom. The same distinction
organizes the classical constructions in Appendix~\ref{app:classical}.

\subsection{Compressed modular comparison}\label{app:compressed-comparison}

Let $\Omega$ be cyclic and separating for $\mathcal A$, and set
$\mathcal H_B=\overline{\mathcal B\Omega}$ for a unital inclusion $\mathcal B\subset\mathcal A$.
The restricted state is faithful, and its modular operator $\Delta_B$ acts on
$\mathcal H_B$. With $V:\mathcal H_B\hookrightarrow\mathcal H_A$ and $e_B=VV^*$,
define the {\it compressed comparison}
\begin{equation}
 C_{B\subset A}(t)=V^*\Delta_A^{-it}V\Delta_B^{it},
 \qquad t\in\mathbb R.
\end{equation}
This is a strongly continuous contraction on $\mathcal H_B$ fixing $\Omega$.
For $e_B=1$ it is the characteristic operator $D_{B\subset A}(t)$.
If an $\Omega$-preserving expectation exists, modular intertwining gives
$C_{B\subset A}(t)=1_{\mathcal H_B}$.

To display the compression algebraically, extend the intrinsic unitary to
$U_B(t)=V\Delta_B^{it}V^*+(1-e_B)$ on $\mathcal H_A$ and set
$u_t=\Delta_A^{-it}U_B(t)$, $\widehat C(t)=e_Bu_te_B=VC(t)V^*$.
Writing $\widetilde\sigma_B^{-t}=\operatorname{Ad}U_B(-t)$ gives
\begin{align}
 \widehat C(t+s)&=\widehat C(t)\widetilde\sigma_B^{-t}(\widehat C(s))
                  +\mathcal L(t,s),\nn\\
 \mathcal L(t,s)&=e_Bu_t(1-e_B)\widetilde\sigma_B^{-t}(u_s)e_B.
\end{align}
The second term records passage through the orthogonal complement. This algebraic
identity does not itself specify an open-system master equation.
On vectors $\psi\in\operatorname{Dom}K_B$ with $V\psi\in\operatorname{Dom}K_A$,
\begin{equation}
 -i\left.\partial_t C_{B\subset A}(t)\psi\right|_{t=0}
 =(V^*K_AV-K_B)\psi,\qquad K_C=-\log\Delta_C.
\end{equation}
This compares modular generators on the stated domain, rather than asserting
a strong derivative on the full Hilbert space or a dissipative generator.
The contraction is a Hilbert-space operator. Conjugation
$a\mapsto C(t)^*aC(t)$ is CP on $B(\mathcal H_B)$ but is unital only if
$C(t)^*C(t)=1$; preservation of $\mathcal B$ would require a further condition.
An interpretation as a quantum channel therefore needs additional structure.

\bibliographystyle{unsrt}
\bibliography{modularscattering}

\end{document}
)